% Options for packages loaded elsewhere
\PassOptionsToPackage{unicode}{hyperref}
\PassOptionsToPackage{hyphens}{url}
\PassOptionsToPackage{dvipsnames,svgnames,x11names}{xcolor}
\documentclass[
  10pt,
]{article}
\usepackage{xcolor}
\usepackage[margin=1cm,top=1cm,bottom=1cm,left=1cm,right=1cm,includeheadfoot]{geometry}
\usepackage{amsmath,amssymb}
\setcounter{secnumdepth}{3}
\usepackage{iftex}
\ifPDFTeX
  \usepackage[T1]{fontenc}
  \usepackage[utf8]{inputenc}
  \usepackage{textcomp} % provide euro and other symbols
\else % if luatex or xetex
  \usepackage{unicode-math} % this also loads fontspec
  \defaultfontfeatures{Scale=MatchLowercase}
  \defaultfontfeatures[\rmfamily]{Ligatures=TeX,Scale=1}
\fi
\usepackage{lmodern}
\ifPDFTeX\else
  % xetex/luatex font selection
  \setmainfont[]{DejaVu Serif}
  \setmonofont[]{DejaVu Sans Mono}
\fi
% Use upquote if available, for straight quotes in verbatim environments
\IfFileExists{upquote.sty}{\usepackage{upquote}}{}
\IfFileExists{microtype.sty}{% use microtype if available
  \usepackage[]{microtype}
  \UseMicrotypeSet[protrusion]{basicmath} % disable protrusion for tt fonts
}{}
\usepackage{setspace}
\makeatletter
\@ifundefined{KOMAClassName}{% if non-KOMA class
  \IfFileExists{parskip.sty}{%
    \usepackage{parskip}
  }{% else
    \setlength{\parindent}{0pt}
    \setlength{\parskip}{6pt plus 2pt minus 1pt}}
}{% if KOMA class
  \KOMAoptions{parskip=half}}
\makeatother
\usepackage{listings}
\newcommand{\passthrough}[1]{#1}
\lstset{defaultdialect=[5.3]Lua}
\lstset{defaultdialect=[x86masm]Assembler}
\usepackage{longtable,booktabs,array}
\usepackage{caption}
\captionsetup[table]{skip=6pt}
\newcounter{none} % for unnumbered tables
\usepackage{calc} % for calculating minipage widths
% Correct order of tables after \paragraph or \subparagraph
\usepackage{etoolbox}
\makeatletter
\patchcmd\longtable{\par}{\if@noskipsec\mbox{}\fi\par}{}{}
\makeatother
\usepackage{graphicx}
\makeatletter
\newsavebox\pandoc@box
\newcommand*\pandocbounded[1]{% scales image to fit in text height/width
  \sbox\pandoc@box{#1}%
  \Gscale@div\@tempa{\textheight}{\dimexpr\ht\pandoc@box+\dp\pandoc@box\relax}%
  \Gscale@div\@tempb{\linewidth}{\wd\pandoc@box}%
  \ifdim\@tempb\p@<\@tempa\p@\let\@tempa\@tempb\fi% select the smaller of both
  \ifdim\@tempa\p@<\p@\scalebox{\@tempa}{\usebox\pandoc@box}%
  \else\usebox{\pandoc@box}%
  \fi%
}
% Set default figure placement to htbp
\def\fps@figure{htbp}
\makeatother
\setlength{\emergencystretch}{3em} % prevent overfull lines
\providecommand{\tightlist}{%
  \setlength{\itemsep}{0pt}\setlength{\parskip}{0pt}}
% Enable graphics inclusion and ensure figure numbering works
\usepackage{graphicx}
\renewcommand{\figurename}{Figure}

% Configure fonts for Unicode support with fallbacks
\usepackage{newunicodechar}
\newunicodechar{⁴}{\textsuperscript{4}}
\newunicodechar{₄}{\textsubscript{4}}

% Math notation macros
\newcommand{\norm}[1]{\lVert #1\rVert}
\newcommand{\abs}[1]{\lvert #1\rvert}

% Enhanced code block styling for better contrast and readability
\usepackage{fancyvrb}
\usepackage{xcolor}
\usepackage{listings}

% Define custom colors for code blocks
\definecolor{codebg}{RGB}{245, 245, 245}      % Light gray background
\definecolor{codeborder}{RGB}{200, 200, 200}  % Medium gray border
\definecolor{codefg}{RGB}{50, 50, 50}         % Dark gray text

% Configure Verbatim environment for inline code
\DefineVerbatimEnvironment{Verbatim}{Verbatim}{%
    fontsize=\small,
    frame=single,
    framerule=0.5pt,
    framesep=3pt,
    rulecolor=\color{codeborder},
    bgcolor=\color{codebg},
    fgcolor=\color{codefg}
}

% Configure code block styling
\DefineVerbatimEnvironment{Highlighting}{Verbatim}{%
    fontsize=\footnotesize,
    frame=single,
    framerule=0.5pt,
    framesep=5pt,
    rulecolor=\color{codeborder},
    bgcolor=\color{codebg},
    fgcolor=\color{codefg}
}

% Style inline code with \texttt
\renewcommand{\texttt}[1]{%
    \colorbox{codebg}{\color{codefg}\ttfamily #1}%
}

% Configure listings package for code blocks
\lstset{
    backgroundcolor=\color{codebg},
    basicstyle=\footnotesize\ttfamily\color{codefg},
    breakatwhitespace=false,
    breaklines=true,
    captionpos=b,
    commentstyle=\color{codefg},
    deletekeywords={...},
    escapeinside={\%*}{*)},
    extendedchars=true,
    frame=single,
    framerule=0.5pt,
    framesep=5pt,
    keepspaces=true,
    keywordstyle=\color{codefg},
    language=Python,
    morekeywords={*,...},
    numbers=left,
    numbersep=5pt,
    numberstyle=\tiny\color{codefg},
    rulecolor=\color{codeborder},
    showspaces=false,
    showstringspaces=false,
    showtabs=false,
    stepnumber=1,
    stringstyle=\color{codefg},
    tabsize=2,
    title=\lstname
}

% Override any Pandoc default lstset configurations
\AtBeginDocument{
    \lstset{
        backgroundcolor=\color{codebg},
        basicstyle=\footnotesize\ttfamily\color{codefg},
        frame=single,
        framerule=0.5pt,
        framesep=5pt,
        rulecolor=\color{codeborder},
        numbers=left,
        numbersep=5pt,
        numberstyle=\tiny\color{codefg}
    }
}

% Configure hyperref colors consistently
\AtBeginDocument{
% Override pandoc's hidelinks setting with consistent options
\hypersetup{
    colorlinks=true,
    allcolors=red,
    linkcolor=red,
    urlcolor=red,
    citecolor=red,
    filecolor=red,
    menucolor=red,
    linktoc=all
}
}

% Simple page break support for document structure
\usepackage{bookmark}
\IfFileExists{xurl.sty}{\usepackage{xurl}}{} % add URL line breaks if available
\urlstyle{same}
% fallback for those not using the hyperref driver hyperxmp:
\makeatletter
\@ifundefined{xmpquote}{\newcommand{\xmpquote}[1]{#1}}{}
\makeatother
\hypersetup{
  colorlinks=true,
  linkcolor={red},
  filecolor={red},
  citecolor={red},
  urlcolor={red},
  pdfcreator={LaTeX via pandoc}}

\title{QuadMath: An Analytical Review of 4D and Quadray Coordinates}
\author{Daniel Ari Friedman\\ ORCID: 0000-0001-6232-9096\\ Email: daniel@activeinference.institute\\ DOI: 10.5281/zenodo.16887791}
\date{October 07, 2026}

\begin{document}
\maketitle

{
\hypersetup{linkcolor=black}
\setcounter{tocdepth}{3}
\tableofcontents
}
\setstretch{1.0}
\section{Introduction}\label{introduction}

\subsection{Abstract}\label{abstract}

We review a unified analytical framework for four-dimensional (4D)
modeling with Quadray coordinates, synthesizing geometric foundations,
optimization on tetrahedral lattices, and information geometry. Building
on R. Buckminster Fuller's
\href{https://en.wikipedia.org/wiki/Synergetics_(Fuller)}{Synergetics}
and Kirby Urner's computational implementations in the
\href{https://github.com/4dsolutions}{4dsolutions ecosystem} (Python,
Rust, Clojure, and POV-Ray), we show that integer lattice constraints
force the volume of any tetrahedron with lattice vertices to be an
integer multiple of the unit tetrahedron's volume --- a quantization
into discrete ``energy levels'' that regularizes optimization on the
tetrahedral lattice. We adapt the
\href{https://en.wikipedia.org/wiki/Nelder\%E2\%80\%93Mead_method}{Nelder--Mead
method} to the Quadray lattice, define
\href{https://en.wikipedia.org/wiki/Fisher_information}{Fisher
information} in Quadray parameter space, and analyze optimization as
geodesic motion on an information manifold via the
\href{https://en.wikipedia.org/wiki/Natural_gradient}{natural gradient}.
We distinguish three 4D namespaces --- Coxeter.4D (Euclidean E⁴),
Einstein.4D (Minkowski spacetime), and Fuller.4D (synergetics/Quadrays)
--- and map each to the analytical tools it supports. The result is a
cohesive, interpretable approach for robust, geometry-grounded
computation in 4D. All source code for the manuscript is available at
\href{https://github.com/docxology/quadmath}{QuadMath}.

\textbf{Keywords}: Quadray coordinates, 4D geometry, tetrahedral
lattice, integer volume quantization, information geometry,
optimization, synergetics, active inference.

\subsection{Overview}\label{overview}

Quadray coordinates provide a tetrahedral basis for modeling space and
computation, standing in contrast to Cartesian cubic frameworks.
Originating in Buckminster Fuller's Synergetics, Quadray coordinates
replace right-angle orthonormal assumptions with 60-degree coordination
and a unit tetrahedron of volume 1. This reframing yields striking
integer relationships among common polyhedra and provides a natural
account of space via close-packed spheres and the isotropic vector
matrix (IVM).

This paper unifies three threads:

\begin{itemize}
\tightlist
\item
  \textbf{Foundations}: Quadray coordinates and their relation to 4D
  modeling more generally, with explicit namespace usage (Coxeter.4D,
  Einstein.4D, Fuller.4D) to maintain clarity.
\item
  \textbf{Optimization framework}: Leverages integer volume quantization
  on tetrahedral lattices to achieve robust, discrete convergence.
\item
  \textbf{Information geometry}: Tools (e.g., Fisher Information,
  free-energy minimization) for interpreting optimization as geodesic
  motion on statistical manifolds.
\end{itemize}

\subsection{4D Namespace Framework}\label{d-namespace-framework}

In this synthetic review, we distinguish three internal meanings of
``4D,'' following a dot-notation that avoids cross-domain confusion. For
comprehensive details, see \href{02_4d_namespaces.md}{Section 2: 4D
Namespaces}.

\begin{itemize}
\tightlist
\item
  \textbf{Coxeter.4D} --- four-dimensional Euclidean space (E⁴) with
  mutually orthogonal axes and a positive-definite metric: the setting
  of classical polytope theory, and explicitly not spacetime (Coxeter,
  Regular Polytopes, Dover ed., p.~119). Lattice packings in four
  dimensions align with the treatment in Conway \& Sloane's
  \href{https://link.springer.com/book/10.1007/978-1-4757-6568-7}{Sphere
  Packings, Lattices and Groups}.
\item
  \textbf{Einstein.4D} --- Minkowski spacetime: three space dimensions
  plus time, with an indefinite metric of signature \((-,+,+,+)\); the
  arena of relativistic physics, distinct from Euclidean E⁴.
\item
  \textbf{Fuller.4D} --- synergetics' tetrahedral accounting of space:
  Quadray coordinates (four non-negative coordinates with at least one
  zero after normalization) on the Isotropic Vector Matrix (IVM) = Cubic
  Close Packing (CCP) = Face-Centered Cubic (FCC) lattice, with the
  regular tetrahedron as the natural unit container; it emphasizes
  angle/shape relations independent of time and energy.
\end{itemize}

\subsection{Contributions}\label{contributions}

The paper makes the following key contributions:

\begin{itemize}
\tightlist
\item
  \textbf{Namespaces mapping}: Coxeter.4D (Euclidean E⁴), Einstein.4D
  (Minkowski spacetime), and Fuller.4D (Quadrays/IVM) → analytical tools
  and examples.
\item
  \textbf{Quadray-adapted Nelder--Mead}: Integer-lattice normalization
  and volume-level tracking.
\item
  \textbf{Equations and methods}: Comprehensive supplement with guidance
  for high-precision computation using
  \passthrough{\lstinline!libquadmath!}.
\item
  \textbf{Discrete optimizer}: Integer-valued variational descent over
  the IVM (\passthrough{\lstinline!discrete\_ivm\_descent!}) with
  animation tooling, connecting lattice geometry to
  information-theoretic objectives.
\end{itemize}

\subsection{Manuscript Structure}\label{manuscript-structure}

\begin{itemize}
\tightlist
\item
  \textbf{Introduction}: motivates Quadrays, clarifies 4D namespaces,
  and summarizes contributions.
\item
  \textbf{Methods}: details coordinate conventions, exact tetravolumes,
  conversions, and lattice-aware optimization methods (Nelder--Mead and
  discrete IVM descent).
\item
  \textbf{Results}: empirical comparisons and demonstrations are shown
  inline and saved under \passthrough{\lstinline!quadmath/output/!}
  (PNG/CSV/NPZ/MP4) for reproducibility.
\item
  \textbf{Discussion}: interprets results, limitations, and
  implications; outlines future work.
\item
  \textbf{Appendices}: equations, free-energy background, and a
  consolidated symbols/glossary with an auto-generated API index.
\end{itemize}

\subsection{Companion Code and Tests}\label{companion-code-and-tests}

The manuscript is accompanied by a fully-tested Python codebase under
\passthrough{\lstinline!src/!} with unit tests under
\passthrough{\lstinline!tests/!}. Key artifacts used throughout the
paper:

\begin{itemize}
\tightlist
\item
  \textbf{Quadray APIs}:
  \passthrough{\lstinline!src/quadmath/core/quadray.py!}
  (\passthrough{\lstinline!Quadray!},
  \passthrough{\lstinline!integer\_tetra\_volume!},
  \passthrough{\lstinline!ace\_tetravolume\_5x5!}).
\item
  \textbf{Determinant utilities}:
  \passthrough{\lstinline!src/quadmath/core/linalg\_utils.py!}
  (\passthrough{\lstinline!bareiss\_determinant\_int!}).
\item
  \textbf{Length-based volume}:
  \passthrough{\lstinline!src/quadmath/core/cayley\_menger.py!}
  (\passthrough{\lstinline!tetra\_volume\_cayley\_menger!},
  \passthrough{\lstinline!ivm\_tetra\_volume\_cayley\_menger!}).
\item
  \textbf{XYZ conversion}:
  \passthrough{\lstinline!src/quadmath/lattice/conversions.py!}
  (\passthrough{\lstinline!urner\_embedding!},
  \passthrough{\lstinline!quadray\_to\_xyz!}).
\item
  \textbf{Examples}:
  \passthrough{\lstinline!src/quadmath/core/examples.py!}
  (\passthrough{\lstinline!example\_ivm\_neighbors!},
  \passthrough{\lstinline!example\_volume!},
  \passthrough{\lstinline!example\_optimize!}).
\end{itemize}

For comprehensive background resources, computational implementations,
and related work, see the \href{07_resources.md}{Resources} section.

\subsection{Reproducibility and Data
Availability}\label{reproducibility-and-data-availability}

\begin{itemize}
\tightlist
\item
  The manuscript Markdown and code to generate the PDF are available on
  the project repository (\passthrough{\lstinline!QuadMath!} on GitHub,
  \passthrough{\lstinline!@docxology!} username). See the repository
  home page for source, figures, and scripts:
  \href{https://github.com/docxology/QuadMath}{QuadMath repository}. The
  repository is also archived on
  \href{https://zenodo.org/records/16887791}{Zenodo record 16887791}
  with DOI 10.5281/zenodo.16887791.
\item
  The manuscript is licensed under the Apache License 2.0. See the
  \href{../../LICENSE}{LICENSE} file for details.
\item
  The manuscript is accompanied by a fully-tested Python codebase under
  \passthrough{\lstinline!src/!} with unit tests under
  \passthrough{\lstinline!tests/!}, complemented by extensive
  cross-validation against Kirby Urner's reference implementations in
  the \href{https://github.com/4dsolutions}{4dsolutions ecosystem}. See
  the \href{07_resources.md}{Resources} section for comprehensive
  details on computational implementations and validation.
\item
  All figures referenced in the manuscript are generated by scripts
  under \passthrough{\lstinline!quadmath/scripts/!} and saved to
  \passthrough{\lstinline!quadmath/output/!} with lightweight CSV/NPZ
  alongside images.
\item
  Tests accompany all methods under \passthrough{\lstinline!src/!} and
  enforce 100\% coverage for \passthrough{\lstinline!src/!}.
\item
  Symbols and notation are standardized across sections; see
  \href{10_symbols_glossary.md}{Appendix: Symbols and Glossary} for a
  consolidated table of variables and constants used throughout.
  Equation labels (e.g., Eq. \eqref{eq:lattice_det} and Eq.
  \eqref{eq:fim}) and figure labels are automatically numbered by LaTeX
  for consistent cross-referencing.
\item
  The manuscript is a work in progress and will be updated as the
  project progresses. There may be errors and missing references, check
  all methods and equations for consistency.
\end{itemize}

\subsection{Graphical Abstract}\label{graphical-abstract}

\textbf{Panel A} shows the four Quadray axes (A, B, C, D) under the
default symmetric embedding with a tetrahedron wireframe. \textbf{Panel
B} shows the same vertices as close-packed spheres (IVM = CCP = FCC)
sized to kiss along the tetrahedron edges.

\begin{figure}
\centering
\pandocbounded{\includegraphics[keepaspectratio,alt={Quadray coordinate system overview (graphical abstract). Panel A: The four Quadray axes --- rays A, B, C, D from the origin to the vertices of a reference regular tetrahedron under the default symmetric embedding --- drawn as colored arrows (A=blue, B=orange, C=green, D=red) with axis labels at the vertex endpoints and a light gray wireframe joining the four vertices. Notice that the axes are four canonical directions (spokes) separated by 60° angles, not orthogonal Cartesian dimensions --- the defining ``directions, not dimensions'' character of Fuller.4D. Panel B: The same four vertices shown as close-packed spheres, each colored to match its Panel A axis and with radius equal to half the minimum edge length, so neighboring spheres kiss exactly along the tetrahedron edges; a light black wireframe marks those edges. This is the local building block of the Isotropic Vector Matrix (IVM) = Cubic Close Packing (CCP) = Face-Centered Cubic (FCC) correspondence; extending the kissing arrangement across the lattice yields the ``twelve around one'' coordination central to synergetics. What to notice: one and the same tetrahedral skeleton underlies both the direction-based Quadray coordinate system and the densest sphere packing in 3D, which is why integer Quadray coordinates and integer sphere counts (and hence integer volumes) live on a single lattice. Generated by quadmath/scripts/graphical\_abstract\_quadray.py (output at quadmath/output/figures/graphical\_abstract\_quadray.png).}]{../output/figures/graphical_abstract_quadray.png}}
\caption{\textbf{Quadray coordinate system overview (graphical
abstract)}. \textbf{Panel A}: The four Quadray axes --- rays A, B, C, D
from the origin to the vertices of a reference regular tetrahedron under
the default symmetric embedding --- drawn as colored arrows (A=blue,
B=orange, C=green, D=red) with axis labels at the vertex endpoints and a
light gray wireframe joining the four vertices. Notice that the axes are
four canonical directions (spokes) separated by 60° angles, not
orthogonal Cartesian dimensions --- the defining ``directions, not
dimensions'' character of Fuller.4D. \textbf{Panel B}: The same four
vertices shown as close-packed spheres, each colored to match its Panel
A axis and with radius equal to half the minimum edge length, so
neighboring spheres kiss exactly along the tetrahedron edges; a light
black wireframe marks those edges. This is the local building block of
the Isotropic Vector Matrix (IVM) = Cubic Close Packing (CCP) =
Face-Centered Cubic (FCC) correspondence; extending the kissing
arrangement across the lattice yields the ``twelve around one''
coordination central to synergetics. What to notice: one and the same
tetrahedral skeleton underlies both the direction-based Quadray
coordinate system and the densest sphere packing in 3D, which is why
integer Quadray coordinates and integer sphere counts (and hence integer
volumes) live on a single lattice. Generated by
\protect\passthrough{\lstinline!quadmath/scripts/graphical\_abstract\_quadray.py!}
(output at
\protect\passthrough{\lstinline!quadmath/output/figures/graphical\_abstract\_quadray.png!}).}
\end{figure}

\newpage

\section{4D Namespaces: Coxeter.4D, Einstein.4D,
Fuller.4D}\label{d-namespaces-coxeter.4d-einstein.4d-fuller.4d}

This section provides the definitive reference for the three 4D
frameworks used throughout this manuscript. Each namespace represents a
distinct mathematical framework with specific applications in our
Quadray-based computational system.

\subsection{Coxeter.4D (Euclidean
E⁴)}\label{coxeter.4d-euclidean-eux2074}

\textbf{Definition}: Four-dimensional Euclidean space E⁴ with mutually
orthogonal axes and a positive-definite metric --- the native setting of
classical regular polytopes, and, per Coxeter (Regular Polytopes, Dover
ed., p.~119), explicitly not spacetime. Lattice/packing discussions
connect to Conway \& Sloane's systematic treatment of higher-dimensional
sphere packings and lattices
(\href{https://link.springer.com/book/10.1007/978-1-4757-6568-7}{Sphere
Packings, Lattices and Groups (Springer)}).

\textbf{Usage}: Embed Quadray configurations or compare alternative
parameterizations when a strictly Euclidean 4D setting is desired.

\textbf{Simplexes}: Simplex structures extend naturally to 4D and beyond
(e.g., pentachora).

\textbf{Mathematical context}: This framework is appropriate for
standard Euclidean geometry, including the Cayley--Menger determinant
for computing volumes from edge lengths.

\subsection{Einstein.4D (Relativistic
spacetime)}\label{einstein.4d-relativistic-spacetime}

\textbf{Definition}: Minkowski spacetime --- three space dimensions
joined to one time dimension by an indefinite metric of signature
\((-,+,+,+)\) (mostly-plus convention) --- the geometric arena of
special relativity; the sign of the squared line element classifies
separations as timelike, lightlike, or spacelike.

\textbf{Spacetime}: Minkowski metric signature.

\textbf{Line element} (mostly-plus convention; see
\href{https://en.wikipedia.org/wiki/Minkowski_space}{Minkowski space}):
see Eq. \eqref{eq:minkowski_line_element} in the equations appendix.

\textbf{Optimization analogy}: Metric-aware geodesics generalize to
information geometry where the Fisher metric replaces the physical
metric. See
\href{https://en.wikipedia.org/wiki/Fisher_information}{Fisher
information} and
\href{https://en.wikipedia.org/wiki/Natural_gradient}{natural gradient}.

\textbf{Important note}: This namespace is used ONLY as a
metric/geodesic analogy when discussing information geometry. Physical
constants G, c, Λ do not appear in Quadray lattice methods and should
not be mixed with IVM unit conventions.

\subsection{Fuller.4D (Synergetics /
Quadrays)}\label{fuller.4d-synergetics-quadrays}

\textbf{Definition}: The synergetic account of space built on Quadray
coordinates --- four non-negative components (A, B, C, D) giving
directions from the center of a reference regular tetrahedron to its
vertices --- together with the IVM = CCP = FCC lattice correspondence,
the regular tetrahedron as the unit of volume, and exact integer
tetravolumes for lattice tetrahedra; it tracks shape and angle relations
among containers, independent of time and energy.

\textbf{Basis}: Four non-negative components A,B,C,D with at least one
zero post-normalization, treated as a vector (direction and magnitude),
not merely a point. Overview:
\href{https://en.wikipedia.org/wiki/Quadray_coordinates}{Quadray
coordinates}.

\textbf{Geometry}: Tetrahedral; unit tetrahedron volume = 1; integer
lattice aligns with close-packed spheres (IVM). Background:
\href{https://en.wikipedia.org/wiki/Synergetics_(Fuller)}{Synergetics}.

\textbf{Distances}: Computed via appropriate projective normalization;
edges align with tetrahedral axes. The IVM = CCP = FCC shortcut allows
working in 3D embeddings for visualization while preserving the
underlying Fuller.4D tetrahedral accounting.

\textbf{Implementation heritage}: Extensive computational validation
through Kirby Urner's \href{https://github.com/4dsolutions}{4dsolutions
ecosystem}. See the \href{07_resources.md}{Resources} section for
comprehensive details on computational implementations and educational
materials.

\subsubsection{Directions, not dimensions (language and
models)}\label{directions-not-dimensions-language-and-models}

\textbf{Vector-first framing}: Treat Quadrays as four canonical
directions (``spokes'' to the vertices of a regular tetrahedron from its
center), not as four orthogonal dimensions. The methane molecule (CH₄)
and caltrop shape are helpful mental models.

\textbf{Origins outside Synergetics}: Quadrays did not originate with
Fuller; we adopt the coordinate system within the IVM context. See
\href{https://en.wikipedia.org/wiki/Quadray_coordinates}{Quadray
coordinates}.

\textbf{Language games}: Quadrays and Cartesian are parallel vector
languages on the same Euclidean container; teaching them together avoids
oscillating between ``points now, vectors later.''

\subsubsection{Figures}\label{figures}

\begin{figure}
\centering
\pandocbounded{\includegraphics[keepaspectratio,alt={IVM neighbors and coordination patterns (2×2 panel layout). Panel A: The twelve nearest IVM neighbors plotted as blue points in 3D space under the default embedding, showing the positions corresponding to permutations of the Quadray integer coordinates \{2,1,1,0\}. These points form the vertices of a cuboctahedron (vector equilibrium) centered at the origin with uniform radial distances. Panel B: The same neighbor points with radial edges (light lines) connecting each neighbor to the central origin, emphasizing the spoke-like radial symmetry and equal distances from center to shell. Panel C: Twelve-around-one close-packed spheres configuration where each neighbor position hosts a sphere with radius chosen so neighboring spheres kiss along cuboctahedron edges, illustrating the fundamental CCP/FCC/IVM correspondence. The central gray sphere represents the ``one'' in Fuller's ``twelve around one'' motif. Panel D: Adjacency graph showing strut connections (solid lines) between touching neighbor spheres, revealing the cuboctahedron's edge structure, plus light radial cables to the origin representing a stylized tensegrity interpretation of the vector equilibrium geometry.}]{../output/figures/ivm_neighbors_edges.png}}
\caption{\textbf{IVM neighbors and coordination patterns (2×2 panel
layout)}. \textbf{Panel A}: The twelve nearest IVM neighbors plotted as
blue points in 3D space under the default embedding, showing the
positions corresponding to permutations of the Quadray integer
coordinates \{2,1,1,0\}. These points form the vertices of a
cuboctahedron (vector equilibrium) centered at the origin with uniform
radial distances. \textbf{Panel B}: The same neighbor points with radial
edges (light lines) connecting each neighbor to the central origin,
emphasizing the spoke-like radial symmetry and equal distances from
center to shell. \textbf{Panel C}: Twelve-around-one close-packed
spheres configuration where each neighbor position hosts a sphere with
radius chosen so neighboring spheres kiss along cuboctahedron edges,
illustrating the fundamental CCP/FCC/IVM correspondence. The central
gray sphere represents the ``one'' in Fuller's ``twelve around one''
motif. \textbf{Panel D}: Adjacency graph showing strut connections
(solid lines) between touching neighbor spheres, revealing the
cuboctahedron's edge structure, plus light radial cables to the origin
representing a stylized tensegrity interpretation of the vector
equilibrium geometry.}
\end{figure}

\begin{figure}
\centering
\pandocbounded{\includegraphics[keepaspectratio,alt={Random Quadray point clouds under different embeddings (3-panel comparison). Each panel shows 200 randomly sampled integer Quadray coordinates with components in \{0,1,2,3,4,5\} projected to 3D space using different embedding matrices. Left panel (Default embedding): Points (blue) under the default symmetric embedding matrix showing the natural tetrahedral-symmetric distribution of normalized Quadrays in 3D space. Center panel (Scaled embedding, 0.75×): The same Quadray points (orange) under a uniformly scaled version of the default embedding, demonstrating how the point cloud structure scales proportionally while preserving relative geometries. Right panel (Urner embedding): The same points (purple) projected through the canonical Urner embedding matrix, illustrating how different linear mappings from Fuller.4D to Coxeter.4D (3D slice) affect the spatial distribution while preserving the underlying discrete lattice relationships. This comparison demonstrates the flexibility in choosing embeddings for visualization and analysis while maintaining the fundamental Quadray coordinate relationships.}]{../output/figures/quadray_clouds.png}}
\caption{\textbf{Random Quadray point clouds under different embeddings
(3-panel comparison)}. Each panel shows 200 randomly sampled integer
Quadray coordinates with components in \{0,1,2,3,4,5\} projected to 3D
space using different embedding matrices. \textbf{Left panel (Default
embedding)}: Points (blue) under the default symmetric embedding matrix
showing the natural tetrahedral-symmetric distribution of normalized
Quadrays in 3D space. \textbf{Center panel (Scaled embedding, 0.75×)}:
The same Quadray points (orange) under a uniformly scaled version of the
default embedding, demonstrating how the point cloud structure scales
proportionally while preserving relative geometries. \textbf{Right panel
(Urner embedding)}: The same points (purple) projected through the
canonical Urner embedding matrix, illustrating how different linear
mappings from Fuller.4D to Coxeter.4D (3D slice) affect the spatial
distribution while preserving the underlying discrete lattice
relationships. This comparison demonstrates the flexibility in choosing
embeddings for visualization and analysis while maintaining the
fundamental Quadray coordinate relationships.}
\end{figure}

In the previous figure, we show the twelve nearest IVM neighbors with
coordination patterns and vector equilibrium geometry; the current
figure illustrates random Quadray clouds under several embeddings.

\textbf{Vector equilibrium (cuboctahedron)}: The shell formed by the 12
nearest IVM neighbors is the cuboctahedron, also called the vector
equilibrium in synergetics. All 12 vertices are equidistant from the
origin with equal edge lengths, modeling a balanced local packing. This
geometry underlies the ``twelve around one'' close-packing motif and
appears in tensegrity discussions as a canonical balanced structure. See
background:
\href{https://en.wikipedia.org/wiki/Cuboctahedron}{Cuboctahedron (vector
equilibrium)} and synergetics references. Computational demonstrations
include related visualizations in the 4dsolutions ecosystem. See the
\href{07_resources.md}{Resources} section for comprehensive details.

\subsubsection{Clarifying remarks}\label{clarifying-remarks}

``A time machine is not a tesseract.''
\href{https://groups.io/g/synergeo/topic/my_take_on_close_pack/114531919}{KU
on synergeo} The tesseract is a Euclidean 4D object (Coxeter.4D), while
Minkowski spacetime (Einstein.4D) is indefinite and not Euclidean;
conflating the two leads to category errors. Fuller.4D, in turn, is a
tetrahedral, mereological framing of ordinary space emphasizing
shape/angle relations and IVM quantization. Each namespace carries
distinct assumptions and should be used accordingly in analysis.

\subsection{Practical usage guide}\label{practical-usage-guide}

\begin{itemize}
\tightlist
\item
  Use \textbf{Fuller.4D} when working with Quadrays, integer
  tetravolumes, and IVM neighbors (native lattice calculations).
\item
  Use \textbf{Coxeter.4D} for Euclidean length-based formulas,
  higher-dimensional polytopes, or comparisons in E⁴ (including
  Cayley--Menger).
\item
  Use \textbf{Einstein.4D} as a metric analogy when discussing geodesics
  or time-evolution; do not mix with synergetic unit conventions.
\end{itemize}

\newpage

\section{Quadray Analytical Details and
Methods}\label{quadray-analytical-details-and-methods}

\subsection{Overview}\label{overview-1}

This section provides detailed analytical methods for working with
Quadray coordinates, including coordinate conventions, volume
calculations, and optimization approaches. We emphasize the distinction
between different 4D frameworks and provide practical computational
methods.

\subsection{Mathematical Foundations}\label{mathematical-foundations}

The methods presented here rest on three interconnected mathematical
frameworks. For comprehensive definitions, see
\href{02_4d_namespaces.md}{Section 2: 4D Namespaces}.

\subsubsection{Framework Distinctions}\label{framework-distinctions}

\begin{itemize}
\tightlist
\item
  \textbf{Coxeter.4D}: Euclidean 4D geometry with standard metric tensor
  and volume forms
\item
  \textbf{Einstein.4D}: Minkowski spacetime with indefinite metric for
  information geometry analogies\\
\item
  \textbf{Fuller.4D}: Synergetics/Quadray coordinates with integer
  lattice constraints and IVM unit conventions
\end{itemize}

\subsubsection{Key Mathematical
Principles}\label{key-mathematical-principles}

\begin{itemize}
\tightlist
\item
  \textbf{Rational volume quantization}: Lattice constraints quantize
  tetrahedral volumes to exact quarter-unit (1/4-grain) rationals in IVM
  units --- integral for unit-tetra tilings, e.g.~\(\tfrac14\) for the
  primitive tetrahedron
\item
  \textbf{Coordinate system bridges}: Linear transformations between
  Fuller.4D and Coxeter.4D preserve geometric relationships
\item
  \textbf{Information geometry}: Fisher metric provides Riemannian
  structure for optimization on parameter manifolds
\item
  \textbf{Exact arithmetic}: Bareiss algorithm ensures determinant
  calculations remain exact for integer inputs
\end{itemize}

\subsubsection{Implementation Strategy}\label{implementation-strategy}

All mathematical concepts are implemented with: - \textbf{Exact
arithmetic} where possible (integer determinants, symbolic computation)
- \textbf{Numerical stability} for floating-point operations (ridge
regularization, condition number monitoring) - \textbf{Cross-validation}
between different formulations (Ace 5×5 vs Cayley--Menger, native vs
bridging approaches) - \textbf{Comprehensive testing} ensuring
mathematical correctness across edge cases

\subsection{Framework Integration}\label{framework-integration}

The three 4D frameworks serve distinct but complementary roles in our
implementation:

\begin{itemize}
\item
  \textbf{Coxeter.4D provides the container}: Euclidean 4D space serves
  as the mathematical foundation for volume calculations, distance
  metrics, and geometric transformations. When we compute volumes via
  Cayley--Menger determinants or XYZ coordinate determinants, we operate
  in this framework.
\item
  \textbf{Einstein.4D provides the analogy}: The Minkowski metric
  structure inspires our information geometry approach, where the Fisher
  information matrix acts as a Riemannian metric on parameter space.
  Natural gradient descent follows geodesics on this information
  manifold, analogous to how particles follow geodesics in spacetime.
\item
  \textbf{Fuller.4D provides the constraints}: The Quadray coordinate
  system and IVM lattice impose integer constraints that enable exact
  arithmetic and discrete optimization. The synergetics unit conventions
  (regular tetrahedron volume = 1) create a quantized geometry where
  volumes are exact rationals on the 1/4 grid --- integral for
  tetrahedra tiling unit IVM tetras, \(\tfrac14\)-grain otherwise.
\end{itemize}

This multi-framework approach allows us to: 1. Use standard Euclidean
methods for volume calculations where relevant or already in use
(Coxeter.4D) 2. Apply information geometry principles for geodesic
optimization (Einstein.4D analogy)\\
3. Maintain exact arithmetic in the Isotropic Vector Matrix (IVM)
setting through integer lattice constraints (Fuller.4D) 4. Bridge and
swap among frameworks via coordinate transformations and unit
conversions

\subsection{Fuller.4D Coordinates and
Normalization}\label{fuller.4d-coordinates-and-normalization}

\begin{itemize}
\tightlist
\item
  Quadray vector q = (a,b,c,d), a,b,c,d ≥ 0, with at least one
  coordinate zero under normalization.
\item
  Projective normalization can add/subtract (k,k,k,k) without changing
  direction; choose k to enforce non-negativity and one zero minimum.
\item
  Isotropic Vector Matrix (IVM): integer quadrays describe CCP sphere
  centers; the 12 permutations of \{2,1,1,0\} form the cuboctahedron
  (vector equilibrium).

  \begin{itemize}
  \tightlist
  \item
    Integer-coordinate models: assigning unit IVM tetravolume to the
    regular tetrahedron yields integer coordinates for several familiar
    polyhedra (inverse tetrahedron, cube, octahedron, rhombic
    dodecahedron, cuboctahedron) when expressed as linear combinations
    of the four quadray basis vectors. See overview:
    \href{https://en.wikipedia.org/wiki/Quadray_coordinates}{Quadray
    coordinates}.
  \end{itemize}
\end{itemize}

\subsection{Conversions and Vector Operations: Quadray ↔ Cartesian
(Fuller.4D ↔
Coxeter.4D/XYZ)}\label{conversions-and-vector-operations-quadray-cartesian-fuller.4d-coxeter.4dxyz}

\begin{itemize}
\tightlist
\item
  \textbf{Embedding conventions} determine the linear maps between
  Quadray (Fuller.4D) and Cartesian XYZ (a 3D slice or embedding aligned
  with Coxeter.4D conventions).
\item
  \textbf{References}: Urner provides practical conversion write-ups and
  matrices; see:

  \begin{itemize}
  \tightlist
  \item
    Quadrays and XYZ:
    \href{https://www.grunch.net/synergetics/quadxyz.html}{Urner --
    Quadrays and XYZ}
  \item
    Introduction with examples:
    \href{https://www.grunch.net/synergetics/quadintro.html}{Urner --
    Quadray intro}
  \end{itemize}
\item
  \textbf{Implementation}: choose a fixed tetrahedral embedding;
  construct a 3×4 matrix M that maps (a,b,c,d) to (x,y,z), respecting
  A,B,C,D directions to tetra vertices. The inverse map can be defined
  up to projective normalization (adding (k,k,k,k)). When comparing
  volumes, use the \passthrough{\lstinline!S3=\\sqrt\{9/8\}!} scale to
  convert XYZ (Euclidean) volumes to IVM (Fuller.4D) units.
\item
  \textbf{Vector view}: treat \passthrough{\lstinline!q!} as a vector
  with magnitude and direction; define dot products and norms by pushing
  to XYZ via \passthrough{\lstinline!M!}.
\end{itemize}

\subsubsection{Integer-coordinate constructions (compact derivation
box)}\label{integer-coordinate-constructions-compact-derivation-box}

\begin{itemize}
\tightlist
\item
  Under the synergetics convention (unit regular tetrahedron has
  tetravolume 1), many familiar solids admit Quadray integer
  coordinates. For example, the octahedron at the same edge length has
  tetravolume 4, and its vertices can be formed as integer linear
  combinations of the four axes A,B,C,D subject to the Quadray
  normalization rule.
\item
  The cuboctahedron (vector equilibrium) arises as the shell of the 12
  nearest IVM neighbors given by the permutations of \((2,1,1,0)\). The
  rhombic dodecahedron (tetravolume 6) is the Voronoi cell of the
  FCC/CCP packing centered at the origin under the same embedding.
\item
  See the following figure for a schematic summary of these
  relationships.
\end{itemize}

{\def\LTcaptype{none} % do not increment counter
\begin{longtable}[]{@{}
  >{\raggedright\arraybackslash}p{(\linewidth - 4\tabcolsep) * \real{0.3333}}
  >{\raggedright\arraybackslash}p{(\linewidth - 4\tabcolsep) * \real{0.3333}}
  >{\raggedright\arraybackslash}p{(\linewidth - 4\tabcolsep) * \real{0.3333}}@{}}
\toprule\noalign{}
\begin{minipage}[b]{\linewidth}\raggedright
Object
\end{minipage} & \begin{minipage}[b]{\linewidth}\raggedright
Quadray construction (sketch)
\end{minipage} & \begin{minipage}[b]{\linewidth}\raggedright
IVM volume
\end{minipage} \\
\midrule\noalign{}
\endhead
\bottomrule\noalign{}
\endlastfoot
Regular tetrahedron & Vertices \passthrough{\lstinline!o=(0,0,0,0)!},
\passthrough{\lstinline!p=(2,1,0,1)!},
\passthrough{\lstinline!q=(2,1,1,0)!},
\passthrough{\lstinline!r=(2,0,1,1)!} & 1 \\
Cube (same edge) & Union of 3 mutually orthogonal rhombic belts wrapped
on the tetra frame; edges tracked by XYZ embedding; compare the
following figure & 3 \\
Octahedron (same edge) & Convex hull of mid-edges of the tetra frame
(pairwise axis sums normalized) & 4 \\
Rhombic dodecahedron & Voronoi cell of FCC/CCP packing at origin (dual
to cuboctahedron) & 6 \\
Cuboctahedron (vector equilibrium) & Shell of the 12 nearest IVM
neighbors: permutations of \passthrough{\lstinline!(2,1,1,0)!} & 20 \\
Truncated octahedron & Archimedean solid with 6 square and 8 hexagonal
faces; space-filling tiling & 20 \\
\end{longtable}
}

Small coordinate examples (subset):

\begin{itemize}
\tightlist
\item
  Cuboctahedron neighbors (representatives):
  \passthrough{\lstinline!(2,1,1,0)!},
  \passthrough{\lstinline!(2,1,0,1)!},
  \passthrough{\lstinline!(2,0,1,1)!},
  \passthrough{\lstinline!(1,2,1,0)!}; the full shell is all distinct
  permutations.
\item
  Tetrahedron:
  \passthrough{\lstinline![(0,0,0,0), (2,1,0,1), (2,1,1,0), (2,0,1,1)]!}.
\end{itemize}

Short scripts:

\begin{lstlisting}[language=bash]
python3 quadmath/scripts/polyhedra_quadray_constructions.py
\end{lstlisting}

Programmatic check (neighbors, equal radii, adjacency):

\begin{lstlisting}[language=Python]
import numpy as np
from quadmath.core.examples import example_cuboctahedron_vertices_xyz

xyz = np.array(example_cuboctahedron_vertices_xyz())
r = np.linalg.norm(xyz[0])
assert np.allclose(np.linalg.norm(xyz, axis=1), r)

# Touching neighbors have separation 2r
touch = []
for i in range(len(xyz)):
    for j in range(i+1, len(xyz)):
        d = np.linalg.norm(xyz[i] - xyz[j])
        if abs(d - 2*r) / (2*r) < 0.05:
            touch.append((i, j))
assert len(touch) > 0
\end{lstlisting}

\subsubsection{Example vertex lists and volume checks
(illustrative)}\label{example-vertex-lists-and-volume-checks-illustrative}

The following snippets show executable decompositions for standard
synergetics volumes. Each tetra volume is computed via
\passthrough{\lstinline!ace\_tetravolume\_5x5!} and summed.

Octahedron (V = 4): the octahedron with vertices \(\pm 2\hat e_i\) (edge
\(2\sqrt2\), matching the unit IVM tetra) decomposes into eight
origin-apex orthant tetras, each of volume \(\tfrac12\):

\begin{lstlisting}[language=Python]
from quadmath.core.quadray import Quadray, ace_tetravolume_5x5

o = Quadray(0,0,0,0)
axes = [
    Quadray(1,0,0,1), Quadray(0,1,1,0),   # +x, -x
    Quadray(1,1,0,0), Quadray(0,0,1,1),   # +y, -y
    Quadray(1,0,1,0), Quadray(0,1,0,1),   # +z, -z
]
V_oct = sum(
    ace_tetravolume_5x5(o, x, y, z)
    for x in axes[0:2] for y in axes[2:4] for z in axes[4:6]
)  # 8 * 1/2 = 4
\end{lstlisting}

Cube (V = 3): the cube with XYZ vertices \((\pm1,\pm1,\pm1)\) (edge 2)
decomposes into the inscribed tetra on four alternating cube vertices (V
= 1) plus four corner tetras (V = \(\tfrac12\) each):

\begin{lstlisting}[language=Python]
from quadmath.core.quadray import Quadray, ace_tetravolume_5x5

inner = ace_tetravolume_5x5(
    Quadray(1,0,0,0), Quadray(0,1,0,0), Quadray(0,0,1,0), Quadray(0,0,0,1),
)  # 1
corner = ace_tetravolume_5x5(
    Quadray(0,1,1,1), Quadray(0,0,0,1), Quadray(0,1,0,0), Quadray(0,0,1,0),
)  # 1/2
V_cube = inner + 4 * corner  # 1 + 4*(1/2) = 3
\end{lstlisting}

Notes.

\begin{itemize}
\tightlist
\item
  The octahedron and cube decompositions above execute exactly as shown
  (verified against \passthrough{\lstinline!ace\_tetravolume\_5x5!}).
  Other equivalent tilings are possible.
\item
  Volumes are invariant to adding \((k,k,k,k)\) to each vertex of a
  tetra (projective normalization), which the 5×5 determinant respects.
\end{itemize}

\subsection{Integer Volume Quantization}\label{sec:integer_volume}

For a tetrahedron with vertices P₀..P₃ in the Quadray integer lattice
(Fuller.4D), see Eq. \eqref{eq:lattice_det} in the equations appendix.

\begin{itemize}
\tightlist
\item
  With integer quadray coordinates the projected determinant is an exact
  integer, and the IVM tetravolume is the exact rational \(|{\det}|/4\)
  (see Eq. \eqref{eq:gdj}), implemented with
  \passthrough{\lstinline!fractions.Fraction!} in
  \passthrough{\lstinline!integer\_tetra\_volume!} and
  \passthrough{\lstinline!ace\_tetravolume\_5x5!}. Volumes are integral
  for tetrahedra that tile into unit IVM tetras (determinant divisible
  by 4) and fractional otherwise: the primitive tetrahedron spanned by
  \((0,0,0,0)\), \((1,0,0,0)\), \((0,1,0,0)\), \((0,0,1,0)\) has
  determinant 1 and exact volume \(\tfrac{1}{4}\).
\item
  Unit conventions: the unit IVM tetrahedron (origin plus three IVM
  neighbor moves, permutations of \((2,1,1,0)\)) has determinant 4 and
  volume exactly 1 (synergetics).
\end{itemize}

Notes.

\begin{itemize}
\tightlist
\item
  \(P_0,\ldots,P_3\) are tetrahedron vertices in Quadray coordinates.
\item
  \(V_{xyz}\) in Eq. \eqref{eq:lattice_det} is the Euclidean (XYZ)
  volume; converting it to IVM tetra-units requires the synergetics
  scale factor \(S3=\sqrt{9/8}\) in sphere-radius units (Eq.
  \eqref{eq:xyz_det}). For lattice tetrahedra, the direct quadray
  formula Eq. \eqref{eq:gdj} yields the IVM volume exactly, without unit
  conversion.
\item
  Background and variations are discussed under Tetrahedron volume
  formulas:
  \href{https://en.wikipedia.org/wiki/Tetrahedron\#Volume}{Tetrahedron
  -- volume}.
\end{itemize}

Tom Ace 5×5 determinant (tetravolume directly from quadrays), see Eq.
\eqref{eq:ace5x5} in the equations appendix.

This returns the same exact tetravolumes as
\passthrough{\lstinline!integer\_tetra\_volume!} for every lattice
tetrahedron: \(|\det_{5\times5}|\) always equals the magnitude of the
projected \(3\times3\) determinant, so both implementations agree
exactly (both return \passthrough{\lstinline!fractions.Fraction!} values
\(|\det|/4\)).

Notes.

\begin{itemize}
\tightlist
\item
  Rows correspond to the Quadray 4-tuples of the four vertices with a
  final affine column of ones; the last row enforces projective
  normalization.
\item
  The factor \(\tfrac{1}{4}\) returns tetravolumes in IVM units
  consistent with synergetics. See also
  \href{https://en.wikipedia.org/wiki/Quadray_coordinates}{Quadray
  coordinates}.
\end{itemize}

Equivalently, define the 5×5 matrix of quadray coordinates augmented
with an affine 1 as shown in Eq. \eqref{eq:ace5x5_expanded} in the
equations appendix.

Points vs vectors: subtracting points is shorthand for forming edge
vectors. We treat quadray 4-tuples as vectors from the origin;
differences like \((P_1-P_0)\) mean ``edge vectors,'' avoiding ambiguity
between ``points'' and ``vectors.''

Equivalently via Cayley--Menger determinant (Coxeter.4D/Euclidean
lengths)
(\href{https://en.wikipedia.org/wiki/Cayley\%E2\%80\%93Menger_determinant}{Cayley--Menger
determinant}), see Eq. \eqref{eq:cayley_menger} in the equations
appendix.

References:
\href{https://en.wikipedia.org/wiki/Cayley\%E2\%80\%93Menger_determinant}{Cayley--Menger
determinant}, lattice tetrahedra discussions in geometry texts; see also
\href{https://en.wikipedia.org/wiki/Tetrahedron\#Volume}{Tetrahedron --
volume}. Code: \passthrough{\lstinline!integer\_tetra\_volume!},
\passthrough{\lstinline!ace\_tetravolume\_5x5!}.

Notes.

\begin{itemize}
\tightlist
\item
  \textbf{Pairwise distances}: \(d_{ij}\) are Euclidean distances
  between vertices \(P_i\) and \(P_j\).
\item
  \textbf{Length-only formulation}: Cayley--Menger provides a
  length-only formula for simplex volumes, here specialized to
  tetrahedra; see the canonical reference above.
\end{itemize}

\begin{longtable}[]{@{}ll@{}}
\caption{Polyhedra tetravolumes in IVM units (edge length equal to the
unit tetra edge).}\label{tbl:polyhedra_volumes}\tabularnewline
\toprule\noalign{}
Polyhedron (edge = tetra edge) & Volume (tetra-units) \\
\midrule\noalign{}
\endfirsthead
\toprule\noalign{}
Polyhedron (edge = tetra edge) & Volume (tetra-units) \\
\midrule\noalign{}
\endhead
\bottomrule\noalign{}
\endlastfoot
Regular Tetrahedron & 1 \\
Cube & 3 \\
Octahedron & 4 \\
Rhombic Dodecahedron & 6 \\
Cuboctahedron (Vector Equilibrium) & 20 \\
Truncated Octahedron & 20 \\
\end{longtable}

\subsection{Distances and Metrics}\label{distances-and-metrics}

Distance definitions depend on the chosen embedding and normalization.
For cross-references to information geometry, see
\href{08_equations_appendix.md\#eq:fim}{Eq. (FIM)} and
\href{08_equations_appendix.md\#eq:natgrad}{natural gradient} in the
Equations appendix.

\subsection{XYZ determinant and S3 conversion}\label{sec:xyz_conversion}

Given XYZ coordinates of tetrahedron vertices (x\_i, y\_i, z\_i), the
Euclidean volume is computed as shown in Eq. \eqref{eq:xyz_det} in the
equations appendix.

Convert to IVM units via \(V_{ivm} = S3 \cdot V_{xyz}\) with
\(S3=\sqrt{9/8}\). See background discussion under
\href{https://en.wikipedia.org/wiki/Tetrahedron\#Volume}{Tetrahedron --
volume}.

\subsubsection{Alternative Volume
Formulas}\label{alternative-volume-formulas}

\begin{itemize}
\item
  \textbf{Piero della Francesca (PdF) formula}, which consumes edge
  lengths and returns Euclidean volumes. See Eq. \eqref{eq:pdf} in the
  equations appendix.

  Convert to IVM units via \(V_{ivm} = S3 \cdot V_{xyz}\) with
  \(S3=\sqrt{9/8}\). See background discussion under
  \href{https://en.wikipedia.org/wiki/Tetrahedron\#Volume}{Tetrahedron
  -- volume}.
\item
  \textbf{Gerald de Jong (GdJ) formula}, which natively returns
  tetravolumes. See Eq. \eqref{eq:gdj} in the equations appendix.

  In Quadray coordinates, one convenient native form uses edge-vector
  differences and an integer-preserving determinant (agreeing with Ace
  5×5):

  where each column is formed from Quadray component differences of
  \(P_1-P_0\), \(P_2-P_0\), \(P_3-P_0\) projected via
  \(\pi(P) = (a - d,\, b - d,\, c - d)\) to the synergetics 3D slice
  (matching \passthrough{\lstinline!integer\_tetra\_volume!}); integer
  arithmetic is exact and the factor \(\tfrac{1}{4}\) produces IVM
  tetravolumes. See de Jong's Quadray notes and Urner's implementations
  for derivations
  (\href{https://en.wikipedia.org/wiki/Quadray_coordinates}{Quadray
  coordinates}).
\item
  Euclidean embedding distance via appropriate linear map from quadray
  to R³.
\item
  Information geometry metric: Fisher Information Matrix (FIM)

  \begin{itemize}
  \tightlist
  \item
    \(\mathrm{FIM}[i,j] = \mathbb{E}\big[\, \partial_{\theta_i} \log p(x;\theta)\,\partial_{\theta_j} \log p(x;\theta)\,\big]\)
  \item
    Acts as Riemannian metric; natural gradient uses FIM⁻¹ ∇θ L. See
    \href{https://en.wikipedia.org/wiki/Fisher_information}{Fisher
    information}.
  \end{itemize}
\end{itemize}

\subsection{Fisher Geometry in Quadray
Space}\label{fisher-geometry-in-quadray-space}

\begin{itemize}
\tightlist
\item
  Symmetries of quadray lattices often induce near block-diagonal FIM.
\item
  Determinant and spectrum characterize conditioning and information
  concentration.
\end{itemize}

\subsection{Practical Methods}\label{practical-methods}

\subsection{Tetravolumes with Quadrays}\label{sec:tetravolumes_quadrays}

\begin{itemize}
\item
  The tetravolume of a tetrahedron with vertices given as Quadrays
  \passthrough{\lstinline!a,b,c,d!} can be computed directly from their
  4-tuples via the Tom Ace 5×5 determinant; see Eq. \eqref{eq:ace5x5}
  for the canonical form.
\item
  Unit regular tetrahedron from origin: with
  \passthrough{\lstinline!o=(0,0,0,0)!},
  \passthrough{\lstinline!p=(2,1,0,1)!},
  \passthrough{\lstinline!q=(2,1,1,0)!},
  \passthrough{\lstinline!r=(2,0,1,1)!}, we have
  \passthrough{\lstinline!V\_ivm(o,p,q,r)=1!}. Doubling each vector
  scales volume by 8, as expected.
\item
  Equivalent length-based formulas agree with the 5×5 determinant:

  \begin{itemize}
  \tightlist
  \item
    Cayley--Menger:
    \(288\,V^2 = \det\begin{pmatrix}0&1&1&1&1\\1&0&d_{01}^2&d_{02}^2&d_{03}^2\\1&d_{10}^2&0&d_{12}^2&d_{13}^2\\1&d_{20}^2&d_{21}^2&0&d_{23}^2\\1&d_{30}^2&d_{31}^2&d_{32}^2&0\end{pmatrix}\).
  \item
    Piero della Francesca (PdF) Heron-like formula (converted to IVM via
    \(S3 = \sqrt{9/8}\)).
  \end{itemize}

  Let edge lengths meeting at a vertex be \(a,b,c\), and the opposite
  edges be \(d,e,f\). The Euclidean volume satisfies

  Convert to IVM units via \(V_{ivm} = S3 \cdot V_{xyz}\) with
  \(S3=\sqrt{9/8}\). See background discussion under
  \href{https://en.wikipedia.org/wiki/Tetrahedron\#Volume}{Tetrahedron
  -- volume}.
\item
  Gerald de Jong (GdJ) formula, which natively returns tetravolumes.

  In Quadray coordinates, one convenient native form uses edge-vector
  differences and an integer-preserving determinant (agreeing with Ace
  5×5):

  where each column is formed from Quadray component differences of
  \(P_1-P_0\), \(P_2-P_0\), \(P_3-P_0\) projected via
  \(\pi(P) = (a - d,\, b - d,\, c - d)\) to the synergetics 3D slice
  (matching \passthrough{\lstinline!integer\_tetra\_volume!}); integer
  arithmetic is exact and the factor \(\tfrac{1}{4}\) produces IVM
  tetravolumes. See de Jong's Quadray notes and Urner's implementations
  for derivations
  (\href{https://en.wikipedia.org/wiki/Quadray_coordinates}{Quadray
  coordinates}).
\end{itemize}

\subsubsection{Bridging and native tetravolume
formulas}\label{bridging-and-native-tetravolume-formulas}

\begin{itemize}
\tightlist
\item
  \textbf{Lengths (bridging)}: PdF and Cayley--Menger (CM) consume
  Cartesian lengths (XYZ) and produce Euclidean volumes; convert to IVM
  units via \(S3 = \sqrt{9/8}\).
\item
  \textbf{Quadray-native}: Gerald de Jong (GdJ) returns IVM tetravolumes
  directly (no XYZ bridge). Tom Ace's 5×5 coordinate formula is likewise
  native IVM. All agree numerically with CM+S3 on shared cases.
\end{itemize}

References and discussion: \href{https://flic.kr/p/2rn22en}{Urner --
Flickr diagram}. For computational implementations and educational
materials, see the \href{07_resources.md}{Resources} section.

Figure: automated comparison (native Ace 5×5 vs CM+S3) across small
examples (see script \passthrough{\lstinline!sympy\_formalisms.py!}).
The figure and source CSV/NPZ are in
\passthrough{\lstinline!quadmath/output/!}.

\begin{figure}
\centering
\pandocbounded{\includegraphics[keepaspectratio,alt={Validation of bridging vs native tetravolume formulations across canonical examples. This bar chart compares IVM tetravolumes computed via two independent methods: the ``bridging'' approach using Cayley--Menger determinants on Euclidean edge lengths converted to IVM units via the synergetics factor S3=\textbackslash sqrt\{9/8\}, versus the ``native'' approach using Tom Ace's 5×5 determinant formula that operates directly on Quadray coordinates without XYZ intermediates. Test cases: Unit tetrahedron (V=1), 2× edge scaling (V=8), mixed coordinate tetrahedron, centered tetrahedron (V=3), and large mixed tetrahedron, all using integer Quadray coordinates. Results: The overlapping bars demonstrate numerical agreement at machine precision between the length-based Coxeter.4D approach (Cayley--Menger + S3 conversion) and the coordinate-based Fuller.4D approach (Ace 5×5), confirming the mathematical equivalence of these formulations under synergetics unit conventions. Raw numerical data saved as bridging\_vs\_native.csv for reproducibility and further analysis.}]{../output/figures/bridging_vs_native.png}}
\caption{\textbf{Validation of bridging vs native tetravolume
formulations across canonical examples}. This bar chart compares IVM
tetravolumes computed via two independent methods: the ``bridging''
approach using Cayley--Menger determinants on Euclidean edge lengths
converted to IVM units via the synergetics factor \(S3=\sqrt{9/8}\),
versus the ``native'' approach using Tom Ace's 5×5 determinant formula
that operates directly on Quadray coordinates without XYZ intermediates.
\textbf{Test cases}: Unit tetrahedron (V=1), 2× edge scaling (V=8),
mixed coordinate tetrahedron, centered tetrahedron (V=3), and large
mixed tetrahedron, all using integer Quadray coordinates.
\textbf{Results}: The overlapping bars demonstrate numerical agreement
at machine precision between the length-based Coxeter.4D approach
(Cayley--Menger + S3 conversion) and the coordinate-based Fuller.4D
approach (Ace 5×5), confirming the mathematical equivalence of these
formulations under synergetics unit conventions. Raw numerical data
saved as \protect\passthrough{\lstinline!bridging\_vs\_native.csv!} for
reproducibility and further analysis.}
\end{figure}

\begin{figure}
\centering
\pandocbounded{\includegraphics[keepaspectratio,alt={Tetrahedron volume scaling relationships: Euclidean vs IVM unit conventions. This plot demonstrates the mathematical relationship between edge length scaling and tetravolume under both Euclidean (XYZ) and IVM (synergetics) unit conventions. X-axis: Edge length scaling factor (0 to 5.0). Y-axis: Tetrahedron volume in respective units. Blue line (Euclidean): Volume scales as the cube of edge length, following the standard V = \textbackslash frac\{\textbackslash sqrt\{2\}\}\{12\} \textbackslash cdot L\^{}3 relationship for regular tetrahedra. Orange line (IVM): Volume scales as the cube of edge length but in IVM tetra-units, following V\_\{ivm\} = \textbackslash frac\{1\}\{8\} \textbackslash cdot L\^{}3 where the regular tetrahedron with unit edge has volume 1/8. Key insight: The ratio between these two scaling laws is the synergetics factor S3 = \textbackslash sqrt\{9/8\} \textbackslash approx 1.06066, which converts between Euclidean and IVM volume conventions. Mathematical foundation: This scaling relationship demonstrates how both conventions preserve the cubic scaling relationship, but with different fundamental units reflecting the different geometric assumptions of Coxeter.4D (Euclidean) versus Fuller.4D (synergetics) frameworks. The plot provides the theoretical foundation for understanding volume conversions and scaling behavior in the IVM system.}]{../output/figures/volumes_scale_plot.png}}
\caption{\textbf{Tetrahedron volume scaling relationships: Euclidean vs
IVM unit conventions}. This plot demonstrates the mathematical
relationship between edge length scaling and tetravolume under both
Euclidean (XYZ) and IVM (synergetics) unit conventions. \textbf{X-axis}:
Edge length scaling factor (0 to 5.0). \textbf{Y-axis}: Tetrahedron
volume in respective units. \textbf{Blue line (Euclidean)}: Volume
scales as the cube of edge length, following the standard
\(V = \frac{\sqrt{2}}{12} \cdot L^3\) relationship for regular
tetrahedra. \textbf{Orange line (IVM)}: Volume scales as the cube of
edge length but in IVM tetra-units, following
\(V_{ivm} = \frac{1}{8} \cdot L^3\) where the regular tetrahedron with
unit edge has volume 1/8. \textbf{Key insight}: The ratio between these
two scaling laws is the synergetics factor
\(S3 = \sqrt{9/8} \approx 1.06066\), which converts between Euclidean
and IVM volume conventions. \textbf{Mathematical foundation}: This
scaling relationship demonstrates how both conventions preserve the
cubic scaling relationship, but with different fundamental units
reflecting the different geometric assumptions of Coxeter.4D (Euclidean)
versus Fuller.4D (synergetics) frameworks. The plot provides the
theoretical foundation for understanding volume conversions and scaling
behavior in the IVM system.}
\end{figure}

\begin{figure}
\centering
\pandocbounded{\includegraphics[keepaspectratio,alt={Synergetic polyhedra volume relationships in the Quadray/IVM framework (comprehensive visualization). This figure combines 3D polyhedra visualizations with an extended network diagram showing integer volume relationships among key synergetic polyhedra. Left panel (3D visualizations): Color-coded polyhedra including regular tetrahedron (V=1, fundamental unit), cube (V=3), octahedron (V=4), rhombic dodecahedron (V=6), cuboctahedron (V=20), and truncated octahedron (V=20), all constructed with consistent edge lengths and proper geometric faces. Right panel (network diagram): Extended volume relationships showing fundamental shapes (V=1,3,4,6), complex constructions (V=20), and scaling relationships (2× edge length → 8× volume). Additional polyhedra: Includes truncated octahedron (V=20) and scaled variants demonstrating the ``third power'' volume scaling law V ∝ L³ in IVM units. Geometric constructions: Edge-union relationships, truncation operations, dual polyhedra, and Voronoi cell constructions. Fuller.4D significance: These integer volume ratios reflect the quantized nature of space-filling in synergetics, where the regular tetrahedron provides a natural unit container and other polyhedra emerge as integer multiples, supporting discrete geometric computation and exact lattice-based optimization methods. All constructions respect the IVM unit convention where the regular tetrahedron has tetravolume 1.}]{../output/figures/polyhedra_quadray_constructions.png}}
\caption{\textbf{Synergetic polyhedra volume relationships in the
Quadray/IVM framework (comprehensive visualization)}. This figure
combines 3D polyhedra visualizations with an extended network diagram
showing integer volume relationships among key synergetic polyhedra.
\textbf{Left panel (3D visualizations)}: Color-coded polyhedra including
regular tetrahedron (V=1, fundamental unit), cube (V=3), octahedron
(V=4), rhombic dodecahedron (V=6), cuboctahedron (V=20), and truncated
octahedron (V=20), all constructed with consistent edge lengths and
proper geometric faces. \textbf{Right panel (network diagram)}: Extended
volume relationships showing fundamental shapes (V=1,3,4,6), complex
constructions (V=20), and scaling relationships (2× edge length → 8×
volume). \textbf{Additional polyhedra}: Includes truncated octahedron
(V=20) and scaled variants demonstrating the ``third power'' volume
scaling law V ∝ L³ in IVM units. \textbf{Geometric constructions}:
Edge-union relationships, truncation operations, dual polyhedra, and
Voronoi cell constructions. \textbf{Fuller.4D significance}: These
integer volume ratios reflect the quantized nature of space-filling in
synergetics, where the regular tetrahedron provides a natural unit
container and other polyhedra emerge as integer multiples, supporting
discrete geometric computation and exact lattice-based optimization
methods. All constructions respect the IVM unit convention where the
regular tetrahedron has tetravolume 1.}
\end{figure}

\subsubsection{Short Python snippets}\label{short-python-snippets}

\begin{lstlisting}[language=Python]
from quadmath.core.quadray import Quadray, ace_tetravolume_5x5

o = Quadray(0,0,0,0)
p = Quadray(2,1,0,1)
q = Quadray(2,1,1,0)
r = Quadray(2,0,1,1)
assert ace_tetravolume_5x5(o,p,q,r) == 1  # unit IVM tetra
\end{lstlisting}

\begin{lstlisting}[language=Python]
import numpy as np
from quadmath.core.cayley_menger import ivm_tetra_volume_cayley_menger

# Example: regular tetrahedron with edge length 1 (XYZ units)
d2 = np.ones((4,4)) - np.eye(4)  # squared distances
V_ivm = ivm_tetra_volume_cayley_menger(d2)   # = 1/8 in IVM tetra-units
\end{lstlisting}

\begin{lstlisting}[language=Python]
# SymPy implementation of Tom Ace 5×5 (symbolic determinant)
from sympy import Matrix

def qvolume(q0, q1, q2, q3):
    M = Matrix([
        q0 + (1,),
        q1 + (1,),
        q2 + (1,),
        q3 + (1,),
        [1, 1, 1, 1, 0],
    ])
    return abs(M.det()) / 4
\end{lstlisting}

\begin{lstlisting}[language=Python]
# Symbolic variant with SymPy (exact radicals)
from sympy import Matrix, sqrt, simplify
from quadmath.core.symbolic import cayley_menger_volume_symbolic, convert_xyz_volume_to_ivm_symbolic

d2 = Matrix([[0,1,1,1],[1,0,1,1],[1,1,0,1],[1,1,1,0]])
V_xyz_sym = cayley_menger_volume_symbolic(d2)      # sqrt(2)/12
V_ivm_sym = simplify(convert_xyz_volume_to_ivm_symbolic(V_xyz_sym))  # 1/8
\end{lstlisting}

\subsubsection{Random tetrahedra in the IVM (integer
volumes)}\label{random-tetrahedra-in-the-ivm-integer-volumes}

\begin{itemize}
\tightlist
\item
  The 12 CCP directions are the permutations of \((2,1,1,0)\). Random
  walks on this move set generate integer-coordinate Quadrays; resulting
  tetrahedra have integer tetravolumes.
\end{itemize}

\begin{lstlisting}[language=Python]
from itertools import permutations
from random import choice
from quadmath.core.quadray import Quadray, ace_tetravolume_5x5

moves = [Quadray(*p) for p in set(permutations((2,1,1,0)))]

def random_walk(start: Quadray, steps: int) -> Quadray:
    cur = start
    for _ in range(steps):
        m = choice(moves)
        cur = Quadray(cur.a+m.a, cur.b+m.b, cur.c+m.c, cur.d+m.d)
    return cur

A = random_walk(Quadray(0,0,0,0), 1000)
B = random_walk(Quadray(0,0,0,0), 1000)
C = random_walk(Quadray(0,0,0,0), 1000)
D = random_walk(Quadray(0,0,0,0), 1000)
V = ace_tetravolume_5x5(A,B,C,D)            # exact IVM volume as a Fraction
\end{lstlisting}

\subsubsection{Algebraic precision}\label{algebraic-precision}

\begin{itemize}
\tightlist
\item
  Determinants via floating-point introduce rounding noise. For exact
  arithmetic, use the
  \href{https://en.wikipedia.org/wiki/Bareiss_algorithm}{Bareiss
  algorithm} (already used by
  \passthrough{\lstinline!ace\_tetravolume\_5x5!}) or symbolic engines
  (e.g., \passthrough{\lstinline!sympy!}). For random-walk examples with
  integer quadray inputs, volumes are exact rationals
  (\passthrough{\lstinline!fractions.Fraction!}): integral when the
  tetrahedron decomposes into unit IVM tetras, fractional (in quarters)
  otherwise.
\item
  When computing via XYZ determinants, high-precision floats (e.g.,
  \passthrough{\lstinline!gmpy2.mpfr!}) or symbolic matrices avoid
  vestigial errors; round at the end if the underlying result is known
  to be integral.
\end{itemize}

\subsubsection{XYZ determinant and the S3
conversion}\label{xyz-determinant-and-the-s3-conversion}

\begin{itemize}
\tightlist
\item
  Using XYZ coordinates of the four vertices: see Eq. \eqref{eq:xyz_det}
  for the determinant form and the S3 conversion to IVM units.
\end{itemize}

\subsubsection{D\^{}3 vs R\^{}3: 60° ``closing the lid'' vs orthogonal
``cubing''}\label{d3-vs-r3-60-closing-the-lid-vs-orthogonal-cubing}

\begin{itemize}
\tightlist
\item
  \textbf{IVM (D\^{}3) heuristic}: From a 60--60--60 corner, three
  non-negative edge lengths \(A,B,C\) along quadray directions enclose a
  tetrahedron by ``closing the lid.'' In synergetics, the tetravolume
  scales as the simple product \(ABC\) under IVM conventions (unit
  regular tetra has volume 1). By contrast, in the orthogonal (R\^{}3)
  habit, one constructs a full parallelepiped (12 edges); the tetra
  occupies one-sixth of the triple product of edge vectors. The IVM path
  is more direct for tetrahedra.
\item
  \textbf{Pedagogical note}: Adopt a vector-first approach. Differences
  like \((P_i-P_0)\) denote edge vectors; Quadrays and Cartesian can be
  taught in parallel as vector languages on the same Euclidean
  container.
\end{itemize}

Reference notebook with worked examples and code: See the
\href{07_resources.md}{Resources} section for comprehensive educational
materials and computational implementations.

See implementation:
\passthrough{\lstinline!tetra\_volume\_cayley\_menger!}.

\begin{itemize}
\tightlist
\item
  Lattice projection: round to nearest integer quadray; renormalize to
  maintain non-negativity and a minimal zero.
\end{itemize}

\subsection{Practical Implementation
Workflow}\label{practical-implementation-workflow}

The methods described in this document follow a unified workflow that
ensures mathematical correctness, computational efficiency, and
reproducibility:

\subsubsection{1. Mathematical
Formulation}\label{mathematical-formulation}

\begin{itemize}
\tightlist
\item
  \textbf{Theoretical foundations}: Establish mathematical relationships
  between frameworks (Coxeter.4D ↔ Fuller.4D ↔ Einstein.4D)
\item
  \textbf{Coordinate transformations}: Define linear maps between
  coordinate systems with exact arithmetic
\item
  \textbf{Volume formulations}: Implement multiple approaches (Ace 5×5,
  Cayley--Menger, symbolic) for cross-validation
\end{itemize}

\subsubsection{2. Implementation
Strategy}\label{implementation-strategy-1}

\begin{itemize}
\tightlist
\item
  \textbf{Exact arithmetic}: Use Bareiss algorithm for integer
  determinants, symbolic computation for exact results
\item
  \textbf{Numerical stability}: Implement ridge regularization,
  condition number monitoring, and error handling
\item
  \textbf{Cross-validation}: Ensure different formulations produce
  identical results on test cases
\end{itemize}

\subsubsection{3. Testing and Validation}\label{testing-and-validation}

\begin{itemize}
\tightlist
\item
  \textbf{Unit tests}: 100\% coverage of all mathematical functions with
  edge case handling
\item
  \textbf{Integration tests}: Validate coordinate transformations and
  volume calculations across frameworks
\item
  \textbf{Numerical validation}: Compare floating-point implementations
  with exact symbolic results
\end{itemize}

\subsubsection{4. Documentation and
Reproducibility}\label{documentation-and-reproducibility}

\begin{itemize}
\tightlist
\item
  \textbf{Code-documentation sync}: All mathematical concepts have
  corresponding implementations in \passthrough{\lstinline!src/!}
\item
  \textbf{Figure generation}: Scripts in
  \passthrough{\lstinline!quadmath/scripts/!} generate all figures from
  source code
\item
  \textbf{Data export}: CSV/NPZ files accompany all figures for complete
  reproducibility
\end{itemize}

\subsubsection{5. Optimization and
Extension}\label{optimization-and-extension}

\begin{itemize}
\tightlist
\item
  \textbf{Lattice constraints}: Integer volume quantization enables
  discrete optimization methods
\item
  \textbf{Information geometry}: Fisher metric guides natural gradient
  descent on parameter manifolds
\item
  \textbf{Active inference}: Free energy minimization drives both
  perception and action updates
\end{itemize}

This workflow ensures that mathematical theory, computational
implementation, and practical application remain coherent and verifiable
throughout the development process.

\subsection{Code methods (anchors)}\label{code-methods-anchors}

\subsubsection{Core Quadray operations}\label{code:core_quadray}

\paragraph{\texorpdfstring{\texttt{Quadray}}{Quadray}}\label{code:Quadray}

Source: \passthrough{\lstinline!src/quadmath/core/quadray.py!} ---
Quadray vector class with non-negative components and at least one zero
(Fuller.4D).

\paragraph{\texorpdfstring{\texttt{DEFAULT\_EMBEDDING}}{DEFAULT\_EMBEDDING}}\label{code:DEFAULT_EMBEDDING}

Source: \passthrough{\lstinline!src/quadmath/core/quadray.py!} ---
canonical 3×4 symmetric embedding matrix for Quadray to XYZ conversion.

\paragraph{\texorpdfstring{\texttt{to\_xyz}}{to\_xyz}}\label{code:to_xyz}

Source: \passthrough{\lstinline!src/quadmath/core/quadray.py!} --- map
quadray to R³ via a 3×4 embedding matrix (Fuller.4D → Coxeter.4D slice).

\paragraph{\texorpdfstring{\texttt{magnitude}}{magnitude}}\label{code:magnitude}

Source: \passthrough{\lstinline!src/quadmath/core/quadray.py!} ---
return the Euclidean magnitude of \passthrough{\lstinline!q!} under the
given embedding (vector norm).

\paragraph{\texorpdfstring{\texttt{dot}}{dot}}\label{code:dot}

Source: \passthrough{\lstinline!src/quadmath/core/quadray.py!} ---
return Euclidean dot product \textless q1,q2\textgreater{} under the
given embedding.

\subsubsection{Quaternion operations}\label{code:quaternion_operations}

\paragraph{\texorpdfstring{\texttt{qmul}}{qmul}}\label{code:qmul}

Source: \passthrough{\lstinline!src/quadmath/core/quadray.py!} ---
Hamilton product of two quaternions in scalar-first \((w, x, y, z)\)
order (\(q = w + x\,i + y\,j + z\,k\)), an independent representation
from the Quadray lattice components \((a, b, c, d)\) and from the XYZ
triples produced by \passthrough{\lstinline!to\_xyz!}.

\paragraph{\texorpdfstring{\texttt{qconjugate}}{qconjugate}}\label{code:qconjugate}

Source: \passthrough{\lstinline!src/quadmath/core/quadray.py!} ---
conjugate \((w, -x, -y, -z)\) of a quaternion; for a unit quaternion the
conjugate is the multiplicative inverse, so
\passthrough{\lstinline!q * qconjugate(q)!} equals \((1, 0, 0, 0)\) up
to floating-point error.

\paragraph{\texorpdfstring{\texttt{qrotate}}{qrotate}}\label{code:qrotate}

Source: \passthrough{\lstinline!src/quadmath/core/quadray.py!} ---
Rodrigues rotation of a 3-vector by a unit quaternion via
\(v' = q\,(0, v)\,q^*\), composed with \passthrough{\lstinline!qmul!}
and \passthrough{\lstinline!qconjugate!}; \passthrough{\lstinline!q!}
must be unit (within 1e-9) and the rotation magnitude it encodes,
\(2\,\mathrm{atan2}(\lVert (x, y, z) \rVert, w)\), must match
\(|\mathrm{angle}|\) modulo \(2\pi\) within 1e-9, so a mismatched (q,
angle) pair raises instead of silently rotating by the wrong amount.

\paragraph{\texorpdfstring{\texttt{slerp}}{slerp}}\label{code:slerp}

Source: \passthrough{\lstinline!src/quadmath/core/quadray.py!} ---
shortest-arc spherical linear interpolation between unit quaternions:
when \(\langle q_a, q_b\rangle < 0\) the second quaternion is negated
first (\(q\) and \(-q\) encode the same rotation) so the path always
takes the shorter arc on the rotation sphere; nearly parallel inputs
(including exactly antipodal pairs after sign alignment) fall back to
normalized lerp to avoid dividing by \(\sin\theta \approx 0\), and the
endpoints are exact. The site path traced by these interpolations is
visualized by \passthrough{\lstinline!plot\_slerp\_path!} in
\passthrough{\lstinline!src/quadmath/viz/plots.py!} (figure below).

\paragraph{\texorpdfstring{\texttt{rotate\_about\_axis}}{rotate\_about\_axis}}\label{code:rotate_about_axis}

Source: \passthrough{\lstinline!src/quadmath/core/quadray.py!} ---
axis-angle convenience wrapper: builds the unit quaternion
\(q = (\cos(\theta/2),\, \sin(\theta/2)\,\hat{u})\) from the normalized
rotation axis (any non-zero scale accepted; the zero vector is rejected)
and delegates to \passthrough{\lstinline!qrotate!}, reducing the angle
to \((-\pi, \pi]\) first so rotation is right-handed about the axis and
\passthrough{\lstinline!qrotate!}'s encoded-angle validation stays
exact.

The property checks in
\passthrough{\lstinline!src/quadmath/validate/validate.py!} ---
normalization, conjugate-inverse, double cover, slerp midpoint, and
associativity --- consume these operations over quaternion sequences via
\passthrough{\lstinline!run\_validation!}, which collects the resulting
deterministic validation reports.

\begin{figure}
\centering
\pandocbounded{\includegraphics[keepaspectratio,alt={Shortest-arc slerp geodesic from the identity to a \textbackslash pi/2 rotation about the z axis. The path traced by the canonical (2, 0, 0, 0) quadray axis point (embedded at (2, 2, 2) under DEFAULT\_EMBEDDING) as the rotation quaternion is interpolated from the identity (1, 0, 0, 0) to (\textbackslash cos\textbackslash frac\{\textbackslash pi\}\{4\}, 0, 0, \textbackslash sin\textbackslash frac\{\textbackslash pi\}\{4\}) --- a right-handed \textbackslash pi/2 rotation about the z axis --- over 16 samples of t \textbackslash in {[}0, 1{]}, rendered by plot\_slerp\_path in src/quadmath/viz/plots.py via the script quadmath/scripts/quaternion\_gallery.py. Because slerp negates the second quaternion whenever \textbackslash langle q\_0, q\_1 \textbackslash rangle \textless{} 0, the interpolation always follows the shorter arc of the geodesic on the rotation sphere: the rotation angle grows linearly from 0 to \textbackslash pi/2, so the site sweeps the quarter circle of radius 2\textbackslash sqrt\{2\} about the z axis at constant height z = 2 (blue line), from the start marker at t = 0 (green) to the end marker at t = 1 (red). Reproduce with uv run python quadmath/scripts/quaternion\_gallery.py.}]{../output/figures/quaternion_slerp_path.png}}
\caption{\textbf{Shortest-arc slerp geodesic from the identity to a
\(\pi/2\) rotation about the \(z\) axis.} The path traced by the
canonical \((2, 0, 0, 0)\) quadray axis point (embedded at \((2, 2, 2)\)
under \protect\passthrough{\lstinline!DEFAULT\_EMBEDDING!}) as the
rotation quaternion is interpolated from the identity \((1, 0, 0, 0)\)
to \((\cos\frac{\pi}{4}, 0, 0, \sin\frac{\pi}{4})\) --- a right-handed
\(\pi/2\) rotation about the \(z\) axis --- over 16 samples of
\(t \in [0, 1]\), rendered by
\protect\passthrough{\lstinline!plot\_slerp\_path!} in
\protect\passthrough{\lstinline!src/quadmath/viz/plots.py!} via the
script
\protect\passthrough{\lstinline!quadmath/scripts/quaternion\_gallery.py!}.
Because \protect\passthrough{\lstinline!slerp!} negates the second
quaternion whenever \(\langle q_0, q_1 \rangle < 0\), the interpolation
always follows the shorter arc of the geodesic on the rotation sphere:
the rotation angle grows linearly from \(0\) to \(\pi/2\), so the site
sweeps the quarter circle of radius \(2\sqrt{2}\) about the \(z\) axis
at constant height \(z = 2\) (blue line), from the start marker at
\(t = 0\) (green) to the end marker at \(t = 1\) (red). Reproduce with
\protect\passthrough{\lstinline!uv run python quadmath/scripts/quaternion\_gallery.py!}.}
\end{figure}

\subsubsection{Quaternion metrics}\label{code:quaternion_metrics}

\paragraph{\texorpdfstring{\texttt{angle\_error}}{angle\_error}}\label{code:angle_error}

Source: \passthrough{\lstinline!src/quadmath/core/metrics.py!} ---
geodesic rotation angle \(2\arccos(|\langle q_1, q_2\rangle|)\) between
two quaternions, in radians and confined to \([0, \pi]\) by clamping the
acos argument; the metric is symmetric in its arguments and invariant
under \(q \to -q\) (the quaternion double cover), with inputs normalized
internally so non-unit but non-zero quaternions are accepted.

\paragraph{\texorpdfstring{\texttt{quat\_log\_euclidean\_dispersion}}{quat\_log\_euclidean\_dispersion}}\label{code:quat_log_euclidean_dispersion}

Source: \passthrough{\lstinline!src/quadmath/core/metrics.py!} ---
root-mean-square chordal dispersion of a quaternion set about its
normalized component-wise mean,
\(\sqrt{\mathrm{mean}_i\,\|q_i - m\|^2}\) in \(\mathbb{R}^4\), with
signs aligned to the first quaternion before averaging because \(q\) and
\(-q\) encode the same rotation; returns 0 for a single quaternion and
raises if the sign-aligned mean vanishes.

\subsubsection{Volume calculations}\label{code:volume_calculations}

\paragraph{\texorpdfstring{\texttt{integer\_tetra\_volume}}{integer\_tetra\_volume}}\label{code:integer_tetra_volume}

Source: \passthrough{\lstinline!src/quadmath/core/quadray.py!} --- exact
projected \(3\times3\) determinant over 4, i.e.~the IVM tetravolume
\(|\det|/4\) as a \passthrough{\lstinline!fractions.Fraction!}.

\paragraph{\texorpdfstring{\texttt{ace\_tetravolume\_5x5}}{ace\_tetravolume\_5x5}}\label{code:ace_tetravolume_5x5}

Source: \passthrough{\lstinline!src/quadmath/core/quadray.py!} --- Tom
Ace 5×5 determinant in IVM units.

\paragraph{\texorpdfstring{\texttt{tetra\_volume\_cayley\_menger}}{tetra\_volume\_cayley\_menger}}\label{code:tetra_volume_cayley_menger}

Source: \passthrough{\lstinline!src/quadmath/core/cayley\_menger.py!}
--- length-based formula (XYZ units).

\paragraph{\texorpdfstring{\texttt{ivm\_tetra\_volume\_cayley\_menger}}{ivm\_tetra\_volume\_cayley\_menger}}\label{code:ivm_tetra_volume_cayley_menger}

Source: \passthrough{\lstinline!src/quadmath/core/cayley\_menger.py!}
--- Cayley--Menger volume converted to IVM units.

\subsubsection{Coordinate
conversions}\label{code:coordinate_conversions}

\paragraph{\texorpdfstring{\texttt{urner\_embedding}}{urner\_embedding}}\label{code:urner_embedding}

Source: \passthrough{\lstinline!src/quadmath/lattice/conversions.py!}
--- canonical XYZ embedding.

\paragraph{\texorpdfstring{\texttt{quadray\_to\_xyz}}{quadray\_to\_xyz}}\label{code:quadray_to_xyz}

Source: \passthrough{\lstinline!src/quadmath/lattice/conversions.py!}
--- apply embedding matrix to map Quadray to XYZ.

\subsubsection{Linear algebra utilities}\label{code:linear_algebra}

\paragraph{\texorpdfstring{\texttt{bareiss\_determinant\_int}}{bareiss\_determinant\_int}}\label{code:bareiss_determinant_int}

Source: \passthrough{\lstinline!src/quadmath/core/linalg\_utils.py!} ---
exact integer Bareiss determinant.

\subsubsection{Optimization methods}\label{code:optimization}

\paragraph{\texorpdfstring{\texttt{nelder\_mead\_quadray}}{nelder\_mead\_quadray}}\label{code:nelder_mead_quadray}

Source:
\passthrough{\lstinline!src/quadmath/optimize/nelder\_mead\_quadray.py!}
--- Nelder--Mead optimization adapted to the integer quadray lattice.

\paragraph{\texorpdfstring{\texttt{discrete\_ivm\_descent}}{discrete\_ivm\_descent}}\label{code:discrete_ivm_descent}

Source:
\passthrough{\lstinline!src/quadmath/optimize/discrete\_variational.py!}
--- greedy integer-valued descent over the IVM using canonical neighbor
moves; returns a \passthrough{\lstinline!DiscretePath!} with visited
Quadrays and objective values.

\paragraph{\texorpdfstring{\texttt{neighbor\_moves\_ivm}}{neighbor\_moves\_ivm}}\label{code:neighbor_moves_ivm}

Source:
\passthrough{\lstinline!src/quadmath/optimize/discrete\_variational.py!}
--- return the 12 canonical IVM neighbor moves as Quadray deltas.

\paragraph{\texorpdfstring{\texttt{apply\_move}}{apply\_move}}\label{code:apply_move}

Source:
\passthrough{\lstinline!src/quadmath/optimize/discrete\_variational.py!}
--- apply a lattice move and normalize to the canonical representative.

\subsubsection{Information geometry
methods}\label{code:information_geometry}

\paragraph{\texorpdfstring{\texttt{fisher\_information\_matrix}}{fisher\_information\_matrix}}\label{code:fisher_information_matrix}

Source: \passthrough{\lstinline!src/quadmath/inference/information.py!}
--- empirical outer-product estimator.

\paragraph{\texorpdfstring{\texttt{fisher\_information\_quadray}}{fisher\_information\_quadray}}\label{code:fisher_information_quadray}

Source: \passthrough{\lstinline!src/quadmath/inference/information.py!}
--- compute Fisher information matrix in both Cartesian and Quadray
coordinates.

\paragraph{\texorpdfstring{\texttt{natural\_gradient\_step}}{natural\_gradient\_step}}\label{code:natural_gradient_step}

Source: \passthrough{\lstinline!src/quadmath/inference/information.py!}
--- damped inverse-Fisher step.

\paragraph{\texorpdfstring{\texttt{free\_energy}}{free\_energy}}\label{code:free_energy}

Source: \passthrough{\lstinline!src/quadmath/inference/information.py!}
--- discrete-state variational free energy.

\paragraph{\texorpdfstring{\texttt{expected\_free\_energy}}{expected\_free\_energy}}\label{code:expected_free_energy}

Source: \passthrough{\lstinline!src/quadmath/inference/information.py!}
--- canonical expected free energy \(G\) for Active Inference (epistemic
KL + ambiguity − pragmatic; see the equations appendix).

\paragraph{\texorpdfstring{\texttt{active\_inference\_step}}{active\_inference\_step}}\label{code:active_inference_step}

Source: \passthrough{\lstinline!src/quadmath/inference/information.py!}
--- joint perception-action update step in Active Inference.

\paragraph{\texorpdfstring{\texttt{information\_geometric\_distance}}{information\_geometric\_distance}}\label{code:information_geometric_distance}

Source: \passthrough{\lstinline!src/quadmath/inference/information.py!}
--- compute information-geometric distance between two points.

\paragraph{\texorpdfstring{\texttt{perception\_update}}{perception\_update}}\label{code:perception_update}

Source: \passthrough{\lstinline!src/quadmath/inference/information.py!}
--- continuous-time perception update: dμ/dt = D μ - dF/dμ.

\paragraph{\texorpdfstring{\texttt{action\_update}}{action\_update}}\label{code:action_update}

Source: \passthrough{\lstinline!src/quadmath/inference/information.py!}
--- continuous-time action update: da/dt = -dF/da.

\paragraph{\texorpdfstring{\texttt{finite\_difference\_gradient}}{finite\_difference\_gradient}}\label{code:finite_difference_gradient}

Source: \passthrough{\lstinline!src/quadmath/inference/information.py!}
--- compute numerical gradient of a scalar function via central
differences.

\subsubsection{Metrics and analysis}\label{code:metrics}

\paragraph{\texorpdfstring{\texttt{shannon\_entropy}}{shannon\_entropy}}\label{code:shannon_entropy}

Source: \passthrough{\lstinline!src/quadmath/core/metrics.py!} ---
Shannon entropy H(p) for a discrete distribution.

\paragraph{\texorpdfstring{\texttt{information\_length}}{information\_length}}\label{code:information_length}

Source: \passthrough{\lstinline!src/quadmath/core/metrics.py!} --- path
length in information space via gradient-weighted arc length.

\paragraph{\texorpdfstring{\texttt{fim\_eigenspectrum}}{fim\_eigenspectrum}}\label{code:fim_eigenspectrum}

Source: \passthrough{\lstinline!src/quadmath/core/metrics.py!} ---
eigen-decomposition of a Fisher information matrix.

\paragraph{\texorpdfstring{\texttt{fisher\_condition\_number}}{fisher\_condition\_number}}\label{code:fisher_condition_number}

Source: \passthrough{\lstinline!src/quadmath/core/metrics.py!} ---
compute the condition number of the Fisher information matrix.

\paragraph{\texorpdfstring{\texttt{fisher\_curvature\_analysis}}{fisher\_curvature\_analysis}}\label{code:fisher_curvature_analysis}

Source: \passthrough{\lstinline!src/quadmath/core/metrics.py!} ---
comprehensive analysis of Fisher information matrix curvature.

\paragraph{\texorpdfstring{\texttt{fisher\_quadray\_comparison}}{fisher\_quadray\_comparison}}\label{code:fisher_quadray_comparison}

Source: \passthrough{\lstinline!src/quadmath/core/metrics.py!} ---
compare Fisher information matrices between coordinate systems.

\subsubsection{Examples and utilities}\label{code:examples}

\paragraph{\texorpdfstring{\texttt{example\_ivm\_neighbors}}{example\_ivm\_neighbors}}\label{code:example_ivm_neighbors}

Source: \passthrough{\lstinline!src/quadmath/core/examples.py!} ---
return the 12 nearest IVM neighbors as permutations of \{2,1,1,0\}
(neighbor-move-set role).

\paragraph{\texorpdfstring{\texttt{example\_cuboctahedron\_neighbors}}{example\_cuboctahedron\_neighbors}}\label{code:example_cuboctahedron_neighbors}

Source: \passthrough{\lstinline!src/quadmath/core/examples.py!} ---
return twelve-around-one IVM neighbors (vector-equilibrium-shell role).
Same set as \passthrough{\lstinline!example\_ivm\_neighbors!}; both
public names are documented API sharing one sorted, deterministic
implementation.

\paragraph{\texorpdfstring{\texttt{example\_cuboctahedron\_vertices\_xyz}}{example\_cuboctahedron\_vertices\_xyz}}\label{code:example_cuboctahedron_vertices_xyz}

Source: \passthrough{\lstinline!src/quadmath/core/examples.py!} ---
return XYZ coordinates for the twelve-around-one neighbors.

\paragraph{\texorpdfstring{\texttt{example\_partition\_tetra\_volume}}{example\_partition\_tetra\_volume}}\label{code:example_partition_tetra_volume}

Source: \passthrough{\lstinline!src/quadmath/core/examples.py!} ---
construct a tetrahedron from the four-fold partition and return the
exact IVM tetravolume (\passthrough{\lstinline!fractions.Fraction!}).

\subsubsection{Symbolic computation}\label{code:symbolic}

\paragraph{\texorpdfstring{\texttt{cayley\_menger\_volume\_symbolic}}{cayley\_menger\_volume\_symbolic}}\label{code:cayley_menger_volume_symbolic}

Source: \passthrough{\lstinline!src/quadmath/core/symbolic.py!} ---
return symbolic Euclidean tetrahedron volume from squared distances.

\paragraph{\texorpdfstring{\texttt{convert\_xyz\_volume\_to\_ivm\_symbolic}}{convert\_xyz\_volume\_to\_ivm\_symbolic}}\label{code:convert_xyz_volume_to_ivm_symbolic}

Source: \passthrough{\lstinline!src/quadmath/core/symbolic.py!} ---
convert a symbolic Euclidean volume to IVM tetravolume via S3.

\subsubsection{Visualization and animation}\label{code:visualization}

\paragraph{\texorpdfstring{\texttt{animate\_discrete\_path}}{animate\_discrete\_path}}\label{code:animate_discrete_path}

Source: \passthrough{\lstinline!src/quadmath/viz/visualize.py!} ---
animate a \passthrough{\lstinline!DiscretePath!} to MP4; saves CSV/NPZ
trajectory to \passthrough{\lstinline!quadmath/output/!}.

\paragraph{\texorpdfstring{\texttt{plot\_ivm\_neighbors}}{plot\_ivm\_neighbors}}\label{code:plot_ivm_neighbors}

Source: \passthrough{\lstinline!src/quadmath/viz/visualize.py!} ---
scatter the 12 IVM neighbor points in 3D.

\paragraph{\texorpdfstring{\texttt{plot\_partition\_tetrahedron}}{plot\_partition\_tetrahedron}}\label{code:plot_partition_tetrahedron}

Source: \passthrough{\lstinline!src/quadmath/viz/visualize.py!} --- plot
the four-fold partition as a labeled tetrahedron in 3D.

\paragraph{\texorpdfstring{\texttt{animate\_simplex}}{animate\_simplex}}\label{code:animate_simplex}

Source: \passthrough{\lstinline!src/quadmath/viz/visualize.py!} ---
animate simplex evolution across iterations.

\paragraph{\texorpdfstring{\texttt{plot\_simplex\_trace}}{plot\_simplex\_trace}}\label{code:plot_simplex_trace}

Source: \passthrough{\lstinline!src/quadmath/viz/visualize.py!} --- plot
per-iteration diagnostics for Nelder--Mead.

\subsubsection{Path and file utilities}\label{code:utilities}

\paragraph{\texorpdfstring{\texttt{get\_repo\_root}}{get\_repo\_root}}\label{code:get_repo_root}

Source: \passthrough{\lstinline!src/quadmath/paths.py!} ---
heuristically find repository root by walking up from start.

\paragraph{\texorpdfstring{\texttt{get\_output\_dir}}{get\_output\_dir}}\label{code:get_output_dir}

Source: \passthrough{\lstinline!src/quadmath/paths.py!} --- return
\passthrough{\lstinline!quadmath/output!} path at the repo root and
ensure it exists.

\paragraph{\texorpdfstring{\texttt{get\_data\_dir}}{get\_data\_dir}}\label{code:get_data_dir}

Source: \passthrough{\lstinline!src/quadmath/paths.py!} --- return
\passthrough{\lstinline!quadmath/output/data!} path and ensure it
exists.

\paragraph{\texorpdfstring{\texttt{get\_figure\_dir}}{get\_figure\_dir}}\label{code:get_figure_dir}

Source: \passthrough{\lstinline!src/quadmath/paths.py!} --- return
\passthrough{\lstinline!quadmath/output/figures!} path and ensure it
exists.

\subsubsection{Additional Nelder--Mead
components}\label{code:nelder_mead_components}

\paragraph{\texorpdfstring{\texttt{SimplexState}}{SimplexState}}\label{code:SimplexState}

Source:
\passthrough{\lstinline!src/quadmath/optimize/nelder\_mead\_quadray.py!}
--- optimization trajectory state containing vertices, values, volume,
and history.

\paragraph{\texorpdfstring{\texttt{order\_simplex}}{order\_simplex}}\label{code:order_simplex}

Source:
\passthrough{\lstinline!src/quadmath/optimize/nelder\_mead\_quadray.py!}
--- sort vertices by objective value ascending and return paired lists.

\paragraph{\texorpdfstring{\texttt{centroid\_excluding}}{centroid\_excluding}}\label{code:centroid_excluding}

Source:
\passthrough{\lstinline!src/quadmath/optimize/nelder\_mead\_quadray.py!}
--- integer centroid of three vertices, excluding the specified index.

\paragraph{\texorpdfstring{\texttt{project\_to\_lattice}}{project\_to\_lattice}}\label{code:project_to_lattice}

Source:
\passthrough{\lstinline!src/quadmath/optimize/nelder\_mead\_quadray.py!}
--- project a quadray to the canonical lattice representative via
normalize.

\paragraph{\texorpdfstring{\texttt{compute\_volume}}{compute\_volume}}\label{code:compute_volume}

Source:
\passthrough{\lstinline!src/quadmath/optimize/nelder\_mead\_quadray.py!}
--- exact IVM tetra-volume
(\passthrough{\lstinline!fractions.Fraction!}, \(|\det|/4\)) from the
first four vertices.

\subsubsection{Discrete variational
components}\label{code:discrete_variational}

\paragraph{\texorpdfstring{\texttt{DiscretePath}}{DiscretePath}}\label{code:DiscretePath}

Source:
\passthrough{\lstinline!src/quadmath/optimize/discrete\_variational.py!}
--- optimization trajectory on the integer quadray lattice.

\subsubsection{Glossary generation}\label{code:glossary}

\paragraph{\texorpdfstring{\texttt{build\_api\_index}}{build\_api\_index}}\label{code:build_api_index}

Source: \passthrough{\lstinline!src/quadmath/tools/glossary\_gen.py!}
--- build API index from source directory.

\paragraph{\texorpdfstring{\texttt{generate\_markdown\_table}}{generate\_markdown\_table}}\label{code:generate_markdown_table}

Source: \passthrough{\lstinline!src/quadmath/tools/glossary\_gen.py!}
--- generate markdown table from API entries.

\paragraph{\texorpdfstring{\texttt{inject\_between\_markers}}{inject\_between\_markers}}\label{code:inject_between_markers}

Source: \passthrough{\lstinline!src/quadmath/tools/glossary\_gen.py!}
--- inject payload between markers in markdown text.

\subsubsection{Geometry utilities}\label{code:geometry}

\paragraph{\texorpdfstring{\texttt{minkowski\_interval}}{minkowski\_interval}}\label{code:minkowski_interval}

Source: \passthrough{\lstinline!src/quadmath/core/geometry.py!} ---
return the Minkowski interval squared ds² (Einstein.4D).

Relevant tests (\passthrough{\lstinline!tests/!}):

\begin{itemize}
\tightlist
\item
  \passthrough{\lstinline!test\_quadray.py!} (unit IVM tetra, primitive
  tetra = 1/4, exact scaled volume, Ace vs.~integer agreement)
\item
  \passthrough{\lstinline!test\_quadray\_cov.py!} (Ace determinant basic
  check)
\item
  \passthrough{\lstinline!test\_cayley\_menger.py!} (regular tetra
  volume in XYZ units)
\item
  \passthrough{\lstinline!test\_linalg\_utils.py!} (Bareiss determinant
  behavior)
\item
  \passthrough{\lstinline!test\_examples.py!},
  \passthrough{\lstinline!test\_examples\_cov.py!} (neighbors, examples)
\item
  \passthrough{\lstinline!test\_metrics.py!},
  \passthrough{\lstinline!test\_metrics\_cov.py!},
  \passthrough{\lstinline!test\_information.py!},
  \passthrough{\lstinline!test\_paths.py!},
  \passthrough{\lstinline!test\_paths\_cov.py!}
\end{itemize}

\subsubsection{Comprehensive test
coverage}\label{comprehensive-test-coverage}

The test suite provides 100\% coverage of all source modules with the
following focus areas:

\paragraph{Core functionality tests}\label{core-functionality-tests}

\begin{itemize}
\tightlist
\item
  \textbf{\passthrough{\lstinline!test\_quadray.py!}}: Quadray class
  operations, normalization, volume calculations, vector operations
\item
  \textbf{\passthrough{\lstinline!test\_cayley\_menger.py!}}:
  Cayley--Menger determinant validation, IVM conversion accuracy
\item
  \textbf{\passthrough{\lstinline!test\_linalg\_utils.py!}}: Bareiss
  algorithm correctness, edge cases, numerical stability
\end{itemize}

\paragraph{Optimization and variational
methods}\label{optimization-and-variational-methods}

\begin{itemize}
\tightlist
\item
  \textbf{\passthrough{\lstinline!test\_nelder\_mead\_visual.py!}}:
  Nelder--Mead adaptation to quadray lattice, simplex evolution
\item
  \textbf{\passthrough{\lstinline!test\_discrete\_variational.py!}}: IVM
  neighbor moves, discrete descent algorithms, path tracking
\item
  \textbf{\passthrough{\lstinline!test\_active\_inference.py!}}: Free
  energy minimization, perception-action updates
\end{itemize}

\paragraph{Information geometry and
metrics}\label{information-geometry-and-metrics}

\begin{itemize}
\tightlist
\item
  \textbf{\passthrough{\lstinline!test\_information.py!}}: Fisher
  information matrix computation, natural gradient steps
\item
  \textbf{\passthrough{\lstinline!test\_metrics.py!}}: FIM
  eigendecomposition, curvature analysis, coordinate system comparisons
\item
  \textbf{\passthrough{\lstinline!test\_information\_cov.py!}}: Extended
  coverage of information geometry functions
\end{itemize}

\paragraph{Examples and utilities}\label{examples-and-utilities}

\begin{itemize}
\tightlist
\item
  \textbf{\passthrough{\lstinline!test\_examples.py!}}: IVM neighbor
  constructions, polyhedra volume relationships
\item
  \textbf{\passthrough{\lstinline!test\_examples\_cov.py!}}: Extended
  example function coverage
\item
  \textbf{\passthrough{\lstinline!test\_paths.py!}}: Repository
  structure utilities, output directory management
\end{itemize}

\paragraph{Visualization and symbolic
computation}\label{visualization-and-symbolic-computation}

\begin{itemize}
\tightlist
\item
  \textbf{\passthrough{\lstinline!test\_visualize.py!}}: Plotting
  functions, animation generation, data export
\item
  \textbf{\passthrough{\lstinline!test\_visualize\_cov.py!}}: Extended
  visualization coverage, edge case handling
\item
  \textbf{\passthrough{\lstinline!test\_symbolic.py!}}: SymPy symbolic
  volume calculations, exact arithmetic
\item
  \textbf{\passthrough{\lstinline!test\_symbolic\_cov.py!}}: Symbolic
  computation error handling
\item
  \textbf{\passthrough{\lstinline!test\_sympy\_formalisms.py!}}:
  End-to-end symbolic workflow validation
\end{itemize}

\paragraph{API and documentation
generation}\label{api-and-documentation-generation}

\begin{itemize}
\tightlist
\item
  \textbf{\passthrough{\lstinline!test\_glossary\_gen.py!}}: Automatic
  API index generation, markdown table formatting
\item
  \textbf{\passthrough{\lstinline!test\_glossary\_gen\_cov.py!}}:
  Extended glossary generation coverage
\end{itemize}

\subsection{Reproducibility checklist}\label{reproducibility-checklist}

\begin{itemize}
\tightlist
\item
  \textbf{Complete implementation coverage}: All formulas and methods
  discussed in the paper are implemented in
  \passthrough{\lstinline!src/!} modules and verified by comprehensive
  test suites in \passthrough{\lstinline!tests/!}.
\item
  \textbf{Exact arithmetic for integer inputs}: Determinants are
  computed using the Bareiss algorithm for exact integer arithmetic;
  floating-point paths are used only where appropriate and results are
  converted (e.g., via S3) as specified.
\item
  \textbf{Deterministic random experiments}: Random-walk experiments use
  fixed seeds and produce integer volumes; Ace 5×5 determinant agrees
  with length-based methods across all test cases.
\item
  \textbf{Volume tracking and convergence}: Integer simplex volume
  monitoring detects convergence plateaus; face/edge analyses interpret
  sensitivity along edges and enable subspace searches across faces.
\item
  \textbf{Test-driven development}: All source code follows TDD
  principles with 100\% coverage requirements enforced via
  \passthrough{\lstinline!.coveragerc!} configuration.
\item
  \textbf{Figure and data generation}: All figures referenced in
  documentation are generated by scripts in
  \passthrough{\lstinline!quadmath/scripts/!} that import from
  \passthrough{\lstinline!src/!} modules, ensuring code-documentation
  coherence.
\item
  \textbf{Cross-platform compatibility}: Headless matplotlib backend
  (MPLBACKEND=Agg) ensures CI compatibility; deterministic RNG seeds
  guarantee reproducible outputs.
\item
  \textbf{Dependency management}: All dependencies managed through
  \passthrough{\lstinline!uv!} and
  \passthrough{\lstinline!pyproject.toml!} with exact version pinning
  via \passthrough{\lstinline!uv.lock!}.
\item
  \textbf{Path resolution}: No hardcoded paths; all output directories
  resolved via \passthrough{\lstinline!paths.py!} utilities for
  cross-platform compatibility.
\item
  \textbf{Data export formats}: Figures saved alongside CSV/NPZ data for
  complete reproducibility; all generation scripts print output paths
  for manifest collection.
\end{itemize}

All source code, tests, and documentation are available in the
docxology/QuadMath repository, ensuring complete transparency and
reproducibility of the methods described herein.

\newpage

\section{Optimization in 4D}\label{optimization-in-4d}

\subsection{Overview}\label{overview-2}

This section describes optimization methods adapted to the integer
Quadray lattice: a discrete Nelder--Mead simplex method with lattice
projection, a greedy 12-neighbor descent, and Fisher-metric
(natural-gradient) guidance for continuous parameter spaces. The integer
lattice supplies natural quantization of both steps and simplex volumes,
which the convergence criteria below exploit; higher-dimensional
extensions are developed in Section 5
(\passthrough{\lstinline!05\_extensions.md!}).

\subsection{Nelder--Mead on Integer
Lattice}\label{neldermead-on-integer-lattice}

\begin{itemize}
\tightlist
\item
  \textbf{Adaptation}: the four standard simplex operations ---
  reflection (\(\alpha\)), expansion (\(\gamma\)), contraction
  (\(\rho\)), and shrink (\(\sigma\)) --- applied to a simplex of four
  lattice points, with the objective \(f\) evaluated on the embedded
  coordinates \((x, y, z)\) of each vertex.
\item
  \textbf{Projection}: candidate moves are formed with per-component
  integer truncation of the scaled offsets (the centroid of the best
  three uses floor division), then mapped back to the canonical lattice
  representative by projective normalization, so the entire trajectory
  stays on the lattice.
\item
  \textbf{Volume tracking}: monitor the exact IVM tetravolume (absolute
  quadray determinant divided by 4, kept as an exact Fraction; Section
  3) as a convergence diagnostic; discrete steps create stable volume
  plateaus.
\item
  \textbf{Degenerate-simplex restart}: a zero-volume
  (collinear/coplanar) simplex cannot leave its own affine subspace,
  since every NM move is an affine combination of the vertices. When the
  simplex collapses (zero volume with spread below tolerance), the
  algorithm probes \(\pm 1\) and \(\pm 2\) lattice steps along each
  quadray axis from the best vertex; if any probe improves the best
  value, it re-seeds a full-volume simplex along the three most
  promising axes (a CVP-style restart, consuming one iteration),
  otherwise the collapse is accepted as convergence and the run
  terminates.
\end{itemize}

\subsubsection{Parameters}\label{parameters}

\begin{itemize}
\tightlist
\item
  \textbf{Reflection} \(\alpha \approx 1\)
\item
  \textbf{Expansion} \(\gamma \approx 2\)
\item
  \textbf{Contraction} \(\rho \approx 0.5\)
\item
  \textbf{Shrink} \(\sigma \approx 0.5\)
\end{itemize}

References: original Nelder--Mead method and common parameterizations in
optimization texts and survey articles; see overview:
\href{https://en.wikipedia.org/wiki/Nelder\%E2\%80\%93Mead_method}{Nelder--Mead
method}.

\subsection{Volume-Level Dynamics}\label{volume-level-dynamics}

\begin{itemize}
\tightlist
\item
  Simplex tetravolume decreases in discrete steps
  (\(\lvert\det\rvert/4\) values such as \(0.5\), \(0.25\), \(0\) in IVM
  units), producing plateaus (``energy levels'') between accepted moves.
\item
  Termination: a collapsed simplex (zero volume) with objective spread
  below tolerance \(\tau\) is checked by axis probes (\(\pm 1\),
  \(\pm 2\) lattice steps along each quadray axis from the best vertex);
  if no probe improves the best value, the run terminates, and if one
  does, a CVP-style restart re-seeds a full-volume simplex along the
  most promising axes. A degenerate simplex (zero volume) with spread
  still above \(\tau\) simply continues --- its affine moves remain
  confined to the collapse subspace while the spread shrinks toward
  collapse.
\item
  Monitoring: record best/worst objective, spread, and exact IVM volume
  at each iteration (see \passthrough{\lstinline!simplex\_trace.png!}
  below).
\end{itemize}

\subsection{Quadray Lattice Optimization
Pseudocode}\label{code:nelder_mead_on_integer_lattice}

\begin{lstlisting}
while not converged:
  order vertices by objective
  centroid of best three
  propose reflected (then possibly expanded/contracted) point
  accept per standard tests; else shrink toward best
  if simplex volume is zero and spread is below tolerance:
    probe +/-1 and +/-2 lattice steps along each quadray axis from best
    if any probe improves best: restart (full-volume simplex along 3 best axes); else stop
  update integer volume and function spread trackers
\end{lstlisting}

\subsubsection{Figures}\label{figures-1}

\begin{figure}
\centering
\pandocbounded{\includegraphics[keepaspectratio,alt={Per-iteration diagnostics for discrete Nelder--Mead on the integer Quadray lattice. Time series of the 13 recorded states (iterations 0--12, i.e.~12 optimization iterations after the initial simplex), plotted by plot\_simplex\_trace (src/quadmath/viz/visualize.py) from the run in quadmath/scripts/simplex\_animation.py for the penalized objective f(q) = (x-2)\^{}2 + (y-2)\^{}2 + (z-2)\^{}2 + 0.1\textbackslash,\textbackslash lvert x+y+z\textbackslash rvert, plus +5 in the quadrant x\textless0 \textbackslash wedge y\textless0 and +10 when any of \textbackslash lvert x\textbackslash rvert, \textbackslash lvert y\textbackslash rvert, \textbackslash lvert z\textbackslash rvert exceeds 4, evaluated at the embedded Cartesian coordinates (x, y, z) from to\_xyz with DEFAULT\_EMBEDDING. Left axis (dimensionless objective units): best value \textbackslash min\_i f(v\_i) (green), worst value \textbackslash max\_i f(v\_i) (faint red), and spread \textbackslash max\_i f(v\_i) - \textbackslash min\_i f(v\_i) (orange, dashed); only the best value is monotone --- the restart at iteration 9 raises the worst value from 3.5 to 4.4. Right axis (blue step line): exact tetravolume in IVM units (\textbackslash lvert\textbackslash det\textbackslash rvert/4, i.e.~0.5, 0.25, 0). The best value descends in plateaus (11.1 on 0--2, 4.8 on 3--6, 3.5 on 7--8, 0.6 on 9--12); the spread reaches zero at iterations 8 and 12, and the iteration-8 collapse triggers the CVP-style restart visible at iteration 9 (spread 0 \textbackslash to 3.8, volume 0 \textbackslash to 0.25). Regenerate with uv run python quadmath/scripts/simplex\_animation.py; raw data in quadmath/output/data/simplex\_trace.csv/.npz.}]{../output/figures/simplex_trace.png}}
\caption{\textbf{Per-iteration diagnostics for discrete Nelder--Mead on
the integer Quadray lattice}. Time series of the 13 recorded states
(iterations 0--12, i.e.~12 optimization iterations after the initial
simplex), plotted by
\protect\passthrough{\lstinline!plot\_simplex\_trace!}
(\protect\passthrough{\lstinline!src/quadmath/viz/visualize.py!}) from
the run in
\protect\passthrough{\lstinline!quadmath/scripts/simplex\_animation.py!}
for the penalized objective
\(f(q) = (x-2)^2 + (y-2)^2 + (z-2)^2 + 0.1\,\lvert x+y+z\rvert\), plus
\(+5\) in the quadrant \(x<0 \wedge y<0\) and \(+10\) when any of
\(\lvert x\rvert, \lvert y\rvert, \lvert z\rvert\) exceeds 4, evaluated
at the embedded Cartesian coordinates \((x, y, z)\) from
\protect\passthrough{\lstinline!to\_xyz!} with
\protect\passthrough{\lstinline!DEFAULT\_EMBEDDING!}. \textbf{Left axis}
(dimensionless objective units): best value \(\min_i f(v_i)\) (green),
worst value \(\max_i f(v_i)\) (faint red), and spread
\(\max_i f(v_i) - \min_i f(v_i)\) (orange, dashed); only the best value
is monotone --- the restart at iteration 9 raises the worst value from
\(3.5\) to \(4.4\). \textbf{Right axis} (blue step line): exact
tetravolume in IVM units (\(\lvert\det\rvert/4\), i.e.~\(0.5\),
\(0.25\), \(0\)). The best value descends in plateaus (\(11.1\) on
\(0\)--\(2\), \(4.8\) on \(3\)--\(6\), \(3.5\) on \(7\)--\(8\), \(0.6\)
on \(9\)--\(12\)); the spread reaches zero at iterations 8 and 12, and
the iteration-8 collapse triggers the CVP-style restart visible at
iteration 9 (spread \(0 \to 3.8\), volume \(0 \to 0.25\)). Regenerate
with
\protect\passthrough{\lstinline!uv run python quadmath/scripts/simplex\_animation.py!};
raw data in
\protect\passthrough{\lstinline!quadmath/output/data/simplex\_trace.csv!}/\protect\passthrough{\lstinline!.npz!}.}
\end{figure}

The 2×2 panel below shows the simplex itself at iterations 0, 3, 6, and
9; the plateau structure of the trace above appears as iterations where
the simplex repositions without improving its best value.

\begin{figure}
\centering
\pandocbounded{\includegraphics[keepaspectratio,alt={Nelder--Mead simplex evolution on the integer Quadray lattice (2×2 panel). The four-vertex simplex (vertices as red spheres, the six tetrahedral edges as blue lines) at iterations 0, 3, 6, and 9 of the same run as the trace figure, plotted in the embedded (x, y, z) space produced by to\_xyz with DEFAULT\_EMBEDDING; each axis spans {[}-6, 6{]}, and each panel title reports that iteration's best objective value and spread. Top-left (iteration 0): the widely dispersed initial simplex (Quadray(5,0,0,0), Quadray(4,1,0,0), Quadray(0,4,1,0), Quadray(1,1,1,0)). Top-right (iteration 3): contraction toward the main basin. Bottom-left (iteration 6): near-collapsed simplex (spread 1.3) just before the degenerate restart. Bottom-right (iteration 9): the re-seeded simplex contracting again (spread 3.8); the run terminates at iteration 11 with all four vertices coinciding at Quadray(2,0,0,0), i.e.~embedded (2, 2, 2). Generated by quadmath/scripts/simplex\_animation.py.}]{../output/figures/simplex_final.png}}
\caption{\textbf{Nelder--Mead simplex evolution on the integer Quadray
lattice (2×2 panel)}. The four-vertex simplex (vertices as red spheres,
the six tetrahedral edges as blue lines) at iterations 0, 3, 6, and 9 of
the same run as the trace figure, plotted in the embedded \((x, y, z)\)
space produced by \protect\passthrough{\lstinline!to\_xyz!} with
\protect\passthrough{\lstinline!DEFAULT\_EMBEDDING!}; each axis spans
\([-6, 6]\), and each panel title reports that iteration's best
objective value and spread. \textbf{Top-left (iteration 0)}: the widely
dispersed initial simplex
(\protect\passthrough{\lstinline!Quadray(5,0,0,0)!},
\protect\passthrough{\lstinline!Quadray(4,1,0,0)!},
\protect\passthrough{\lstinline!Quadray(0,4,1,0)!},
\protect\passthrough{\lstinline!Quadray(1,1,1,0)!}). \textbf{Top-right
(iteration 3)}: contraction toward the main basin. \textbf{Bottom-left
(iteration 6)}: near-collapsed simplex (spread \(1.3\)) just before the
degenerate restart. \textbf{Bottom-right (iteration 9)}: the re-seeded
simplex contracting again (spread \(3.8\)); the run terminates at
iteration 11 with all four vertices coinciding at
\protect\passthrough{\lstinline!Quadray(2,0,0,0)!}, i.e.~embedded
\((2, 2, 2)\). Generated by
\protect\passthrough{\lstinline!quadmath/scripts/simplex\_animation.py!}.}
\end{figure}

\begin{figure}
\centering
\pandocbounded{\includegraphics[keepaspectratio,alt={Complete simplex vertex trajectories (3D). Full paths of the four simplex vertices across all 12 recorded iterations of the same run, in the same embedded (x, y, z) axes ({[}-6, 6{]} per direction) as the panel above. Each vertex trace uses a fixed color and marker (red circle, blue square, green triangle, orange diamond); large black-edged markers flag iterations 0, 4, and 8. The traces show coordinated contraction toward the converged point (2, 2, 2), with the degenerate-simplex restart visible as the outward jump near iteration 8. The script also draws a black star labeled ``Converged (0,0,0)'' at the embedded origin; the vertices themselves converge to (2, 2, 2), so read the star as a decorative marker, not the optimum. Generated by quadmath/scripts/simplex\_animation.py.}]{../output/figures/simplex_trace_visualization.png}}
\caption{\textbf{Complete simplex vertex trajectories (3D)}. Full paths
of the four simplex vertices across all 12 recorded iterations of the
same run, in the same embedded \((x, y, z)\) axes (\([-6, 6]\) per
direction) as the panel above. Each vertex trace uses a fixed color and
marker (red circle, blue square, green triangle, orange diamond); large
black-edged markers flag iterations 0, 4, and 8. The traces show
coordinated contraction toward the converged point \((2, 2, 2)\), with
the degenerate-simplex restart visible as the outward jump near
iteration 8. The script also draws a black star labeled ``Converged
(0,0,0)'' at the embedded origin; the vertices themselves converge to
\((2, 2, 2)\), so read the star as a decorative marker, not the optimum.
Generated by
\protect\passthrough{\lstinline!quadmath/scripts/simplex\_animation.py!}.}
\end{figure}

Raw artifacts: the full trajectory animation
\passthrough{\lstinline!simplex\_animation.mp4!} and per-frame vertices
(\passthrough{\lstinline!simplex\_animation\_vertices.csv!}/\passthrough{\lstinline!.npz!})
are available in \passthrough{\lstinline!quadmath/output/!}.

\subsection{Discrete Lattice Descent (Information-Theoretic
Variant)}\label{discrete-lattice-descent-information-theoretic-variant}

\begin{itemize}
\tightlist
\item
  Integer-valued greedy descent over the IVM: from the current lattice
  point \(q\), evaluate \(f\) at the 12 nearest neighbors (the distinct
  permutations of \((2, 1, 1, 0)\) in quadray offsets), then move to the
  minimizing neighbor.
\item
  Each accepted move strictly decreases \(f\), so the objective is
  monotone along the path; the walk terminates at a local minimum of
  \(f\) over the IVM adjacency, i.e.~when no neighbor improves on the
  current value.
\item
  The objective may be geometric (e.g.~Euclidean distance in an
  embedding) or information-theoretic (e.g.~a local free-energy proxy).
\item
  API: \passthrough{\lstinline!discrete\_ivm\_descent!} in
  \passthrough{\lstinline!src/quadmath/optimize/discrete\_variational.py!}.
  Animation helper: \passthrough{\lstinline!animate\_discrete\_path!} in
  \passthrough{\lstinline!src/quadmath/viz/visualize.py!}.
\end{itemize}

Short snippet (paper reproducibility):

\begin{lstlisting}[language=Python]
from quadmath.core.quadray import Quadray, DEFAULT_EMBEDDING, to_xyz
from quadmath.optimize.discrete_variational import discrete_ivm_descent
from quadmath.viz.visualize import animate_discrete_path

def f(q: Quadray) -> float:
    x, y, z = to_xyz(q, DEFAULT_EMBEDDING)
    return (x - 0.5)**2 + (y + 0.2)**2 + (z - 0.1)**2

path = discrete_ivm_descent(f, Quadray(6,0,0,0))
animate_discrete_path(path)
\end{lstlisting}

\subsection{Convergence and
Robustness}\label{convergence-and-robustness}

\begin{itemize}
\tightlist
\item
  Discrete steps reduce numerical drift; improved stability
  vs.~unconstrained Cartesian.
\item
  Natural regularization from volume quantization; fewer wasted
  evaluations.
\item
  Compatible with Gauss--Newton/Natural Gradient guidance using FIM for
  metric-aware steps (Amari, natural gradient).
\end{itemize}

\subsection{Information-Geometric View (Einstein.4D analogy in metric
form)}\label{information-geometric-view-einstein.4d-analogy-in-metric-form}

The Fisher Information Matrix (FIM) provides a fundamental bridge
between the three 4D frameworks, establishing a Riemannian metric on
parameter space that guides optimization through information geometry.
This section demonstrates how the FIM connects Coxeter.4D (Euclidean
parameter space), Einstein.4D (information-geometric flows), and
Fuller.4D (tetrahedral structure) in a unified optimization framework.

\subsubsection{Fisher Information as Riemannian
Metric}\label{fisher-information-as-riemannian-metric}

The empirical Fisher Information Matrix \(F_{ij}\) quantifies the local
curvature of the log-likelihood surface around parameter estimates,
providing a natural metric for parameter space geometry. This
fundamental concept in information geometry establishes a Riemannian
structure on the statistical manifold, where distances and angles are
measured according to the intrinsic geometry of the probability
distributions rather than the extrinsic Euclidean geometry of the
parameter space.

For a model with parameter vector \(\theta\) and per-sample scores
\(\partial_{\theta_i} \log p(x_n; \theta)\), the empirical FIM is the
average outer product of score functions,
\(F_{i,j} = \frac{1}{N} \sum_{n=1}^{N} \partial_{\theta_i} \log p(x_n; \theta)\, \partial_{\theta_j} \log p(x_n; \theta)\)
(Eq. \eqref{eq:fim_empirical} in the equations appendix). Diagonal
entries quantify parameter sensitivity and off-diagonal entries pairwise
parameter interactions. The running example below keeps the parameter
names \(\mathbf{w} = (w_0, w_1, w_2)\) used by
\passthrough{\lstinline!information\_demo.py!}; the two notations
coincide.

The Fisher Information Matrix serves as the natural metric tensor
\(g_{ij} = F_{ij}\) on the statistical manifold, replacing the Euclidean
metric \(\delta_{ij}\) with a data-dependent metric that reflects the
actual curvature structure of the objective function. This geometric
interpretation enables the application of differential geometry concepts
to optimization problems, where geodesics (locally distance-minimizing
paths) follow the natural gradient direction \(F^{-1}\nabla L\) rather
than the standard gradient \(\nabla L\).

The theoretical foundation of this approach stems from the work of
\href{https://en.wikipedia.org/wiki/Cram\%C3\%A9r\%E2\%80\%93Rao_bound}{Rao
(1945)} and \href{https://en.wikipedia.org/wiki/Shun-ichi_Amari}{Amari
(1985)}, who established information geometry as a framework for
analyzing statistical models through differential geometry. The FIM
naturally arises as the Hessian of the Kullback-Leibler divergence
between nearby probability distributions, making it the canonical choice
for measuring distances on the statistical manifold.

In the context of optimization, the FIM provides several key advantages:

\begin{enumerate}
\def\labelenumi{\arabic{enumi}.}
\item
  \textbf{Invariance to parameterization}: The natural gradient
  \(F^{-1}\nabla L\) is invariant to smooth, invertible parameter
  transformations, unlike the standard gradient which depends on the
  choice of coordinate system.
\item
  \textbf{Optimal step sizing}: The FIM automatically determines
  appropriate step sizes in different parameter directions, scaling
  updates according to local curvature.
\item
  \textbf{Geometric consistency}: Optimization follows geodesics on the
  statistical manifold, respecting the intrinsic geometry of the
  parameter space rather than imposing an artificial Euclidean
  structure.
\end{enumerate}

This geometric approach to optimization is particularly powerful in the
context of the 4D frameworks, where it provides a unified mathematical
language for describing optimization dynamics across different geometric
paradigms.

\subsubsection{4D Framework Integration through Fisher
Information}\label{d-framework-integration-through-fisher-information}

\textbf{Coxeter.4D (Euclidean)}: In standard Euclidean parameter space,
the metric tensor is simply \(\delta_{ij}\), providing uniform scaling
in all directions. The FIM \(F_{ij}\) generalizes this to capture the
actual curvature structure of the objective function.

\textbf{Einstein.4D (Minkowski analogy)}: the Fisher metric
\(g_{ij} = F_{ij}\) replaces the flat metric, and geodesic motion on the
statistical manifold corresponds to the natural gradient update
\(\theta \leftarrow \theta - \eta\, F(\theta)^{-1} \nabla_\theta L(\theta)\)
(Eq. \eqref{eq:natural_gradient} in the equations appendix) --- steepest
descent measured in the Fisher metric rather than straight-line motion
in parameter space.

\textbf{Fuller.4D (Synergetics)}: The tetrahedral structure of Quadray
coordinates naturally encodes the four-fold partition of optimization
problems, while the FIM provides the metric structure for efficient
navigation through this space. The discrete nature of the IVM lattice
creates natural quantization effects that can be exploited for
computational efficiency.

\subsubsection{Comprehensive Fisher Information
Analysis}\label{comprehensive-fisher-information-analysis}

Two figures summarize the empirical Fisher analysis of the regression
example introduced below: the matrix structure and its eigenspectrum.
Both are generated by
\passthrough{\lstinline!quadmath/scripts/information\_demo.py!} and
interpreted through the three 4D frameworks.

\begin{figure}
\centering
\pandocbounded{\includegraphics[keepaspectratio,alt={Empirical Fisher Information Matrix with 4D framework context (three panels; generated by quadmath/scripts/information\_demo.py with seed default\_rng(0)). The model is a 3-parameter linear regression y = \textbackslash mathbf\{x\}\^{}\textbackslash top \textbackslash mathbf\{w\} + \textbackslash varepsilon on N = 200 samples with standard-Gaussian features and noise \textbackslash varepsilon \textbackslash sim \textbackslash mathcal\{N\}(0, 0.1\^{}2), generated at \textbackslash mathbf\{w\}\_\{\textbackslash text\{true\}\} = (1.0, -2.0, 0.5) and evaluated at the deliberately misspecified estimate \textbackslash mathbf\{w\}\_\{\textbackslash text\{est\}\} = (0.3, -1.2, 0.0). Left panel: the data points with the 1-D slice y(x) = w\_0 + w\_1 x + w\_2 x\^{}2 of both parameter vectors (green solid: true; red dashed: estimate) and the MSE annotation. Center panel: the 3 \textbackslash times 3 matrix of per-sample squared-loss gradients 2\textbackslash,x\_\{ni\}\textbackslash,r\_n (a Gauss--Newton/Gram surrogate of the FIM; see the estimator note below) as an annotated heatmap, with diagonal entries \textbackslash approx (9.39, 12.66, 5.60) and off-diagonal entries \textbackslash approx (-4.09, 2.05, -3.34); units are squared-loss gradient products. Right panel: schematic tetrahedron relating Coxeter.4D (Euclidean parameter space, metric \textbackslash delta\_\{ij\}), Einstein.4D (Fisher metric replacing the flat metric), and Fuller.4D (tetrahedral/IVM structure).}]{../output/figures/fisher_information_matrix.png}}
\caption{\textbf{Empirical Fisher Information Matrix with 4D framework
context} (three panels; generated by
\protect\passthrough{\lstinline!quadmath/scripts/information\_demo.py!}
with seed \protect\passthrough{\lstinline!default\_rng(0)!}). The model
is a 3-parameter linear regression
\(y = \mathbf{x}^\top \mathbf{w} + \varepsilon\) on \(N = 200\) samples
with standard-Gaussian features and noise
\(\varepsilon \sim \mathcal{N}(0, 0.1^2)\), generated at
\(\mathbf{w}_{\text{true}} = (1.0, -2.0, 0.5)\) and evaluated at the
deliberately misspecified estimate
\(\mathbf{w}_{\text{est}} = (0.3, -1.2, 0.0)\). \textbf{Left panel}: the
data points with the 1-D slice \(y(x) = w_0 + w_1 x + w_2 x^2\) of both
parameter vectors (green solid: true; red dashed: estimate) and the MSE
annotation. \textbf{Center panel}: the \(3 \times 3\) matrix of
per-sample squared-loss gradients \(2\,x_{ni}\,r_n\) (a
Gauss--Newton/Gram surrogate of the FIM; see the estimator note below)
as an annotated heatmap, with diagonal entries
\(\approx (9.39, 12.66, 5.60)\) and off-diagonal entries
\(\approx (-4.09, 2.05, -3.34)\); units are squared-loss gradient
products. \textbf{Right panel}: schematic tetrahedron relating
Coxeter.4D (Euclidean parameter space, metric \(\delta_{ij}\)),
Einstein.4D (Fisher metric replacing the flat metric), and Fuller.4D
(tetrahedral/IVM structure).}
\end{figure}

Note on the estimator: the matrix in the figure above is computed from
per-sample \textbf{squared-loss} gradients
(\passthrough{\lstinline!2·x\_i·r\_i!}) at a misspecified
\passthrough{\lstinline!w\_est!}
(\passthrough{\lstinline!information\_demo.py!}), i.e.~a
Gauss--Newton/Gram surrogate. It coincides with the empirical FIM of Eq.
\eqref{eq:fim_empirical} only for true score functions --- exactly at
\(w_{\text{true}}\) in expectation under Gaussian noise --- so treat the
displayed matrix as a curvature-scale illustration rather than a model
FIM estimate.

\begin{figure}
\centering
\pandocbounded{\includegraphics[keepaspectratio,alt={Fisher Information eigenspectrum and parameter-space curvature (three panels; generated by quadmath/scripts/information\_demo.py from the same seeded run as the matrix figure above). Left panel: bar chart of the eigenvalues of the empirical FIM, sorted descending and annotated: \textbackslash lambda\_0 \textbackslash approx 16.80, \textbackslash lambda\_1 \textbackslash approx 6.63, \textbackslash lambda\_2 \textbackslash approx 4.22 (units: squared-loss gradient products) --- the principal curvature scales of the loss surface. Center panel: text summary of the curvature metrics: condition number \textbackslash lambda\_\{\textbackslash max\}/\textbackslash lambda\_\{\textbackslash min\} \textbackslash approx 3.98 (anisotropy), anisotropy index \textbackslash approx 0.59, and total curvature (trace of F) \textbackslash approx 27.65, with per-direction and 4D-framework interpretation. Right panel: a parameter-space tetrahedron with one vertex at the origin and three vertices along the eigenvector directions scaled by \textbackslash sqrt\{\textbackslash lambda\_i\}, so the tetrahedron's shape encodes the anisotropy; edge colors mark eigenvalue rank, and vertices are labeled with their eigenvalues. Large eigenvalues mark directions of rapid objective change (where the natural gradient takes small steps); small eigenvalues mark flat directions (larger steps).}]{../output/figures/fisher_information_eigenspectrum.png}}
\caption{\textbf{Fisher Information eigenspectrum and parameter-space
curvature} (three panels; generated by
\protect\passthrough{\lstinline!quadmath/scripts/information\_demo.py!}
from the same seeded run as the matrix figure above). \textbf{Left
panel}: bar chart of the eigenvalues of the empirical FIM, sorted
descending and annotated: \(\lambda_0 \approx 16.80\),
\(\lambda_1 \approx 6.63\), \(\lambda_2 \approx 4.22\) (units:
squared-loss gradient products) --- the principal curvature scales of
the loss surface. \textbf{Center panel}: text summary of the curvature
metrics: condition number \(\lambda_{\max}/\lambda_{\min} \approx 3.98\)
(anisotropy), anisotropy index \(\approx 0.59\), and total curvature
(trace of \(F\)) \(\approx 27.65\), with per-direction and 4D-framework
interpretation. \textbf{Right panel}: a parameter-space tetrahedron with
one vertex at the origin and three vertices along the eigenvector
directions scaled by \(\sqrt{\lambda_i}\), so the tetrahedron's shape
encodes the anisotropy; edge colors mark eigenvalue rank, and vertices
are labeled with their eigenvalues. Large eigenvalues mark directions of
rapid objective change (where the natural gradient takes small steps);
small eigenvalues mark flat directions (larger steps).}
\end{figure}

\subsubsection{Natural Gradient Descent: Geodesic Motion on Information
Manifold}\label{natural-gradient-descent-geodesic-motion-on-information-manifold}

The Fisher Information Matrix enables natural gradient descent, which
implements geodesic motion on the information manifold. Unlike standard
gradient descent that follows straight lines in parameter space, natural
gradient descent follows curved paths that respect the intrinsic
geometry defined by the FIM.

The natural gradient update used throughout this manuscript is Eq.
\eqref{eq:natural_gradient} in the equations appendix,
\(\theta \leftarrow \theta - \eta\, F(\theta)^{-1} \nabla_\theta L(\theta)\),
where \(\eta\) is the step size, \(F\) the empirical FIM of Eq.
\eqref{eq:fim_empirical}, and \(\nabla_\theta L\) the gradient of the
loss. Because \(F\) is the metric tensor \(g_{ij} = F_{ij}\) of the
statistical manifold, this update is steepest descent measured in the
Fisher metric --- geodesic motion rather than straight-line motion in
parameter space.

The theoretical foundation of natural gradient descent was established
by \href{https://en.wikipedia.org/wiki/Natural_gradient}{Amari (1998)}
in the context of information geometry. The key insight is that the
natural gradient \(F^{-1}\nabla L\) is the steepest descent direction
when distances are measured using the Fisher metric rather than the
Euclidean metric. This makes natural gradient descent invariant to
smooth, invertible parameter transformations, a property that standard
gradient descent lacks.

In the context of the 4D frameworks, natural gradient descent provides a
unified approach to optimization that respects the intrinsic geometry of
each framework:

\begin{itemize}
\tightlist
\item
  \textbf{Coxeter.4D}: The natural gradient respects the actual
  curvature structure of the objective function rather than imposing
  artificial Euclidean geometry.
\item
  \textbf{Einstein.4D}: The Fisher metric replaces the spacetime metric,
  creating geodesic flows that follow the intrinsic geometry of the
  parameter space.
\item
  \textbf{Fuller.4D}: The tetrahedral structure provides natural
  coordinate systems where the FIM can exhibit beneficial structural
  properties.
\end{itemize}

The efficiency of natural gradient descent comes from its ability to
automatically adapt step sizes to local curvature. In directions of high
curvature (large eigenvalues of \(F\)), the natural gradient takes
smaller steps, while in directions of low curvature (small eigenvalues),
it takes larger steps. This anisotropic scaling leads to faster
convergence and better numerical stability compared to standard gradient
descent.

\begin{figure}
\centering
\pandocbounded{\includegraphics[keepaspectratio,alt={Natural gradient trajectory on a quadratic bowl (generated by quadmath/scripts/information\_demo.py). The objective is the quadratic L(\textbackslash mathbf\{w\}) = \textbackslash frac\{1\}\{2\}(\textbackslash mathbf\{w\} - \textbackslash mathbf\{w\}\_\{\textbackslash text\{true\}\})\^{}\textbackslash top A\textbackslash, (\textbackslash mathbf\{w\} - \textbackslash mathbf\{w\}\_\{\textbackslash text\{true\}\}) with \textbackslash mathbf\{w\}\_\{\textbackslash text\{true\}\} = (1, -2, 0.5) and A the positive-definite matrix with diagonal (3, 2, 1) and off-diagonal A\_\{12\} = 0.5; the metric used for the steps is the empirical FIM F of the regression example above (plus a 10\^{}\{-3\} ridge for invertibility), not A. Starting from (2, 2, 2), the blue line with markers shows 20 updates of the form \textbackslash Delta\textbackslash mathbf\{w\} = -0.5\textbackslash, F\^{}\{-1\} A\textbackslash,(\textbackslash mathbf\{w\} - \textbackslash mathbf\{w\}\_\{\textbackslash text\{true\}\}) --- the 3-parameter trajectory (w\_0, w\_1, w\_2) projected onto the (w\_0, w\_1) plane (parameter units). Green circle: start; red circle: the final iterate (0.70, -1.40) after 20 steps --- a fixed-metric preconditioned run stopped at its step budget, not yet at \textbackslash mathbf\{w\}\_\{\textbackslash text\{true\}\} = (1, -2, 0.5). The anisotropic F\^{}\{-1\} preconditioning is visible as unequal progress along the two plotted coordinates.}]{../output/figures/natural_gradient_path.png}}
\caption{\textbf{Natural gradient trajectory on a quadratic bowl}
(generated by
\protect\passthrough{\lstinline!quadmath/scripts/information\_demo.py!}).
The objective is the quadratic
\(L(\mathbf{w}) = \frac{1}{2}(\mathbf{w} - \mathbf{w}_{\text{true}})^\top A\, (\mathbf{w} - \mathbf{w}_{\text{true}})\)
with \(\mathbf{w}_{\text{true}} = (1, -2, 0.5)\) and \(A\) the
positive-definite matrix with diagonal \((3, 2, 1)\) and off-diagonal
\(A_{12} = 0.5\); the metric used for the steps is the empirical FIM
\(F\) of the regression example above (plus a \(10^{-3}\) ridge for
invertibility), not \(A\). Starting from \((2, 2, 2)\), the blue line
with markers shows 20 updates of the form
\(\Delta\mathbf{w} = -0.5\, F^{-1} A\,(\mathbf{w} - \mathbf{w}_{\text{true}})\)
--- the 3-parameter trajectory \((w_0, w_1, w_2)\) projected onto the
\((w_0, w_1)\) plane (parameter units). Green circle: start; red circle:
the final iterate \((0.70, -1.40)\) after 20 steps --- a fixed-metric
preconditioned run stopped at its step budget, not yet at
\(\mathbf{w}_{\text{true}} = (1, -2, 0.5)\). The anisotropic \(F^{-1}\)
preconditioning is visible as unequal progress along the two plotted
coordinates.}
\end{figure}

\subsubsection{Quadray-Specific
Considerations}\label{quadray-specific-considerations}

Under Quadray parameterizations, the FIM often exhibits block-structured
and symmetric patterns that simplify matrix inversion for
natural-gradient steps. This structural regularity arises from the
tetrahedral symmetry of the IVM lattice and can be exploited for
computational efficiency.

The discrete nature of the IVM lattice also influences the FIM
structure, as parameter updates are constrained to integer coordinate
positions. This creates a natural regularization effect that can improve
optimization stability and convergence.

\subsubsection{Variational Free Energy and Active Inference
Integration}\label{variational-free-energy-and-active-inference-integration}

The Fisher Information framework naturally extends to variational
inference and active inference, where the free energy principle guides
both perception and action through information-geometric optimization.

\begin{figure}
\centering
\pandocbounded{\includegraphics[keepaspectratio,alt={Variational free energy for a 2-state toy model (generated by quadmath/scripts/information\_demo.py). The curve shows the free energy of Eq. , \textbackslash mathcal\{F\} = -\textbackslash log P(o\textbackslash mid s) + \textbackslash mathrm\{KL\}\textbackslash big{[}Q(s)\textbackslash,\textbackslash big\textbackslash Vert\textbackslash,P(s)\textbackslash big{]}, as a function of the variational parameter q = Q(\textbackslash text\{state\}=0) over the grid q \textbackslash in (0, 1) (200 points), for likelihood P(o\textbackslash mid s) = (0.7, 0.3), uniform prior P(s) = (0.5, 0.5), and variational family Q(s) = (q, 1 - q). X-axis: q (dimensionless probability). Y-axis: \textbackslash mathcal\{F\} in nats. The minimum (red marker) is \textbackslash mathcal\{F\} \textbackslash approx \textbackslash ln 2 \textbackslash approx 0.693 at q = 0.7, where Q coincides with the likelihood --- the standard variational result that the optimal Q in the unconstrained family is the posterior itself. In the 4D reading, minimizing \textbackslash mathcal\{F\} is geodesic motion under the Fisher metric on the variational manifold (Einstein.4D analogy).}]{../output/figures/free_energy_curve.png}}
\caption{\textbf{Variational free energy for a 2-state toy model}
(generated by
\protect\passthrough{\lstinline!quadmath/scripts/information\_demo.py!}).
The curve shows the free energy of Eq. \eqref{eq:free_energy},
\(\mathcal{F} = -\log P(o\mid s) + \mathrm{KL}\big[Q(s)\,\big\Vert\,P(s)\big]\),
as a function of the variational parameter \(q = Q(\text{state}=0)\)
over the grid \(q \in (0, 1)\) (200 points), for likelihood
\(P(o\mid s) = (0.7, 0.3)\), uniform prior \(P(s) = (0.5, 0.5)\), and
variational family \(Q(s) = (q, 1 - q)\). \textbf{X-axis}: \(q\)
(dimensionless probability). \textbf{Y-axis}: \(\mathcal{F}\) in nats.
The minimum (red marker) is \(\mathcal{F} \approx \ln 2 \approx 0.693\)
at \(q = 0.7\), where \(Q\) coincides with the likelihood --- the
standard variational result that the optimal \(Q\) in the unconstrained
family is the posterior itself. In the 4D reading, minimizing
\(\mathcal{F}\) is geodesic motion under the Fisher metric on the
variational manifold (Einstein.4D analogy).}
\end{figure}

For the full Active Inference treatment --- expected free energy,
perception and action updates, and further 4D natural-gradient
visualizations --- see \href{09_free_energy_active_inference.md}{Section
9: Free Energy and Active Inference}.

\subsection{Multi-Objective and Higher-Dimensional
Notes}\label{multi-objective-and-higher-dimensional-notes}

The extension of these ideas --- simplex faces as Pareto trade-off
surfaces, integer volume as a solution-diversity measure, and
higher-simplex volume decompositions --- is developed in Section 5
(\passthrough{\lstinline!05\_extensions.md!}).

\subsection{External Validation and Computational
Context}\label{external-validation-and-computational-context}

The methods above complement the computational framework in Kirby
Urner's \href{https://github.com/4dsolutions}{4dsolutions ecosystem};
implementation and educational context are collected in the
\href{07_resources.md}{Resources} section.

\subsection{Results}\label{results}

On the penalized quadratic of
\passthrough{\lstinline!quadmath/scripts/simplex\_animation.py!}, the
discrete Nelder--Mead reaches the lattice point
\passthrough{\lstinline!Quadray(2,0,0,0)!} (embedded \((2, 2, 2)\), best
objective \(0.6\)) in 12 recorded iterations, with the best value
descending through the plateaus \(11.1 \to 4.8 \to 3.5 \to 0.6\) and one
degenerate-simplex restart mid-run; see the simplex figures above and
the artifacts in \passthrough{\lstinline!quadmath/output/!}.

\newpage

\section{Extensions of 4D and
Quadrays}\label{extensions-of-4d-and-quadrays}

Here we review some extensions of the Quadray 4D framework, including
multi-objective optimization, machine learning, computer graphics and
GPU acceleration, active inference, complex systems, pedagogy, and
implementations, with an emphasis on cognitive security.

\subsection{Multi-Objective
Optimization}\label{multi-objective-optimization}

\begin{itemize}
\tightlist
\item
  A candidate \(\mathbf{q}\) Pareto-dominates \(\mathbf{q}'\) when
  \(f_k(\mathbf{q}) \le f_k(\mathbf{q}')\) for every objective \(k\),
  with strict inequality for at least one; the Pareto front is the set
  of undominated candidates. Simplex faces of a candidate vertex set
  encode these trade-offs, each face selecting one non-dominated
  combination.
\item
  Pareto front exploration via tetrahedral traversal: walking
  edge-adjacent simplices over the candidate set enumerates neighboring
  non-dominated combinations without a continuous relaxation.
\item
  Integer tetravolume of the solution simplex serves as a discrete
  diversity measure: larger volume corresponds to more mutually
  separated trade-off points.
\end{itemize}

\subsection{Machine Learning and
Robustness}\label{machine-learning-and-robustness}

\begin{itemize}
\tightlist
\item
  \textbf{Geometric regularization}: Quadray-constrained
  weights/topologies yield structural priors and improved stability.
\item
  \textbf{Adversarial robustness}: Discrete lattice projection reduces
  vulnerability to gradient-based adversarial perturbations by limiting
  directions.
\item
  \textbf{Ensembles}: Tetrahedral vertex voting and consensus improve
  robustness.
\end{itemize}

References: see
\href{https://en.wikipedia.org/wiki/Fisher_information}{Fisher
information},
\href{https://en.wikipedia.org/wiki/Natural_gradient}{Natural gradient},
and quadray conversion notes by Urner for embedding choices.

\subsection{Computer Graphics and GPU
Acceleration}\label{computer-graphics-and-gpu-acceleration}

\begin{itemize}
\tightlist
\item
  \textbf{Quadray visualization acceleration}: GPU-accelerated rendering
  of tetrahedral coordinate systems enables real-time exploration of 4D
  geometric structures. The parallel nature of GPU architectures
  naturally maps to the four-basis vector representation of quadrays,
  allowing simultaneous computation of vertex positions, edge
  connections, and face tessellations across thousands of tetrahedra.
\item
  \textbf{Integer arithmetic optimization}: GPU compute shaders excel at
  integer-based volume calculations and determinant computations using
  the Bareiss algorithm. The discrete lattice structure of quadray
  coordinates benefits from parallel integer arithmetic units, achieving
  significant speedups over CPU implementations for large-scale
  geometric computations.
\item
  \textbf{Dynamic programming acceleration}: GPU-accelerated dynamic
  programming algorithms leverage CUDA Dynamic Parallelism for adaptive
  parallel computation of recursive geometric algorithms. This approach
  enables efficient handling of varying computational workloads in
  tetrahedral decomposition and optimization problems, as demonstrated
  in applications like the Mandelbrot set computation where dynamic
  parallelism manages computational complexity effectively.
\item
  \textbf{Parallel geometric algorithms}: Implementation of
  GPU-optimized versions of algorithms like QuickHull for convex hull
  computation in quadray space achieves substantial performance
  improvements. The tetrahedral lattice structure naturally supports
  parallel prefix sum operations and efficient neighbor queries,
  enabling real-time visualization of complex 4D geometric
  transformations.
\item
  \textbf{Memory bandwidth optimization}: The structured memory access
  patterns of quadray coordinates align well with GPU memory
  hierarchies, enabling efficient coalesced memory access for
  large-scale geometric datasets. This optimization is particularly
  beneficial for applications requiring real-time rendering of complex
  polyhedral structures and dynamic tessellations.
\end{itemize}

References: GPU-accelerated geometry processing techniques
(\href{https://arxiv.org/abs/1501.04706?utm_source=openai}{arxiv.org}),
CUDA Dynamic Parallelism for adaptive computation
(\href{https://developer.nvidia.com/blog/introduction-cuda-dynamic-parallelism/?utm_source=openai}{developer.nvidia.com}),
and parallel scan algorithms for optimization
(\href{https://developer.nvidia.com/gpugems/gpugems3/part-vi-gpu-computing?utm_source=openai}{developer.nvidia.com}).

\subsection{Active Inference and Free
Energy}\label{active-inference-and-free-energy}

\begin{itemize}
\tightlist
\item
  Free energy
  \(\mathcal{F} = -\log P(o\mid s) + \mathrm{KL}[Q(s)\,\|\,P(s)]\) (see
  Eq. \eqref{eq:free_energy} in the equations appendix); background:
  \href{https://en.wikipedia.org/wiki/Free_energy_principle}{Free energy
  principle} and overviews connecting to predictive coding and control.
\item
  Belief updates follow steepest descent in Fisher geometry using the
  natural gradient (see Eq. \eqref{eq:natural_gradient} in the equations
  appendix); quadray constraints improve stability/interpretability.
\item
  Links to metabolic efficiency and biologically plausible computation.
\item
  For the full treatment, see
  \passthrough{\lstinline!09\_free\_energy\_active\_inference.md!}
  (Appendix: The Free Energy Principle and Active Inference).
\end{itemize}

\subsection{Complex Systems and Collective
Intelligence}\label{complex-systems-and-collective-intelligence}

\begin{itemize}
\tightlist
\item
  Tetrahedral interaction patterns support distributed consensus and
  emergent behavior.
\item
  Resource allocation and network flows benefit from geometric
  constraints.
\item
  \textbf{Cognitive security}: Applying cognitive security can safeguard
  distributed consensus mechanisms from manipulation, preserving the
  reliability of emergent behaviors in complex systems. Incorporating
  cognitive security measures can protect the integrity of belief
  updates and decision-making processes, ensuring that actions are based
  on accurate and unmanipulated information.
\end{itemize}

\subsection{Geospatial Intelligence and the World
Game}\label{geospatial-intelligence-and-the-world-game}

\begin{itemize}
\tightlist
\item
  \textbf{Spatial data integration}: Quadray tetrahedral frameworks
  provide natural tessellations for geospatial data analysis, where the
  Dymaxion projection's minimal distortion aligns with Fuller's World
  Game objectives of holistic global perspective. The tetrahedral
  lattice supports efficient spatial indexing and neighbor queries for
  distributed geospatial intelligence operations.
\item
  \textbf{Resource allocation optimization}: The World Game's goal of
  ``making the world work for 100\% of humanity'' translates to
  multi-objective optimization problems where tetrahedral simplex faces
  encode trade-offs between population centers, resource distribution,
  and ecological constraints. Integer volume quantization ensures
  discrete, interpretable solutions for global resource allocation.
\item
  \textbf{Cognitive security in distributed sensing}: Geospatial
  intelligence networks benefit from tetrahedral consensus mechanisms
  that resist manipulation of spatial data streams. The geometric
  constraints of Fuller.4D provide natural validation frameworks for
  detecting anomalous spatial patterns and maintaining data integrity
  across distributed sensor networks.
\item
  \textbf{Tetrahedral tessellations for global modeling}: The World
  Game's emphasis on interconnected global systems maps naturally to
  tetrahedral decompositions of the Dymaxion projection, where each
  tetrahedron represents a coherent region for local optimization while
  maintaining global connectivity through shared faces and edges.
\end{itemize}

\subsection{Quadrays, Synergetics (Fuller.4D), and William
Blake}\label{quadrays-synergetics-fuller.4d-and-william-blake}

\begin{itemize}
\tightlist
\item
  Quadrays (tetrahedral coordinates) instantiate Fuller's Synergetics
  emphasis on the tetrahedron as a structural primitive; in this
  manuscript's terminology this corresponds to Fuller.4D. Tetrahedral
  frames support part--whole reasoning and efficient decompositions used
  throughout.
\item
  William Blake's ``fourfold vision'' (single, twofold, threefold,
  fourfold) provides a historical metaphor for multiscale perception and
  inference. Read through Fisher geometry and natural gradient dynamics,
  it parallels multilayer predictive processing and counterfactual
  simulation. For background, see a concise overview of Blake's
  visionary psycho‑topographies in British Art Studies
  (\href{https://www.britishartstudies.ac.uk/index/article-index/visionary-sense-of-london/article-category/cover-collaboration}{visionary
  art analysis}) and the Active Inference Institute's MathArt Stream \#8
  (\href{https://zenodo.org/records/13711302}{Active Inference \&
  Blake}).
\item
  Juxtaposing Blake and Fuller foregrounds ``comprehensivity'': holistic
  design and sensemaking via geometric primitives. Context:
  (\href{https://zenodo.org/records/7519132}{Fuller \& Blake: Lives in
  Juxtaposition}) and pedagogical antecedents in experimental design
  education at Black Mountain College
  (\href{https://commons.princeton.edu/eng574-s23/wp-content/uploads/sites/348/2023/03/Diaz-The-Experimenters-Chance-and-Design-at-Black-Mountain-College.pdf}{Diaz,
  Chance and Design at Black Mountain College -- PDF}).
\item
  Implications for Quadray practice: four‑facet summaries of
  models/trajectories, tetrahedral consensus in ensembles, and
  stigmergic annotation patterns for cognitive security and distributed
  sensemaking.
\end{itemize}

\subsection{Pedagogy and
Implementations}\label{pedagogy-and-implementations}

Kirby Urner's comprehensive
\href{https://github.com/4dsolutions}{4dsolutions ecosystem} provides
extensive educational resources and cross-platform implementations for
Quadray computation and visualization. For comprehensive details on
educational frameworks, cross-language implementations, historical
context, and community development, see the
\href{07_resources.md}{Resources} section.

\subsection{Higher Dimensions and
Decompositions}\label{higher-dimensions-and-decompositions}

\begin{itemize}
\tightlist
\item
  Decompose an \(n\)-simplex (\(n > 3\)) into tetrahedra and sum their
  signed integer volumes; the sum extends the quantized tetravolume of
  Section 3 to higher-dimensional content.
\item
  Tessellations support parallel and distributed implementations (cells
  processed independently).
\end{itemize}

\subsection{Limitations and Future
Work}\label{limitations-and-future-work}

\begin{itemize}
\tightlist
\item
  Benchmark breadth: extend beyond convex/quadratic toys to real tasks
  (registration, robust regression, control) with ablations.
\item
  Distance sensitivity: compare embeddings and their effect on optimizer
  trajectories; document recommended defaults.
\item
  Hybrid schemes: study schedules that interleave continuous proposals
  with lattice projection.
\end{itemize}

\newpage

\section{Discussion}\label{discussion}

Quadray geometry (Fuller.4D) offers an interpretable, quantized view of
geometry, topology, information, and optimization. The claim this
discussion defends is that two mechanisms carry most of the weight.
First, integer volumes enforce discrete dynamics, acting as a structural
prior that regularizes optimization, prevents numerical fragility, and
enables exact integer-based accelerated methods (Sections 3--5 present
the evidence). Second, information geometry supplies the right
optimization language for the synergetic tradition: updates proceed not
through arbitrary parameter-space moves in continuous space, but along
geodesics defined by information content (see Eq. \eqref{eq:fim} and Eq.
\eqref{eq:natural_gradient} in the equations appendix; overview:
\href{https://en.wikipedia.org/wiki/Natural_gradient}{Natural
gradient}). The subsections below state what follows from combining the
two mechanisms, what remains unestablished, and where external
validation stands.

\subsection{Scope and Limitations}\label{scope-and-limitations}

\begin{itemize}
\tightlist
\item
  \textbf{Embeddings and distances}: distance calculations are
  embedding-dependent --- the quadray-to-Euclidean map must be selected
  and reported deliberately (Section 5 catalogs the comparison agenda);
  no single embedding is privileged.
\item
  \textbf{Hybrid strategies}: some problems need both worlds ---
  continuous steps interleaved with periodic lattice projection balance
  curvature-aware efficiency against discrete robustness.
\item
  \textbf{Benchmarking}: the structural-prior benefits above remain
  hypotheses until benchmarked per domain; Section 5 lists the breadth
  and ablation gaps.
\end{itemize}

In practical analysis and simulation, numerical precision matters.
Integer-volume reasoning is exact in theory, but empirical evaluation
(e.g., determinants, Fisher Information, geodesics) can benefit from
high-precision arithmetic when double precision is insufficient; the
High-Precision Arithmetic Note in the equations appendix covers quad
precision (\passthrough{\lstinline!libquadmath!},
\passthrough{\lstinline!\_\_float128!}) and symbolic (SymPy) evaluation
routes.

\subsection{Fisher Information and
Curvature}\label{fisher-information-and-curvature}

Claim: the Fisher Information Matrix (FIM) defines a Riemannian metric
on parameter space whose eigenspectrum is a readable map of the loss
surface --- large eigenvalues of \passthrough{\lstinline!F!} mark
sensitive (high-curvature) directions, small eigenvalues mark sloppy
ones. Evidence: the eigenspectrum and curvature panels of Section 4 (the
``Fisher Information eigenspectrum and parameter-space curvature''
figure) display these scales for the empirical estimate of Eq.
\eqref{eq:fim_empirical} in the equations appendix. Implication:
curvature-aware steps using Eq. \eqref{eq:natural_gradient} in the
equations appendix adaptively scale updates by the inverse metric,
improving conditioning relative to vanilla gradient descent. Background:
\href{https://en.wikipedia.org/wiki/Fisher_information}{Fisher
information}.

A curious connection unites geodesics in information geometry, the
physical principle of least action, and Buckminster Fuller's tensegrity
geodesic domes (Fuller.4D). On statistical manifolds, geodesics are
shortest paths under the Fisher metric, and natural-gradient flows trace
paths that motivate gradient-weighted path-length heuristics (the
\passthrough{\lstinline!information\_length!} proxy in
\passthrough{\lstinline!metrics.py!}) constrained by curvature (Eqs.
\eqref{eq:fim}, \eqref{eq:natural_gradient} in the equations appendix).
In tensegrity domes, geodesic lines on triangulated spherical shells
distribute stress nearly uniformly while the network balances continuous
tension with discontinuous compression, attaining maximal stiffness with
minimal material. Both systems exemplify constraint-balanced minimalism:
an extremal path emerges by trading off cost (action or information
length) against structure (metric curvature or tensegrity
compatibility). The shared economy---optimal routing through low-cost
directions---links geodesic shells in architecture to geodesic flows in
parameter spaces; see background on tensegrity/geodesic domes online.

\subsection{Quadray Coordinates and 4D Structure (Fuller.4D vs
Coxeter.4D vs
Einstein.4D)}\label{quadray-coordinates-and-4d-structure-fuller.4d-vs-coxeter.4d-vs-einstein.4d}

Quadray coordinates provide a tetrahedral basis with projective
normalization, aligning with close-packed sphere centers (IVM);
overview:
\href{https://en.wikipedia.org/wiki/Quadray_coordinates}{Quadray
coordinates} and synergetics background. The optimization-relevant
consequence: symmetries common in quadray parameterizations often yield
near block-diagonal structure in \passthrough{\lstinline!F!},
simplifying inversion and preconditioning (Section 4 develops this
observation; the appendix's namespaces summary fixes the definitions).
We stress the boundaries because category errors here are the
framework's main failure mode: (i) Fuller.4D for lattice and integer
volumes, (ii) Coxeter.4D for Euclidean embeddings, lengths, and simplex
families, (iii) Einstein.4D for metric analogies only --- not for
interpreting synergetic tetravolumes.

\subsection{Integrating FIM with Quadray
Models}\label{integrating-fim-with-quadray-models}

Applying the FIM within quadray-parameterized models ties statistical
curvature to tetrahedral structure. Practical takeaways:

\begin{itemize}
\tightlist
\item
  Use \passthrough{\lstinline!fisher\_information\_matrix!} to estimate
  \passthrough{\lstinline!F!} from per-sample gradients; inspect
  principal directions via \passthrough{\lstinline!fim\_eigenspectrum!}.
\item
  Exploit block patterns induced by quadray symmetries to stabilize
  metric inverses and reduce compute.
\item
  Combine integer-lattice projection with natural-gradient steps to
  balance discrete robustness and curvature-aware efficiency.
\item
  Purely discrete alternatives (e.g.,
  \passthrough{\lstinline!discrete\_ivm\_descent!}) provide monotone
  integer-valued descent when gradients are unreliable; hybrid schemes
  can interleave discrete steps with curvature-aware continuous
  proposals.
\end{itemize}

\subsection{Implications for Optimization and
Estimation}\label{implications-for-optimization-and-estimation}

\subsubsection{Clarifications on ``frequency/time''
dimensions}\label{clarifications-on-frequencytime-dimensions}

\begin{itemize}
\tightlist
\item
  Fuller's discussions often treat frequency/energy as an additional
  organizing dimension distinct from Euclidean coordinates. In our
  manuscript, we keep the shape/angle relations (Fuller.4D) separate
  from time/energy bookkeeping; when temporal evolution is needed, we
  use explicit trajectories and metric analogies (Einstein.4D) without
  conflating with Euclidean 4D objects (Coxeter.4D). This separation
  avoids category errors while preserving the intended interpretability.
\end{itemize}

\subsubsection{On distance-based tetravolume formulas
(clarification)}\label{on-distance-based-tetravolume-formulas-clarification}

\begin{itemize}
\tightlist
\item
  The bridging-vs-native dichotomy of Section 3 reduces to a practical
  rule: when inputs are edge lengths, compute with PdF or Cayley--Menger
  in Euclidean length space and convert via the S3 factor (Eqs.
  \eqref{eq:pdf}, \eqref{eq:cayley_menger} in the equations appendix);
  when inputs are integer quadrays, use the native determinants of Tom
  Ace or Gerald de Jong (Eqs. \eqref{eq:ace5x5}, \eqref{eq:gdj}), which
  return IVM tetravolumes directly with no XYZ intermediates. All routes
  agree numerically on shared cases, so the choice is driven by input
  type and desired exactness, not by the answer.
\end{itemize}

\subsubsection{Symbolic analysis (bridging vs native) (Results
linkage)}\label{symbolic-analysis-bridging-vs-native-results-linkage}

\begin{itemize}
\item
  Evidence: exact (SymPy) comparisons confirm that CM+S3 and Ace 5×5
  produce identical IVM tetravolumes on canonical small integer-quadray
  examples --- see the bridging vs native validation figure in Section 3
  and the manifest \passthrough{\lstinline!sympy\_symbolics.txt!}
  alongside \passthrough{\lstinline!bridging\_vs\_native.csv!} in
  \passthrough{\lstinline!quadmath/output/!}. Implication: the two
  computational cultures are interchangeable on the lattice, so
  exactness can be chosen per call site.
\item
  Curvature-aware optimizers: Kronecker-factored approximations (K-FAC)
  leverage structure in \passthrough{\lstinline!F!} to accelerate
  training and improve stability; see
  \href{https://arxiv.org/abs/1503.05671}{K-FAC (arXiv:1503.05671)}.
  Similar ideas apply when quadray structure induces separable blocks.
\item
  Model selection: eigenvalue spread of \passthrough{\lstinline!F!}
  provides a lens on parameter identifiability; near-zero modes suggest
  redundancies or over-parameterization.
\item
  Robust computation: lattice normalization in quadray space yields
  discrete plateaus that complement FIM-based scaling for numerically
  stable trajectories.
\end{itemize}

\subsection{Community Ecosystem and
Validation}\label{community-ecosystem-and-validation}

The computational ecosystem around Quadrays and synergetic geometry ---
chiefly the 4dsolutions organization (Section 7) --- is treated here as
an external validation surface: its cross-language implementations
independently verify algorithmic correctness against this manuscript's
codebase, and its educational materials document applications across
environments this paper does not attempt to survey. The claim worth
stating explicitly: agreement between independent implementations is
evidence for the mathematics, not merely for the code.

\newpage

\section{Resources}\label{resources}

This section provides comprehensive resources for learning about and
working with Quadrays, synergetics, and the computational methods
discussed in this manuscript.

\subsection{Core Concepts and
Background}\label{core-concepts-and-background}

\subsubsection{Information Geometry and
Optimization}\label{information-geometry-and-optimization}

\begin{itemize}
\tightlist
\item
  \textbf{Fisher information}:
  \href{https://en.wikipedia.org/wiki/Fisher_information}{Fisher
  information (reference)} --- see also Eq. \eqref{eq:fim} in the
  equations appendix
\item
  \textbf{Natural gradient}:
  \href{https://en.wikipedia.org/wiki/Natural_gradient}{Natural gradient
  (reference)} --- see also Eq. \eqref{eq:natural_gradient} in the
  equations appendix
\end{itemize}

\subsubsection{Active Inference and Free
Energy}\label{active-inference-and-free-energy-1}

\begin{itemize}
\tightlist
\item
  \textbf{Active Inference Institute}:
  \href{https://welcome.activeinference.institute/}{Welcome to Active
  Inference Institute}
\item
  \textbf{Comprehensive review}:
  \href{https://discovery.ucl.ac.uk/id/eprint/10176959/1/1-s2.0-S1571064523001094-main.pdf}{Active
  Inference --- recent review (UCL Discovery, 2023)}
\end{itemize}

\subsubsection{Mathematical
Foundations}\label{mathematical-foundations-1}

\begin{itemize}
\tightlist
\item
  \textbf{Tetrahedron volume formulas}: length-based
  \href{https://en.wikipedia.org/wiki/Cayley\%E2\%80\%93Menger_determinant}{Cayley--Menger
  determinant} and determinant-based expressions on vertex coordinates
  (see
  \href{https://en.wikipedia.org/wiki/Tetrahedron\#Volume}{Tetrahedron
  -- volume})
\item
  \textbf{Exact determinants}:
  \href{https://en.wikipedia.org/wiki/Bareiss_algorithm}{Bareiss
  algorithm}, used in our integer tetravolume implementations
\item
  \textbf{Optimization baseline}: the
  \href{https://en.wikipedia.org/wiki/Nelder\%E2\%80\%93Mead_method}{Nelder--Mead
  method}, adapted here to the Quadray lattice
\end{itemize}

\subsection{Quadrays and Synergetics (Core Starting
Points)}\label{quadrays-and-synergetics-core-starting-points}

\subsubsection{Introductory Materials}\label{introductory-materials}

\begin{itemize}
\tightlist
\item
  \textbf{Quadray coordinates (intro and conversions)}:
  \href{https://www.grunch.net/synergetics/quadintro.html}{Urner --
  Quadray intro},
  \href{https://www.grunch.net/synergetics/quadxyz.html}{Urner --
  Quadrays and XYZ}
\item
  \textbf{Quadrays and the Philosophy of Mathematics}:
  \href{https://www.grunch.net/synergetics/quadphil.html}{Urner --
  Quadrays and the Philosophy of Mathematics}
\item
  \textbf{Synergetics background and IVM}:
  \href{https://en.wikipedia.org/wiki/Synergetics_(Fuller)}{Synergetics
  (Fuller, overview)}
\item
  \textbf{Quadray coordinates overview}:
  \href{https://en.wikipedia.org/wiki/Quadray_coordinates}{Quadray
  coordinates (reference)}
\end{itemize}

\subsubsection{Historical and Background
Materials}\label{historical-and-background-materials}

\begin{itemize}
\tightlist
\item
  \textbf{RW Gray projects --- Synergetics text}:
  \href{http://www.rwgrayprojects.com/synergetics/s00/p0000.html}{rwgrayprojects.com
  (synergetics)}
\item
  \textbf{Fuller FAQ}:
  \href{https://www.cjfearnley.com/fuller-faq.pdf}{C. J. Fearnley's
  Fuller FAQ}
\item
  \textbf{Synergetics resource list}:
  \href{https://www.cjfearnley.com/fuller-faq-2.html}{C. J. Fearnley's
  resource page}
\item
  \textbf{Wikieducator}:
  \href{https://wikieducator.org/Synergetics}{Synergetics hub}
\item
  \textbf{Quadray animation}:
  \href{https://commons.wikimedia.org/wiki/File:Quadray.gif}{Quadray.gif
  (Wikimedia Commons)}
\item
  \textbf{Fuller Institute}:
  \href{https://www.bfi.org/about-fuller/big-ideas/synergetics/}{BFI ---
  Big Ideas: Synergetics}
\end{itemize}

\subsection{4dsolutions Ecosystem: Comprehensive Computational
Framework}\label{dsolutions-ecosystem-comprehensive-computational-framework}

The \href{https://github.com/4dsolutions}{4dsolutions organization}
provides the most extensive computational framework for Quadrays and
synergetic geometry, spanning 29+ repositories with implementations
across multiple programming languages.

\subsubsection{Core Computational
Modules}\label{core-computational-modules}

\paragraph{Primary Python Libraries}\label{primary-python-libraries}

\begin{itemize}
\tightlist
\item
  \textbf{Math for Wisdom (m4w)}:
  \href{https://github.com/4dsolutions/m4w}{m4w (repo)}

  \begin{itemize}
  \tightlist
  \item
    \textbf{Quadray vectors and conversions}:
    \href{https://github.com/4dsolutions/m4w/blob/main/qrays.py}{\passthrough{\lstinline!qrays.py!}
    (Qvector, SymPy-aware)}
  \item
    \textbf{Synergetic tetravolumes and modules}:
    \href{https://github.com/4dsolutions/m4w/blob/main/tetravolume.py}{\passthrough{\lstinline!tetravolume.py!}
    with PdF-CM vs native IVM and BEAST algorithms}
  \end{itemize}
\end{itemize}

\paragraph{Cross-Language Validation}\label{cross-language-validation}

\begin{itemize}
\tightlist
\item
  \textbf{Rust implementation}:
  \href{https://github.com/4dsolutions/rusty_rays}{rusty\_rays}
  (performance-oriented)

  \begin{itemize}
  \tightlist
  \item
    Sources:
    \href{https://github.com/4dsolutions/rusty_rays/blob/master/src/lib.rs}{Rust
    library implementation},
    \href{https://github.com/4dsolutions/rusty_rays/blob/master/src/main.rs}{Rust
    command-line interface}
  \end{itemize}
\item
  \textbf{Clojure implementation}:
  \href{https://github.com/4dsolutions/synmods}{synmods} (functional
  paradigm)

  \begin{itemize}
  \tightlist
  \item
    Sources:
    \href{https://github.com/4dsolutions/synmods/blob/master/qrays.clj}{\passthrough{\lstinline!qrays.clj!}},
    \href{https://github.com/4dsolutions/synmods/blob/master/ramping_up.clj}{\passthrough{\lstinline!ramping\_up.clj!}}
  \end{itemize}
\end{itemize}

\subsubsection{Primary Hub: School\_of\_Tomorrow (Python +
Notebooks)}\label{primary-hub-school_of_tomorrow-python-notebooks}

\textbf{Repository}:
\href{https://github.com/4dsolutions/School_of_Tomorrow}{School\_of\_Tomorrow}

\paragraph{Core Modules}\label{core-modules}

\begin{itemize}
\tightlist
\item
  \textbf{\passthrough{\lstinline!qrays.py!}}: Quadray implementation
  with normalization, conversions, and vector ops
  (\href{https://github.com/4dsolutions/School_of_Tomorrow/blob/master/qrays.py}{source})
\item
  \textbf{\passthrough{\lstinline!quadcraft.py!}}: POV-Ray scenes for
  CCP/IVM arrangements, animations, and tutorials
  (\href{https://github.com/4dsolutions/School_of_Tomorrow/blob/master/quadcraft.py}{source})
\item
  \textbf{\passthrough{\lstinline!flextegrity.py!}}: Polyhedron
  framework, concentric hierarchy, POV-Ray export
  (\href{https://github.com/4dsolutions/School_of_Tomorrow/blob/master/flextegrity.py}{source})
\item
  \textbf{Additional modules}: \passthrough{\lstinline!polyhedra.py!},
  \passthrough{\lstinline!identities.py!},
  \passthrough{\lstinline!smod\_play.py!} (synergetic modules)
\end{itemize}

\paragraph{Key Notebooks}\label{key-notebooks}

\begin{itemize}
\tightlist
\item
  \textbf{\passthrough{\lstinline!Qvolume.ipynb!}}: Tom Ace 5×5
  determinant with random-walk demonstrations
  (\href{https://github.com/4dsolutions/School_of_Tomorrow/blob/master/Qvolume.ipynb}{source})
\item
  \textbf{\passthrough{\lstinline!VolumeTalk.ipynb!}}: Comparative
  analysis of bridging vs native tetravolume formulations
  (\href{https://github.com/4dsolutions/School_of_Tomorrow/blob/master/VolumeTalk.ipynb}{source})
\item
  \textbf{\passthrough{\lstinline!QuadCraft\_Project.ipynb!}}: 1,255
  lines of interactive CCP navigation and visualization tutorials
  (\href{https://github.com/4dsolutions/School_of_Tomorrow/blob/master/QuadCraft_Project.ipynb}{source})
\item
  \textbf{Additional notebooks}:
  \passthrough{\lstinline!TetraBook.ipynb!},
  \passthrough{\lstinline!CascadianSynergetics.ipynb!},
  \passthrough{\lstinline!Rendering\_IVM.ipynb!},
  \passthrough{\lstinline!SphereVolumes.ipynb!} (visual and curricular
  materials)
\end{itemize}

\subsubsection{Additional Repositories}\label{additional-repositories}

\paragraph{Tetravolumes (Algorithms and
Pedagogy)}\label{tetravolumes-algorithms-and-pedagogy}

\begin{itemize}
\tightlist
\item
  \textbf{Repository}:
  \href{https://github.com/4dsolutions/tetravolumes}{tetravolumes}
\item
  \textbf{Code}:
  \href{https://github.com/4dsolutions/tetravolumes/blob/master/tetravolume.py}{\passthrough{\lstinline!tetravolume.py!}}
\item
  \textbf{Notebooks}:
  \href{https://raw.githubusercontent.com/4dsolutions/tetravolumes/refs/heads/master/Atoms\%20R\%20Us.ipynb}{Atoms
  R Us.ipynb},
  \href{https://raw.githubusercontent.com/4dsolutions/tetravolumes/refs/heads/master/Computing\%20Volumes.ipynb}{Computing
  Volumes.ipynb}
\end{itemize}

\paragraph{Visualization and
Rendering}\label{visualization-and-rendering}

\begin{itemize}
\tightlist
\item
  \textbf{BookCovers}: VPython for interactive educational animations
  (\href{https://github.com/4dsolutions/BookCovers}{repo})

  \begin{itemize}
  \tightlist
  \item
    Examples:
    \href{https://github.com/4dsolutions/BookCovers/blob/master/bookdemo.py}{\passthrough{\lstinline!bookdemo.py!}},
    \href{https://github.com/4dsolutions/BookCovers/blob/master/stickworks.py}{\passthrough{\lstinline!stickworks.py!}},
    \href{https://github.com/4dsolutions/BookCovers/blob/master/tetravolumes.py}{\passthrough{\lstinline!tetravolumes.py!}}
  \end{itemize}
\end{itemize}

\subsubsection{Educational Framework and
Curricula}\label{educational-framework-and-curricula}

\paragraph{Oregon Curriculum Network
(OCN)}\label{oregon-curriculum-network-ocn}

\begin{itemize}
\tightlist
\item
  \textbf{OCN portal}: \href{http://www.4dsolutions.net/ocn/}{OCN
  portal}
\item
  \textbf{Python for Everyone}:
  \href{http://www.4dsolutions.net/ocn/pymath.html}{pymath page}
\end{itemize}

\paragraph{Historical Documentation}\label{historical-documentation}

\begin{itemize}
\tightlist
\item
  \textbf{Python5 notebooks}:
  \href{https://raw.githubusercontent.com/4dsolutions/Python5/master/Polyhedrons\%20101.ipynb}{Polyhedrons
  101.ipynb}
\item
  \textbf{Historical variants}: \passthrough{\lstinline!qrays.py!} also
  appears in
  \href{https://github.com/4dsolutions/Python5/blob/master/qrays.py}{Python5
  (archive)}
\item
  \textbf{Python edu-sig archives}:
  \href{https://mail.python.org/pipermail/edu-sig/2000-May/000498.html}{Python
  edu-sig archives} tracing 25+ years of development
\end{itemize}

\subsubsection{Media and Publications}\label{media-and-publications}

\begin{itemize}
\tightlist
\item
  \textbf{YouTube demonstrations}:
  \href{https://www.youtube.com/watch?v=g14mu4uWD4E}{Synergetics talk
  1}, \href{https://www.youtube.com/watch?v=i9oij02oje0}{Synergetics
  talk 2},
  \href{https://www.youtube.com/watch?v=D0M1h_gjA_w}{Additional}
\item
  \textbf{Academia profile}:
  \href{https://princeton.academia.edu/kirbyurner}{Kirby Urner at
  Academia.edu}
\end{itemize}

\subsection{Community Discussions and Collaborative
Platforms}\label{community-discussions-and-collaborative-platforms}

\subsubsection{Active Platforms}\label{active-platforms}

\begin{itemize}
\tightlist
\item
  \textbf{Math4Wisdom Knowledge Engineering}:
  \href{https://coda.io/d/_d0SvdI3KSto/Knowledge-Engineering_suxu39sp}{Collaborative
  platform} with various art, resources, and cross-reference materials
\item
  \textbf{synergeo discussion archive}:
  \href{https://groups.io/g/synergeo/topics}{Groups.io platform} with
  ongoing community discussions and technical exchanges
\end{itemize}

\subsubsection{Historical Archives}\label{historical-archives}

\begin{itemize}
\tightlist
\item
  \textbf{GeodesicHelp threads}:
  \href{https://groups.google.com/g/GeodesicHelp/}{GeodesicHelp
  computations archive (Google Groups)} documenting computational
  approaches and problem-solving techniques
\end{itemize}

\subsection{Related Projects and
Applications}\label{related-projects-and-applications}

\subsubsection{Tetrahedral Voxel
Engines}\label{tetrahedral-voxel-engines}

\begin{itemize}
\tightlist
\item
  \textbf{QuadCraft}:
  \href{https://github.com/docxology/quadcraft/}{Tetrahedral voxel
  engine using Quadrays}
\end{itemize}

\subsubsection{Academic Publications}\label{academic-publications}

\begin{itemize}
\tightlist
\item
  \textbf{Flextegrity}:
  \href{https://www.academia.edu/44531954/Generating_the_Flextegrity_Lattice}{Generating
  the Flextegrity Lattice (academia.edu)}
\end{itemize}

\subsubsection{Context and Integration}\label{context-and-integration}

These materials popularize the IVM/CCP/FCC framing of space, integer
tetravolumes, and projective Quadray normalization. They inform the
methods in this paper and complement the \passthrough{\lstinline!src/!}
implementations (see \passthrough{\lstinline!quadray.py!},
\passthrough{\lstinline!cayley\_menger.py!},
\passthrough{\lstinline!linalg\_utils.py!}).

The ecosystem provides extensive validation, pedagogical context, and
practical implementations that complement and extend the methods
developed in this manuscript. Cross-language implementations serve as
independent verification of algorithmic correctness while educational
materials demonstrate practical applications across diverse
computational environments.

\newpage

\section{Equations and Math Supplement
(Appendix)}\label{equations-and-math-supplement-appendix}

Conventions for this appendix: vertices are \(P_0,\ldots,P_3\) with
Quadray components \((a_i, b_i, c_i, d_i)\); volumes are \(V_{xyz}\)
(Euclidean, cubic length units) and \(V_{ivm}\) (synergetics/IVM
tetravolume units, unit regular tetrahedron \(V_{ivm} = 1\)), related by
\(V_{ivm} = S3\,V_{xyz}\) with \(S3=\sqrt{9/8}\) (Section 3); \(M\)
denotes the Quadray-to-XYZ embedding matrix (Sections 3 and 14),
\(\lVert\cdot\rVert_2\) the Euclidean norm, and
\(\mathrm{KL}\big[Q\,\|\,P\big]\) the Kullback--Leibler divergence of
\(Q\) relative to \(P\) (Section 10). Three symbol overloads are scoped
by context and flagged where they occur: \(c\) is a Quadray component
(and a PdF edge length) in the volume formulas but the speed of light in
the Minkowski line element; \(d\) is the fourth Quadray component in the
coordinate formulas but an edge length in the length-based formulas and
the distance function \(d(\cdot,\cdot)\) of the embedding section; and
\(P\) is a vertex label \(P_i\) in the volume sections but a probability
distribution in the free-energy sections.

\subsection{Volume of a Tetrahedron
(Lattice)}\label{volume-of-a-tetrahedron-lattice}

\begin{equation}\label{eq:lattice_det}
V_{xyz} = \tfrac{1}{6}\,\left|\det\,[\,P_1 - P_0,\; P_2 - P_0,\; P_3 - P_0\,]\right|
\end{equation}

Notes.

\begin{itemize}
\tightlist
\item
  \(P_0,\ldots,P_3\) are Cartesian (XYZ) vertex coordinates, in length
  units; each bracketed column is an edge vector, the determinant is the
  volume of the parallelepiped they span, and the \(1/6\) factor
  converts it to the tetrahedron volume \(V_{xyz}\) in cubic length
  units. This is the coordinate (difference) form of the homogeneous-row
  determinant of Eq. \eqref{eq:xyz_det}; for IVM units directly from
  Quadray coordinates, see the native formula of Eq. \eqref{eq:gdj}.
\end{itemize}

Tom Ace 5×5 tetravolume (IVM units):

\begin{equation}\label{eq:ace5x5}
V_{ivm} = \tfrac{1}{4} \left| \det \begin{pmatrix}
 a_0 & b_0 & c_0 & d_0 & 1 \\
 a_1 & b_1 & c_1 & d_1 & 1 \\
 a_2 & b_2 & c_2 & d_2 & 1 \\
 a_3 & b_3 & c_3 & d_3 & 1 \\
  1 & 1 & 1 & 1 & 0
\end{pmatrix} \right|
\end{equation}

Notes.

\begin{itemize}
\tightlist
\item
  Row \(i\) lists the four Quadray components \((a_i, b_i, c_i, d_i)\)
  of vertex \(P_i\), augmented with an affine 1; the last row
  \((1,1,1,1,0)\) encodes the projective normalization constraint, so
  the determinant is invariant to adding \((t,t,t,t)\) to every vertex
  (Section 3). Division by 4 returns the IVM tetravolume \(V_{ivm}\);
  for integer quadrays the determinant is an exact integer (Bareiss
  algorithm), computed as a
  \passthrough{\lstinline!fractions.Fraction!}. This matrix is
  identical, row for row, to the named matrix \(\mathsf{Q}\) of Eq.
  \eqref{eq:ace5x5_expanded}.
\end{itemize}

\subsection{Expanded Ace 5×5 Matrix}\label{expanded-ace-55-matrix}

The Ace 5×5 matrix of Eq. \eqref{eq:ace5x5} written out explicitly and
named \(\mathsf{Q}\):

\begin{equation}\label{eq:ace5x5_expanded}
\mathsf{Q}(P_0,P_1,P_2,P_3) = \begin{bmatrix}
 a_0 & b_0 & c_0 & d_0 & 1 \\
 a_1 & b_1 & c_1 & d_1 & 1 \\
 a_2 & b_2 & c_2 & d_2 & 1 \\
 a_3 & b_3 & c_3 & d_3 & 1 \\
1 & 1 & 1 & 1 & 0
\end{bmatrix}, \qquad V_{ivm} = \tfrac{1}{4}\,\big|\det \mathsf{Q}(P_0,\ldots,P_3)\big|
\end{equation}

Notes.

\begin{itemize}
\tightlist
\item
  \textbf{Matrix structure}: row \(i\) holds the four Quadray components
  \((a_i, b_i, c_i, d_i)\) of vertex \(P_i\), plus the affine coordinate
  1 --- literally the same rows as Eq. \eqref{eq:ace5x5}.
\item
  \textbf{Last row}: \((1,1,1,1,0)\) enforces the projective
  normalization constraint (Section 3).
\item
  \textbf{Volume computation}:
  \(V_{ivm} = \tfrac{1}{4}\,\big|\det \mathsf{Q}\big|\) in IVM units,
  exact for integer quadrays.
\item
  \textbf{Notation}: the Ace matrix is written \(\mathsf{Q}\); the
  symbol \(M\) is reserved for the Quadray-to-XYZ embedding matrix
  (Sections 3 and 14).
\end{itemize}

XYZ determinant volume and S3 conversion:

\begin{equation}\label{eq:xyz_det}
V_{xyz} = \tfrac{1}{6} \left| \det \begin{pmatrix}
 x_0 & y_0 & z_0 & 1 \\
 x_1 & y_1 & z_1 & 1 \\
 x_2 & y_2 & z_2 & 1 \\
 x_3 & y_3 & z_3 & 1 \\
\end{pmatrix} \right|, \qquad V_{ivm} = S3\, V_{xyz},\quad S3=\sqrt{\tfrac{9}{8}}
\end{equation}

Notes.

\begin{itemize}
\tightlist
\item
  Homogeneous-row determinant in Cartesian coordinates: each row is
  \((x_i, y_i, z_i, 1)\) for vertex \(P_i\), with coordinates in length
  units; the absolute determinant divided by 6 is the Euclidean volume
  \(V_{xyz}\) in cubic length units. This is the homogeneous form of Eq.
  \eqref{eq:lattice_det}; conversion to IVM units uses
  \(V_{ivm} = S3\,V_{xyz}\) with \(S3=\sqrt{9/8}\) as used throughout.
\end{itemize}

\subsection{Cayley-Menger Determinant
(Coxeter.4D)}\label{cayley-menger-determinant-coxeter.4d}

For tetrahedron volume from edge lengths (Coxeter.4D approach):

\begin{equation}\label{eq:cayley_menger}
288\,V_{xyz}^2 = \det\begin{pmatrix}
  0 & 1 & 1 & 1 & 1 \\
  1 & 0 & d_{01}^2 & d_{02}^2 & d_{03}^2 \\
  1 & d_{10}^2 & 0 & d_{12}^2 & d_{13}^2 \\
  1 & d_{20}^2 & d_{21}^2 & 0 & d_{23}^2 \\
  1 & d_{30}^2 & d_{31}^2 & d_{32}^2 & 0
\end{pmatrix}
\end{equation}

Notes.

\begin{itemize}
\tightlist
\item
  \textbf{Pairwise distances}: \(d_{ij}\) is the Euclidean distance
  between vertices \(P_i\) and \(P_j\) in length units; the matrix
  stores the squared distances \(d_{ij}^2\), so the formula is
  length-only --- no coordinates enter (Coxeter.4D).
\item
  \textbf{Length-only formulation}: Cayley--Menger provides a
  length-only formula for simplex volumes, here specialized to
  tetrahedra.
\item
  \textbf{Conversion to IVM}: \(V_{xyz}\) is the Euclidean volume in
  cubic length units; use \(V_{ivm} = S3\,V_{xyz}\) with
  \(S3=\sqrt{9/8}\) (Eq. \eqref{eq:xyz_det}); the PdF formula of Eq.
  \eqref{eq:pdf} consumes the same lengths in closed algebraic form.
\end{itemize}

\subsection{Piero della Francesca Formula
(PdF)}\label{piero-della-francesca-formula-pdf}

For tetrahedron volume from edge lengths meeting at a vertex:

\begin{equation}\label{eq:pdf}
144\,V_{xyz}^2 = 4 a^2 b^2 c^2 - a^2\,(b^2 + c^2 - f^2)^2 - b^2\,(c^2 + a^2 - e^2)^2 - c^2\,(a^2 + b^2 - d^2)^2 + (b^2 + c^2 - f^2)(c^2 + a^2 - e^2)(a^2 + b^2 - d^2)
\end{equation}

Notes.

\begin{itemize}
\tightlist
\item
  \textbf{Edge lengths}: \(a, b, c\) are the three edges meeting at the
  apex vertex \(P_0\) --- \(a = d_{01}\), \(b = d_{02}\), \(c = d_{03}\)
  --- and \(d, e, f\) are the respective opposite edges, \(d = d_{23}\),
  \(e = d_{13}\), \(f = d_{12}\), in length units (the same \(d_{ij}\)
  as Eq. \eqref{eq:cayley_menger}).
\item
  \textbf{Conversion to IVM}: \(V_{xyz}\) is the Euclidean volume in
  cubic length units (the prefactor 144 is the standard Heron-like PdF
  normalization); use \(V_{ivm} = S3\,V_{xyz}\) with \(S3=\sqrt{9/8}\)
  (Eq. \eqref{eq:xyz_det}).
\end{itemize}

\subsection{Gerald de Jong Formula
(GdJ)}\label{gerald-de-jong-formula-gdj}

Native Quadray formula for tetrahedron volume:

\begin{equation}\label{eq:gdj}
V_{ivm} = \tfrac{1}{4}\,\left|\det\big[\, \pi(P_1) - \pi(P_0),\; \pi(P_2) - \pi(P_0),\; \pi(P_3) - \pi(P_0) \,\big]\right|, \qquad \pi(P_i) = (\,a_i - d_i,\; b_i - d_i,\; c_i - d_i\,)
\end{equation}

Notes.

\begin{itemize}
\tightlist
\item
  \textbf{Projection}: \(\pi\) maps a Quadray vertex \(P_i\) with
  components \((a_i, b_i, c_i, d_i)\) to \(\mathbb{R}^3\) by subtracting
  the fourth component from the first three; each bracketed column is a
  projected edge vector. Because \(\pi\) kills the direction
  \((1,1,1,1)\), it is invariant under \(q \to q + t\,(1,1,1,1)\) ---
  the projective fiber of Eq. \eqref{eq:conv-fiber} (Section 14) --- so
  the determinant is unchanged by projective normalization of the
  vertices. When all four vertices share the same fourth component (in
  particular \(d_i = 0\)), \(\pi(P_i)\) reduces to \((a_i, b_i, c_i)\).
\item
  \textbf{Native IVM}: no S3 conversion; the factor \(1/4\) returns IVM
  tetravolume directly --- for the unit tetrahedron (origin plus three
  IVM neighbor moves) the determinant is exactly 4, giving
  \(V_{ivm} = 1\) (Section 3).
\item
  \textbf{Exact arithmetic}: integer Quadray coordinates give an integer
  determinant, evaluated exactly by the Bareiss algorithm as a
  \passthrough{\lstinline!fractions.Fraction!}; the code implements Eq.
  \eqref{eq:gdj} as \passthrough{\lstinline!integer\_tetra\_volume!},
  and the magnitude agrees exactly with the Ace determinant of Eq.
  \eqref{eq:ace5x5} on every integer-quadray input.
\end{itemize}

See code:
\href{03_quadray_methods.md\#code:tetra_volume_cayley_menger}{\passthrough{\lstinline!tetra\_volume\_cayley\_menger!}}.
For tetrahedron volume background, see
\href{https://en.wikipedia.org/wiki/Tetrahedron\#Volume}{Tetrahedron --
volume}. Exact integer determinants in code use the
\href{https://en.wikipedia.org/wiki/Bareiss_algorithm}{Bareiss
algorithm}. External validation: these formulas align with
implementations in the 4dsolutions ecosystem. See the
\href{07_resources.md}{Resources} section for comprehensive details.

\subsection{Fisher Information Matrix (FIM)}\label{eq:fim}

Background:
\href{https://en.wikipedia.org/wiki/Fisher_information}{Fisher
information}.

\begin{equation}\label{eq:fim}
F_{i,j} = \mathbb{E}\left[ \frac{\partial \, \log p(x;\theta)}{\partial \theta_i}\, \frac{\partial \, \log p(x;\theta)}{\partial \theta_j} \right]
\end{equation}

Notes.

\begin{itemize}
\tightlist
\item
  \textbf{Symbols}: \(x\) is an observation,
  \(\theta = (\theta_1, \ldots, \theta_m)\) the parameter vector, and
  \(p(x;\theta)\) the likelihood; the expectation is over
  \(x \sim p(\cdot;\theta)\), so \(F\) is a function of \(\theta\). Each
  score \(\partial \log p(x;\theta)/\partial \theta_i\) carries inverse
  parameter units, so \(F_{i,j}\) carries inverse squared parameter
  units per observation.
\item
  \(F\) is symmetric positive semi-definite: large eigenvalues mark
  sensitive (high-curvature) directions, small eigenvalues sloppy ones
  (Section 4 and the Discussion). The empirical estimate of Eq.
  \eqref{eq:fim_empirical} is the per-observation counterpart.
\item
  See code:
  \href{03_quadray_methods.md\#code:fisher_information_matrix}{\passthrough{\lstinline!fisher\_information\_matrix!}}
  in \passthrough{\lstinline!src/quadmath/inference/information.py!} ---
  empirical outer-product estimator. Figure: the empirical estimate is
  shown in the FIM heatmap figure of Section 4.
\end{itemize}

\subsection{Empirical Fisher Information
Matrix}\label{empirical-fisher-information-matrix}

For empirical estimation from data, the Fisher Information Matrix is
computed as:

\begin{equation}\label{eq:fim_empirical}
F_{i,j} = \frac{1}{N} \sum_{n=1}^{N} \frac{\partial \, \log p(x_n;\theta)}{\partial \theta_i}\, \frac{\partial \, \log p(x_n;\theta)}{\partial \theta_j}
\end{equation}

Notes.

\begin{itemize}
\tightlist
\item
  \textbf{Symbols}: \(x_1, \ldots, x_N\) are \(N\) i.i.d. observations
  and \(g_n = \nabla_\theta \log p(x_n;\theta)\) the per-sample score
  vector (same units as Eq. \eqref{eq:fim}); the estimator is
  \(F = \tfrac{1}{N} \sum_n g_n g_n^{\mathsf{T}}\), with a
  \passthrough{\lstinline!normalize!} flag controlling the \(1/N\)
  division.
\item
  Converges to Eq. \eqref{eq:fim} as \(N \to \infty\); used by
  natural-gradient descent (Eq. \eqref{eq:natural_gradient}) and the
  information-geometry applications of Section 4.
\end{itemize}

\subsection{Natural Gradient}\label{eq:natgrad}

Background:
\href{https://en.wikipedia.org/wiki/Natural_gradient}{Natural gradient}
(Amari).

\begin{equation}\label{eq:natural_gradient}
\theta \leftarrow \theta - \eta\, F(\theta)^{-1}\, \nabla_{\theta} L(\theta)
\end{equation}

Explanation.

\begin{itemize}
\tightlist
\item
  \textbf{Symbols}: \(\theta\) the parameter vector, \(L(\theta)\) a
  differentiable loss, \(\nabla_\theta L\) its gradient, \(F(\theta)\)
  the Fisher matrix of Eq. \eqref{eq:fim} --- or its empirical estimate,
  Eq. \eqref{eq:fim_empirical} --- evaluated at \(\theta\), and
  \(\eta > 0\) the step size (learning rate).
\item
  \textbf{Update}: right-preconditioning by the inverse Fisher metric
  gives the steepest-descent direction under the Fisher metric (Amari),
  invariant to smooth invertible reparameterizations of \(\theta\).
\item
  \textbf{Damping}: the implementation
  \passthrough{\lstinline!natural\_gradient\_step!} solves the damped
  system \((F + \lambda I)\,\delta = \nabla_\theta L\) and applies
  \(\theta \leftarrow \theta - \eta\,\delta\), with a small Tikhonov
  ridge \(\lambda\) (default \(10^{-9}\)) keeping the solve stable when
  \(F\) is near-singular.
\end{itemize}

See code:
\href{03_quadray_methods.md\#code:natural_gradient_step}{\passthrough{\lstinline!natural\_gradient\_step!}}
in \passthrough{\lstinline!src/quadmath/inference/information.py!} ---
damped inverse-Fisher step.

\subsection{Free Energy (Active Inference)}\label{eq:free_energy}

\begin{equation}\label{eq:free_energy}
\mathcal{F} = -\log P(o\mid s) + \mathrm{KL}\big[ Q(s)\;\|\; P(s) \big]
\end{equation}

Explanation.

\begin{itemize}
\tightlist
\item
  \textbf{Symbols}: \(o\) the observation, \(s\) the latent state,
  \(Q(s)\) the approximate posterior (recognition distribution), and
  \(P\) the generative model, so \(P(o \mid s)\) is the likelihood and
  \(P(s)\) the prior; \(\mathcal{F}\) is the variational free energy in
  nats (natural logarithm, Section 10). \textbf{Notation flag}: capital
  \(P\) here is a probability distribution, unrelated to the vertex
  labels \(P_i\) of the volume sections.
\item
  \textbf{Partition}: minimizing \(\mathcal{F}\) over \(Q\) trades the
  expected negative log-likelihood \(\mathbb{E}_Q[-\log P(o \mid s)]\)
  (the displayed \(-\log P(o\mid s)\) is read under this
  \(Q\)-expectation, as implemented) against the KL divergence
  \(\mathrm{KL}\big[Q\,\|\,P\big]\) that ties the posterior to the
  prior; see
  \href{https://en.wikipedia.org/wiki/Free_energy_principle}{Free energy
  principle}.
\end{itemize}

See code:
\href{03_quadray_methods.md\#code:free_energy}{\passthrough{\lstinline!free\_energy!}}
in \passthrough{\lstinline!src/quadmath/inference/information.py!} ---
discrete-state variational free energy (inputs are unnormalized
distributions, normalized internally).

\textbf{Note}: The main figures demonstrating natural gradient
trajectories and free energy landscapes are shown in
\href{04_optimization_in_4d.md}{Section 4: Optimization in 4D}. The
appendix focuses on unique figures specific to mathematical formulations
and validation.

\subsection{Expected Free Energy (Active
Inference)}\label{eq:expected_free_energy}

Background:
\href{https://direct.mit.edu/books/oa-monograph/5299/Active-InferenceThe-Free-Energy-Principle-in-Mind}{Active
Inference (Parr, Pezzulo \& Friston, MIT Press, 2022)}.

\begin{equation}\label{eq:expected_free_energy}
G = \mathrm{KL}\big[ Q(s)\;\|\;P(s) \big] \;-\; \mathbb{E}_{q}\big[\log P(o\mid s)\big] \;-\; \log P(o)
\end{equation}

Explanation.

\begin{itemize}
\tightlist
\item
  \textbf{Symbols}: \(Q(s)\) the approximate posterior and \(P(s)\) the
  prior over states (as in Eq. \eqref{eq:free_energy}), \(P(o \mid s)\)
  the likelihood, \(P(o)\) the prior preference over outcomes in nats
  (uniform when omitted), \(H\big[Q\big] = -\sum_s Q(s) \log Q(s)\) the
  Shannon entropy in nats, and \(\mathbb{E}_q\) expectation under \(Q\);
  \(G\) is the expected free energy minimized during action selection.
\item
  \textbf{Epistemic term}: the KL divergence between variational
  posterior and prior over states.
\item
  \textbf{Entropy}: the posterior entropy is already inside the KL term,
  since
  \(\mathrm{KL}[Q\|P] = \mathbb{E}_q[\log Q(s)] - \mathbb{E}_q[\log P(s)] = -H[Q(s)] - \mathbb{E}_q[\log P(s)]\);
  subtracting \(H\) again would count it twice.
\item
  \textbf{Ambiguity}: the negative expected log-likelihood of outcomes
  penalizes noisy observations.
\item
  \textbf{Pragmatic term}: prior preferences \(P(o)\) enter negatively
  (\(-\log P(o)\)), so preferred outcomes lower \(G\); agents minimize
  \(G\) during action selection.
\end{itemize}

See code:
\href{03_quadray_methods.md\#code:expected_free_energy}{\passthrough{\lstinline!expected\_free\_energy!}}
in \passthrough{\lstinline!src/quadmath/inference/information.py!} ---
all four terms of Eq. \eqref{eq:expected_free_energy} as implemented.

\subsection{Quadray Normalization
(Fuller.4D)}\label{quadray-normalization-fuller.4d}

Given a Quadray \(q = (a, b, c, d)\), choose \(k = \min(a, b, c, d)\)
and set \(q' = q - (k, k, k, k)\), enforcing non-negative entries with
at least one zero. The shift is the projective normalization of Section
3 --- \(q\) and \(q'\) represent the same direction, the fiber
formalized in Eq. \eqref{eq:conv-fiber} (Section 14). Here \(k\) is the
normalization offset of the glossary (Section 10); it is unrelated to
the shell-frequency \(k\) of the lattice sections (Sections 11 and 13).

\subsection{Distance (Embedding Sketch; Coxeter.4D
slice)}\label{distance-embedding-sketch-coxeter.4d-slice}

Choose a linear map \(M\) from Quadray space to \(\mathbb{R}^3\) (or
\(\mathbb{R}^4\)) consistent with the tetrahedral axes --- the embedding
matrix of Sections 3 and 14 (canonical Urner family, Eq.
\eqref{eq:conv-embedding}). Then for Quadray points \(p, q\), the
distance is \(d(p, q) = \lVert M\,(p - q) \rVert_2\) in length units;
under the integer Urner family the squared distance is the exact integer
identity of Eq. \eqref{eq:conv-distance}, which the lattice-search layer
ranks with (Section 13).

\subsection{Minkowski Line Element (Einstein.4D
analogy)}\label{minkowski-line-element-einstein.4d-analogy}

\begin{equation}\label{eq:minkowski_line_element}
ds^2 = -c^2\,dt^2 + dx^2 + dy^2 + dz^2
\end{equation}

Background:
\href{https://en.wikipedia.org/wiki/Minkowski_space}{Minkowski space}.

Signature convention \((-,+,+,+)\) (mostly-plus, Section 2): \(c\) is
the speed of light, \((t, x, y, z)\) are spacetime coordinates in time
and length units respectively, and \(ds^2\) has units of length squared.
Here \(c\) is the physical constant, not the Quadray component or PdF
edge length used in the volume sections.

\subsection{High-Precision Arithmetic
Note}\label{high-precision-arithmetic-note}

When evaluating determinants, FIMs, or geodesic distances for sensitive
problems, use quad precision (binary128) via GCC's
\passthrough{\lstinline!libquadmath!}
(\passthrough{\lstinline!\_\_float128!}, functions like
\passthrough{\lstinline!expq!}, \passthrough{\lstinline!sqrtq!}, and
\passthrough{\lstinline!quadmath\_snprintf!}). See
\href{https://gcc.gnu.org/onlinedocs/libquadmath/index.html}{GCC
libquadmath}. Where possible, it is useful to use symbolic math
libraries like SymPy to compute exact values.

\subsubsection{Reproducibility artifacts and external
validation}\label{reproducibility-artifacts-and-external-validation}

\begin{itemize}
\tightlist
\item
  \textbf{This manuscript's artifacts}: Raw data in
  \passthrough{\lstinline!quadmath/output/!} for reproducibility and
  downstream analysis:

  \begin{itemize}
  \tightlist
  \item
    \passthrough{\lstinline!fisher\_information\_matrix.csv!} /
    \passthrough{\lstinline!.npz!}: empirical Fisher matrix and inputs
  \item
    \passthrough{\lstinline!fisher\_information\_eigenvalues.csv!} /
    \passthrough{\lstinline!fisher\_information\_eigensystem.npz!}:
    eigenspectrum and eigenvectors
  \item
    \passthrough{\lstinline!natural\_gradient\_path.png!} with
    \passthrough{\lstinline!natural\_gradient\_path.csv!} /
    \passthrough{\lstinline!.npz!}: projected trajectory and raw
    coordinates
  \item
    \passthrough{\lstinline!ivm\_neighbors\_data.csv!} /
    \passthrough{\lstinline!ivm\_neighbors\_edges\_data.npz!}: neighbor
    coordinates (Quadray and XYZ)
  \item
    \passthrough{\lstinline!polyhedra\_quadray\_constructions.png!}:
    synergetics volume relationships schematic
  \end{itemize}
\item
  \textbf{External validation resources}: The
  \href{https://github.com/4dsolutions}{4dsolutions ecosystem} provides
  extensive cross-validation. See the \href{07_resources.md}{Resources}
  section for comprehensive details on computational implementations and
  validation.
\end{itemize}

\subsection{Namespaces summary
(notation)}\label{namespaces-summary-notation}

\begin{itemize}
\tightlist
\item
  Coxeter.4D: Euclidean E⁴; regular polytopes; not spacetime
  (cf.~Coxeter, Regular Polytopes, Dover ed., p.~119). Connections to
  higher-dimensional lattices and packings as in Conway \& Sloane.
\item
  Einstein.4D: Minkowski spacetime; indefinite metric; used here only as
  a metric analogy when discussing geodesics and information geometry.
\item
  Fuller.4D: Quadrays/IVM; tetrahedral lattice with integer tetravolume;
  unit regular tetrahedron has volume 1; synergetics scale relations
  (e.g., S3).
\end{itemize}

\newpage

\section{Appendix: The Free Energy Principle and Active
Inference}\label{appendix-the-free-energy-principle-and-active-inference}

\subsection{Overview}\label{overview-3}

The Free Energy Principle (FEP) posits that biological systems maintain
their states by minimizing variational free energy, thereby reducing
surprise via prediction and model updating. Active Inference extends
this by casting action selection as inference under prior preferences.
Background: see the concise overview on the
\href{https://en.wikipedia.org/wiki/Free_energy_principle}{Free energy
principle} and the monograph
\href{https://direct.mit.edu/books/oa-monograph/5299/Active-InferenceThe-Free-Energy-Principle-in-Mind}{Active
Inference (MIT Press)}.

This appendix emphasizes relationships among: (i) the four-fold
partition of Active Inference, (ii) Quadrays (Fuller.4D) as a geometric
scaffold for mapping this partition, and (iii) information-geometric
flows (Einstein.4D analogy) that underpin perception--action updates.
For the naming of 4D namespaces used throughout---Coxeter.4D (Euclidean
E4), Einstein.4D (Minkowski spacetime analogy), Fuller.4D
(Synergetics/Quadrays)---see
\passthrough{\lstinline!02\_4d\_namespaces.md!}.

\subsection{Mathematical Formulation and Equation Callouts (Equations
linkage)}\label{mathematical-formulation-and-equation-callouts-equations-linkage}

\begin{itemize}
\item
  Variational free energy (discrete states) --- see Eq.
  \eqref{eq:free_energy} in the equations appendix, implemented by
  \href{03_quadray_methods.md\#code:free_energy}{\passthrough{\lstinline!free\_energy!}}:
  \(\mathcal{F} = -\log P(o \mid s) + \mathrm{KL}[\,Q(s)\,\|\,P(s)\,]\),
  where \(o\) are observations, \(s\) latent states, \(Q(s)\) the
  variational posterior, and \(P(s)\) the prior; minimizing
  \(\mathcal{F}\) drives perception.
\item
  Expected free energy (action selection) --- see Eq.
  \eqref{eq:expected_free_energy} in the equations appendix, implemented
  by
  \href{03_quadray_methods.md\#code:expected_free_energy}{\passthrough{\lstinline!expected\_free\_energy!}}:
  \(G = \mathrm{KL}[\,Q(s)\,\|\,P(s)\,] - \mathbb{E}_{q}[\log P(o \mid s)] - \log P(o)\),
  where the epistemic term is the KL divergence between posterior
  \(Q(s)\) and prior \(P(s)\) over states (its entropy part,
  \(\mathbb{E}_{q}[\log Q(s)] = -H[Q(s)]\), is already included, so
  \(H\) is not subtracted again), the ambiguity term
  \(-\mathbb{E}_{q}[\log P(o \mid s)]\) penalizes expected surprise
  about outcomes \(o\), and the pragmatic preference \(-\log P(o)\)
  lowers \(G\) for preferred outcomes; agents minimize \(G\) when
  selecting actions.
\item
  Fisher Information Matrix (FIM) as metric --- see Eq. \eqref{eq:fim}
  in the equations appendix and
  \href{03_quadray_methods.md\#code:fisher_information_matrix}{\passthrough{\lstinline!fisher\_information\_matrix!}}:
  \(F_{ij} = \mathbb{E}\big[\partial_{\theta_i} \log p(x;\theta)\, \partial_{\theta_j} \log p(x;\theta)\big]\),
  the expected outer product of score functions
  \(\partial_{\theta_i} \log p(x;\theta)\), serving as the Riemannian
  metric on parameter space \(\theta\).
\item
  Natural gradient descent under information geometry --- see Eq.
  \eqref{eq:natural_gradient} in the equations appendix and
  \href{03_quadray_methods.md\#code:natural_gradient_step}{\passthrough{\lstinline!natural\_gradient\_step!}}:
  the update
  \(\theta \leftarrow \theta - \eta\, F^{-1} \nabla_{\theta} L\)
  preconditions the loss gradient \(\nabla_{\theta} L\) by the inverse
  Fisher metric \(F^{-1}\) at step size \(\eta\); the implementation
  solves the damped system
  \((F + \varepsilon I)\,\delta = \nabla_{\theta} L\) and moves by
  \(-\eta\,\delta\). Overview:
  \href{https://en.wikipedia.org/wiki/Natural_gradient}{Natural
  gradient}.
\end{itemize}

Figures: three Active Inference figures follow --- the partition
tetrahedron, the 4D natural-gradient trajectory
(\passthrough{\lstinline!figure\_13\_4d\_trajectory.png!}), and the
free-energy landscape
(\passthrough{\lstinline!figure\_14\_free\_energy\_landscape.png!}).

Discrete variational optimization on the quadray lattice:
\passthrough{\lstinline!discrete\_ivm\_descent!} greedily descends a
free-energy-like objective over IVM moves, yielding integer-valued
trajectories. See the path animation artifact
\passthrough{\lstinline!discrete\_path.mp4!} in
\passthrough{\lstinline!quadmath/output/!}.

\begin{figure}
\centering
\pandocbounded{\includegraphics[keepaspectratio,alt={Active Inference four-fold partition mapped to a Quadray tetrahedron in Fuller.4D. The four vertices of a regular tetrahedron, embedded in XYZ via the quadray embedding to\_xyz (axes range -2 to 2), carry the four partition components: \textbackslash mu internal states (blue), s sensory observations (orange), a actions (green), and \textbackslash psi external causes (red). The gray edges are the six pairwise couplings of the partition --- e.g., \textbackslash mu--s perceptual inference and a--\textbackslash psi control (see the partition section below). Generated by quadmath/scripts/active\_inference\_figures.py (MPLBACKEND=Agg, fixed seed 42).}]{../output/figures/partition_tetrahedron.png}}
\caption{\textbf{Active Inference four-fold partition mapped to a
Quadray tetrahedron in Fuller.4D}. The four vertices of a regular
tetrahedron, embedded in XYZ via the quadray embedding
\protect\passthrough{\lstinline!to\_xyz!} (axes range \(-2\) to \(2\)),
carry the four partition components: \(\mu\) internal states (blue),
\(s\) sensory observations (orange), \(a\) actions (green), and \(\psi\)
external causes (red). The gray edges are the six pairwise couplings of
the partition --- e.g., \(\mu\)--\(s\) perceptual inference and
\(a\)--\(\psi\) control (see the partition section below). Generated by
\protect\passthrough{\lstinline!quadmath/scripts/active\_inference\_figures.py!}
(\protect\passthrough{\lstinline!MPLBACKEND=Agg!}, fixed seed 42).}
\end{figure}

\begin{figure}
\centering
\pandocbounded{\includegraphics[keepaspectratio,alt={4D natural-gradient trajectory with Active Inference context (figure\_13\_4d\_trajectory.png). Top left: the trajectory of natural gradient descent in (\textbackslash mu, a, s) coordinates --- perception weight \textbackslash mu, action weight a, internal-state coordinate s --- from the initial state (green circle) through the converged state (red star) toward the true optimum (blue triangle), converging in 11 steps with final parameter errors below 0.015. Top right: free energy (squared loss, log scale) falling from about 10\^{}\{-1\} to 10\^{}\{-4\} over the 11 optimization steps. Bottom left: evolution of the four partition components \textbackslash mu, a, s, \textbackslash psi against their true values (dashed lines). Bottom middle: step size (linear scale) and gradient norm (log scale) per step. Bottom right: the 4 \textbackslash times 4 Fisher information matrix F\_\{ij\} (entries of order 10\^{}\{-7\}) that preconditions the updates via Eq. . Generated by quadmath/scripts/active\_inference\_figures.py (MPLBACKEND=Agg, fixed seed 42); raw data in figure\_13\_data.npz.}]{../output/figures/figure_13_4d_trajectory.png}}
\caption{\textbf{4D natural-gradient trajectory with Active Inference
context}
(\protect\passthrough{\lstinline!figure\_13\_4d\_trajectory.png!}). Top
left: the trajectory of natural gradient descent in \((\mu, a, s)\)
coordinates --- perception weight \(\mu\), action weight \(a\),
internal-state coordinate \(s\) --- from the initial state (green
circle) through the converged state (red star) toward the true optimum
(blue triangle), converging in 11 steps with final parameter errors
below 0.015. Top right: free energy (squared loss, log scale) falling
from about \(10^{-1}\) to \(10^{-4}\) over the 11 optimization steps.
Bottom left: evolution of the four partition components
\(\mu, a, s, \psi\) against their true values (dashed lines). Bottom
middle: step size (linear scale) and gradient norm (log scale) per step.
Bottom right: the \(4 \times 4\) Fisher information matrix \(F_{ij}\)
(entries of order \(10^{-7}\)) that preconditions the updates via Eq.
\eqref{eq:natural_gradient}. Generated by
\protect\passthrough{\lstinline!quadmath/scripts/active\_inference\_figures.py!}
(\protect\passthrough{\lstinline!MPLBACKEND=Agg!}, fixed seed 42); raw
data in \protect\passthrough{\lstinline!figure\_13\_data.npz!}.}
\end{figure}

\begin{figure}
\centering
\pandocbounded{\includegraphics[keepaspectratio,alt={Free-energy landscape over perception and action parameters (figure\_14\_free\_energy\_landscape.png). Top left: the variational free energy \textbackslash mathcal\{F\} of Eq.  as a surface over the perception parameter q\_1 and action parameter q\_2 (values in nats, about 1.2 to 2.4; color bar at right), with the global minimum starred. Top right: level contours of the same surface. Bottom left: one-dimensional cross-sections \textbackslash mathcal\{F\}(q\_1) at three fixed values of q\_2, showing sensitivity to each parameter. Bottom right: local curvature over (q\_1, q\_2) (Fisher information structure). The text panel summarizes the four-fold partition (\textbackslash mu, s, a, \textbackslash psi), the roles of Fuller.4D/Coxeter.4D/Einstein.4D, and the generative-model objects Q(s), P(s), P(o \textbackslash mid s). Generated by quadmath/scripts/active\_inference\_figures.py (MPLBACKEND=Agg, fixed seed 42); raw data in figure\_14\_data.npz.}]{../output/figures/figure_14_free_energy_landscape.png}}
\caption{\textbf{Free-energy landscape over perception and action
parameters}
(\protect\passthrough{\lstinline!figure\_14\_free\_energy\_landscape.png!}).
Top left: the variational free energy \(\mathcal{F}\) of Eq.
\eqref{eq:free_energy} as a surface over the perception parameter
\(q_1\) and action parameter \(q_2\) (values in nats, about 1.2 to 2.4;
color bar at right), with the global minimum starred. Top right: level
contours of the same surface. Bottom left: one-dimensional
cross-sections \(\mathcal{F}(q_1)\) at three fixed values of \(q_2\),
showing sensitivity to each parameter. Bottom right: local curvature
over \((q_1, q_2)\) (Fisher information structure). The text panel
summarizes the four-fold partition (\(\mu, s, a, \psi\)), the roles of
Fuller.4D/Coxeter.4D/Einstein.4D, and the generative-model objects
\(Q(s)\), \(P(s)\), \(P(o \mid s)\). Generated by
\protect\passthrough{\lstinline!quadmath/scripts/active\_inference\_figures.py!}
(\protect\passthrough{\lstinline!MPLBACKEND=Agg!}, fixed seed 42); raw
data in \protect\passthrough{\lstinline!figure\_14\_data.npz!}.}
\end{figure}

\subsection{Four-Fold Partition and Tetrahedral Mapping (Quadrays;
Fuller.4D)}\label{four-fold-partition-and-tetrahedral-mapping-quadrays-fuller.4d}

Active Inference partitions the agent--environment system into four
coupled states:

\begin{itemize}
\tightlist
\item
  Internal (\(\mu\)) --- agent's internal states
\item
  Sensory (\(s\)) --- observations
\item
  Active (\(a\)) --- actions
\item
  External (\(\psi\)) --- latent environmental causes
\end{itemize}

See, for an overview of this partition and generative process
formulations, the
\href{https://discovery.ucl.ac.uk/id/eprint/10176959/1/1-s2.0-S1571064523001094-main.pdf}{Active
Inference review} and the general entry on
\href{https://en.wikipedia.org/wiki/Active_inference}{Active inference}.

Tetrahedral mapping via Quadrays (Fuller.4D): assign each state to a
vertex of a tetrahedron, using Quadray coordinates
\passthrough{\lstinline!(A,B,C,D)!} with non-negative components and at
least one zero after normalization. One canonical mapping is
\(A \leftrightarrow \mu\) (internal), \(B \leftrightarrow s\) (sensory),
\(C \leftrightarrow a\) (active), \(D \leftrightarrow \psi\) (external).
The edges capture the pairwise couplings (e.g., \(\mu\text{--}s\) for
perceptual inference; \(a\text{--}\psi\) for control). Integer
tetravolume then quantifies the ``coupled capacity'' region spanned by
jointly feasible states in a time slice; see
\passthrough{\lstinline!Quadray!} and tetravolume methods in
\passthrough{\lstinline!03\_quadray\_methods.md!}.

Interpretation note: this Quadray-based mapping is a didactic geometric
scaffold. It is not standard in the Active Inference literature, which
typically develops the four-state partition in probabilistic graphical
terms. Our use highlights structural symmetries and discrete volumetric
quantities available in Fuller.4D, building on the computational
foundations developed in the
\href{https://github.com/4dsolutions}{4dsolutions ecosystem} for
tetrahedral modeling and volume calculations. See the
\href{07_resources.md}{Resources} section for comprehensive details on
the computational implementations.

Code linkage (no snippet): see
\passthrough{\lstinline!example\_partition\_tetra\_volume!} in
\passthrough{\lstinline!src/quadmath/core/examples.py!} and the
partition tetrahedron figure above.

\subsection{How the 4D namespaces relate
here}\label{how-the-4d-namespaces-relate-here}

\begin{itemize}
\tightlist
\item
  Fuller.4D (Quadrays): geometric embedding of the four-state partition
  on a tetrahedron; integer tetravolumes and IVM moves provide discrete
  combinatorial structure.
\item
  Coxeter.4D (Euclidean E4): exact Euclidean measurements (e.g.,
  Cayley--Menger determinants) for tetrahedra underlying volumetric
  comparisons and scale relations.
\item
  Einstein.4D (Minkowski analogy): information-geometric flows (natural
  gradient, metric-aware updates) supply a continuum picture for
  perception--action dynamics.
\end{itemize}

The three roles are complementary: Fuller.4D encodes partition
structure, Coxeter.4D provides exact metric geometry for static
comparisons, and Einstein.4D guides dynamical descent.

\subsection{Joint Optimization in the Tetrahedral Framework (Methods
linkage)}\label{joint-optimization-in-the-tetrahedral-framework-methods-linkage}

\begin{itemize}
\tightlist
\item
  Perception: update \(\mu\) to minimize prediction error on \(s\) under
  the generative model; \passthrough{\lstinline!perception\_update!}
  integrates the flow \(\dot{\mu} = D\,\mu - \partial_{\mu} F\), where
  \(D\) is the model's derivative operator and \(F\) the free energy of
  Eq. \eqref{eq:free_energy}.
\item
  Action: select \(a\) that steers \(\psi\) toward preferred outcomes;
  \passthrough{\lstinline!action\_update!} descends the action gradient,
  \(a \leftarrow a - \eta\,\partial_{a} F\) at step size \(\eta\).
\end{itemize}

Continuous-time flows (Einstein.4D analogy for metric/geodesic
intuition): see \passthrough{\lstinline!perception\_update!} and
\passthrough{\lstinline!action\_update!} in
\passthrough{\lstinline!src/quadmath/inference/information.py!}; their
discrete Quadray counterpart is
\passthrough{\lstinline!discrete\_ivm\_descent!} in
\passthrough{\lstinline!src/quadmath/optimize/discrete\_variational.py!},
described above.

\subsection{Implications for AI and Robust
Computation}\label{implications-for-ai-and-robust-computation}

FEP/Active Inference provide algorithms that unify perception and action
under uncertainty, offering biologically plausible alternatives to
standard RL with adaptive exploration and robust decision-making. See
\href{https://arxiv.org/abs/1907.03876}{applications in AI
(arXiv:1907.03876)}.

\subsection{Code, Reproducibility, and
Cross-References}\label{code-reproducibility-and-cross-references}

-- Equation references:
\href{08_equations_appendix.md\#eq:free_energy}{Eq. (Free Energy)},
\href{08_equations_appendix.md\#eq:fim}{Eq. (FIM)},
\href{08_equations_appendix.md\#eq:natgrad}{Eq. (Natural Gradient)} in
\passthrough{\lstinline!08\_equations\_appendix.md!}. -- Code anchors
(for readers who want to run experiments):
\href{03_quadray_methods.md\#code:free_energy}{\passthrough{\lstinline!free\_energy!}},
\href{03_quadray_methods.md\#code:fisher_information_matrix}{\passthrough{\lstinline!fisher\_information\_matrix!}},
\href{03_quadray_methods.md\#code:natural_gradient_step}{\passthrough{\lstinline!natural\_gradient\_step!}},
\passthrough{\lstinline!perception\_update!},
\passthrough{\lstinline!action\_update!}, and
\passthrough{\lstinline!discrete\_ivm\_descent!} in
\passthrough{\lstinline!src/quadmath/inference/information.py!} and
\passthrough{\lstinline!src/quadmath/optimize/discrete\_variational.py!}.

Demo and figures generated by
\passthrough{\lstinline!quadmath/scripts/information\_demo.py!} and
\passthrough{\lstinline!quadmath/scripts/active\_inference\_figures.py!}
output to \passthrough{\lstinline!quadmath/output/!}:

\begin{itemize}
\tightlist
\item
  \textbf{Active Inference Visualizations (embedded above)}:
  \passthrough{\lstinline!figure\_13\_4d\_trajectory.png!},
  \passthrough{\lstinline!figure\_14\_free\_energy\_landscape.png!},
  \passthrough{\lstinline!partition\_tetrahedron.png!}
\item
  \textbf{Information Geometry Visualizations} (rendered in the
  Optimization in 4D section):
  \passthrough{\lstinline!fisher\_information\_matrix.png!},
  \passthrough{\lstinline!fisher\_information\_eigenspectrum.png!},
  \passthrough{\lstinline!natural\_gradient\_path.png!},
  \passthrough{\lstinline!free\_energy\_curve.png!}
\item
  \textbf{Raw data}: \passthrough{\lstinline!figure\_13\_data.npz!},
  \passthrough{\lstinline!figure\_14\_data.npz!},
  \passthrough{\lstinline!fisher\_information\_matrix.csv!},
  \passthrough{\lstinline!fisher\_information\_matrix.npz!} (F, grads,
  X, y, w\_true, w\_est),
  \passthrough{\lstinline!fisher\_information\_eigenvalues.csv!},
  \passthrough{\lstinline!fisher\_information\_eigensystem.npz!}
\item
  \textbf{External validation}: Cross-reference with volume calculations
  and tetrahedral modeling tools from the
  \href{https://github.com/4dsolutions}{4dsolutions ecosystem}. See the
  \href{07_resources.md}{Resources} section for comprehensive details.
\end{itemize}

\newpage

\section{Appendix: Symbols and
Glossary}\label{appendix-symbols-and-glossary}

This appendix consolidates the symbols, variables, and constants used
throughout the manuscript.

\subsection{Sets and Spaces}\label{sets-and-spaces}

{\def\LTcaptype{none} % do not increment counter
\begin{longtable}[]{@{}ll@{}}
\toprule\noalign{}
Symbol & Name \\
\midrule\noalign{}
\endhead
\bottomrule\noalign{}
\endlastfoot
\(\mathbb{R}^n\) & Euclidean space \\
IVM & Isotropic Vector Matrix \\
Coxeter.4D & Euclidean 4D (E⁴) \\
Einstein.4D & Minkowski spacetime (3+1) \\
Fuller.4D & Synergetics/Quadray tetrahedral space \\
\end{longtable}
}

Descriptions:

\begin{itemize}
\tightlist
\item
  \(\mathbb{R}^n\): \(n\)-dimensional real vector space.
\item
  IVM: Quadray integer lattice (CCP sphere centers).
\item
  Coxeter.4D: Four-dimensional Euclidean geometry (not spacetime); see
  Coxeter, Regular Polytopes (Dover ed., p.~119); related
  lattice/packing background in Conway \& Sloane.
\item
  Einstein.4D: Relativistic spacetime with Minkowski metric.
\item
  Fuller.4D: Quadrays with projective normalization and IVM unit
  conventions.
\end{itemize}

\subsection{Quadray Coordinates and
Geometry}\label{quadray-coordinates-and-geometry}

{\def\LTcaptype{none} % do not increment counter
\begin{longtable}[]{@{}
  >{\raggedright\arraybackslash}p{(\linewidth - 4\tabcolsep) * \real{0.3333}}
  >{\raggedright\arraybackslash}p{(\linewidth - 4\tabcolsep) * \real{0.3333}}
  >{\raggedright\arraybackslash}p{(\linewidth - 4\tabcolsep) * \real{0.3333}}@{}}
\toprule\noalign{}
\begin{minipage}[b]{\linewidth}\raggedright
Symbol
\end{minipage} & \begin{minipage}[b]{\linewidth}\raggedright
Name
\end{minipage} & \begin{minipage}[b]{\linewidth}\raggedright
Description
\end{minipage} \\
\midrule\noalign{}
\endhead
\bottomrule\noalign{}
\endlastfoot
\(q=(a,b,c,d)\) & Quadray point & Non-negative coordinates with at least
one zero after normalization \\
\(A,B,C,D\) & Quadray axes & Canonical tetrahedral axes mapped by the
embedding \\
\(k\) & Normalization offset & \(k=\min(a,b,c,d)\) used to set
\(q' = q - (k,k,k,k)\) \\
\(q'\) & Normalized Quadray & Canonical representative with at least one
zero and non-negative entries \\
\(P_0,\ldots,P_3\) & Tetrahedron vertices & Vertices used in volume
formulas \\
\(d_{ij}\) & Pairwise distances & Distance between vertices \(P_i\) and
\(P_j\) (squared in CM matrix) \\
\(\det(\cdot)\) & Determinant & Determinant of a matrix \\
\(\lvert\cdot\rvert\) & Magnitude & Absolute value (determinant
magnitude) \\
\(V_{ivm}\) & Tetravolume (IVM) & Tetrahedron volume in synergetics/IVM
units; unit regular tetra has \(V_{ivm}=1\) \\
\(V_{xyz}\) & Tetravolume (XYZ) & Euclidean tetrahedron volume \\
\(S3\) & Scale factor & \(S3=\sqrt{9/8}\) with \(V_{ivm} = S3\,V_{xyz}\)
(synergetics unit convention) \\
Coxeter.4D & Namespace & Euclidean E⁴; regular polytopes \\
Einstein.4D & Namespace & Minkowski spacetime (metric analogy only
here) \\
Fuller.4D & Namespace & Quadrays/IVM; integer tetravolume \\
Eq. (lattice\_det) & Lattice determinant & Integer-lattice volume via
3x3 determinant \\
Eq. (ace5x5) & Tom Ace 5x5 & Direct IVM tetravolume from Quadrays \\
Eq. (cayley\_menger) & Cayley--Menger & Length-based formula: 288 V\^{}2
= det(·) \\
\end{longtable}
}

\subsection{Optimization and
Algorithms}\label{optimization-and-algorithms}

{\def\LTcaptype{none} % do not increment counter
\begin{longtable}[]{@{}ll@{}}
\toprule\noalign{}
Symbol & Name \\
\midrule\noalign{}
\endhead
\bottomrule\noalign{}
\endlastfoot
\(\alpha\) & Reflection coefficient \\
\(\gamma\) & Expansion coefficient \\
\(\rho\) & Contraction coefficient \\
\(\sigma\) & Shrink coefficient \\
\(V_{ivm}\) & Integer volume monitor \\
\end{longtable}
}

Descriptions:

\begin{itemize}
\tightlist
\item
  \(\alpha,\gamma,\rho,\sigma\): Nelder--Mead parameters (typical values
  1, 2, 0.5, 0.5).
\item
  \(V_{ivm}\): Tracks simplex volume across iterations.
\end{itemize}

\subsection{Information Theory and
Geometry}\label{information-theory-and-geometry}

{\def\LTcaptype{none} % do not increment counter
\begin{longtable}[]{@{}
  >{\raggedright\arraybackslash}p{(\linewidth - 4\tabcolsep) * \real{0.3333}}
  >{\raggedright\arraybackslash}p{(\linewidth - 4\tabcolsep) * \real{0.3333}}
  >{\raggedright\arraybackslash}p{(\linewidth - 4\tabcolsep) * \real{0.3333}}@{}}
\toprule\noalign{}
\begin{minipage}[b]{\linewidth}\raggedright
Symbol
\end{minipage} & \begin{minipage}[b]{\linewidth}\raggedright
Name
\end{minipage} & \begin{minipage}[b]{\linewidth}\raggedright
Description
\end{minipage} \\
\midrule\noalign{}
\endhead
\bottomrule\noalign{}
\endlastfoot
\(\log\) & Natural logarithm & Logarithm base \(e\) \\
\(\mathbb{E}[\cdot]\) & Expectation & Mean with respect to a
distribution \\
\(F_{ij}\) & Fisher Information Matrix &
\(\mathbb{E}[\partial_{\theta_i}\log p \cdot \partial_{\theta_j}\log p]\);
Eq. \eqref{eq:fim} in the equations appendix \\
\(\mathcal{F}\) & Variational free energy &
\(-\log P(o\mid s) + \mathrm{KL}\big[Q(s)\,\|\,P(s)\big]\); Eq.
\eqref{eq:free_energy} in the equations appendix \\
\(\mathrm{KL}[Q\,\|\,P]\) & Kullback--Leibler divergence &
\(\sum Q\log(Q/P)\); information distance \\
\(\nabla_{\theta} L\) & Natural gradient &
\(F(\theta)^{-1} \nabla_{\theta} L(\theta)\); Eq.
\eqref{eq:natural_gradient} in the equations appendix \\
\(\eta\) & Step size & Learning-rate scalar used in updates \\
\(\theta\) & Parameters & Model parameter vector; indices
\(\theta_i\) \\
\(ds^2\) & Minkowski line element & \(-c^2\,dt^2 + dx^2 + dy^2 + dz^2\);
Eq. \eqref{eq:minkowski_line_element} in the equations appendix \\
\(c\) & Speed of light & Physical constant appearing in Minkowski
metric \\
\end{longtable}
}

\subsection{Embeddings and Distances}\label{embeddings-and-distances}

{\def\LTcaptype{none} % do not increment counter
\begin{longtable}[]{@{}
  >{\raggedright\arraybackslash}p{(\linewidth - 4\tabcolsep) * \real{0.3333}}
  >{\raggedright\arraybackslash}p{(\linewidth - 4\tabcolsep) * \real{0.3333}}
  >{\raggedright\arraybackslash}p{(\linewidth - 4\tabcolsep) * \real{0.3333}}@{}}
\toprule\noalign{}
\begin{minipage}[b]{\linewidth}\raggedright
Symbol
\end{minipage} & \begin{minipage}[b]{\linewidth}\raggedright
Name
\end{minipage} & \begin{minipage}[b]{\linewidth}\raggedright
Description
\end{minipage} \\
\midrule\noalign{}
\endhead
\bottomrule\noalign{}
\endlastfoot
\(M\) & Embedding matrix & Linear map from Quadray to \(\mathbb{R}^3\)
(Urner-style unless noted) \\
\(\lVert\cdot\rVert_2\) & Euclidean norm &
\(\sqrt{x_1^2+\cdots+x_n^2}\) \\
\(R, D\) & Edge scales & Cube edge \(R\) and Quadray edge \(D\) with
\(D=2R\) (common convention) \\
\end{longtable}
}

\subsection{Greek Letters (usage)}\label{greek-letters-usage}

{\def\LTcaptype{none} % do not increment counter
\begin{longtable}[]{@{}
  >{\raggedright\arraybackslash}p{(\linewidth - 4\tabcolsep) * \real{0.3333}}
  >{\raggedright\arraybackslash}p{(\linewidth - 4\tabcolsep) * \real{0.3333}}
  >{\raggedright\arraybackslash}p{(\linewidth - 4\tabcolsep) * \real{0.3333}}@{}}
\toprule\noalign{}
\begin{minipage}[b]{\linewidth}\raggedright
Symbol
\end{minipage} & \begin{minipage}[b]{\linewidth}\raggedright
Name
\end{minipage} & \begin{minipage}[b]{\linewidth}\raggedright
Description
\end{minipage} \\
\midrule\noalign{}
\endhead
\bottomrule\noalign{}
\endlastfoot
\(\alpha,\gamma,\rho,\sigma\) & NM coefficients & Nelder--Mead
parameters (reflection, expansion, contraction, shrink) \\
\(\theta\) & Theta & Parameter vector in models and metrics \\
\(\mu\) & Mu & Internal states (Active Inference) \\
\(\psi\) & Psi & External states (Active Inference) \\
\(\eta\) & Eta & Step size / learning rate \\
\end{longtable}
}

\subsection{Notes (usage and
cross-references)}\label{notes-usage-and-cross-references}

\begin{itemize}
\tightlist
\item
  \textbf{Figures referenced}: In-text references use LaTeX's automatic
  figure numbering for consistent cross-referencing.
\item
  \textbf{Equation references}: Use labels defined in the text (e.g.,
  Eq. \eqref{eq:lattice_det} in the equations appendix).
\item
  \textbf{Namespaces}: We use Coxeter.4D, Einstein.4D, Fuller.4D
  consistently to designate Euclidean E⁴, Minkowski spacetime, and
  Quadray/IVM synergetics, respectively. This avoids conflation of
  Euclidean 4D objects (e.g., tesseracts) with spacetime constructs and
  synergetic tetravolume conventions.
\item
  \textbf{External validation}: Cross-reference implementations from the
  \href{https://github.com/4dsolutions}{4dsolutions ecosystem} for
  algorithmic verification and performance comparison baselines. See the
  \href{07_resources.md}{Resources} section for comprehensive details.
\end{itemize}

\subsection{Polyhedra and Synergetic
Shapes}\label{polyhedra-and-synergetic-shapes}

{\def\LTcaptype{none} % do not increment counter
\begin{longtable}[]{@{}
  >{\raggedright\arraybackslash}p{(\linewidth - 4\tabcolsep) * \real{0.3333}}
  >{\raggedright\arraybackslash}p{(\linewidth - 4\tabcolsep) * \real{0.3333}}
  >{\raggedright\arraybackslash}p{(\linewidth - 4\tabcolsep) * \real{0.3333}}@{}}
\toprule\noalign{}
\begin{minipage}[b]{\linewidth}\raggedright
Symbol
\end{minipage} & \begin{minipage}[b]{\linewidth}\raggedright
Name
\end{minipage} & \begin{minipage}[b]{\linewidth}\raggedright
Description
\end{minipage} \\
\midrule\noalign{}
\endhead
\bottomrule\noalign{}
\endlastfoot
Tetrahedron & Regular tetrahedron & Fundamental unit with V=1 in IVM
units \\
Cube & Regular hexahedron & V=3 in IVM units; orthogonal
space-filling \\
Octahedron & Regular octahedron & V=4 in IVM units; edge-midpoint
construction \\
Rhombic Dodecahedron & 12-faced solid & V=6 in IVM units; Voronoi cell
of FCC packing \\
Cuboctahedron & Vector equilibrium & V=20 in IVM units; shell of 12 IVM
neighbors \\
Truncated Octahedron & Archimedean solid & V=20 in IVM units;
space-filling tiling \\
\end{longtable}
}

\subsection{Acronyms and
abbreviations}\label{acronyms-and-abbreviations}

{\def\LTcaptype{none} % do not increment counter
\begin{longtable}[]{@{}
  >{\raggedright\arraybackslash}p{(\linewidth - 2\tabcolsep) * \real{0.5000}}
  >{\raggedright\arraybackslash}p{(\linewidth - 2\tabcolsep) * \real{0.5000}}@{}}
\toprule\noalign{}
\begin{minipage}[b]{\linewidth}\raggedright
Acronym
\end{minipage} & \begin{minipage}[b]{\linewidth}\raggedright
Meaning
\end{minipage} \\
\midrule\noalign{}
\endhead
\bottomrule\noalign{}
\endlastfoot
CM & Cayley--Menger (determinant-based tetrahedron volume) \\
PdF & Piero della Francesca (Heron-like tetrahedron volume) \\
GdJ & Gerald de Jong (Quadray-native tetravolume expression) \\
K-FAC & Kronecker-Factored Approximate Curvature (optimizer using
structured Fisher) \\
CCP & Cubic Close Packing (same centers as FCC) \\
FCC & Face-Centered Cubic (same centers as CCP) \\
E⁴ & Four-dimensional Euclidean space (Coxeter.4D) \\
NM & Nelder--Mead (simplex optimization algorithm) \\
4dsolutions & Kirby Urner's GitHub organization with extensive Quadray
implementations \\
BEAST & Synergetic modules (B, E, A, S, T) in Fuller's hierarchical
system \\
OCN & Oregon Curriculum Network (educational framework integrating
Quadrays) \\
POV-Ray & Persistence of Vision Raytracer (used in quadcraft.py
visualizations) \\
\end{longtable}
}

\subsection{API Index (auto-generated; Methods
linkage)}\label{api-index-auto-generated-methods-linkage}

The table below enumerates public symbols from
\passthrough{\lstinline!src/!} modules.

{\def\LTcaptype{none} % do not increment counter
\begin{longtable}[]{@{}
  >{\raggedright\arraybackslash}p{(\linewidth - 8\tabcolsep) * \real{0.2000}}
  >{\raggedright\arraybackslash}p{(\linewidth - 8\tabcolsep) * \real{0.2000}}
  >{\raggedright\arraybackslash}p{(\linewidth - 8\tabcolsep) * \real{0.2000}}
  >{\raggedright\arraybackslash}p{(\linewidth - 8\tabcolsep) * \real{0.2000}}
  >{\raggedright\arraybackslash}p{(\linewidth - 8\tabcolsep) * \real{0.2000}}@{}}
\toprule\noalign{}
\begin{minipage}[b]{\linewidth}\raggedright
Module
\end{minipage} & \begin{minipage}[b]{\linewidth}\raggedright
Symbol
\end{minipage} & \begin{minipage}[b]{\linewidth}\raggedright
Kind
\end{minipage} & \begin{minipage}[b]{\linewidth}\raggedright
Signature
\end{minipage} & \begin{minipage}[b]{\linewidth}\raggedright
Summary
\end{minipage} \\
\midrule\noalign{}
\endhead
\bottomrule\noalign{}
\endlastfoot
\passthrough{\lstinline!quadmath.core.cayley\_menger!} &
\passthrough{\lstinline!ivm\_tetra\_volume\_cayley\_menger!} & function
& \passthrough{\lstinline!(d2)!} & Compute IVM tetravolume from squared
distances via Cayley--Menger. \\
\passthrough{\lstinline!quadmath.core.cayley\_menger!} &
\passthrough{\lstinline!squared\_distances\_from\_quadrays!} & function
& \passthrough{\lstinline!(p0, p1, p2, p3, embedding)!} & Build the 4x4
squared-distance matrix from four quadray vertices. \\
\passthrough{\lstinline!quadmath.core.cayley\_menger!} &
\passthrough{\lstinline!tetra\_circumradius!} & function &
\passthrough{\lstinline!(d2)!} & Circumscribed sphere radius of a
tetrahedron from squared distances. \\
\passthrough{\lstinline!quadmath.core.cayley\_menger!} &
\passthrough{\lstinline!tetra\_inradius!} & function &
\passthrough{\lstinline!(d2)!} & Inscribed sphere radius of a
tetrahedron from squared distances. \\
\passthrough{\lstinline!quadmath.core.cayley\_menger!} &
\passthrough{\lstinline!tetra\_volume\_cayley\_menger!} & function &
\passthrough{\lstinline!(d2)!} & Compute Euclidean tetrahedron volume
from squared distances (Coxeter.4D). \\
\passthrough{\lstinline!quadmath.core.examples!} &
\passthrough{\lstinline!example\_cuboctahedron\_neighbors!} & function &
\passthrough{\lstinline!()!} & Return twelve-around-one IVM neighbors
(vector equilibrium shell). \\
\passthrough{\lstinline!quadmath.core.examples!} &
\passthrough{\lstinline!example\_cuboctahedron\_vertices\_xyz!} &
function & \passthrough{\lstinline!()!} & Return XYZ coordinates for the
twelve-around-one neighbors. \\
\passthrough{\lstinline!quadmath.core.examples!} &
\passthrough{\lstinline!example\_ivm\_neighbors!} & function &
\passthrough{\lstinline!()!} & Return the 12 nearest IVM neighbors as
permutations of \{2,1,1,0\} (Fuller.4D). \\
\passthrough{\lstinline!quadmath.core.examples!} &
\passthrough{\lstinline!example\_optimize!} & function &
\passthrough{\lstinline!()!} & Run Nelder--Mead over integer quadrays
for a simple convex objective (Fuller.4D). \\
\passthrough{\lstinline!quadmath.core.examples!} &
\passthrough{\lstinline!example\_partition\_tetra\_volume!} & function &
\passthrough{\lstinline!(mu, s, a, psi)!} & Construct a tetrahedron from
the four-fold partition and return tetravolume (Fuller.4D). \\
\passthrough{\lstinline!quadmath.core.examples!} &
\passthrough{\lstinline!example\_volume!} & function &
\passthrough{\lstinline!()!} & Return the exact IVM tetravolume of the
primitive lattice tetrahedron. \\
\passthrough{\lstinline!quadmath.core.geometry!} &
\passthrough{\lstinline!lorentz\_factor!} & function &
\passthrough{\lstinline!(v, c)!} & Lorentz factor gamma = 1 / sqrt(1 -
v\textsuperscript{2/c}2) (Einstein.4D). \\
\passthrough{\lstinline!quadmath.core.geometry!} &
\passthrough{\lstinline!minkowski\_interval!} & function &
\passthrough{\lstinline!(dt, dx, dy, dz, c)!} & Return the Minkowski
interval squared ds\^{}2 (Einstein.4D). \\
\passthrough{\lstinline!quadmath.core.geometry!} &
\passthrough{\lstinline!proper\_time!} & function &
\passthrough{\lstinline!(dt, dx, dy, dz, c)!} & Proper time elapsed for
a timelike interval (Einstein.4D). \\
\passthrough{\lstinline!quadmath.core.geometry!} &
\passthrough{\lstinline!spacetime\_classify!} & function &
\passthrough{\lstinline!(ds2, tol)!} & Classify a Minkowski interval
squared as timelike, spacelike, or lightlike. \\
\passthrough{\lstinline!quadmath.core.linalg\_utils!} &
\passthrough{\lstinline!bareiss\_determinant\_int!} & function &
\passthrough{\lstinline!(matrix)!} & Compute an exact integer
determinant using the Bareiss algorithm. \\
\passthrough{\lstinline!quadmath.core.linalg\_utils!} &
\passthrough{\lstinline!bareiss\_rank!} & function &
\passthrough{\lstinline!(matrix)!} & Compute the exact integer rank of a
matrix via Bareiss elimination. \\
\passthrough{\lstinline!quadmath.core.linalg\_utils!} &
\passthrough{\lstinline!integer\_adjugate!} & function &
\passthrough{\lstinline!(matrix)!} & Compute the exact integer adjugate
(classical adjoint) of a square matrix. \\
\passthrough{\lstinline!quadmath.core.metrics!} &
\passthrough{\lstinline!angle\_error!} & function &
\passthrough{\lstinline!(q1, q2)!} & Geodesic rotation angle between two
quaternions, in radians. \\
\passthrough{\lstinline!quadmath.core.metrics!} &
\passthrough{\lstinline!fim\_eigenspectrum!} & function &
\passthrough{\lstinline!(F)!} & Eigen-decomposition of a Fisher
information matrix. \\
\passthrough{\lstinline!quadmath.core.metrics!} &
\passthrough{\lstinline!fisher\_condition\_number!} & function &
\passthrough{\lstinline!(F)!} & Compute the condition number of the
Fisher information matrix. \\
\passthrough{\lstinline!quadmath.core.metrics!} &
\passthrough{\lstinline!fisher\_curvature\_analysis!} & function &
\passthrough{\lstinline!(F)!} & Comprehensive analysis of Fisher
information matrix curvature. \\
\passthrough{\lstinline!quadmath.core.metrics!} &
\passthrough{\lstinline!fisher\_quadray\_comparison!} & function &
\passthrough{\lstinline!(F\_cartesian, F\_quadray)!} & Compare Fisher
information matrices between coordinate systems. \\
\passthrough{\lstinline!quadmath.core.metrics!} &
\passthrough{\lstinline!fisher\_rao\_metric!} & function &
\passthrough{\lstinline!(p, q, eps)!} & Fisher--Rao geodesic distance on
the probability simplex. \\
\passthrough{\lstinline!quadmath.core.metrics!} &
\passthrough{\lstinline!information\_length!} & function &
\passthrough{\lstinline!(path\_gradients)!} & Gradient-weighted proxy
for informational path length (NOT the \\
\passthrough{\lstinline!quadmath.core.metrics!} &
\passthrough{\lstinline!jensen\_shannon\_divergence!} & function &
\passthrough{\lstinline!(p, q, eps)!} & Jensen--Shannon divergence
between two discrete distributions. \\
\passthrough{\lstinline!quadmath.core.metrics!} &
\passthrough{\lstinline!kl\_divergence!} & function &
\passthrough{\lstinline!(p, q, eps)!} & Kullback--Leibler divergence
between two discrete distributions. \\
\passthrough{\lstinline!quadmath.core.metrics!} &
\passthrough{\lstinline!quat\_log\_euclidean\_dispersion!} & function &
\passthrough{\lstinline!(quats)!} & Root-mean-square chordal dispersion
of quaternions about their mean. \\
\passthrough{\lstinline!quadmath.core.metrics!} &
\passthrough{\lstinline!shannon\_entropy!} & function &
\passthrough{\lstinline!(p, eps)!} & Shannon entropy H(p) for a discrete
distribution. \\
\passthrough{\lstinline!quadmath.core.quadray!} &
\passthrough{\lstinline!DEFAULT\_EMBEDDING!} & constant & & \\
\passthrough{\lstinline!quadmath.core.quadray!} &
\passthrough{\lstinline!IVM\_NEIGHBOR\_MOVES!} & constant & & \\
\passthrough{\lstinline!quadmath.core.quadray!} &
\passthrough{\lstinline!Quadray!} & class & & Quadray vector with
non-negative components and at least one zero (Fuller.4D). \\
\passthrough{\lstinline!quadmath.core.quadray!} &
\passthrough{\lstinline!ace\_tetravolume\_5x5!} & function &
\passthrough{\lstinline!(p0, p1, p2, p3)!} & Tom Ace 5x5 determinant as
the exact IVM tetra-volume (Fuller.4D). \\
\passthrough{\lstinline!quadmath.core.quadray!} &
\passthrough{\lstinline!angle!} & function &
\passthrough{\lstinline!(q1, q2, q3, embedding)!} & Angle at vertex q2
formed by rays q2-\textgreater q1 and q2-\textgreater q3 (radians). \\
\passthrough{\lstinline!quadmath.core.quadray!} &
\passthrough{\lstinline!centroid!} & function &
\passthrough{\lstinline!(*quads)!} & Component-wise mean of quadray
points, rounded to the nearest lattice point. \\
\passthrough{\lstinline!quadmath.core.quadray!} &
\passthrough{\lstinline!distance!} & function &
\passthrough{\lstinline!(q1, q2, embedding)!} & Euclidean distance
between two quadray points under the given embedding. \\
\passthrough{\lstinline!quadmath.core.quadray!} &
\passthrough{\lstinline!dot!} & function &
\passthrough{\lstinline!(q1, q2, embedding)!} & Return Euclidean dot
product \textless q1,q2\textgreater{} under the given embedding. \\
\passthrough{\lstinline!quadmath.core.quadray!} &
\passthrough{\lstinline!integer\_tetra\_volume!} & function &
\passthrough{\lstinline!(p0, p1, p2, p3)!} & Compute the exact IVM
tetra-volume of a lattice tetrahedron (Fuller.4D). \\
\passthrough{\lstinline!quadmath.core.quadray!} &
\passthrough{\lstinline!magnitude!} & function &
\passthrough{\lstinline!(q, embedding)!} & Return the Euclidean
magnitude of \passthrough{\lstinline!q!} under the given embedding
(vector norm). \\
\passthrough{\lstinline!quadmath.core.quadray!} &
\passthrough{\lstinline!qconjugate!} & function &
\passthrough{\lstinline!(q)!} & Conjugate (w, -x, -y, -z) of a
quaternion in (w, x, y, z) order. \\
\passthrough{\lstinline!quadmath.core.quadray!} &
\passthrough{\lstinline!qmul!} & function &
\passthrough{\lstinline!(a, b)!} & Hamilton product of two
quaternions. \\
\passthrough{\lstinline!quadmath.core.quadray!} &
\passthrough{\lstinline!qrotate!} & function &
\passthrough{\lstinline!(q, v\_xyz, angle)!} & Rotate a 3-vector by a
unit quaternion via Rodrigues (v' = q v q*). \\
\passthrough{\lstinline!quadmath.core.quadray!} &
\passthrough{\lstinline!quadray\_from\_xyz!} & function &
\passthrough{\lstinline!(x, y, z, embedding)!} & Map an R\^{}3 point
back to the quadray lattice via pseudoinverse rounding. \\
\passthrough{\lstinline!quadmath.core.quadray!} &
\passthrough{\lstinline!rotate\_about\_axis!} & function &
\passthrough{\lstinline!(v\_xyz, axis\_xyz, angle)!} & Rotate a 3-vector
about an axis by an angle (axis-angle convenience). \\
\passthrough{\lstinline!quadmath.core.quadray!} &
\passthrough{\lstinline!slerp!} & function &
\passthrough{\lstinline!(qa, qb, t)!} & Shortest-arc spherical linear
interpolation between unit quaternions. \\
\passthrough{\lstinline!quadmath.core.quadray!} &
\passthrough{\lstinline!to\_xyz!} & function &
\passthrough{\lstinline!(q, embedding)!} & Map quadray to R\^{}3 via a
3x4 embedding matrix (Fuller.4D -\textgreater{} Coxeter.4D slice). \\
\passthrough{\lstinline!quadmath.core.symbolic!} &
\passthrough{\lstinline!cayley\_menger\_volume\_symbolic!} & function &
\passthrough{\lstinline!(d2)!} & Return symbolic Euclidean tetrahedron
volume from squared distances. \\
\passthrough{\lstinline!quadmath.core.symbolic!} &
\passthrough{\lstinline!convert\_xyz\_volume\_to\_ivm\_symbolic!} &
function & \passthrough{\lstinline!(V\_xyz)!} & Convert a symbolic
Euclidean volume to IVM tetravolume via S3. \\
\passthrough{\lstinline!quadmath.inference.information!} &
\passthrough{\lstinline!action\_update!} & function &
\passthrough{\lstinline!(action, free\_energy\_fn, step\_size, epsilon)!}
& Continuous-time action update: da/dt = - dF/da. \\
\passthrough{\lstinline!quadmath.inference.information!} &
\passthrough{\lstinline!active\_inference\_step!} & function &
\passthrough{\lstinline!(mu, action, free\_energy\_fn, derivative\_operator, step\_size, epsilon)!}
& Joint perception-action update step in Active Inference. \\
\passthrough{\lstinline!quadmath.inference.information!} &
\passthrough{\lstinline!expected\_free\_energy!} & function &
\passthrough{\lstinline!(log\_p\_o\_given\_s, q, p, log\_p\_o)!} &
Expected free energy for Active Inference with prior preferences. \\
\passthrough{\lstinline!quadmath.inference.information!} &
\passthrough{\lstinline!finite\_difference\_gradient!} & function &
\passthrough{\lstinline!(function, x, epsilon)!} & Compute numerical
gradient of a scalar function via central differences. \\
\passthrough{\lstinline!quadmath.inference.information!} &
\passthrough{\lstinline!fisher\_information\_matrix!} & function &
\passthrough{\lstinline!(gradients, normalize)!} & Estimate the Fisher
information matrix via sample gradients. \\
\passthrough{\lstinline!quadmath.inference.information!} &
\passthrough{\lstinline!fisher\_information\_quadray!} & function &
\passthrough{\lstinline!(gradients, embedding\_matrix)!} & Compute
Fisher information matrix in both Cartesian and Quadray coordinates. \\
\passthrough{\lstinline!quadmath.inference.information!} &
\passthrough{\lstinline!free\_energy!} & function &
\passthrough{\lstinline!(log\_p\_o\_given\_s, q, p)!} & Variational free
energy for discrete latent states. \\
\passthrough{\lstinline!quadmath.inference.information!} &
\passthrough{\lstinline!information\_gain!} & function &
\passthrough{\lstinline!(prior, posterior, eps)!} & Information gain
(Bayesian surprise) between prior and posterior. \\
\passthrough{\lstinline!quadmath.inference.information!} &
\passthrough{\lstinline!information\_geometric\_distance!} & function &
\passthrough{\lstinline!(F, x1, x2)!} & Compute information-geometric
distance between two points. \\
\passthrough{\lstinline!quadmath.inference.information!} &
\passthrough{\lstinline!mutual\_information!} & function &
\passthrough{\lstinline!(p\_joint, eps)!} & Mutual information I(X; Y)
from a joint probability matrix. \\
\passthrough{\lstinline!quadmath.inference.information!} &
\passthrough{\lstinline!natural\_gradient\_step!} & function &
\passthrough{\lstinline!(gradient, fisher, step\_size, ridge)!} &
Compute a natural gradient step using a damped inverse Fisher. \\
\passthrough{\lstinline!quadmath.inference.information!} &
\passthrough{\lstinline!perception\_update!} & function &
\passthrough{\lstinline!(mu, derivative\_operator, free\_energy\_fn, step\_size, epsilon)!}
& Continuous-time perception update: dmu/dt = D mu - dF/dmu. \\
\passthrough{\lstinline!quadmath.lattice.conversions!} &
\passthrough{\lstinline!embedding\_basis!} & function &
\passthrough{\lstinline!(M)!} & Return the four embedding COLUMNS as a
(4,3) basis matrix (Fuller.4D axes). \\
\passthrough{\lstinline!quadmath.lattice.conversions!} &
\passthrough{\lstinline!quadray\_roundtrip!} & function &
\passthrough{\lstinline!(q, M)!} & Assert the exact round-trip identity
recon == q for q -\textgreater{} XYZ -\textgreater{} quadray. \\
\passthrough{\lstinline!quadmath.lattice.conversions!} &
\passthrough{\lstinline!quadray\_to\_xyz!} & function &
\passthrough{\lstinline!(q, M)!} & Map a
\passthrough{\lstinline!Quadray!} to Cartesian XYZ via a 3x4 embedding
matrix (Fuller.4D -\textgreater{} Coxeter.4D slice). \\
\passthrough{\lstinline!quadmath.lattice.conversions!} &
\passthrough{\lstinline!urner\_embedding!} & function &
\passthrough{\lstinline!(scale)!} & Return a 3x4 Urner-style symmetric
embedding matrix (Fuller.4D -\textgreater{} Coxeter.4D slice). \\
\passthrough{\lstinline!quadmath.lattice.conversions!} &
\passthrough{\lstinline!xyz\_to\_quadray\_canonical!} & function &
\passthrough{\lstinline!(xyz, M)!} & Recover the canonical integer
quadray for an XYZ point in the embedding image, EXACTLY. \\
\passthrough{\lstinline!quadmath.lattice.ivm\_dynamics!} &
\passthrough{\lstinline!DynamicsParams!} & class & & Parameters of a
discrete IVM lattice dynamics run. \\
\passthrough{\lstinline!quadmath.lattice.ivm\_dynamics!} &
\passthrough{\lstinline!FitResult!} & class & & Outcome of gradient-free
coupling identification. \\
\passthrough{\lstinline!quadmath.lattice.ivm\_dynamics!} &
\passthrough{\lstinline!IVMLattice!} & class & & Finite IVM lattice ball
with neighbor adjacency and diffusion operators. \\
\passthrough{\lstinline!quadmath.lattice.ivm\_dynamics!} &
\passthrough{\lstinline!Trajectory!} & class & & Deterministic
simulation record. \\
\passthrough{\lstinline!quadmath.lattice.ivm\_dynamics!} &
\passthrough{\lstinline!ball\_sites!} & function &
\passthrough{\lstinline!(radius)!} & Enumerate canonical quadray sites
with squared IVM radius \textless= radius**2. \\
\passthrough{\lstinline!quadmath.lattice.ivm\_dynamics!} &
\passthrough{\lstinline!fit\_trajectory!} & function &
\passthrough{\lstinline!(observed, grid, lattice, kind, refine\_rounds)!}
& Identify the coupling alpha from an observed trajectory
(gradient-free). \\
\passthrough{\lstinline!quadmath.lattice.ivm\_dynamics!} &
\passthrough{\lstinline!heat\_step!} & function &
\passthrough{\lstinline!(u, lattice, alpha)!} & One heat-diffusion step
u \textless- (1 - alpha) u + alpha S u. \\
\passthrough{\lstinline!quadmath.lattice.ivm\_dynamics!} &
\passthrough{\lstinline!is\_nonincreasing!} & function &
\passthrough{\lstinline!(values, tol)!} & True iff no consecutive pair
of values increases by more than tol. \\
\passthrough{\lstinline!quadmath.lattice.ivm\_dynamics!} &
\passthrough{\lstinline!majority\_step!} & function &
\passthrough{\lstinline!(u, lattice, alpha)!} & One rounded-averaging
(``majority'') step on integer states. \\
\passthrough{\lstinline!quadmath.lattice.ivm\_dynamics!} &
\passthrough{\lstinline!make\_lattice!} & function &
\passthrough{\lstinline!(radius)!} & Build the finite IVM lattice ball
of the given radius. \\
\passthrough{\lstinline!quadmath.lattice.ivm\_dynamics!} &
\passthrough{\lstinline!neighbor\_shifts!} & function &
\passthrough{\lstinline!()!} & Return the 12 canonical IVM neighbor
shifts as quadray deltas. \\
\passthrough{\lstinline!quadmath.lattice.ivm\_dynamics!} &
\passthrough{\lstinline!render\_dynamics\_demo!} & function &
\passthrough{\lstinline!(output\_path, radius=, seed=, alpha=, horizon=, fit\_alpha=, grid\_count=, refine\_rounds=)!}
& Render the multi-snapshot IVM dynamics demo figure; return its
path. \\
\passthrough{\lstinline!quadmath.lattice.ivm\_dynamics!} &
\passthrough{\lstinline!simulate!} & function &
\passthrough{\lstinline!(T, params, lattice, u0)!} & Simulate T updates;
deterministic given \passthrough{\lstinline!params!} (fixed seed). \\
\passthrough{\lstinline!quadmath.lattice.ivm\_dynamics!} &
\passthrough{\lstinline!site\_radius\_sq!} & function &
\passthrough{\lstinline!(q)!} & Integer squared IVM radius of a quadray
site (squared embedding norm). \\
\passthrough{\lstinline!quadmath.lattice.ivm\_dynamics!} &
\passthrough{\lstinline!step!} & function &
\passthrough{\lstinline!(u, lattice, params)!} & Apply one update of the
dynamics in \passthrough{\lstinline!params!} to the field. \\
\passthrough{\lstinline!quadmath.lattice.ivm\_dynamics!} &
\passthrough{\lstinline!sum\_of\_squares!} & function &
\passthrough{\lstinline!(u)!} & Sum of squares of a field --- the
observable of the heat lemma. \\
\passthrough{\lstinline!quadmath.lattice.ivm\_field!} &
\passthrough{\lstinline!IVMField!} & class & & Scalar field over an IVM
lattice ball, stored on a deterministic site index. \\
\passthrough{\lstinline!quadmath.lattice.ivm\_field!} &
\passthrough{\lstinline!IVM\_NEIGHBOR\_STEPS!} & constant & & \\
\passthrough{\lstinline!quadmath.lattice.ivm\_field!} &
\passthrough{\lstinline!TetrahedronFit!} & class & & Result of
:func:\passthrough{\lstinline!fit\_geometry!}. \\
\passthrough{\lstinline!quadmath.lattice.ivm\_field!} &
\passthrough{\lstinline!fit\_geometry!} & function &
\passthrough{\lstinline!(points, labels, embedding)!} & Least-squares
recovery of a tetrahedron's orientation+scale from noisy 3D points. \\
\passthrough{\lstinline!quadmath.lattice.ivm\_field!} &
\passthrough{\lstinline!is\_ivm\_site!} & function &
\passthrough{\lstinline!(q)!} & Return True iff
\passthrough{\lstinline!q!} (after normalization) is an IVM lattice
site. \\
\passthrough{\lstinline!quadmath.lattice.ivm\_field!} &
\passthrough{\lstinline!quadray\_shell\_norm!} & function &
\passthrough{\lstinline!(q)!} & Return the IVM shell norm of
\passthrough{\lstinline!q!}: an even integer equal to
\passthrough{\lstinline!2k!}. \\
\passthrough{\lstinline!quadmath.lattice.ivm\_field!} &
\passthrough{\lstinline!shell\_ball\_sites!} & function &
\passthrough{\lstinline!(radius)!} & Enumerate the IVM lattice ball of
the given shell radius, deterministically. \\
\passthrough{\lstinline!quadmath.lattice.ivm\_field!} &
\passthrough{\lstinline!shell\_cardinalities!} & function &
\passthrough{\lstinline!(max\_shell)!} & Cardinalities of shells
\passthrough{\lstinline!0 .. max\_shell!} (cuboctahedral numbers). \\
\passthrough{\lstinline!quadmath.lattice.ivm\_field!} &
\passthrough{\lstinline!shell\_sites!} & function &
\passthrough{\lstinline!(k)!} & Enumerate all IVM lattice sites with
quadray shell norm \passthrough{\lstinline!2k!}. \\
\passthrough{\lstinline!quadmath.lattice.lattice\_search!} &
\passthrough{\lstinline!nearest!} & function &
\passthrough{\lstinline!(site, R, k)!} & Return the
\passthrough{\lstinline!k!} nearest IVM lattice sites within radius
\passthrough{\lstinline!R!}. \\
\passthrough{\lstinline!quadmath.lattice.lattice\_search!} &
\passthrough{\lstinline!squared\_distance!} & function &
\passthrough{\lstinline!(p, sites)!} & Exact lattice squared distances
from \passthrough{\lstinline!p!} to rows of
\passthrough{\lstinline!sites!}. \\
\passthrough{\lstinline!quadmath.lattice.lattice\_search!} &
\passthrough{\lstinline!within\_radius!} & function &
\passthrough{\lstinline!(site, R)!} & Return all IVM lattice sites
within Euclidean radius \passthrough{\lstinline!R!} of
\passthrough{\lstinline!site!}. \\
\passthrough{\lstinline!quadmath.lattice.omni\_numbering!} &
\passthrough{\lstinline!MAX\_SHELL!} & constant & & \\
\passthrough{\lstinline!quadmath.lattice.omni\_numbering!} &
\passthrough{\lstinline!NEIGHBOR\_MOVES!} & constant & & \\
\passthrough{\lstinline!quadmath.lattice.omni\_numbering!} &
\passthrough{\lstinline!clear\_shell\_cache!} & function &
\passthrough{\lstinline!()!} & Drop memoized shell enumerations so the
next call recomputes them. \\
\passthrough{\lstinline!quadmath.lattice.omni\_numbering!} &
\passthrough{\lstinline!cumulative\_count!} & function &
\passthrough{\lstinline!(k)!} & Return the total number of IVM sites
through shell \passthrough{\lstinline!k!} (inclusive). \\
\passthrough{\lstinline!quadmath.lattice.omni\_numbering!} &
\passthrough{\lstinline!generate\_shell!} & function &
\passthrough{\lstinline!(k)!} & Generate the sites of shell
\passthrough{\lstinline!k!} of the omnidirectional close packing. \\
\passthrough{\lstinline!quadmath.lattice.omni\_numbering!} &
\passthrough{\lstinline!shell\_count!} & function &
\passthrough{\lstinline!(k)!} & Return the number of IVM sites on shell
\passthrough{\lstinline!k!} of the close packing. \\
\passthrough{\lstinline!quadmath.lattice.omni\_numbering!} &
\passthrough{\lstinline!site\_at\_index!} & function &
\passthrough{\lstinline!(index, max\_shell)!} & Return the IVM site at a
canonical global index. \\
\passthrough{\lstinline!quadmath.lattice.omni\_numbering!} &
\passthrough{\lstinline!site\_index!} & function &
\passthrough{\lstinline!(site, max\_shell)!} & Return the canonical
global index of an IVM site, or -1 if absent. \\
\passthrough{\lstinline!quadmath.lattice.omni\_numbering!} &
\passthrough{\lstinline!sites\_through\_shell!} & function &
\passthrough{\lstinline!(max\_shell)!} & Return all IVM sites through
shell \passthrough{\lstinline!max\_shell!} in canonical order. \\
\passthrough{\lstinline!quadmath.learn.learning\_eval!} &
\passthrough{\lstinline!CrossValidationResult!} & class & & K-fold
cross-validation table for the IVM field learner. \\
\passthrough{\lstinline!quadmath.learn.learning\_eval!} &
\passthrough{\lstinline!GradientDescentTrainer!} & class & & Full-batch
gradient-descent fit of a linear model on standardized features. \\
\passthrough{\lstinline!quadmath.learn.learning\_eval!} &
\passthrough{\lstinline!LearningCurveResult!} & class & & Data-coverage
curve for the IVM field learner. \\
\passthrough{\lstinline!quadmath.learn.learning\_eval!} &
\passthrough{\lstinline!RidgeSiteFit!} & class & & Closed-form ridge fit
of a single linear site model. \\
\passthrough{\lstinline!quadmath.learn.learning\_eval!} &
\passthrough{\lstinline!TrajectorySplitResult!} & class & & Outcome of a
temporal train/test evaluation of dynamics identification. \\
\passthrough{\lstinline!quadmath.learn.learning\_eval!} &
\passthrough{\lstinline!cross\_validate\_field!} & function &
\passthrough{\lstinline!(field\_values, sites, k, lam\_grid, seed, radius=, kernel\_width=)!}
& Cross-validate the Laplacian-regularized field learner over k
folds. \\
\passthrough{\lstinline!quadmath.learn.learning\_eval!} &
\passthrough{\lstinline!enclosing\_radius!} & function &
\passthrough{\lstinline!(sites)!} & Return the smallest IVM ball radius
containing every given site. \\
\passthrough{\lstinline!quadmath.learn.learning\_eval!} &
\passthrough{\lstinline!kfold\_site\_splits!} & function &
\passthrough{\lstinline!(observed\_sites, k, seed)!} & Partition
observed lattice sites into \passthrough{\lstinline!k!} seeded folds. \\
\passthrough{\lstinline!quadmath.learn.learning\_eval!} &
\passthrough{\lstinline!learning\_curve!} & function &
\passthrough{\lstinline!(values, sites, train\_fracs, seed, radius=, lam=, kernel\_width=)!}
& Trace held-out MSE as a function of the fraction of observed sites. \\
\passthrough{\lstinline!quadmath.learn.learning\_eval!} &
\passthrough{\lstinline!ridge\_site\_fit!} & function &
\passthrough{\lstinline!(features, values, lam)!} & Fit a closed-form
ridge regression with intercept on tabular rows. \\
\passthrough{\lstinline!quadmath.learn.learning\_eval!} &
\passthrough{\lstinline!three\_way\_split!} & function &
\passthrough{\lstinline!(n, train\_frac, val\_frac, seed)!} & Split
\passthrough{\lstinline!range(n)!} into seeded, disjoint train/val/test
index lists. \\
\passthrough{\lstinline!quadmath.learn.learning\_eval!} &
\passthrough{\lstinline!trajectory\_train\_test!} & function &
\passthrough{\lstinline!(observed, split\_frac, lattice, kind=, grid=, refine\_rounds=)!}
& Identify dynamics parameters on a training prefix, score the held-out
suffix. \\
\passthrough{\lstinline!quadmath.optimize.discrete\_variational!} &
\passthrough{\lstinline!DiscretePath!} & class & & Optimization
trajectory on the integer quadray lattice. \\
\passthrough{\lstinline!quadmath.optimize.discrete\_variational!} &
\passthrough{\lstinline!apply\_move!} & function &
\passthrough{\lstinline!(q, delta)!} & Apply a lattice move and
normalize to the canonical representative. \\
\passthrough{\lstinline!quadmath.optimize.discrete\_variational!} &
\passthrough{\lstinline!discrete\_ivm\_descent!} & function &
\passthrough{\lstinline!(objective, start, moves=, max\_iter=, on\_step=)!}
& Greedy discrete descent over the quadray integer lattice. \\
\passthrough{\lstinline!quadmath.optimize.discrete\_variational!} &
\passthrough{\lstinline!neighbor\_moves\_ivm!} & function &
\passthrough{\lstinline!()!} & Return the 12 canonical IVM neighbor
moves as Quadray deltas. \\
\passthrough{\lstinline!quadmath.optimize.nelder\_mead\_quadray!} &
\passthrough{\lstinline!SimplexState!} & class & & \\
\passthrough{\lstinline!quadmath.optimize.nelder\_mead\_quadray!} &
\passthrough{\lstinline!centroid\_excluding!} & function &
\passthrough{\lstinline!(vertices, exclude\_idx)!} & Integer centroid of
three vertices, excluding the specified index. \\
\passthrough{\lstinline!quadmath.optimize.nelder\_mead\_quadray!} &
\passthrough{\lstinline!compute\_volume!} & function &
\passthrough{\lstinline!(vertices)!} & Exact IVM tetra-volume (a
Fraction, absolute determinant divided by 4) of the first four
vertices. \\
\passthrough{\lstinline!quadmath.optimize.nelder\_mead\_quadray!} &
\passthrough{\lstinline!nelder\_mead\_quadray!} & function &
\passthrough{\lstinline!(f, initial\_vertices, alpha, gamma, rho, sigma, max\_iter, tol, on\_step)!}
& Nelder--Mead on the integer quadray lattice. \\
\passthrough{\lstinline!quadmath.optimize.nelder\_mead\_quadray!} &
\passthrough{\lstinline!order\_simplex!} & function &
\passthrough{\lstinline!(vertices, f)!} & Sort vertices by objective
value ascending and return paired lists. \\
\passthrough{\lstinline!quadmath.optimize.nelder\_mead\_quadray!} &
\passthrough{\lstinline!project\_to\_lattice!} & function &
\passthrough{\lstinline!(q)!} & Project a quadray to the canonical
lattice representative via normalize. \\
\passthrough{\lstinline!quadmath.paths!} &
\passthrough{\lstinline!get\_data\_dir!} & function &
\passthrough{\lstinline!()!} & Return
\passthrough{\lstinline!quadmath/output/data!} path and ensure it
exists. \\
\passthrough{\lstinline!quadmath.paths!} &
\passthrough{\lstinline!get\_figure\_dir!} & function &
\passthrough{\lstinline!()!} & Return
\passthrough{\lstinline!quadmath/output/figures!} path and ensure it
exists. \\
\passthrough{\lstinline!quadmath.paths!} &
\passthrough{\lstinline!get\_output\_dir!} & function &
\passthrough{\lstinline!()!} & Return
\passthrough{\lstinline!quadmath/output!} path at the repo root and
ensure it exists. \\
\passthrough{\lstinline!quadmath.paths!} &
\passthrough{\lstinline!get\_repo\_root!} & function &
\passthrough{\lstinline!(start)!} & Heuristically find repository root
by walking up from \passthrough{\lstinline!start!}. \\
\passthrough{\lstinline!quadmath.pipeline!} &
\passthrough{\lstinline!FieldLearner!} & class & & Laplacian-regularized
field learner over an IVM ball; satisfies
:class:\passthrough{\lstinline!Fittable!}. \\
\passthrough{\lstinline!quadmath.pipeline!} &
\passthrough{\lstinline!FieldModel!} & class & & Anything that predicts
a scalar field value at a lattice site. \\
\passthrough{\lstinline!quadmath.pipeline!} &
\passthrough{\lstinline!Fittable!} & class & & Anything that can fit a
field model from data, then predict. \\
\passthrough{\lstinline!quadmath.pipeline!} &
\passthrough{\lstinline!LatticeBall!} & class & & Immutable view of an
IVM lattice ball; satisfies
:class:\passthrough{\lstinline!LatticeSource!}. \\
\passthrough{\lstinline!quadmath.pipeline!} &
\passthrough{\lstinline!LatticeSource!} & class & & Anything that can
enumerate IVM lattice sites (structural). \\
\passthrough{\lstinline!quadmath.pipeline!} &
\passthrough{\lstinline!Pipeline!} & class & & Immutable sequence of
:class:\passthrough{\lstinline!Step!} objects; monoid-style
composition. \\
\passthrough{\lstinline!quadmath.pipeline!} &
\passthrough{\lstinline!Step!} & class & & A named, typed callable unit
of a :class:\passthrough{\lstinline!Pipeline!}. \\
\passthrough{\lstinline!quadmath.pipeline!} &
\passthrough{\lstinline!dynamics\_step!} & function &
\passthrough{\lstinline!(params, T)!} & Step simulating the discrete IVM
dynamics in \passthrough{\lstinline!params!} (ignores input). \\
\passthrough{\lstinline!quadmath.pipeline!} &
\passthrough{\lstinline!learn\_step!} & function &
\passthrough{\lstinline!(lam, seed, radius)!} & Step fitting an
:class:\passthrough{\lstinline!IVMField!} by Laplacian-regularized
learning. \\
\passthrough{\lstinline!quadmath.pipeline!} &
\passthrough{\lstinline!sites\_step!} & function &
\passthrough{\lstinline!(radius)!} & Step producing the IVM ball sites
of shell \passthrough{\lstinline!radius!} (ignores input). \\
\passthrough{\lstinline!quadmath.stats.benchmarks!} &
\passthrough{\lstinline!BENCH\_DEFAULTS!} & constant & & \\
\passthrough{\lstinline!quadmath.stats.benchmarks!} &
\passthrough{\lstinline!BenchRow!} & class & & One timed benchmark
result. \\
\passthrough{\lstinline!quadmath.stats.benchmarks!} &
\passthrough{\lstinline!bench\_conversions!} & function &
\passthrough{\lstinline!(n, trials)!} & Benchmark quadray/XYZ
conversions over \passthrough{\lstinline!n!} deterministic samples. \\
\passthrough{\lstinline!quadmath.stats.benchmarks!} &
\passthrough{\lstinline!bench\_field\_fit!} & function &
\passthrough{\lstinline!(n\_sites, trials)!} & Benchmark
\passthrough{\lstinline!IVMField.learn!} on a synthetic field with a
fixed seed. \\
\passthrough{\lstinline!quadmath.stats.benchmarks!} &
\passthrough{\lstinline!bench\_lattice\_search!} & function &
\passthrough{\lstinline!(n\_sites, queries, trials)!} & Benchmark
nearest-site queries through the
\passthrough{\lstinline!lattice\_search!} ball index. \\
\passthrough{\lstinline!quadmath.stats.benchmarks!} &
\passthrough{\lstinline!bench\_shell\_enumeration!} & function &
\passthrough{\lstinline!(k\_max, trials)!} & Benchmark shell enumeration
through shell \passthrough{\lstinline!k\_max!}. \\
\passthrough{\lstinline!quadmath.stats.benchmarks!} &
\passthrough{\lstinline!run\_all!} & function &
\passthrough{\lstinline!()!} & Run every benchmark with the module-level
:\url{data:\%60BENCH_DEFAULTS}`. \\
\passthrough{\lstinline!quadmath.stats.benchmarks!} &
\passthrough{\lstinline!summary\_table!} & function &
\passthrough{\lstinline!(rows)!} & Render aligned fixed-width ASCII rows
as a table. \\
\passthrough{\lstinline!quadmath.stats.benchmarks!} &
\passthrough{\lstinline!time\_callable!} & function &
\passthrough{\lstinline!(fn, trials=, warmup=)!} & Time
\passthrough{\lstinline!fn!} with
\passthrough{\lstinline!time.perf\_counter!} and return per-trial wall
seconds. \\
\passthrough{\lstinline!quadmath.stats.statistics!} &
\passthrough{\lstinline!benjamini\_hochberg!} & function &
\passthrough{\lstinline!(pvals)!} & Benjamini-Hochberg FDR-adjusted
p-values (step-up procedure). \\
\passthrough{\lstinline!quadmath.stats.statistics!} &
\passthrough{\lstinline!bootstrap\_ci!} & function &
\passthrough{\lstinline!(x, stat, iters=, seed=, alpha=)!} & Percentile
bootstrap confidence interval for \passthrough{\lstinline!stat!} on
\passthrough{\lstinline!x!}. \\
\passthrough{\lstinline!quadmath.stats.statistics!} &
\passthrough{\lstinline!cohens\_d!} & function &
\passthrough{\lstinline!(a, b)!} & Pooled-standard-deviation Cohen's d
between two samples. \\
\passthrough{\lstinline!quadmath.stats.statistics!} &
\passthrough{\lstinline!jackknife\_ci!} & function &
\passthrough{\lstinline!(x, stat, alpha)!} & Leave-one-out jackknife
interval and bias estimate for \passthrough{\lstinline!stat!}. \\
\passthrough{\lstinline!quadmath.stats.statistics!} &
\passthrough{\lstinline!p\_adjust\_bonferroni!} & function &
\passthrough{\lstinline!(pvals)!} & Bonferroni-adjusted p-values,
elementwise \passthrough{\lstinline!min(1, p * m)!}. \\
\passthrough{\lstinline!quadmath.stats.statistics!} &
\passthrough{\lstinline!permutation\_test!} & function &
\passthrough{\lstinline!(a, b, iters=, seed=, alternative=)!} & Pooled
permutation test on the difference of sample means. \\
\passthrough{\lstinline!quadmath.stats.statistics!} &
\passthrough{\lstinline!rotation\_stats!} & function &
\passthrough{\lstinline!(angles)!} & Circular statistics for angles
given in radians. \\
\passthrough{\lstinline!quadmath.stats.statistics!} &
\passthrough{\lstinline!scaling\_fit!} & function &
\passthrough{\lstinline!(sizes, times)!} & Power-law (log-log linear)
fit of runtimes against input sizes. \\
\passthrough{\lstinline!quadmath.stats.statistics!} &
\passthrough{\lstinline!summarize!} & function &
\passthrough{\lstinline!(x)!} & Descriptive summary of a sample. \\
\passthrough{\lstinline!quadmath.stats.statistics!} &
\passthrough{\lstinline!welch\_t\_test!} & function &
\passthrough{\lstinline!(a, b, alternative)!} & Welch's unequal-variance
two-sample t test. \\
\passthrough{\lstinline!quadmath.tools.atomic\_write!} &
\passthrough{\lstinline!atomic\_open!} & function &
\passthrough{\lstinline!(path, newline=)!} & Yield a text handle whose
contents replace \passthrough{\lstinline!path!} only if the block
succeeds. \\
\passthrough{\lstinline!quadmath.tools.atomic\_write!} &
\passthrough{\lstinline!atomic\_savez!} & function &
\passthrough{\lstinline!(**arrays)!} & Write
\passthrough{\lstinline!arrays!} as an \passthrough{\lstinline!.npz!}
archive at \passthrough{\lstinline!path!} atomically and
byte-reproducibly. \\
\passthrough{\lstinline!quadmath.tools.atomic\_write!} &
\passthrough{\lstinline!atomic\_write\_text!} & function &
\passthrough{\lstinline!(path, text, newline=)!} & Write
\passthrough{\lstinline!text!} to \passthrough{\lstinline!path!}
atomically (see :func:\passthrough{\lstinline!atomic\_open!}). \\
\passthrough{\lstinline!quadmath.tools.glossary\_gen!} &
\passthrough{\lstinline!ApiEntry!} & class & & \\
\passthrough{\lstinline!quadmath.tools.glossary\_gen!} &
\passthrough{\lstinline!build\_api\_index!} & function &
\passthrough{\lstinline!(src\_dir)!} & \\
\passthrough{\lstinline!quadmath.tools.glossary\_gen!} &
\passthrough{\lstinline!generate\_markdown\_table!} & function &
\passthrough{\lstinline!(entries)!} & \\
\passthrough{\lstinline!quadmath.tools.glossary\_gen!} &
\passthrough{\lstinline!inject\_between\_markers!} & function &
\passthrough{\lstinline!(markdown\_text, begin, end, payload)!} & \\
\passthrough{\lstinline!quadmath.validate.validate!} &
\passthrough{\lstinline!DEFAULT\_CHECKS!} & constant & & \\
\passthrough{\lstinline!quadmath.validate.validate!} &
\passthrough{\lstinline!DEFAULT\_TOLERANCE!} & constant & & \\
\passthrough{\lstinline!quadmath.validate.validate!} &
\passthrough{\lstinline!NOTES!} & constant & & \\
\passthrough{\lstinline!quadmath.validate.validate!} &
\passthrough{\lstinline!SKIPPED\_CHECKS!} & constant & & \\
\passthrough{\lstinline!quadmath.validate.validate!} &
\passthrough{\lstinline!ValidationReport!} & class & & Immutable outcome
of a single validation check. \\
\passthrough{\lstinline!quadmath.validate.validate!} &
\passthrough{\lstinline!check\_associativity!} & function &
\passthrough{\lstinline!(a, b, c, tol)!} & Verify (a\emph{b)}c ==
a\emph{(b}c) within \passthrough{\lstinline!tol!} via the core Hamilton
product. \\
\passthrough{\lstinline!quadmath.validate.validate!} &
\passthrough{\lstinline!check\_conjugate\_inverse!} & function &
\passthrough{\lstinline!(q, tol)!} & Verify q * conj(q) equals the
identity (1, 0, 0, 0) within \passthrough{\lstinline!tol!}. \\
\passthrough{\lstinline!quadmath.validate.validate!} &
\passthrough{\lstinline!check\_double\_cover!} & function &
\passthrough{\lstinline!(q1, q2, tol)!} & Verify the SO(3) homomorphism
R(q1*q2) == R(q1) R(q2) within \passthrough{\lstinline!tol!}. \\
\passthrough{\lstinline!quadmath.validate.validate!} &
\passthrough{\lstinline!check\_normalization!} & function &
\passthrough{\lstinline!(q, tol)!} & Verify the quaternion norm is
within \passthrough{\lstinline!tol!} of 1 (norm of (a, b, c, d)). \\
\passthrough{\lstinline!quadmath.validate.validate!} &
\passthrough{\lstinline!check\_slerp\_midpoint!} & function &
\passthrough{\lstinline!(q0, q1, tol)!} & Verify the shortest-arc slerp
midpoint lies on the geodesic of (q0, q1). \\
\passthrough{\lstinline!quadmath.validate.validate!} &
\passthrough{\lstinline!run\_validation!} & function &
\passthrough{\lstinline!(quaternions, checks)!} & Run checks over
\passthrough{\lstinline!quaternions!} and collect deterministic
reports. \\
\passthrough{\lstinline!quadmath.viz.\_common!} &
\passthrough{\lstinline!atomic\_target!} & function &
\passthrough{\lstinline!(path)!} & Yield a temporary sibling of
\passthrough{\lstinline!path!}; move it onto
\passthrough{\lstinline!path!} only on success. \\
\passthrough{\lstinline!quadmath.viz.\_common!} &
\passthrough{\lstinline!embedding\_array!} & function &
\passthrough{\lstinline!(embedding)!} & Return a 3x4 embedding as a
float array. \\
\passthrough{\lstinline!quadmath.viz.\_common!} &
\passthrough{\lstinline!encoded\_angle!} & function &
\passthrough{\lstinline!(quat)!} & Return the rotation magnitude a unit
quaternion \passthrough{\lstinline!(w, x, y, z)!} encodes. \\
\passthrough{\lstinline!quadmath.viz.\_common!} &
\passthrough{\lstinline!figure\_scope!} & function &
\passthrough{\lstinline!(*args, **kwargs)!} & Yield a new figure and
close it on exit, including when a save fails. \\
\passthrough{\lstinline!quadmath.viz.\_common!} &
\passthrough{\lstinline!mp4\_writer!} & function &
\passthrough{\lstinline!(fps)!} & Return an ffmpeg writer whose MP4
bytes repeat for identical frames. \\
\passthrough{\lstinline!quadmath.viz.\_common!} &
\passthrough{\lstinline!resolve\_output\_path!} & function &
\passthrough{\lstinline!(path, figure\_dir)!} & Return the location to
write \passthrough{\lstinline!path!} under the viz output policy. \\
\passthrough{\lstinline!quadmath.viz.\_common!} &
\passthrough{\lstinline!save\_figure!} & function &
\passthrough{\lstinline!(fig, path, **kwargs)!} & Write
\passthrough{\lstinline!fig!} to \passthrough{\lstinline!path!}
atomically; keyword arguments go to
\passthrough{\lstinline!savefig!}. \\
\passthrough{\lstinline!quadmath.viz.\_common!} &
\passthrough{\lstinline!set\_axes\_equal!} & function &
\passthrough{\lstinline!(ax)!} & Scale a 3D axes box to the data spans
so equal lengths match on every axis. \\
\passthrough{\lstinline!quadmath.viz.animations!} &
\passthrough{\lstinline!Frame!} & class & & A single animation frame. \\
\passthrough{\lstinline!quadmath.viz.animations!} &
\passthrough{\lstinline!GRID\_SIZE!} & constant & & \\
\passthrough{\lstinline!quadmath.viz.animations!} &
\passthrough{\lstinline!diffusion\_frames!} & function &
\passthrough{\lstinline!(n\_steps, seed)!} & Render explicit heat
diffusion on the IVM radius-3 ball adjacency. \\
\passthrough{\lstinline!quadmath.viz.animations!} &
\passthrough{\lstinline!frames\_strip!} & function &
\passthrough{\lstinline!(frames, out\_path, labels, save)!} & Render
frames as a single-row matplotlib strip, one panel per frame. \\
\passthrough{\lstinline!quadmath.viz.animations!} &
\passthrough{\lstinline!frames\_to\_gif!} & function &
\passthrough{\lstinline!(frames, out\_path, fps, scale)!} & Assemble
frames into an animated GIF at \passthrough{\lstinline!out\_path!}. \\
\passthrough{\lstinline!quadmath.viz.animations!} &
\passthrough{\lstinline!lattice\_frames!} & function &
\passthrough{\lstinline!(shells, n)!} & Render a pulsing IVM lattice
ball. \\
\passthrough{\lstinline!quadmath.viz.animations!} &
\passthrough{\lstinline!simplex\_frames!} & function &
\passthrough{\lstinline!(q0, q1, n)!} & Render a quaternion-slerp
rotation of the IVM radius-1 ball. \\
\passthrough{\lstinline!quadmath.viz.plots!} &
\passthrough{\lstinline!plot\_error\_histogram!} & function &
\passthrough{\lstinline!(errors, bins, save, out\_path)!} & Plot a
histogram of error values with a dashed vertical mean line. \\
\passthrough{\lstinline!quadmath.viz.plots!} &
\passthrough{\lstinline!plot\_lattice\_shell\_3d!} & function &
\passthrough{\lstinline!(k, save, out\_path)!} & Scatter the sites of
one IVM lattice shell in 3D (equal-aspect axes). \\
\passthrough{\lstinline!quadmath.viz.plots!} &
\passthrough{\lstinline!plot\_loss\_history!} & function &
\passthrough{\lstinline!(losses, save, out\_path)!} & Plot a training
loss sequence as a line with markers. \\
\passthrough{\lstinline!quadmath.viz.plots!} &
\passthrough{\lstinline!plot\_shell\_growth!} & function &
\passthrough{\lstinline!(k\_max, save, out\_path)!} & Plot IVM shell
cardinalities (cuboctahedral numbers) versus shell index. \\
\passthrough{\lstinline!quadmath.viz.plots!} &
\passthrough{\lstinline!plot\_slerp\_path!} & function &
\passthrough{\lstinline!(q0, q1, n\_frames, site, save, out\_path)!} &
Trace the 3D path of a fixed lattice site under shortest-arc slerp. \\
\passthrough{\lstinline!quadmath.viz.vis\_lattice!} &
\passthrough{\lstinline!DEFAULT\_PLANE!} & constant & & \\
\passthrough{\lstinline!quadmath.viz.vis\_lattice!} &
\passthrough{\lstinline!GALLERY\_FILES!} & constant & & \\
\passthrough{\lstinline!quadmath.viz.vis\_lattice!} &
\passthrough{\lstinline!dynamics\_strip!} & function &
\passthrough{\lstinline!(axs, trajectory, t\_indices, embedding=, cmap=, titles=)!}
& Render evolution snapshots of a trajectory as a strip of 3D panels. \\
\passthrough{\lstinline!quadmath.viz.vis\_lattice!} &
\passthrough{\lstinline!field\_slice!} & function &
\passthrough{\lstinline!(ax, field, sites, plane, q0=, cmap=, title=, colorbar=)!}
& Heatmap of a scalar IVM field restricted to a lattice plane. \\
\passthrough{\lstinline!quadmath.viz.vis\_lattice!} &
\passthrough{\lstinline!gallery!} & function &
\passthrough{\lstinline!(paths\_out\_dir, seed)!} & Compose the three
lattice-gallery figures deterministically. \\
\passthrough{\lstinline!quadmath.viz.vis\_lattice!} &
\passthrough{\lstinline!shell\_scatter!} & function &
\passthrough{\lstinline!(ax, sites, k, embedding=, color=, size=, axis\_hints=, title=)!}
& Scatter one IVM frequency shell in 3D with tetrahedral axis hints. \\
\passthrough{\lstinline!quadmath.viz.vis\_stats!} &
\passthrough{\lstinline!GALLERY\_FILES!} & constant & & \\
\passthrough{\lstinline!quadmath.viz.vis\_stats!} &
\passthrough{\lstinline!gallery!} & function &
\passthrough{\lstinline!(paths\_out\_dir, seed)!} & Compose the four
statistics-gallery figures deterministically. \\
\passthrough{\lstinline!quadmath.viz.vis\_stats!} &
\passthrough{\lstinline!plot\_ci\_bars!} & function &
\passthrough{\lstinline!(labels, means, lows, highs, ax, title=)!} &
Point estimates with confidence intervals
\passthrough{\lstinline![lows, highs]!} as error bars. \\
\passthrough{\lstinline!quadmath.viz.vis\_stats!} &
\passthrough{\lstinline!plot\_ecdf!} & function &
\passthrough{\lstinline!(values, ax, title=)!} & Empirical cumulative
distribution function as a sorted step plot. \\
\passthrough{\lstinline!quadmath.viz.vis\_stats!} &
\passthrough{\lstinline!plot\_latency\_hist!} & function &
\passthrough{\lstinline!(times, ax, bins=, title=)!} & Histogram of a
latency sample with a dashed vertical mean line. \\
\passthrough{\lstinline!quadmath.viz.vis\_stats!} &
\passthrough{\lstinline!plot\_scaling\_loglog!} & function &
\passthrough{\lstinline!(sizes, times, ax, title=)!} & Log-log scatter
of times versus sizes with the fitted power law. \\
\passthrough{\lstinline!quadmath.viz.visualize!} &
\passthrough{\lstinline!animate\_discrete\_path!} & function &
\passthrough{\lstinline!(path, embedding, save, out\_path)!} & Animate a
point moving along a discrete quadray path. \\
\passthrough{\lstinline!quadmath.viz.visualize!} &
\passthrough{\lstinline!animate\_simplex!} & function &
\passthrough{\lstinline!(vertices\_list, embedding, save, out\_path)!} &
Animate simplex evolution across iterations. \\
\passthrough{\lstinline!quadmath.viz.visualize!} &
\passthrough{\lstinline!plot\_ivm\_neighbors!} & function &
\passthrough{\lstinline!(embedding, save, out\_path)!} & Scatter the 12
IVM neighbor points in 3D. \\
\passthrough{\lstinline!quadmath.viz.visualize!} &
\passthrough{\lstinline!plot\_partition\_tetrahedron!} & function &
\passthrough{\lstinline!(mu, s, a, psi, embedding, save, out\_path)!} &
Plot the four-fold partition as a labeled tetrahedron in 3D. \\
\passthrough{\lstinline!quadmath.viz.visualize!} &
\passthrough{\lstinline!plot\_simplex\_trace!} & function &
\passthrough{\lstinline!(state, save, out\_path)!} & Plot per-iteration
diagnostics for Nelder--Mead. \\
\end{longtable}
}

\newpage

\section{Static IVM Field Learning}\label{static-ivm-field-learning}

\subsection{Overview}\label{overview-4}

This section develops machine learning on a static synergetic geometry:
scalar fields defined over the Isotropic Vector Matrix (IVM) lattice,
learned from sparse noisy observations by Laplacian-regularized
kernel-weighted least squares on the lattice graph. All methods live in
the \passthrough{\lstinline!ivm\_field.py!} module
(\passthrough{\lstinline!quadray.IVMField!},
\passthrough{\lstinline!quadray\_shell\_norm!},
\passthrough{\lstinline!shell\_sites!},
\passthrough{\lstinline!fit\_geometry!}) and reuse the quadray machinery
of \passthrough{\lstinline!quadray.py!} described in
\href{03_quadray_methods.md}{Quadray Methods}. The learner uses
\passthrough{\lstinline!numpy!} only --- no machine-learning frameworks
--- and is fully deterministic for a fixed observation set.

\subsection{The IVM Lattice in Quadray
Coordinates}\label{the-ivm-lattice-in-quadray-coordinates}

Quadray coordinates represent IVM close-packed sphere centers as
non-negative integer quadrays normalized so the minimum component is
zero (see \href{03_quadray_methods.md}{Quadray Methods}). Two facts
organize the lattice. First, normalization adds or subtracts
\((k,k,k,k)\), so the component sum is a projective invariant modulo 4.
Second, the IVM sites are exactly one of the four residue classes:

\begin{equation}
\label{eq:ivm-site-condition}
q \text{ is an IVM site} \quad\Longleftrightarrow\quad \textstyle\sum_i q_i \equiv 0 \pmod{4},
\end{equation}

implemented as \passthrough{\lstinline!is\_ivm\_site()!} in
\passthrough{\lstinline!ivm\_field.py!}. The remaining three cosets are
the octahedral (sum \(\equiv 2\)) and tetrahedral (sum \(\equiv 1, 3\))
voids of the packing --- lattice points of the ambient grid, but not
sphere centers.

The shell norm of a site is the \(L^1\) magnitude of its sum-zero
(Coxeter.4D hyperplane) representative. With \(s = \sum_i q_i\):

\begin{equation}
\label{eq:ivm-shell-norm}
N(q) \;=\; \sum_{i=1}^{4} \left|\, q_i - \tfrac{s}{4} \,\right| \;=\; 2k, \qquad k \in \mathbb{Z}_{\geq 0},
\end{equation}

computed by \passthrough{\lstinline!quadray\_shell\_norm()!}. Shell
\(k\) of the lattice is the set of sites with \(N(q) = 2k\);
\passthrough{\lstinline!shell\_sites(k)!} enumerates it by scanning the
bounding box \([0, 2k]^4\), filtering on the membership condition
\eqref{eq:ivm-site-condition} and the norm \eqref{eq:ivm-shell-norm},
and sorting lexicographically --- a deterministic site order. The shell
cardinalities are the cuboctahedral numbers:

\begin{equation}
\label{eq:ivm-shell-cardinality}
\bigl|\, \{ q : N(q) = 2k \} \,\bigr| \;=\; 10k^2 + 2, \qquad k \geq 1,
\end{equation}

giving the sequence 1, 12, 42, 92, 162, \ldots{} for
\(k = 0, 1, 2, 3, 4\) --- the center plus the cuboctahedral numbers,
with shell 1 the twelve-around-one vector equilibrium of
\href{03_quadray_methods.md}{Quadray Methods}.
\passthrough{\lstinline!shell\_cardinalities()!} counts shells of an
enumerated ball, and the test suite pins the sequence exactly.

\subsection{The Field Model}\label{the-field-model}

\passthrough{\lstinline!IVMField.lattice\_ball(radius)!} stores a scalar
field over the lattice ball \(B_R = \{ q : N(q) \leq 2R \}\) as a
one-dimensional \passthrough{\lstinline!numpy!} array keyed by the
deterministic site index of
\passthrough{\lstinline!shell\_ball\_sites()!} (shell-major,
lexicographic within shell), with a dictionary mapping each normalized
site to its array position. Two lattice-graph ingredients drive
learning. The adjacency structure uses the twelve IVM neighbor moves ---
the permutations of \((2,1,1,0)\) collected in
\passthrough{\lstinline!IVM\_NEIGHBOR\_STEPS!} --- so two sites in the
ball are graph-adjacent exactly when one step of close packing separates
them. The graph Laplacian \(L\) (method
\passthrough{\lstinline!IVMField.\_laplacian()!}) is the symmetric
difference operator on that adjacency:

\begin{equation}
\label{eq:ivm-laplacian}
(L f)_i \;=\; \deg(i)\, f_i \;-\; \sum_{j \sim i} f_j .
\end{equation}

\subsection{Learning: Laplacian-Regularized Kernel-Weighted Least
Squares}\label{learning-laplacian-regularized-kernel-weighted-least-squares}

Observations \(y_j\) arrive at a sampled subset \(\Omega\) of sites.
Graph distances \(d(\cdot,\cdot)\) are hop counts from a multi-source
BFS (\passthrough{\lstinline!\_multi\_source\_distances()!} in
\passthrough{\lstinline!ivm\_field.py!}), and the Gaussian kernel over
hops with width \(\tau\) (\passthrough{\lstinline!kernel\_width!}) is:

\begin{equation}
\label{eq:ivm-kernel}
K(d) \;=\; \exp\!\bigl( -(d/\tau)^2 \bigr), \qquad d(i,j) = \text{hop distance}.
\end{equation}

\passthrough{\lstinline!IVMField.learn()!} minimizes a two-regime
objective over the ball, with Laplacian regularization weight
\(\lambda \geq 0\):

\begin{equation}
\label{eq:ivm-objective}
\min_{f} \;\; \sum_{i \in \Omega} \bigl( f_i - y_i \bigr)^2
\;+\; \sum_{i \notin \Omega} c_i \bigl( f_i - t_i \bigr)^2
\;+\; \lambda \sum_{i \sim j} \bigl( f_i - f_j \bigr)^2 ,
\end{equation}

with three data-fidelity terms: observed sites are pinned to their data
with unit weight (no self-smoothing of real observations); unobserved
sites carry a kernel confidence weight and a Nadaraya--Watson kernel
target,

\begin{equation}
\label{eq:ivm-confidence}
c_i \;=\; \max\bigl( K\bigl(\min_{j \in \Omega} d(i,j)\bigr),\ \varepsilon \bigr),
\qquad
t_i \;=\; \frac{\sum_{j \in \Omega} K(d(i,j))\, y_j}{\sum_{j \in \Omega} K(d(i,j))},
\end{equation}

where \(\varepsilon\) --- the \passthrough{\lstinline!\_WEIGHT\_FLOOR!}
constant, \(10^{-12}\), in \passthrough{\lstinline!ivm\_field.py!} ---
is a numerical floor keeping the normal-equation matrix of
\eqref{eq:ivm-objective} positive definite on the connected lattice
ball. The normal equations of \eqref{eq:ivm-objective} are

\begin{equation}
\label{eq:ivm-normal-equations}
\bigl( C + \lambda L \bigr) f \;=\; C\, t ,
\end{equation}

with \(C = \mathrm{diag}(c_i)\) and \(t\) the target vector;
\passthrough{\lstinline!IVMField.learn()!} solves them with a dense
\passthrough{\lstinline!numpy.linalg.solve!}. The Laplacian term
propagates field structure from observed sites across the lattice graph,
and the confidence weighting degrades gracefully to pure Laplacian
propagation as observation density falls. Because observed rows carry
unit weight, the estimator interpolates exact data as \(\lambda \to 0\)
and denoises noisy data for moderate \(\lambda\) --- both behaviors are
verified numerically in
\passthrough{\lstinline!tests/test\_ivm\_field.py!} with fixed seeds.

Two structural facts justify the estimator. A linear field in the
embedded XYZ coordinates (the \passthrough{\lstinline!to\_xyz()!} image
of \passthrough{\lstinline!quadray.py!}) is harmonic on the IVM graph
--- the twelve neighbor moves sum to zero --- so the harmonic extension
implied by \eqref{eq:ivm-normal-equations} reproduces it exactly from
boundary data (asserted to \(10^{-5}\) in the test suite). And for noisy
smooth fields, the learned field tracks the kernel-weighted local mean,
which averages observation noise across graph neighborhoods:

\begin{equation}
\label{eq:ivm-mse}
\mathrm{MSE}(f) \;=\; \frac{1}{|B_R|} \sum_{i \in B_R} \bigl( f_i - f_i^{\text{truth}} \bigr)^2 ,
\end{equation}

the quantity returned by \passthrough{\lstinline!IVMField.score()!}.

\subsection{Recovering Tetrahedral Geometry from Noisy
Points}\label{recovering-tetrahedral-geometry-from-noisy-points}

\passthrough{\lstinline!fit\_geometry()!} recovers the orientation and
scale of a tetrahedron from noisy 3D point clouds, using the quadray
basis of \passthrough{\lstinline!quadray.py!}. The four canonical vertex
images are \(v_i = \texttt{to\_xyz}(e_i)\), where \(e_i\) are the unit
quadray axes mapped through the 3\$\times\$4 embedding matrix of
\passthrough{\lstinline!quadray.py!}. Given labeled observations
\(p^{(i)}_m \approx R\, v_i + o\) (noisy samples of vertex \(i\) under
rotation \(R\) and offset \(o\)), vertex centroids \(c_i\) are formed
and the linear map \(G\) is fit in closed form:

\begin{equation}
\label{eq:ivm-fit}
G \;=\; \arg\min_{M \in \mathbb{R}^{3\times 3}} \sum_{i=1}^{4}
\bigl\| M v_i - c_i \bigr\|^2 ,
\qquad
\text{scale} \;=\; \operatorname{sign}\bigl(\det G\bigr)\, \bigl|\det G\bigr|^{1/3}.
\end{equation}

The least-squares solve uses
\passthrough{\lstinline!numpy.linalg.lstsq!} on the transpose system, so
noisy multi-sample vertices average out; the residual reported by
\passthrough{\lstinline!TetrahedronFit.residual!} is the RMS
vertex-centroid mismatch. The test suite recovers a rotated, scaled
tetrahedron to \(10^{-6}\) from \(\sigma =
10^{-7}\) noise, verifies the signed scale for reflected (negative
determinant) fits, and checks that doubling the embedding halves the
recovered matrix --- orientation and scale separate cleanly.

\subsection{Demonstration}\label{demonstration}

The script
\passthrough{\lstinline!quadmath/scripts/ivm\_field\_demo.py!} (run with
\passthrough{\lstinline!MPLBACKEND=Agg!}, seed 12) renders the pipeline
on the radius-3 ball (147 lattice sites): the synthetic field (a
harmonic linear part plus a weak quadratic bowl), the noisy observations
at 74 of the 147 sites, and the learned field. The learned
reconstruction beats the raw observation noise (reconstruction MSE
0.0558 against observation-noise MSE 0.0720, printed by the demo):

\begin{figure}
\centering
\pandocbounded{\includegraphics[keepaspectratio,alt={Static IVM field learning on the radius-3 IVM lattice ball (147 sites, shown at their XYZ embedding from the quadray coordinates; axes in embedding units). Panel A: the synthetic ground-truth field --- a harmonic linear part plus a weak quadratic bowl --- with dot color giving field value on the shared color bar (right). Panel B: the noisy observations at 74 of the 147 sites (seed 12). Panel C: the learned field from Laplacian-regularized kernel-weighted least squares (Eqs. --); reconstruction MSE 0.0558 against observation-noise MSE 0.0720. Reproduce with quadmath/scripts/ivm\_field\_demo.py (MPLBACKEND=Agg, seed 12).}]{../output/figures/ivm_field_demo.png}}
\caption{Static IVM field learning on the radius-3 IVM lattice ball (147
sites, shown at their XYZ embedding from the quadray coordinates; axes
in embedding units). Panel A: the synthetic ground-truth field --- a
harmonic linear part plus a weak quadratic bowl --- with dot color
giving field value on the shared color bar (right). Panel B: the noisy
observations at 74 of the 147 sites (seed 12). Panel C: the learned
field from Laplacian-regularized kernel-weighted least squares (Eqs.
\eqref{eq:ivm-objective}--\eqref{eq:ivm-normal-equations});
reconstruction MSE 0.0558 against observation-noise MSE 0.0720.
Reproduce with
\protect\passthrough{\lstinline!quadmath/scripts/ivm\_field\_demo.py!}
(\protect\passthrough{\lstinline!MPLBACKEND=Agg!}, seed 12).}
\end{figure}

\subsection{Cross-References}\label{cross-references}

\begin{itemize}
\tightlist
\item
  Coordinate foundations and the twelve-around-one shell:
  \href{03_quadray_methods.md}{Quadray Methods}
\item
  Optimization methods on tetrahedral lattices:
  \href{04_optimization_in_4d.md}{Optimization in 4D}
\item
  Applications and generalizations: \href{05_extensions.md}{Extensions}
\end{itemize}

\newpage

\section{Dynamics and Learning on the IVM
Lattice}\label{dynamics-and-learning-on-the-ivm-lattice}

\subsection{Overview}\label{overview-5}

Where \href{04_optimization_in_4d.md}{Section 4} descends a \emph{single
point} along the IVM lattice, this section evolves a \emph{field} ---
one scalar per lattice site --- and then asks the inverse question:
given an observed field trajectory, what coupling produced it?
Everything here is implemented in
\passthrough{\lstinline!ivm\_dynamics.py!}
(\passthrough{\lstinline!ball\_sites!},
\passthrough{\lstinline!make\_lattice!},
\passthrough{\lstinline!heat\_step!},
\passthrough{\lstinline!majority\_step!},
\passthrough{\lstinline!step!}, \passthrough{\lstinline!simulate!},
\passthrough{\lstinline!fit\_trajectory!},
\passthrough{\lstinline!sum\_of\_squares!},
\passthrough{\lstinline!is\_nonincreasing!}) on top of the
\passthrough{\lstinline!Quadray!} class and
\passthrough{\lstinline!to\_xyz!} embedding in
\passthrough{\lstinline!quadray.py!}; the demo figure is produced by
\passthrough{\lstinline!quadmath/scripts/ivm\_dynamics\_demo.py!} (which
delegates to \passthrough{\lstinline!render\_dynamics\_demo!} in
\passthrough{\lstinline!src/quadmath/lattice/ivm\_dynamics.py!}, per the
thin-orchestrator contract in
\passthrough{\lstinline!quadmath/scripts/AGENTS.md!}).

\subsection{Lattice sites, shells, and the 12-around-one move
graph}\label{lattice-sites-shells-and-the-12-around-one-move-graph}

A \emph{site} is a canonical quadray representative --- non-negative
integer components with at least one zero, selected by
\passthrough{\lstinline!Quadray.normalize!} (see
\href{03_quadray_methods.md}{Quadray Methods}). The radial observable is
the squared embedding norm
\passthrough{\lstinline!site\_radius\_sq(q)!}; because the shift vector
\((1,1,1,1)\) lies in the kernel of the
\passthrough{\lstinline!DEFAULT\_EMBEDDING!} matrix, the radius is
invariant under quadray normalization.
\passthrough{\lstinline!ball\_sites(R)!} enumerates the canonical sites
with squared radius at most \(R^2\): for a canonical site the squared
radius is at least twice the square of its largest component, so
components in \([0, R]\) cover the ball.

\passthrough{\lstinline!make\_lattice(R)!} joins two sites when one of
the 12 canonical IVM neighbor shifts
(\passthrough{\lstinline!neighbor\_shifts!}, all permutations of
\((2,1,1,0)\), each at squared radius 8 --- the close-packing distance)
maps one site onto the other after renormalization. Normalization never
moves the embedded point, so adjacency is exactly ``embedding difference
is one of the 12 shifts''. Lattice sites are the quadrays whose
coordinate sum is divisible by 4; the remaining residue classes are IVM
voids and are excluded. On the \(R = 3\) ball (13 sites) the graph is a
single connected component: the close-packed \textbf{12-around-one
cluster}, i.e.~the origin (degree 12) and its twelve cuboctahedron
neighbors (degree 5). The update laws below act independently on each
connected component of the ball.

\subsection{Heat update and the monotonicity
lemma}\label{heat-update-and-the-monotonicity-lemma}

The heat (diffusion) update blends a field \(u_t \in \mathbb{R}^{N}\)
--- one value per lattice site, \(N\) the site count --- with its
symmetric-normalized adjacency average \(S = D^{-1/2} A D^{-1/2}\),
where \(A\) is the adjacency matrix of the move graph, \(D\) the
diagonal degree matrix, and rows of \(S\) are zero at isolated sites:

\begin{equation}
\label{eq:ivmdyn-heat}
u_{t+1} \;=\; (1-\alpha)\, u_t \;+\; \alpha\, S\, u_t ,
\qquad \alpha \in [0,1] .
\end{equation}

\textbf{Lemma (heat averaging is non-increasing in sum of squares).} For
every \(\alpha \in [0,1]\), every lattice built by
\passthrough{\lstinline!make\_lattice!}, and every state,

\begin{equation}
\label{eq:ivmdyn-lemma}
\lVert u_{t+1} \rVert_2^2 \;\le\; \lVert u_t \rVert_2^2 .
\end{equation}

\emph{Proof sketch.} \(S\) is symmetric with spectrum in \([-1,1]\) (it
is similar to the random-walk matrix \(D^{-1}A\)). The update operator
\((1-\alpha)I + \alpha S\) is therefore symmetric with eigenvalues in
\([1-2\alpha,\,1] \subseteq [-1,1]\), hence non-expansive in the 2-norm.
\(\square\)

This is exactly the claim the test suite asserts --- \emph{per step},
across seeds and the full coupling range
(\passthrough{\lstinline!test\_heat\_update\_l2\_nonincreasing\_per\_step!}
in \passthrough{\lstinline!tests/unit/lattice/test\_ivm\_dynamics.py!}).
Two scope restrictions are deliberate and tested:

\begin{enumerate}
\def\labelenumi{\arabic{enumi}.}
\tightlist
\item
  The guarantee is specific to the symmetric normalization \(S\). Plain
  row-average diffusion \(P = D^{-1}A\) on a truncated ball is
  \textbf{not} covered (boundary rows make \(P\) non-doubly-stochastic);
  the module neither uses nor claims it.
\item
  In the demo run (\(\alpha = 0.35\), \(T = 40\), seed 7) the sum of
  squares decreases monotonically from \(37.04\) to \(3.93\) while never
  increasing at any single step.
\end{enumerate}

\subsection{Majority update on integer
states}\label{majority-update-on-integer-states}

For integer-valued fields the rounded-averaging (``majority'') update is

\begin{equation}
\label{eq:ivmdyn-majority}
u_{t+1}(s) \;=\; \operatorname{rint}\!\Big( (1-\alpha)\, u_t(s)
\;+\; \alpha\, \tfrac{1}{\deg(s)} \!\!\sum_{n \in N(s)} u_t(n) \Big),
\end{equation}

with \(N(s)\) the neighbor set of site \(s\) and \(\deg(s) = |N(s)|\)
its degree, identity rows at isolated sites, and round-half-to-even ties
(\passthrough{\lstinline!numpy.rint!}), computed by
\passthrough{\lstinline!majority\_step!}. Every new value is a rounded
convex combination of the site and its neighbors, so the extremal
structure is preserved:

\textbf{Proposition (range containment).} Under
\eqref{eq:ivmdyn-majority} the field maximum is non-increasing and the
minimum non-decreasing, at every step, for every \(\alpha \in [0,1]\).

We deliberately make \textbf{no} sum-of-squares claim for this update:
rounding can increase it, and the test suite pins a demonstrating case
(origin at 0 with its twelve packers at 1 maps to thirteen sites at
value 1, raising the sum of squares from 12 to 13 ---
\passthrough{\lstinline!test\_majority\_step\_sum\_of\_squares\_can\_increase!}).
The honest summary: heat is \(L_2\)-monotone, majority is
range-preserving, and neither claim is stretched to cover the other
update.

\subsection{Learning the coupling from an observed
trajectory}\label{learning-the-coupling-from-an-observed-trajectory}

Given observed snapshots \(u_{\mathrm{obs}}(0..T)\) on a known lattice,
the coupling \(\alpha\) is identified gradient-free:
\passthrough{\lstinline!fit\_trajectory!} re-simulates the field from
the observed initial condition for each candidate coupling and minimizes
the trajectory mean squared error over the \(T\) post-initial snapshots
and the \(N\) sites, where \(u_{\alpha}(t)\) is the re-simulated field
under candidate coupling \(\alpha\):

\begin{equation}
\label{eq:ivmdyn-mse}
\mathcal{L}(\alpha) \;=\; \frac{1}{T\,N} \sum_{t=1}^{T}
\big\lVert u_{\alpha}(t) - u_{\mathrm{obs}}(t) \big\rVert_2^2 ,
\end{equation}

scanning the supplied grid, then re-gridding the interval between the
best candidate's grid neighbors \passthrough{\lstinline!refine\_rounds!}
times (grid/coordinate search; ties resolve to the smallest \(\alpha\);
fully deterministic). The result is a global optimum \emph{over the
evaluated candidates only} --- no convergence beyond the refinement
resolution is claimed. In the demo configuration (synthetic trajectory
generated at \(\alpha_{\text{true}} = 0.3\) on the \(R=3\) lattice,
\(T = 40\)), a 21-point grid plus refinement (33 MSE evaluations)
recovers \(\alpha \approx 0.300\) with
\(\mathcal{L} \approx 4.5 \times 10^{-33}\) --- machine precision for a
bitwise-deterministic re-simulation. When the true coupling lies
off-grid, refinement converges to within a grid-spacing fraction of it
(\passthrough{\lstinline!test\_fit\_refinement\_finds\_offgrid\_coupling!}:
0.37 recovered to \(\le 0.0125\) from a coarse
\(\{0, 0.25, 0.5, 0.75, 1\}\) grid).

\subsection{Demo figure}\label{demo-figure}

Generated by
\passthrough{\lstinline!uv run python quadmath/scripts/ivm\_dynamics\_demo.py!}
(\passthrough{\lstinline!MPLBACKEND=Agg!}); all randomness is seeded
(\passthrough{\lstinline!DynamicsParams.seed!}), so the figure is
reproducible byte-for-byte up to PNG encoding.

\begin{figure}
\centering
\pandocbounded{\includegraphics[keepaspectratio,alt={Dynamics and learning on the IVM lattice (R=3 ball, 13 sites). Top row: heat-field snapshots at t = 0, 20, 40 (\textbackslash alpha = 0.35, seed 7) over the embedded quadray sites (XYZ axes in embedding units); dot color encodes field value on a shared symmetric scale, and diffusion homogenizes the field by t = 40. Bottom left: sum-of-squares histories over t = 0 \textbackslash ldots 40 --- the heat curve decreases monotonically (lemma ); the majority curve falls steeply and plateaus, but unlike heat it carries no monotonicity guarantee, since rounding can increase the sum of squares (test\_majority\_step\_sum\_of\_squares\_can\_increase). Bottom middle: trajectory MSE versus coupling \textbackslash alpha on a logarithmic axis over the search grid, with the identified optimum (fit 0.3000) coinciding with the true generating \textbackslash alpha = 0.3 (dotted). Bottom right: the final integer majority field at t = 40. Generated by uv run python quadmath/scripts/ivm\_dynamics\_demo.py (MPLBACKEND=Agg, seeded via DynamicsParams.seed).}]{../output/figures/ivm_dynamics_demo.png}}
\caption{\textbf{Dynamics and learning on the IVM lattice (\(R=3\) ball,
13 sites)}. Top row: heat-field snapshots at \(t = 0, 20, 40\)
(\(\alpha = 0.35\), seed 7) over the embedded quadray sites (XYZ axes in
embedding units); dot color encodes field value on a shared symmetric
scale, and diffusion homogenizes the field by \(t = 40\). Bottom left:
sum-of-squares histories over \(t = 0 \ldots 40\) --- the heat curve
decreases monotonically (lemma \eqref{eq:ivmdyn-lemma}); the majority
curve falls steeply and plateaus, but unlike heat it carries no
monotonicity guarantee, since rounding can increase the sum of squares
(\protect\passthrough{\lstinline!test\_majority\_step\_sum\_of\_squares\_can\_increase!}).
Bottom middle: trajectory MSE versus coupling \(\alpha\) on a
logarithmic axis over the search grid, with the identified optimum
(\protect\passthrough{\lstinline!fit 0.3000!}) coinciding with the true
generating \(\alpha = 0.3\) (dotted). Bottom right: the final integer
majority field at \(t = 40\). Generated by
\protect\passthrough{\lstinline!uv run python quadmath/scripts/ivm\_dynamics\_demo.py!}
(\protect\passthrough{\lstinline!MPLBACKEND=Agg!}, seeded via
\protect\passthrough{\lstinline!DynamicsParams.seed!}).}
\end{figure}

\subsection{Reproducibility and test
contract}\label{reproducibility-and-test-contract}

\begin{itemize}
\tightlist
\item
  Tests:
  \passthrough{\lstinline!PYTEST\_DISABLE\_PLUGIN\_AUTOLOAD=1 uv run coverage run -m pytest tests/unit/lattice/test\_ivm\_dynamics.py -q!}
  then \passthrough{\lstinline!uv run coverage report!} --- 42 tests,
  100\% statement and branch coverage of
  \passthrough{\lstinline!src/quadmath/lattice/ivm\_dynamics.py!}, no
  mocks, all examples real numerics with fixed seeds.
\item
  Markdown:
  \passthrough{\lstinline!uv run python quadmath/scripts/validate\_markdown.py!}.
\item
  The two provable claims (lemma \eqref{eq:ivmdyn-lemma}, the range
  proposition under \eqref{eq:ivmdyn-majority}) are asserted per step;
  the non-claims (row-average diffusion, majority \(L_2\) behavior) are
  pinned by the honesty tests described above rather than left implicit.
\end{itemize}

\subsection{Cross-references}\label{cross-references-1}

\begin{itemize}
\tightlist
\item
  Site geometry and normalization: \href{03_quadray_methods.md}{Quadray
  Methods}.
\item
  Discrete point descent on the same move set:
  \href{04_optimization_in_4d.md}{Optimization in 4D}.
\item
  Field/energy vocabulary (free energy, Fisher information):
  \href{09_free_energy_active_inference.md}{Appendix B}; equation index:
  \href{08_equations_appendix.md}{Appendix A}.
\end{itemize}

\newpage

\section{Lattice Tooling: Omnidirectional Numbering and Nearest-Site
Search}\label{lattice-tooling-omnidirectional-numbering-and-nearest-site-search}

\subsection{Overview}\label{overview-6}

Sections \href{02_4d_namespaces.md}{2} and
\href{03_quadray_methods.md}{3} develop the Quadray/IVM framework
analytically; this section documents the computational tooling that
operates directly on the close-packed lattice. Two modules extend the
analytical core without touching it:

\begin{itemize}
\tightlist
\item
  \passthrough{\lstinline!omni\_numbering!} --- omnidirectional
  close-packing numbering: vectorized enumeration of the IVM shell
  sequence, cumulative counts, and bidirectional site/index mappings.
\item
  \passthrough{\lstinline!lattice\_search!} --- fast
  nearest-lattice-point queries:
  \passthrough{\lstinline!nearest(site, R, k)!} and
  \passthrough{\lstinline!within\_radius(site, R)!} built on a
  precomputed ball index, entirely in NumPy (no scipy).
\end{itemize}

Both modules operate on canonical quadray integer 4-tuples (non-negative
components, at least one zero --- the projective normalization of
\href{03_quadray_methods.md}{Section 3}), and both rely on the same
verified shell invariants.

\subsection{Omnidirectional Close Packing and Frequency
Shells}\label{omnidirectional-close-packing-and-frequency-shells}

In the isotropic vector matrix (IVM), equal spheres pack
omnidirectionally --- closest packing. Starting from a central sphere,
the packing builds up in consecutive \textbf{frequency shells}: the
shell of frequency \(k \in \mathbb{Z}_{\ge 0}\) sits \(k\) radius-length
increments outward of the center, the same four-dimensional buildup
language synergetics uses for the growing vector equilibrium. Writing
\(N_k\) for the population of shell \(k\), the lone center is the whole
of shell 0 (\(N_0 = 1\)), and every shell with \(k \ge 1\) carries
exactly

\begin{equation}\label{eq:lattice-shell-population}
N_k \;=\; 10\,k^{2} \;+\; 2
\end{equation}

sphere centers. The first shell --- the 12 permutations of
\((2, 1, 1, 0)\) --- is the cuboctahedron (vector equilibrium) of the
``twelve around one'' motif; each subsequent shell is the next layer of
the omnidirectional buildup. The cumulative count through frequency
\(k\), written \(C_k\), is the centered-cuboctahedral closed form

\begin{equation}\label{eq:lattice-cumulative}
C_k \;=\; 1 \;+\; \sum_{j=1}^{k} \left(10\,j^{2} + 2\right) \;=\; 1 + 2k + \frac{10\,k\,(k+1)\,(2k+1)}{6}.
\end{equation}

The functions \passthrough{\lstinline!shell\_count(k)!} and
\passthrough{\lstinline!cumulative\_count(k)!} evaluate Eqs.
\eqref{eq:lattice-shell-population} and \eqref{eq:lattice-cumulative}
(scalar or array input), and the test suite cross-validates them against
direct summation.

\subsubsection{Reachability is enumerative, not
congruence-based}\label{reachability-is-enumerative-not-congruence-based}

A normalized integer quadray \(q\) with \(\sum_i q_i \equiv 0 \pmod 4\)
is a candidate IVM point, but not every candidate is reachable from the
origin by the 12 neighbor moves. The module therefore builds the site
set \textbf{layer by layer}: each new shell is produced at once by
adding all 12 moves to the previous frontier, re-normalizing rows, and
removing sites already seen (packed-key membership tests). The result is
cached per depth: \passthrough{\lstinline!sites\_through\_shell(k)!}
returns the \passthrough{\lstinline!(C\_k, 4)!} integer array in
canonical order --- the center first, then each shell in lexicographic
\passthrough{\lstinline!(a, b, c, d)!} order --- and
\passthrough{\lstinline!generate\_shell(k)!} returns one shell. An
independent breadth-first reference over
\passthrough{\lstinline!itertools.permutations!} confirms exact
agreement through frequency 8, including the counts of Eq.
\eqref{eq:lattice-shell-population} on every shell.

\subsubsection{Site/index mapping}\label{siteindex-mapping}

\passthrough{\lstinline!site\_index(site, max\_shell)!} returns the
position of a site in that canonical enumeration (or
\passthrough{\lstinline!-1!} if absent within the depth), and
\passthrough{\lstinline!site\_at\_index(index, max\_shell)!} is its
inverse. Lookups accept any projective representative --- the input is
translated by \passthrough{\lstinline!-(k, k, k, k)!} before matching
--- so \passthrough{\lstinline!(2, 1, 1, 0)!} and
\passthrough{\lstinline!(3, 2, 2, 1)!} map to the same index, while
\passthrough{\lstinline!(0, 1, 1, 2)!}, a different permutation and
hence a different sphere, maps elsewhere. Component overflow is guarded:
inputs whose normalized components reach the 16-bit packed-key width
(\(2^{16}\), the per-component field of the int64 keys) return
\passthrough{\lstinline!-1!} rather than wrapping.

Executable check of the bookends:

\begin{lstlisting}[language=Python]
from quadmath.lattice.omni_numbering import (
    shell_count, cumulative_count, sites_through_shell,
    site_index, site_at_index,
)

assert shell_count(1) == 12 and shell_count(2) == 42      # 10k^2 + 2
assert cumulative_count(2) == 55                          # 1 + 12 + 42
assert site_index(site_at_index(12, 3), 3) == 12          # roundtrip
assert site_index((2, 0, 0, 0), 6) == -1                  # unreachable site
\end{lstlisting}

\subsection{Fast Nearest-Site Queries}\label{fast-nearest-site-queries}

\subsubsection{Exact lattice distances}\label{exact-lattice-distances}

Under \passthrough{\lstinline!quadray.DEFAULT\_EMBEDDING!} the four
embedding columns \(C_j\) satisfy \(C_i \cdot C_j = 4\,\delta_{ij} - 1\)
at unit scale, so the Euclidean squared distance between integer
4-vectors \(p, q\) is the exact integer

\begin{equation}\label{eq:lattice-distance-identity}
d^{2}(p, q) \;=\; 4 \sum_{j} \delta_j^{2} \;-\; \Bigl(\sum_{j} \delta_j\Bigr)^{2},
\qquad \delta = p - q,
\end{equation}

implemented in \passthrough{\lstinline!squared\_distance!} and
cross-validated in the tests against direct embedding coordinates.
Because the identity is exact on integers, ranking and filtering never
suffer floating-point boundary errors: a site is inside the query radius
\(R > 0\) iff \(d^{2} \le R^{2}\) evaluated exactly (in float64 the same
algebra is exact for the magnitudes involved here).

\subsubsection{Truncation bound}\label{truncation-bound}

Every site \(s\) on frequency shell \(g\) satisfies
\(d^{2}(\mathbf{0}, s) \ge 8\,g\) --- verified by exhaustive enumeration
through frequency 8 in the test suite --- and the shell maximum is
exactly \(8\,g^{2}\). Combining this invariant with the triangle
inequality gives the depth at which a query can stop:

\begin{equation}\label{eq:lattice-truncation-bound}
g \;\leq\; \frac{\left(\lVert c \rVert + R\right)^{2}}{8},
\end{equation}

for a query center \(c\) and radius \(R > 0\), where \(\lVert c \rVert\)
is the Euclidean norm of the embedded center: any site within \(R\) of
\(c\) must live on a shell at most this deep.
\passthrough{\lstinline!within\_radius(site, R)!} therefore enumerates
exactly through \(\lfloor (\lVert c \rVert + R)^{2} / 8 \rfloor + 1\)
shells, filters with Eq. \eqref{eq:lattice-distance-identity}, and
returns all hits sorted by squared distance with lexicographic
\passthrough{\lstinline!(a, b, c, d)!} tie-breaking. Requests whose
required depth exceeds the precomputed index depth
(\passthrough{\lstinline!MAX\_SHELL = 32!}, about \(1.1 \times 10^{5}\)
sites --- \(C_{32} = 114465\) by Eq. \eqref{eq:lattice-cumulative})
raise \passthrough{\lstinline!ValueError!} rather than silently
truncating.

\subsubsection{\texorpdfstring{Shell-sweep
\texttt{nearest}}{Shell-sweep nearest}}\label{shell-sweep-nearest}

\passthrough{\lstinline!nearest(site, R, k)!} sweeps shells outward,
ranking each shell's squared distances with
\passthrough{\lstinline!numpy.argsort!}; once
\passthrough{\lstinline!k!} candidates are in hand, the invariant behind
Eq. \eqref{eq:lattice-truncation-bound} supplies an early-stop test ---
no unvisited shell can contain a closer site: a site at depth \(g + 1\)
satisfies \(d^{2}(\mathbf{0}, s) \ge 8\,(g+1)\), so the sweep stops once
\(8\,(g+1) > (\sqrt{d^{2}_{(k)}} + \lVert c \rVert)^{2}\), with
\(d^{2}_{(k)}\) the current k-th best squared distance. The sweep never
under-enumerates (the bound is conservative; ties at the boundary are
kept), and answers always match the brute-force reference over the same
region, as the tests assert on fixed-seed random queries.

\begin{lstlisting}[language=Python]
import numpy as np
from quadmath.lattice.lattice_search import nearest, within_radius

sites, d2 = nearest((0.4, -0.3, 0.9, 0.1), R=2.0, k=5)   # 5 nearest centers
ball, ball_d2 = within_radius((2, 1, 1, 0), R=4.0)       # everything within 4
assert np.all(np.diff(ball_d2) >= 0)                      # ascending distances
\end{lstlisting}

\subsection{API Summary}\label{api-summary}

{\def\LTcaptype{none} % do not increment counter
\begin{longtable}[]{@{}
  >{\raggedright\arraybackslash}p{(\linewidth - 2\tabcolsep) * \real{0.5000}}
  >{\raggedright\arraybackslash}p{(\linewidth - 2\tabcolsep) * \real{0.5000}}@{}}
\toprule\noalign{}
\begin{minipage}[b]{\linewidth}\raggedright
Function
\end{minipage} & \begin{minipage}[b]{\linewidth}\raggedright
Purpose
\end{minipage} \\
\midrule\noalign{}
\endhead
\bottomrule\noalign{}
\endlastfoot
\passthrough{\lstinline!omni\_numbering.shell\_count(k)!} & Shell
population, Eq. \eqref{eq:lattice-shell-population} \\
\passthrough{\lstinline!omni\_numbering.cumulative\_count(k)!} &
Cumulative count through shell \passthrough{\lstinline!k!}, Eq.
\eqref{eq:lattice-cumulative} \\
\passthrough{\lstinline!omni\_numbering.generate\_shell(k)!} & Sites of
one frequency shell (lexicographic) \\
\passthrough{\lstinline!omni\_numbering.sites\_through\_shell(k)!} &
Whole enumeration \passthrough{\lstinline!(C\_k, 4)!} in canonical
order \\
\passthrough{\lstinline!omni\_numbering.site\_index!} /
\passthrough{\lstinline!site\_at\_index!} & Bidirectional site/index
mapping with projective normalization \\
\passthrough{\lstinline!lattice\_search.squared\_distance!} & Exact
squared distances via Eq. \eqref{eq:lattice-distance-identity} \\
\passthrough{\lstinline!lattice\_search.within\_radius(site, R)!} & All
sites within radius, exact and sorted \\
\passthrough{\lstinline!lattice\_search.nearest(site, R, k)!} &
\passthrough{\lstinline!k!} nearest sites within radius, shell-sweep
with early stop \\
\end{longtable}
}

\subsection{Verification}\label{verification}

Both modules ship with 100\% test coverage (branch coverage included):
the enumeration is checked against an independent breadth-first
reference through frequency 8, distance identities against direct
embedding computations, and query results against brute-force filtering
on fixed-seed random centers. Run:

\begin{lstlisting}[language=bash]
PYTEST_DISABLE_PLUGIN_AUTOLOAD=1 uv run coverage run -m pytest -q
uv run coverage report
\end{lstlisting}

\newpage

\section{Conversions \& Specification}\label{conversions-specification}

\subsection{Overview}\label{overview-7}

Sections \href{02_4d_namespaces.md}{2}, \href{03_quadray_methods.md}{3}
and \href{13_lattice_tooling.md}{13} establish the Quadray/IVM framework
and its computational tooling; this section documents the conversion
layer that joins its two coordinate namespaces ---
\passthrough{\lstinline!conversions.py!}, the module that maps between
Fuller.4D quadray coordinates and Coxeter.4D Cartesian coordinates.
Three concerns live here: the embedding itself and its projective fiber,
the Gram identities that make lattice distances exact, and an
\textbf{exact rational inversion} that answers ``which lattice point is
this point?'' with \passthrough{\lstinline!fractions.Fraction!}
arithmetic instead of floating-point rounding.

The full mathematical specification of this layer is the repository-root
file \passthrough{\lstinline!SPEC.md!}; the
\passthrough{\lstinline!tests/test\_spec\_examples.py!} pins its worked
examples numerically. This section is the manuscript projection of that
specification: same definitions, same error taxonomy, same guarantees,
prose instead of source. The analytical foundations it draws on are
\href{03_quadray_methods.md}{Quadray Methods} (normalization, the
\passthrough{\lstinline!Quadray!} class), the shell machinery of
\href{11_ivm_field_learning.md}{IVM Field Learning}, and the
nearest-site search of \href{13_lattice_tooling.md}{Lattice Tooling},
whose exact-distance identity this section re-derives from the
embedding's Gram matrix.

\subsection{The embedding (Fuller.4D to
Coxeter.4D)}\label{the-embedding-fuller.4d-to-coxeter.4d}

The default embedding maps the four quadray axes \((A, B, C, D)\) to the
vertices of a regular tetrahedron in \(\mathbb{R}^3\). With scale factor
\(s\) (uniform; scales all resulting coordinates):

\begin{equation}
\label{eq:conv-embedding}
M \;=\; s \times
\begin{pmatrix}
 1 & -1 & -1 &  1 \\
 1 &  1 & -1 & -1 \\
 1 & -1 &  1 & -1
\end{pmatrix},
\qquad \mathrm{xyz}(q) \;=\; M\,q .
\end{equation}

\passthrough{\lstinline!urner\_embedding(scale=1.0)!} builds exactly
this matrix (each row scaled by \passthrough{\lstinline!scale!}), and
\passthrough{\lstinline!quadray.DEFAULT\_EMBEDDING!} is the same
unscaled matrix stored as a tuple-of-rows; the two agree entry for entry
at scale 1. Every row of \(M\) sums to zero:

\begin{equation}
\label{eq:conv-fiber}
\textstyle\sum_{j} M_{ij} \;=\; 0 \quad \forall\, i,
\qquad \text{hence} \qquad
q \;\sim\; q + t\,\mathbf{1}, \qquad t \in \mathbb{Z},
\end{equation}

where \(\mathbf{1} = (1, 1, 1, 1)\). The vector \(\mathbf{1}\) spans the
kernel of \(M\), so the map is projective:
\passthrough{\lstinline!to\_xyz(q) == to\_xyz(q + t*(1, 1, 1, 1))!} for
every integer \(t\), and quadray normalization (see
\href{03_quadray_methods.md}{Quadray Methods}) simply selects the min-0
representative of that equivalence class. The forward map
\passthrough{\lstinline!quadray\_to\_xyz(q, M=None)!} delegates to
\passthrough{\lstinline!quadray.to\_xyz!}; passing
\passthrough{\lstinline!M=None!} selects
\passthrough{\lstinline!quadray.DEFAULT\_EMBEDDING!}, and an explicit
\passthrough{\lstinline!M!} is shape-validated against the shared
internal helper \passthrough{\lstinline!\_embedding\_rows!} (a
\passthrough{\lstinline!(3, 4)!} array is required,
\passthrough{\lstinline!ValueError!} otherwise) before the product is
taken. Existing callers that pass an explicit embedding --- for example
\passthrough{\lstinline!quadmath/scripts/quadray\_clouds.py!} --- are
unaffected.

\subsection{Gram identities and exact
distances}\label{gram-identities-and-exact-distances}

The embedding matrix satisfies two exact Gram identities. Row-wise, at
scale \(s\):

\begin{equation}
\label{eq:conv-gram-row}
M\,M^{\mathsf{T}} \;=\; 4\,s^{2}\,I_{3},
\end{equation}

and column-wise, with \(C_j\) the four columns of \(M\) and \(J_4\) the
all-ones matrix:

\begin{equation}
\label{eq:conv-gram-col}
M^{\mathsf{T}}M \;=\; 4\,I_{4} \;-\; J_{4},
\qquad \text{i.e.} \qquad
C_i \cdot C_j \;=\; 4\,\delta_{ij} \;-\; 1 \;\; (s = 1).
\end{equation}

\passthrough{\lstinline!embedding\_basis(M=None)!} returns the columns
of the validated embedding as a \passthrough{\lstinline!(4, 3)!} array
\(B\) whose \emph{rows} are the \(C_j\); at scale 1 the identity
\eqref{eq:conv-gram-col} reads \(B\,B^{\mathsf{T}} = 4\,I_4 - J_4\)
exactly (diagonal 3, off-diagonal \(-1\)), and it scales by \(s^2\) for
\passthrough{\lstinline!urner\_embedding(scale)!}. This is the identity
that \passthrough{\lstinline!lattice\_search.squared\_distance!} relies
on: for integer quadrays \(p, q\) with \(\delta = p - q\), the Euclidean
squared distance in embedded coordinates is the exact integer

\begin{equation}
\label{eq:conv-distance}
d^{2}(p, q) \;=\; \sum_{i} (M\delta)_i^{2}
\;=\; 4 \sum_{j} \delta_j^{2} \;-\; \Bigl(\sum_{j} \delta_j\Bigr)^{2},
\end{equation}

the same identity documented as Eq. \eqref{eq:lattice-distance-identity}
in \href{13_lattice_tooling.md}{Lattice Tooling}, where it powers
ranking and filtering; the test suite cross-validates it against float
embeddings and \passthrough{\lstinline!quadray.distance!}. Because the
identity is exact on integers, nearest-site decisions never suffer
floating-point boundary errors.

Executable check of the two identities:

\begin{lstlisting}[language=Python]
import numpy as np
from quadmath.lattice.conversions import urner_embedding, embedding_basis
from quadmath.lattice.lattice_search import squared_distance

M = urner_embedding()
B = embedding_basis()                              # (4, 3): rows are the columns of M
assert np.allclose(M @ M.T, 4.0 * np.eye(3))       # Eq. (eq:conv-gram-row)
assert np.allclose(B @ B.T, 4.0 * np.eye(4) - np.ones((4, 4)))
delta = np.array([2, 1, 1, 0])                     # shell-1 displacement
assert squared_distance(delta, np.zeros((1, 4), dtype=np.int64)) == 8   # Eq. (eq:conv-distance)
\end{lstlisting}

\subsection{Exact canonical inversion}\label{exact-canonical-inversion}

The forward map is projective but not injective: by
\eqref{eq:conv-fiber} the fiber of an embedded point is the full coset
\(\{\,q + t\,\mathbf{1} : t \in \mathbb{Z}\,\}\), and an inverse must
select one representative from it.

\begin{equation}
\label{eq:conv-canonical}
y \;=\; (x_1, x_2, x_3, 0), \qquad
\begin{pmatrix} C_0 & C_1 & C_2 \end{pmatrix}
\begin{pmatrix} x_1 \\ x_2 \\ x_3 \end{pmatrix} \;=\; \mathrm{xyz},
\qquad q^\ast \;=\; \mathrm{normalize}(y).
\end{equation}

\passthrough{\lstinline!xyz\_to\_quadray\_canonical(xyz, M=None)!}
resolves that ambiguity deterministically by exact rational arithmetic
(\passthrough{\lstinline!fractions.Fraction!}; a Python float enters as
its exact binary rational --- no rounding, no tolerance). The algorithm,
in prose: (1) convert \passthrough{\lstinline!xyz!} and
\passthrough{\lstinline!M!} to exact \passthrough{\lstinline!Fraction!}
entries and require every row of \passthrough{\lstinline!M!} to sum to
exactly zero, so that normalization preserves the image point
\eqref{eq:conv-fiber}; (2) because \(C_3 = -(C_0 + C_1 + C_2)\), rank
\(M = 3\) is equivalent to \(\det\,[\,C_0\; C_1\; C_2\,] \neq 0\) (a
\passthrough{\lstinline!ValueError!} otherwise); (3) solve the
\(3 \times 3\) system on columns 0--2 by Cramer's rule for the
particular preimage \(y = (x_1, x_2, x_3, 0)\) of
\eqref{eq:conv-canonical}; (4) the embedded point lies in the image of
the integer lattice \(\{\,M\,q : q \in \mathbb{Z}^{4}\,\}\) if and only
if \(x_1, x_2, x_3\) are all integers --- then every integer preimage is
\(y + t\,\mathbf{1}\), and the min-0 representative
\passthrough{\lstinline!Quadray(x1, x2, x3, 0).normalize()!} is the
unique deterministic tie-break on the fiber \eqref{eq:conv-fiber} (the
same rule as \passthrough{\lstinline!Quadray.normalize!}); a
non-integral preimage raises \passthrough{\lstinline!ValueError!} (the
point is not in the image).

The taxonomy is fail-closed. \passthrough{\lstinline!ValueError!}
covers: \passthrough{\lstinline!xyz!} of length other than 3;
\passthrough{\lstinline!M!} of the wrong shape; any row sum different
from zero; \(\det = 0\) (rank below 3); a non-integral preimage.
\passthrough{\lstinline!TypeError!} covers entries that are not
\passthrough{\lstinline!int!}/\passthrough{\lstinline!float!}/\passthrough{\lstinline!Fraction!}/\passthrough{\lstinline!numbers.Integral!}
--- note that \passthrough{\lstinline!numpy.float64!} (a
\passthrough{\lstinline!float!} subclass) and
\passthrough{\lstinline!numpy.int64!} (integral) are accepted, while
\passthrough{\lstinline!numpy.float32!} is rejected. There is no
floating-point fuzz anywhere: an off-lattice point raises rather than
snapping to a neighbor.

The float-valued inverse
\passthrough{\lstinline!quadray.quadray\_from\_xyz!} is a different,
deliberately fuzzier contract: it applies the pseudoinverse
\(M^{+} = M^{\mathsf{T}}(MM^{\mathsf{T}})^{-1}\) and rounds each
component half-up (\passthrough{\lstinline!floor(v + 0.5)!}, not
banker's rounding --- its docstring argument shows why half-up provably
stays in the \(\mathbf{1}\)-coset of the exact preimage, where
per-component round-half-even can leave it). For points on the lattice
it round-trips exactly; for general \(\mathbb{R}^3\) points it returns
the nearest lattice point in quadray coordinates (component-wise --- not
always nearest in embedded XYZ distance), which is a snapping operation,
not an inverse. The component-wise rule is the chosen contract; an
XYZ-nearest variant would be a separate function that searches
neighboring classes.
\passthrough{\lstinline!quadray\_roundtrip(q, M=None)!} pins the
round-trip contract

\begin{equation}
\label{eq:conv-roundtrip}
\mathrm{from\_xyz}\bigl(\mathrm{to\_xyz}(q)\bigr) \;=\; q ,
\end{equation}

and raises \passthrough{\lstinline!AssertionError!} explicitly (not a
bare \passthrough{\lstinline!assert!}, so the check survives
\passthrough{\lstinline!python -O!}) when it fails. The contract is
exact for normalized integer quadrays and, because
\passthrough{\lstinline!quadray\_roundtrip!} threads the \emph{same}
embedding through both legs, it is \textbf{scale-robust}: for
\(M = c\,M_0\) with any \(c \neq 0\) the pseudoinverse projection is
scale-independent,

\begin{equation}
\label{eq:conv-roundtrip-scale}
\mathrm{pinv}(cM)\,(cM\,q) \;=\; q \;-\; \frac{\sum_i q_i}{4}\,\mathbf{1},
\end{equation}

since \(\mathrm{pinv}(cM) = (1/c)\,\mathrm{pinv}(M)\) and the row space
of \(cM_0\) is that of \(M_0\) (the \((1/c)\cdot c\) cancels). A shell
site embedded at scale \(1/2\) --- the half-integer FCC picture of the
same sites --- round-trips through the same scaled embedding. The single
documented \passthrough{\lstinline!AssertionError!} case is an
\textbf{unnormalized} input: the inverse canonicalizes to the min-0
representative, e.g.~\passthrough{\lstinline!(2, 2, 2, 1)!} comes back
as \passthrough{\lstinline!(1, 1, 1, 0)!}. Embeddings whose rows do not
sum to zero lie outside the Urner family of \eqref{eq:conv-embedding}
and outside this guarantee.

Executable check of exact recovery and round-trips:

\begin{lstlisting}[language=Python]
from fractions import Fraction
from quadmath.lattice.conversions import (
    quadray_to_xyz, xyz_to_quadray_canonical, quadray_roundtrip, urner_embedding,
)
from quadmath.core.quadray import Quadray
from quadmath.lattice.omni_numbering import sites_through_shell

q0 = Quadray(2, 1, 1, 0)                          # cuboctahedron vertex, shell 1
assert xyz_to_quadray_canonical(quadray_to_xyz(q0)) == q0
assert xyz_to_quadray_canonical(
    (Fraction(-1), Fraction(1), Fraction(1))
) == Quadray(1, 1, 1, 0)                          # exact entries, no float fuzz
for site in sites_through_shell(2)[:8]:           # round-trips over enumerated sites
    assert quadray_roundtrip(Quadray(*site)) == Quadray(*site)
assert quadray_roundtrip(q0, urner_embedding(0.5)) == q0   # scale-robust (same rows, both legs)
\end{lstlisting}

\subsection{Lattice context}\label{lattice-context}

The canonical representatives produced by \eqref{eq:conv-canonical} sit
in the IVM site structure of \href{11_ivm_field_learning.md}{IVM Field
Learning}. A normalized quadray whose component sum is
\(\equiv 0 \pmod 4\) is an IVM sphere center
(\passthrough{\lstinline!is\_ivm\_site!}); the other three residue
classes are the octahedral and tetrahedral voids of the packing. The
radial observable is the shell norm

\begin{equation}
\label{eq:conv-shell}
N(q) \;=\; \sum_{i} \Bigl|\, q_i \;-\; \sigma/4 \,\Bigr| \;=\; 2k,
\qquad \sigma \;=\; \sum_{i} q_i ,
\end{equation}

computed by \passthrough{\lstinline!quadray\_shell\_norm!}: shell \(k\)
is the set of sites with \(N(q) = 2k\), populated by \(10k^2 + 2\)
centers (Eq. \eqref{eq:lattice-shell-population} in
\href{13_lattice_tooling.md}{Lattice Tooling}), and
\passthrough{\lstinline!ivm\_field.shell\_sites(k)!} enumerates shell
\(k\) in lexicographic order (1, 12, 42, 92, 162 sites for
\(k = 0..4\)). Conversions, shells, and nearest-site search share one
geometry; Sections \href{11_ivm_field_learning.md}{11} and
\href{13_lattice_tooling.md}{13} detail the field learner and the query
machinery respectively.

\subsection{API Summary}\label{api-summary-1}

{\def\LTcaptype{none} % do not increment counter
\begin{longtable}[]{@{}
  >{\raggedright\arraybackslash}p{(\linewidth - 2\tabcolsep) * \real{0.5000}}
  >{\raggedright\arraybackslash}p{(\linewidth - 2\tabcolsep) * \real{0.5000}}@{}}
\toprule\noalign{}
\begin{minipage}[b]{\linewidth}\raggedright
Function
\end{minipage} & \begin{minipage}[b]{\linewidth}\raggedright
Purpose
\end{minipage} \\
\midrule\noalign{}
\endhead
\bottomrule\noalign{}
\endlastfoot
\passthrough{\lstinline!conversions.urner\_embedding(scale=1.0)!} & The
\((3, 4)\) embedding matrix of Eq. \eqref{eq:conv-embedding} \\
\passthrough{\lstinline!conversions.quadray\_to\_xyz(q, M=None)!} &
Forward map via \passthrough{\lstinline!quadray.to\_xyz!};
\passthrough{\lstinline!M=None!} =
\passthrough{\lstinline!quadray.DEFAULT\_EMBEDDING!} \\
\passthrough{\lstinline!conversions.xyz\_to\_quadray\_canonical(xyz, M=None)!}
& Exact rational canonical inverse, Eqs. \eqref{eq:conv-fiber} and
\eqref{eq:conv-canonical} \\
\passthrough{\lstinline!conversions.quadray\_roundtrip(q, M=None)!} &
Asserts the round-trip identity \eqref{eq:conv-roundtrip}, scale-robust
per \eqref{eq:conv-roundtrip-scale} \\
\passthrough{\lstinline!conversions.embedding\_basis(M=None)!} &
\passthrough{\lstinline!(4, 3)!} array of embedding columns; Gram
identity \eqref{eq:conv-gram-col} \\
\passthrough{\lstinline!quadray.to\_xyz(q, embedding)!} & Underlying
forward product \(\mathrm{xyz} = M\,q\) \\
\passthrough{\lstinline!quadray.quadray\_from\_xyz(x, y, z, embedding)!}
& Float inverse: pseudoinverse + round-half-up (snapping, not exact
inversion) \\
\passthrough{\lstinline!lattice\_search.squared\_distance(p, sites)!} &
Exact squared distances via Eq. \eqref{eq:conv-distance} \\
\end{longtable}
}

\subsection{Verification}\label{verification-1}

\passthrough{\lstinline!tests/test\_conversions.py!} covers the layer
end to end: round-trips of \eqref{eq:conv-roundtrip} over shells 0--4
(309 enumerated sites, center included) and at scaled embeddings per
\eqref{eq:conv-roundtrip-scale}, exact canonical recovery of every site
through \eqref{eq:conv-canonical}, the Gram identities
\eqref{eq:conv-gram-row} and \eqref{eq:conv-gram-col} with
\passthrough{\lstinline!embedding\_basis!}, the distance identity
\eqref{eq:conv-distance} cross-checked against float embeddings and
\passthrough{\lstinline!quadray.distance!}, and the fail-closed error
taxonomy (\passthrough{\lstinline!ValueError!} for off-image points, bad
shapes, zero determinants and non-zero row sums;
\passthrough{\lstinline!TypeError!} for inexact-foreign entry types such
as \passthrough{\lstinline!numpy.float32!}). The repository-root
\passthrough{\lstinline!SPEC.md!} (the full mathematical specification)
and \passthrough{\lstinline!tests/test\_spec\_examples.py!} pin the
worked examples of this section numerically as they land. Run the suite:

\begin{lstlisting}[language=bash]
PYTEST_DISABLE_PLUGIN_AUTOLOAD=1 uv run coverage run -m pytest -q
uv run coverage report
\end{lstlisting}

\newpage

\section{Learning and Evaluation on the IVM
Lattice}\label{learning-and-evaluation-on-the-ivm-lattice}

\subsection{Overview}\label{overview-8}

Learning on a lattice is easy to fake. A learner that merely
interpolates its observations through the graph can look excellent on
the very data it was fit to, and a dynamics model identified on early
snapshots can silently diverge on late ones. This section develops the
training/testing methodology used to evaluate the two IVM learners
honestly --- the static field learner of
\href{11_ivm_field_learning.md}{Static IVM Field Learning} and the
dynamics identifier of \href{12_ivm_dynamics.md}{IVM Lattice Dynamics}.
All utilities live in the \passthrough{\lstinline!learning\_eval.py!}
module (\passthrough{\lstinline!kfold\_site\_splits()!},
\passthrough{\lstinline!cross\_validate\_field()!},
\passthrough{\lstinline!trajectory\_train\_test()!},
\passthrough{\lstinline!learning\_curve()!},
\passthrough{\lstinline!enclosing\_radius()!}), use
\passthrough{\lstinline!numpy!} only, and are fully deterministic for
fixed seeds. Every behavioral claim below is pinned in
\passthrough{\lstinline!tests/test\_learning\_eval.py!} with fixed RNG
seeds and real numerics.

\subsection{Why held-out evaluation is harder on a
lattice}\label{why-held-out-evaluation-is-harder-on-a-lattice}

On independent data, random resampling is benign. On a lattice it is
not: field values at neighboring sites are strongly correlated, and the
learners of \passthrough{\lstinline!ivm\_field.py!} exploit exactly that
correlation --- the kernel-weighted target of the normal equations
\eqref{eq:ivm-normal-equations} predicts an unobserved site as a
weighted average of its graph neighbors. A random split that leaves
graph neighbors on both sides therefore lets the learner ``cheat'': the
held-out site is predicted from observations one hop away, and the
apparent skill measures the smoothness of the field under
\eqref{eq:ivm-objective}, not generalization to unseen territory. The
autocorrelation that makes lattice learning work is the same
autocorrelation that inflates its evaluation:

\begin{equation}
\label{eq:learn-correlation}
\operatorname{Cov}\bigl(f_i, f_j\bigr) \;\approx\; \sigma_f^{2}\,
K\bigl(d(i,j)\bigr),
\qquad d(i,j) \;=\; \text{hop distance on the IVM graph},
\end{equation}

with \(K\) the Gaussian hop kernel of \eqref{eq:ivm-kernel} and
\(\sigma_f^2\) the field variance. Because \(K\) decays with hop
distance, the leakage is local --- but the IVM ball is small, and a
random split leaks everywhere.

The methodology below therefore always holds out \emph{structure}, not
just rows: whole folds of sites for the field learner (never shared
between training and scoring), and a temporal suffix of snapshots for
dynamics identification (never shuffled). The pinned numerical result
worth internalizing: on the radius-2 ball (55 sites, 80\% observed), the
noise-free held-out error of the kernel-weighted interpolator is
dominated not by observation noise but by \emph{coverage geometry} ---
which sites happen to be observed --- so per-fold scores vary
substantially even at fixed sample size, and honest evaluation must
average over folds.

\subsection{K-fold site splits}\label{k-fold-site-splits}

\passthrough{\lstinline!kfold\_site\_splits(observed\_sites, k, seed)!}
partitions observed sites into \(k\) folds by a seeded
\passthrough{\lstinline!numpy.random.default\_rng!} permutation cut with
\passthrough{\lstinline!numpy.array\_split!} (the first \(n \bmod k\)
folds receive one extra site). Each split is a (train, test) pair; folds
are pairwise disjoint and jointly exhaustive, and site order within each
list follows the original normalized index order, so the splits are a
pure function of \((sites, k, seed)\). With \(s\) the seeded permutation
and \(F_j\) fold \(j\)'s index set, the held-out score of a fit
\(\hat{f}^{(-F_j)}_\lambda\) that never saw \(F_j\) is

\begin{equation}
\label{eq:learn-fold}
\widehat{M}(\lambda, F_j) \;=\; \frac{1}{|F_j|}
\sum_{i \in F_j}
\bigl( \hat{f}^{(-F_j)}_\lambda(i) - y_i \bigr)^{2},
\end{equation}

the quantity \passthrough{\lstinline!IVMField.score()!} returns.
Duplicate sites are rejected after normalization (two projective
representatives of one lattice point would sit on both sides of a
split), as is \(k < 2\) or \(k > n\).

\subsection{Cross-validation and honest model
selection}\label{cross-validation-and-honest-model-selection}

\passthrough{\lstinline!cross\_validate\_field(field\_values, sites, k, lam\_grid, seed)!}
runs \eqref{eq:learn-fold} for every fold and every regularization
candidate of \passthrough{\lstinline!lam\_grid!}, fitting each fold from
scratch on the training sites only --- a fresh
\passthrough{\lstinline!IVMField.lattice\_ball!} per fit, via
\passthrough{\lstinline!IVMField.learn()!} with the objective
\eqref{eq:ivm-objective}. The result is the full table
\(\widehat{M}(\lambda, F_j)\) with per-candidate mean and standard
deviation over folds, and the selected strength

\begin{equation}
\label{eq:learn-lambda-star}
\lambda^{*} \;=\; \arg\min_{\lambda \in \mathcal{G}}
\frac{1}{k} \sum_{j=1}^{k} \widehat{M}(\lambda, F_j),
\end{equation}

resolved to the earliest candidate on ties. Selection uses validation
folds only; the returned \passthrough{\lstinline!refit\_field!} is then
refit with \(\lambda^{*}\) on \emph{all} observed sites, matching the
deployment condition where every observation is available. Separating
the two is what makes the selection honest: choosing \(\lambda\) on the
same data the final model fits is how interpolation gets mistaken for
skill.

Two regimes are pinned numerically on synthetic fields. With exact
observations of a harmonic (linear in embedded XYZ) field at 80\%
coverage, cross-validation selects \(\lambda^{*} = 0\): the
kernel-weighted interpolator generalizes best, and Laplacian smoothing
only biases the fit away from exact data --- the mean held-out MSE is
strictly increasing in \(\lambda\) across a grid spanning two decades.
With noisy observations of a smooth non-harmonic field, the selection
flips to \(\lambda^{*} > 0\): interpolation carries observation noise
into the fit, and moderate regularization wins on the held-out folds.
The margin is honest about the learner's limits --- because observed
rows are pinned to their data, the noise propagates through the graph
regardless of \(\lambda\), so the achievable gain is set by how much
neighbor averaging the kernel already performs, not by the regularizer
alone.

\subsection{Temporal splits for dynamics
identification}\label{temporal-splits-for-dynamics-identification}

\passthrough{\lstinline!trajectory\_train\_test(observed, split\_frac, lattice)!}
evaluates \passthrough{\lstinline!fit\_trajectory()!}-style
identification of the coupling \(\alpha\) (\href{12_ivm_dynamics.md}{IVM
Lattice Dynamics}). The snapshot matrix (\(T+1\) rows, row 0 the initial
condition) is split \textbf{temporally}: the first
nearest-integer-rounded fraction of rows trains, the trailing rows are
held out, and no shuffling is ever applied --- a shuffled split
interpolates between adjacent snapshots of a deterministic trajectory
and hides exactly the forecast error being measured. The split point is
clamped so training keeps at least two rows (initial condition plus one
snapshot) and testing keeps at least one. Identification uses the
training prefix only, with candidate grid and refinement passed through
to \passthrough{\lstinline!fit\_trajectory()!}. Two held-out errors are
reported. The multi-step continuation error re-simulates from the
observed initial condition with the identified \(\hat{\alpha}\) and
scores only the held-out rows:

\begin{equation}
\label{eq:learn-horizon}
M_{\mathrm{test}}(\hat{\alpha}) \;=\; \frac{1}{(T - t_s)\,N}
\sum_{t=t_s+1}^{T} \big\lVert u_{\hat{\alpha}}(t) - u_{\mathrm{obs}}(t)
\big\rVert^{2} ,
\end{equation}

where \(t_s\) is the last training row. This is the stringent metric:
for any model that is not exactly right, error compounds with horizon,
so the held-out tail punishes misspecification that early rows forgive.
The one-step-ahead error is teacher-forced --- each held-out transition
is predicted a single step from the \emph{true} previous state, the
first transition starting from the last training row:

\begin{equation}
\label{eq:learn-onestep}
M_{1}(\hat{\alpha}) \;=\; \frac{1}{(T - t_s)\,N}
\sum_{t=t_s}^{T-1} \big\lVert
\mathrm{step}\bigl(u_{\mathrm{obs}}(t);\, \hat{\alpha}\bigr)
- u_{\mathrm{obs}}(t+1) \big\rVert^{2},
\end{equation}

isolating local prediction error from error accumulation. The contrast
between \eqref{eq:learn-horizon} and \eqref{eq:learn-onestep} is
diagnostic: a correct model drives both to zero (pinned exactly: a heat
trajectory with \(\alpha = 0.3\) in the candidate grid yields train,
continuation, and one-step errors of exactly \(0\)), while a
misspecified model class --- a heat model identified on
majority-dynamics data --- shows a modest training error, a held-out
continuation error more than twice as large (horizon compounding), and a
one-step error five times smaller than the continuation error (local
prediction is easier than long-horizon forecasting). Note the direction
of the misspecification: \passthrough{\lstinline!majority\_step()!}
requires integer states, so the float-continuum heat model can be fit to
integer data, but not the reverse.

\subsection{Learning curves}\label{learning-curves}

\passthrough{\lstinline!learning\_curve(values, sites, train\_fracs, seed)!}
traces the classic data-coverage curve on a lattice. One seeded
permutation of the observed sites is drawn; for each requested fraction
\(p\) the first \(m(p) = \operatorname{round}(p \cdot n)\) sites of that
permutation (clamped to \([1, n-1]\)) form the training subset and the
remainder the held-out complement:

\begin{equation}
\label{eq:learn-curve}
M(p) \;=\; \frac{1}{n - m(p)} \sum_{i \notin P_{m(p)}}
\bigl( \hat{f}_{P_{m(p)}}(i) - y_i \bigr)^{2},
\qquad P_{m} \;=\; \text{first } m \text{ entries of the permutation},
\end{equation}

reported together with the in-sample score of the same fit. The subsets
are \emph{nested} by construction --- growing \(p\) never removes an
earlier observation --- so the curve is a proper coverage curve rather
than a family of independent splits. The pinned example (radius-2 ball,
quadratic bowl observed with noise \(\sigma = 0.4\) everywhere,
\(\lambda = 0.05\)) shows held-out MSE falling by more than a factor of
two from 10\% to 85\% coverage, with non-monotone wiggles in between:
with few observations the held-out complement is large and its geometry
varies between fractions, so the curve is monotone-ish, not monotone ---
the final point sits far below the first, and that is the asserted
invariant. The curve also separates interpolation from generalization
visibly: the in-sample MSE stays below \(0.15\) at every coverage level
while the held-out MSE never drops below \(0.7\) --- the pinned
observations are fit closely at all sizes, and the difficulty is
entirely on unseen sites.

\subsection{Three-way splitting and trainable site
fits}\label{sec:learn_three_way_site_fits}

The lattice surfaces above all share one shape: a seeded split of sites
or snapshots, a fit on the training side, and scores on the held-out
side. The same shape is available for plain tabular work in the same
module (\passthrough{\lstinline!src/quadmath/learn/learning\_eval.py!}),
under the same discipline --- all split randomness confined to the
seeded permutation of the split, and fits that are pure deterministic
functions of their inputs.

\passthrough{\lstinline!three\_way\_split(n, train\_frac=0.6, val\_frac=0.2, seed=0)!}
partitions \passthrough{\lstinline!range(n)!} into disjoint
train/validation/test index lists. One seeded
\passthrough{\lstinline!numpy.random.default\_rng!} permutation of
\passthrough{\lstinline!range(n)!} is drawn and cut into three
contiguous blocks: the first \(m_{train} =
\operatorname{round}(\mathrm{train\_frac} \cdot n)\) indices train, the
next \(m_{val} = \operatorname{round}(\mathrm{val\_frac} \cdot n)\)
validate, and the remainder test. Each returned list is sorted, the
three are pairwise disjoint with union exactly
\passthrough{\lstinline!range(n)!}, and the split is a deterministic
function of \((n, \mathrm{train\_frac}, \mathrm{val\_frac}, seed)\) ---
the same deterministic-permutation contract as the k-fold site splits
above, applied to a plain index range rather than observed lattice
sites. Rounding is unguarded at the tails: for small \(n\) a trailing
block may come back empty (\(n = 3\) under the default fractions leaves
the test list empty), and a \passthrough{\lstinline!ValueError!} is
raised when either fraction is not strictly positive or their sum
reaches \(1\) --- a split with no held-out rows at all cannot be honest.

\passthrough{\lstinline!ridge\_site\_fit(features, values, lam=1e-3)!}
fits a single linear site model in closed form. With \(X_c\) the
mean-centered design matrix, \(y_c\) the mean-centered target, and
\(\bar{x}\), \(\bar{y}\) the feature-column means and target mean, the
ridge normal equations are solved exactly for the slope vector \(w\):

\begin{equation}
\label{eq:learn-ridge}
\bigl( X_c^{\!\top} X_c + \lambda I \bigr)\, w \;=\; X_c^{\!\top} y_c,
\qquad b \;=\; \bar{y} - \bar{x}^{\!\top} w,
\end{equation} with \(b\) the intercept recovered from the centered
solve; \(\lambda \geq 0\) penalizes the centered coefficient norm,
shrinking \(w\) toward zero, and \(\lambda = 0\) recovers the exact
least-squares solve. The returned \passthrough{\lstinline!RidgeSiteFit!}
carries \passthrough{\lstinline!coefficients!},
\passthrough{\lstinline!intercept!}, and the in-sample
\passthrough{\lstinline!train\_mse!}, and the fit is a deterministic
function of \passthrough{\lstinline!(features, values, lam)!} --- a
\passthrough{\lstinline!numpy.linalg.solve!}, not an iterative
optimizer, so repeated calls on the same inputs reproduce it bit for
bit. Validation is explicit: a negative \(\lambda\), a non-2-D design
matrix, a non-1-D target, disagreeing sample counts, or zero rows each
raise before any linear algebra runs.

\passthrough{\lstinline!GradientDescentTrainer(lr=0.05, max\_iters=300, tol=1e-9)!}
is the iterative counterpart, useful where the closed form is
deliberately being compared against: a full-batch gradient-descent fit
of the same linear model. \passthrough{\lstinline!fit!} standardizes the
feature columns internally (per-column zero mean and unit standard
deviation, with constant columns passing through at standard deviation
\(1\)), then runs at most \passthrough{\lstinline!max\_iters!}
full-batch updates on the mean-squared-error loss,

\begin{equation}
\label{eq:learn-gd}
\theta \;\leftarrow\; \theta \;-\; \frac{2\,\eta}{m}\,
\begin{bmatrix} X_s^{\!\top} e \\ \mathbf{1}^{\!\top} e \end{bmatrix},
\qquad e \;=\; X_s w + b - y,
\end{equation}

where \(X_s\) is the standardized design, \(m\) the row count, \(\eta\)
the learning rate, \(w\) and \(b\) the coefficient vector and bias, and
\(y\) the target. The MSE at the top of each executed iteration is
appended to \passthrough{\lstinline!loss\_history!} --- one float per
iteration, recorded before the update it describes --- and training
stops early with \passthrough{\lstinline!converged\_ = True!} as soon as
two consecutive losses differ by less than
\passthrough{\lstinline!tol!}. The fit exposes
\passthrough{\lstinline!coef\_!} expressed against the standardized
features and \passthrough{\lstinline!intercept\_!} in original target
units, together with the stored \passthrough{\lstinline!mean\_!} and
\passthrough{\lstinline!std\_!} it used, so
\passthrough{\lstinline!predict!} re-applies the recorded
standardization before the linear map and thus scores consistently
across the train, validation, and test blocks of a
\passthrough{\lstinline!three\_way\_split!}. Like
\passthrough{\lstinline!ridge\_site\_fit!} the trainer draws no
randomness of its own --- no shuffling, no stochastic gradient --- so
its entire loss history is deterministic for fixed inputs: the same rows
and hyperparameters replay the same loss curve to the last float, and
the per-iteration record doubles as the divergence/plateau diagnostic
that a scalar final loss cannot provide.

\begin{figure}
\centering
\pandocbounded{\includegraphics[keepaspectratio,alt={Convergence of the full-batch gradient-descent trainer. Per-iteration training loss (mean squared error, linear y-axis) of GradientDescentTrainer on a deterministic synthetic linear problem: a 40 \textbackslash times 3 standard-normal design matrix with true coefficients (1.5, -2.0, 0.75), intercept 0.5, and Gaussian target noise \textbackslash sigma = 0.05 (seed 0). Fitted with lr = 0.05, max\_iters = 300, tol = 1e-9; the loss falls from 6.23 at iteration 0 to the noise-floor scale \textbackslash sigma\^{}2 = 2.5\textbackslash times 10\^{}\{-3\}, flagged as converged after 126 iterations. Regenerate with quadmath/scripts/learning\_gallery.py.}]{../output/figures/learn_loss_history.png}}
\caption{\textbf{Convergence of the full-batch gradient-descent
trainer.} Per-iteration training loss (mean squared error, linear
\(y\)-axis) of \protect\passthrough{\lstinline!GradientDescentTrainer!}
on a deterministic synthetic linear problem: a \(40 \times 3\)
standard-normal design matrix with true coefficients
\((1.5, -2.0, 0.75)\), intercept \(0.5\), and Gaussian target noise
\(\sigma = 0.05\) (seed \(0\)). Fitted with
\protect\passthrough{\lstinline!lr = 0.05!},
\protect\passthrough{\lstinline!max\_iters = 300!},
\protect\passthrough{\lstinline!tol = 1e-9!}; the loss falls from
\(6.23\) at iteration \(0\) to the noise-floor scale
\(\sigma^2 = 2.5\times 10^{-3}\), flagged as converged after \(126\)
iterations. Regenerate with
\protect\passthrough{\lstinline!quadmath/scripts/learning\_gallery.py!}.}
\end{figure}

\subsection{Verification}\label{verification-2}

\passthrough{\lstinline!tests/test\_learning\_eval.py!} pins all of the
above with fixed seeds and no mocks: disjoint-exhaustive folds with
exact fold sizes, seed-determinism and seed-sensitivity of the splits,
the \(\lambda^{*}=0\) and \(\lambda^{*}>0\) regimes with their table
identities (mean and standard deviation recomputed from the fold table;
the refit reproduced by hand), exact-zero errors for the correctly
specified dynamics against a \(> 2\times\) continuation blow-up for the
misspecified one, manual recomputation of the one-step-ahead error,
fraction clamping at both extremes, and the learning curve's
final-below-first invariant. The module carries 100\% statement and
branch coverage, and the validation rules of every entry point (distinct
normalized sites, \(k \in
[2, n]\), fractions in the open unit interval, non-empty grids, minimum
row counts) are each exercised.

\subsection{Cross-References}\label{cross-references-2}

\begin{itemize}
\tightlist
\item
  Coordinate foundations and the twelve-around-one shell:
  \href{03_quadray_methods.md}{Quadray Methods}
\item
  The learner being evaluated: \href{11_ivm_field_learning.md}{Static
  IVM Field Learning}
\item
  The dynamics being identified: \href{12_ivm_dynamics.md}{IVM Lattice
  Dynamics}
\item
  Enumeration and search infrastructure:
  \href{13_lattice_tooling.md}{Lattice Tooling}
\end{itemize}

\newpage

\section{Lattice Visualization
Gallery}\label{lattice-visualization-gallery}

\subsection{Overview}\label{overview-9}

This section is the figure surface of the lattice layer: every rendering
primitive used by the IVM chapters lives in
\passthrough{\lstinline!vis\_lattice.py!}
(\passthrough{\lstinline!shell\_scatter!},
\passthrough{\lstinline!field\_slice!},
\passthrough{\lstinline!dynamics\_strip!},
\passthrough{\lstinline!gallery!}). Each primitive receives its
matplotlib axes explicitly, so panels compose inside caller-owned
figures; only \passthrough{\lstinline!gallery!} creates figures (three,
written as PNGs). Nothing renders at import time, and all randomness is
drawn from a seeded \passthrough{\lstinline!numpy.random.default\_rng!},
so the gallery is deterministic. The command-line entry is
\passthrough{\lstinline!quadmath/scripts/lattice\_gallery.py!} --- a
thin orchestrator that sets a headless backend and a fixed seed,
delegates to \passthrough{\lstinline!gallery!} (thin-orchestrator
contract of \passthrough{\lstinline!quadmath/scripts/AGENTS.md!}), and
prints each written path on its own line: those stdout lines are the
\passthrough{\lstinline!make\_all\_figures!} manifest contract.

\subsection{Frequency shells}\label{frequency-shells}

\passthrough{\lstinline!shell\_scatter(ax, sites, k)!} embeds the sites
of shell \passthrough{\lstinline!k!} (from
:func:\passthrough{\lstinline!omni\_numbering.generate\_shell!}, which
enumerates the \(10k^2 + 2\) sites of shell \(k\); see
\href{13_lattice_tooling.md}{Lattice Tooling}) with
\passthrough{\lstinline!to\_xyz!} under
\passthrough{\lstinline!DEFAULT\_EMBEDDING!} and scatters them in 3D.
With \passthrough{\lstinline!axis\_hints!} it overlays four dashed rays
from the origin toward the embedding matrix columns --- the images of
the quadray basis directions \(A, B, C, D\), which are the vertices of a
regular tetrahedron (pairwise dot product \(-1\), edge length
\(2\sqrt{2}\); see \href{03_quadray_methods.md}{Quadray Methods}). The
hints make the tetrahedral axis frame of the lattice readable in print.

\begin{figure}
\centering
\pandocbounded{\includegraphics[keepaspectratio,alt={Frequency shells of the omnidirectional close packing. Shell 1 (12 sites) and shell 2 (42 sites) of the IVM lattice embedded via DEFAULT\_EMBEDDING and scattered by shell\_scatter as 3D point clouds in the x, y, z embedding coordinates; dashed gray rays mark the four tetrahedral quadray axes, labeled A, B, C, D, pointing toward the embedding-matrix columns. Reproduced with uv run python quadmath/scripts/lattice\_gallery.py (fixed seed 12).}]{../output/figures/vis_gallery_shell.png}}
\caption{\textbf{Frequency shells of the omnidirectional close packing.}
Shell 1 (12 sites) and shell 2 (42 sites) of the IVM lattice embedded
via \protect\passthrough{\lstinline!DEFAULT\_EMBEDDING!} and scattered
by \protect\passthrough{\lstinline!shell\_scatter!} as 3D point clouds
in the \(x, y, z\) embedding coordinates; dashed gray rays mark the four
tetrahedral quadray axes, labeled A, B, C, D, pointing toward the
embedding-matrix columns. Reproduced with
\protect\passthrough{\lstinline!uv run python quadmath/scripts/lattice\_gallery.py!}
(fixed seed 12).}
\end{figure}

\subsection{Lattice-plane slices}\label{lattice-plane-slices}

\passthrough{\lstinline!field\_slice(ax, field, sites, plane)!} renders
a scalar \passthrough{\lstinline!ivm\_field.IVMField!} as a heatmap
restricted to a lattice plane. A plane is parametrized by an origin
\(q_0\) and two independent integer translation vectors \(u\), \(v\) ---
by default two of the twelve IVM neighbor moves, both permutations of
\((2,1,1,0)\) --- and its points are

\begin{equation}
\label{eq:vis-plane}
q(i, j) \;=\; \operatorname{normalize}\!\big(q_0 + i\,u + j\,v\big),
\qquad (i, j) \in \mathbb{Z}^2 ,
\end{equation}

with \passthrough{\lstinline!normalize!} the quadray canonical
representative. Because each neighbor move has component sum \(4\), the
sum stays divisible by \(4\) along the whole plane, so every \(q(i,j)\)
is again a lattice site. Two coordinate facts make the rendering exact:
the embedding is linear, so the xyz image is the planar set
\(t(q_0) + i\,t(u) + j\,t(v)\); and the shift \((1,1,1,1)\) lies in the
embedding kernel, so plane membership is tested in integer quadray space
--- a site lies on the plane of \eqref{eq:vis-plane} iff
\(q - q_0 = i\,u + j\,v + m\,(1,1,1,1)\) for integers \(i, j, m\).
\passthrough{\lstinline!field\_slice!} decides membership by an exact
least-squares solve followed by an integer verification, so no site is
ever misassigned by floating point. Sites off the plane are ignored,
plane cells outside the field's ball are masked, and the heatmap axes
are the integer plane indices \(i\) (steps along \(u\)) and \(j\) (steps
along \(v\)).

The gallery figure learns a field on the radius-3 ball (147 sites) from
noisy observations of a linear synthetic truth (25\% of sites observed,
noise level \(0.1\), regularization \passthrough{\lstinline!lam=0.01!},
kernel width \passthrough{\lstinline!0.5!}) and slices it through the
default plane \(\big((2,1,1,0), (1,2,1,0)\big)\); the learned surface
follows the planar lattice, with masked tiles where the ball ends. Field
learning itself is the subject of \href{11_ivm_field_learning.md}{Static
IVM Field Learning}.

\begin{figure}
\centering
\pandocbounded{\includegraphics[keepaspectratio,alt={Learned scalar field sliced along a lattice plane. Heatmap of a Laplacian-regularized IVMField over the radius-3 ball, restricted by field\_slice to the plane spanned by the neighbor moves u = (2,1,1,0) and v = (1,2,1,0): a viridis heatmap whose axes are the integer plane indices i (steps along u) and j (steps along v), with the colorbar reporting the field value and light-gray masked tiles marking plane cells outside the ball. Reproduced with uv run python quadmath/scripts/lattice\_gallery.py (fixed seed 12).}]{../output/figures/vis_gallery_field.png}}
\caption{\textbf{Learned scalar field sliced along a lattice plane.}
Heatmap of a Laplacian-regularized
\protect\passthrough{\lstinline!IVMField!} over the radius-3 ball,
restricted by \protect\passthrough{\lstinline!field\_slice!} to the
plane spanned by the neighbor moves u = (2,1,1,0) and v = (1,2,1,0): a
\protect\passthrough{\lstinline!viridis!} heatmap whose axes are the
integer plane indices \(i\) (steps along u) and \(j\) (steps along v),
with the colorbar reporting the field value and light-gray masked tiles
marking plane cells outside the ball. Reproduced with
\protect\passthrough{\lstinline!uv run python quadmath/scripts/lattice\_gallery.py!}
(fixed seed 12).}
\end{figure}

\subsection{Dynamics strips}\label{dynamics-strips}

\passthrough{\lstinline!dynamics\_strip(axs, trajectory, t\_indices)!}
renders selected snapshots of an
\passthrough{\lstinline!ivm\_dynamics.simulate!} trajectory (see
\href{12_ivm_dynamics.md}{Dynamics and Learning on the IVM Lattice}) as
a strip of 3D scatter panels over the embedded lattice. All panels share
one symmetric color scale \([-v_{\max}, v_{\max}]\) with
\(v_{\max} = \max_t |u_t|\) over the selected snapshots, so evolution is
directly comparable across panels. The gallery strip shows heat
diffusion (\(\alpha = 0.45\), horizon \(T = 20\)) on the radius-3
lattice (13 sites) at \(t = 0, 10, 20\): the seeded random initial field
relaxes toward its lattice average --- the visually flat mid panel and
right panel are the \(L_2\)-monotone decay of the heat lemma in action.

\begin{figure}
\centering
\pandocbounded{\includegraphics[keepaspectratio,alt={Heat-diffusion evolution strip. Snapshots of a seeded heat run at t = 0, 10, 20 on the radius-3 IVM lattice (13 sites), rendered by dynamics\_strip as 3D scatter panels on one shared symmetric coolwarm scale {[}-v\_\{\textbackslash max\}, v\_\{\textbackslash max\}{]} with v\_\{\textbackslash max\} = \textbackslash max\_t \textbar u\_t\textbar{} over the selected snapshots; the field visibly relaxes as the sum of squares decays monotonically. Reproduced with uv run python quadmath/scripts/lattice\_gallery.py (fixed seed 12).}]{../output/figures/vis_gallery_dynamics.png}}
\caption{\textbf{Heat-diffusion evolution strip.} Snapshots of a seeded
heat run at t = 0, 10, 20 on the radius-3 IVM lattice (13 sites),
rendered by \protect\passthrough{\lstinline!dynamics\_strip!} as 3D
scatter panels on one shared symmetric
\protect\passthrough{\lstinline!coolwarm!} scale
\([-v_{\max}, v_{\max}]\) with \(v_{\max} = \max_t |u_t|\) over the
selected snapshots; the field visibly relaxes as the sum of squares
decays monotonically. Reproduced with
\protect\passthrough{\lstinline!uv run python quadmath/scripts/lattice\_gallery.py!}
(fixed seed 12).}
\end{figure}

\subsection{Deterministic animations and gallery
plots}\label{sec:animation_frames}

Where the primitives above compose panels inside caller-owned figures,
the module \passthrough{\lstinline!src/quadmath/viz/animations.py!}
renders frame-based animations that own their pixels: each renderer
returns a list of \passthrough{\lstinline!Frame!} objects --- the frozen
\passthrough{\lstinline!dataclass!} \passthrough{\lstinline!Frame!} in
that file pairs a 2D \passthrough{\lstinline!numpy!} array with a
human-readable \passthrough{\lstinline!title!}, and its
\passthrough{\lstinline!\_\_post\_init\_\_!} validation admits only
\passthrough{\lstinline!uint8!} arrays (raw gray levels) or float arrays
with every value in \([0, 1]\), rejecting anything that is not 2D, is of
another dtype, or drifts outside the unit interval. No matplotlib is
involved at frame level and no RNG or wall-clock input exists (except
the explicit \passthrough{\lstinline!seed!} of
\passthrough{\lstinline!diffusion\_frames!}), so identical arguments
yield byte-identical frames. All three renderers share one fixed-camera
convention: sites are embedded in \(R^3\) via
\passthrough{\lstinline!DEFAULT\_EMBEDDING!} and viewed by an
orthographic camera looking down \(+z\) --- each point \((x,y,z)\)
projects to \((x,y)\) on a \passthrough{\lstinline!GRID\_SIZE!} \(=48\)
square grid whose extent is \([-8, 8]\) (module constant
\passthrough{\lstinline!\_HALF\_EXTENT!}, sized to cover a radius-3 IVM
ball of embedded norm \(6\) with margin even while pulsing), the
vertical axis flipped so \(+y\) points up, colliding points resolved by
maximum brightness, and out-of-window points clamped to the border
pixel.

\passthrough{\lstinline!simplex\_frames(q0, q1, n=16)!} animates a
quaternion-slerp rotation of the radius-1 ball: the two unit quaternions
\(q_0\), \(q_1\) (validated to four components of unit norm) are
interpolated on the shortest arc --- sign-flipping \(q_1\) when their
dot product is negative, using the Shoemaker spherical formula
\(q(t) = \sin((1-t)\theta)\,q_0 / \sin\theta + \sin(t\theta)\,q_1 / \sin\theta\),
and falling back to normalized linear interpolation when the inputs are
nearly parallel --- and at each \(t = i/(n-1)\) the 13 lattice sites of
\passthrough{\lstinline!ball\_sites!} radius 1 are rotated by \(q(t)\),
projected onto the 48×48 grid, and shaded by radial shell, with origin
brightness \(1.0\) and the outermost shell never dimmer than \(0.25\) so
nothing renders invisible.

\passthrough{\lstinline!lattice\_frames(shells=3, n=12)!} renders a
pulsing IVM lattice ball: the radius-\passthrough{\lstinline!shells!}
ball from \passthrough{\lstinline!ball\_sites!} (147 sites at the
default radius 3) has its embedded coordinates scaled per frame by the
deterministic pulse \(1 + 0.25\sin(2\pi i/n)\) over one full
growth-and-contraction cycle, while brightness per site stays the fixed
radial-shell shading --- brightest at the center, decaying linearly in
shell index, outermost shell at \(0.25\) --- so the pulse animates
geometry, not shading.

\passthrough{\lstinline!diffusion\_frames(n\_steps=12, seed=0)!} seeds
one-hot heat on the radius-3 ball and diffuses it over the IVM neighbor
graph: the 147 ball sites become a graph in which two sites are adjacent
iff their difference normalizes to one of the twelve
\passthrough{\lstinline!IVM\_NEIGHBOR\_STEPS!}, a
\passthrough{\lstinline!numpy.random.default\_rng(seed)!} draw selects
the source site deterministically, and each explicit update applies
\(u \leftarrow u + \alpha\,(\text{mean of graph neighbors} - u)\) with
\(\alpha = 0.25\) (stable below \(0.5\)) and values clipped at zero.
Each frame is normalized by dividing by its (always positive) maximum
heat, giving float values in \([0, 1]\); because heat only spreads to
graph neighbors, the set of heated sites grows monotonically across
frames.

The companion module \passthrough{\lstinline!src/quadmath/viz/plots.py!}
supplies standalone figure builders for diagnostics that do not belong
to any gallery: \passthrough{\lstinline!plot\_loss\_history!} draws a
training loss sequence as a marker line over its iteration index,
\passthrough{\lstinline!plot\_shell\_growth!} plots the IVM shell
cardinalities --- the cuboctahedral numbers \(10k^2+2\) produced by
\passthrough{\lstinline!shell\_cardinalities!} in
\passthrough{\lstinline!src/quadmath/lattice/ivm\_field.py!} --- versus
shell index, \passthrough{\lstinline!plot\_error\_histogram!} histograms
error values with a dashed vertical mean line, and
\passthrough{\lstinline!plot\_lattice\_shell\_3d!} scatters the sites of
one shell (\passthrough{\lstinline!shell\_sites!}, embedded via
\passthrough{\lstinline!DEFAULT\_EMBEDDING!}) in 3D with equal-aspect
axes. Unlike the axes-composing primitives of
\passthrough{\lstinline!vis\_lattice!}, each owns a fixed-size figure
(6.4×4.8 inches at 160 dpi), applies a fixed style, and on
\passthrough{\lstinline!save=True!} writes its PNG into
\passthrough{\lstinline!quadmath/output/figures/!} via
\passthrough{\lstinline!quadmath.paths.get\_figure\_dir!}, returning the
output path (empty string when the caller keeps the figure). All four
are input-agnostic and seed-free: inputs are plain sequences or shell
indices, styling is fixed, and no randomness or wall-clock time enters,
so runs stay deterministic headless
(\passthrough{\lstinline!MPLBACKEND=Agg!}).

\passthrough{\lstinline!frames\_to\_gif(frames, out\_path, fps=8, scale=8)!}
assembles a frame sequence into an animated GIF: float frames are
quantized to \passthrough{\lstinline!uint8!} gray levels by rounding
\(255\,v\), each frame is upscaled by the integer factor
\passthrough{\lstinline!scale!} with nearest-neighbor resampling, and
the sequence is saved with per-frame duration \(1000\,/\,\texttt{fps}\)
milliseconds and infinite looping. PIL is imported lazily inside the
function (Pillow stays optional at import time), and because the PIL GIF
encoder is deterministic and embeds no timestamps, identical
\passthrough{\lstinline!(frames, fps, scale)!} inputs render
byte-identical GIF files --- the same byte-stability contract the PNG
gallery asserts.

\passthrough{\lstinline!frames\_strip(frames, out\_path=None, labels=None, save=True)!}
renders the print still-counterpart of
\passthrough{\lstinline!frames\_to\_gif!}: a single-row matplotlib strip
with one grayscale \passthrough{\lstinline!imshow!} panel per
\passthrough{\lstinline!Frame!}, figure width scaling with the panel
count while the height stays one panel plus the title band, each panel
pinned to the unit brightness interval
(\passthrough{\lstinline!vmin = 0!}, \passthrough{\lstinline!vmax = 1!};
\passthrough{\lstinline!uint8!} frames are first mapped by \(v/255\), so
both \passthrough{\lstinline!Frame!} dtypes share the GIF's gray-level
convention), per-panel titles taken from
\passthrough{\lstinline!labels!} or, by default, from each frame's own
\passthrough{\lstinline!title!}, and axes switched off with zero
inter-panel spacing so consecutive panels butt together. Matplotlib is
imported lazily inside the function (the module stays matplotlib-free at
import time), and saving follows the \passthrough{\lstinline!plots.py!}
convention: a bare \passthrough{\lstinline!out\_path!} name is written
into \passthrough{\lstinline!quadmath/output/figures/!} via
\passthrough{\lstinline!quadmath.paths.get\_figure\_dir!}, a path
carrying a directory component is used verbatim, and the returned string
is the written path (\passthrough{\lstinline!""!} when
\passthrough{\lstinline!save!} is off or no name is given). An empty
\passthrough{\lstinline!frames!} sequence or a
\passthrough{\lstinline!labels!} length mismatch raises
\passthrough{\lstinline!ValueError!}; the render has no RNG and no
wall-clock input, so identical inputs give byte-identical PNGs --- the
test suite asserts the byte stability, the full \([0, 1]\) grayscale
excursion, the uint8/float equivalence, and the panel-count scaling.

\begin{figure}
\centering
\pandocbounded{\includegraphics[keepaspectratio,alt={Three animation stills in one strip. Single-row frames\_strip composition of one evenly spaced still from each 6-frame animation: left, the pulsing radius-2 IVM ball sampled at the quarter-pulse fraction (lattice\_frames, radial-shell brightness, brightest at the center); center, the quaternion-slerp rotation of the radius-1 ball in the midpoint region of the arc (t = 0.4 on the way from the identity to a 90-degree rotation about z, simplex\_frames); right, the final state of seeded heat diffusion on the radius-3 ball (diffusion\_frames, brightness normalized by the frame's maximum heat). All panels are 48×48 orthographic projections under the fixed +z camera with brightness in {[}0, 1{]}; reproduced with uv run python quadmath/scripts/animation\_stills.py.}]{../output/figures/animation_frames_strip.png}}
\caption{\textbf{Three animation stills in one strip.} Single-row
\protect\passthrough{\lstinline!frames\_strip!} composition of one
evenly spaced still from each 6-frame animation: left, the pulsing
radius-2 IVM ball sampled at the quarter-pulse fraction
(\protect\passthrough{\lstinline!lattice\_frames!}, radial-shell
brightness, brightest at the center); center, the quaternion-slerp
rotation of the radius-1 ball in the midpoint region of the arc
(\(t = 0.4\) on the way from the identity to a 90-degree rotation about
\(z\), \protect\passthrough{\lstinline!simplex\_frames!}); right, the
final state of seeded heat diffusion on the radius-3 ball
(\protect\passthrough{\lstinline!diffusion\_frames!}, brightness
normalized by the frame's maximum heat). All panels are 48×48
orthographic projections under the fixed \(+z\) camera with brightness
in \([0, 1]\); reproduced with
\protect\passthrough{\lstinline!uv run python quadmath/scripts/animation\_stills.py!}.}
\end{figure}

The command-line surface is
\passthrough{\lstinline!quadmath/scripts/animation\_gallery.py!} --- a
thin orchestrator in the
\passthrough{\lstinline!quadmath/scripts/AGENTS.md!} contract that sets
the headless \passthrough{\lstinline!Agg!} backend and renders three
GIFs via \passthrough{\lstinline!frames\_to\_gif!}:
\passthrough{\lstinline!animation\_lattice.gif!} (pulsing radius-3 ball,
12 frames), \passthrough{\lstinline!animation\_simplex.gif!} (slerp from
the identity to a 90-degree rotation about the \(x\) axis, 16 frames),
and \passthrough{\lstinline!animation\_diffusion.gif!} (12 diffusion
steps from seed 0), writing each into
\passthrough{\lstinline!quadmath/output/figures/!} and printing each
path on its own line. The script is registered in
\passthrough{\lstinline!quadmath/scripts/make\_all\_figures.py!}
alongside the other galleries, so the manifest contract of the Overview
applies unchanged: only path lines on stdout, byte-identical re-renders
for identical inputs. The companion
\passthrough{\lstinline!quadmath/scripts/animation\_stills.py!} follows
the same thin-orchestrator contract for print stills: it renders the
three 6-frame sequences of the module ---
\passthrough{\lstinline!lattice\_frames!} at radius 2,
\passthrough{\lstinline!simplex\_frames!} from the identity to a
90-degree rotation about \(z\) with the endpoint quaternion derived from
\passthrough{\lstinline!rotate\_about\_axis!}, and
\passthrough{\lstinline!diffusion\_frames!} at seed 0 --- samples each
at an evenly spaced fraction (quarter-pulse, slerp midpoint region,
final diffusion step), composes the three stills with
\passthrough{\lstinline!frames\_strip!} into
\passthrough{\lstinline!animation\_frames\_strip.png!}, and prints the
written path on its own line.

\subsection{Reproducibility and test
contract}\label{reproducibility-and-test-contract-1}

\begin{itemize}
\tightlist
\item
  \passthrough{\lstinline!gallery(paths\_out\_dir, seed=12)!} writes
  exactly the three files
  \passthrough{\lstinline!vis\_gallery\_shell.png!},
  \passthrough{\lstinline!vis\_gallery\_field.png!},
  \passthrough{\lstinline!vis\_gallery\_dynamics.png!} (module constant
  \passthrough{\lstinline!GALLERY\_FILES!}, in that order) into
  \passthrough{\lstinline!paths\_out\_dir!}, creating it when missing,
  and returns their paths. Every panel is fully seeded, and the PNGs
  carry no timestamp metadata: the test suite asserts byte-identical
  re-renders for a fixed seed.
\item
  Tests: \passthrough{\lstinline!tests/test\_vis\_lattice.py!} --- 19
  tests, no mocks, deterministic assertions without pixel diffs (artist
  placement on caller-provided axes, exact field values at known lattice
  cells, byte-level reproducibility). Scope with
  \passthrough{\lstinline!PYTEST\_DISABLE\_PLUGIN\_AUTOLOAD=1 uv run coverage run -m pytest tests/test\_vis\_lattice.py -q!}
  and \passthrough{\lstinline!uv run coverage report!} --- 100\%
  statement and branch coverage of
  \passthrough{\lstinline!src/quadmath/viz/vis\_lattice.py!}.
\item
  Tests: \passthrough{\lstinline!tests/unit/viz/test\_animations.py!}
  covers the frame module (\passthrough{\lstinline!Frame!} validation,
  the three renderers, \passthrough{\lstinline!frames\_to\_gif!}) and
  the strip renderer --- \passthrough{\lstinline!frames\_strip!} label
  and emptiness validation, byte-identical re-renders in temporary
  directories, uint8/float equivalence, the full grayscale excursion,
  and panel-count scaling --- headless via
  \passthrough{\lstinline!MPLBACKEND=Agg!} (set in
  \passthrough{\lstinline!tests/conftest.py!}).
\item
  Figure regeneration:
  \passthrough{\lstinline!uv run python quadmath/scripts/lattice\_gallery.py!}
  (the script sets \passthrough{\lstinline!MPLBACKEND=Agg!} itself);
  stdout is exactly the three written paths.
\item
  Still-strip regeneration:
  \passthrough{\lstinline!uv run python quadmath/scripts/animation\_stills.py!}
  (the script sets \passthrough{\lstinline!MPLBACKEND=Agg!} itself);
  stdout is exactly the one written strip path,
  \passthrough{\lstinline!quadmath/output/figures/animation\_frames\_strip.png!}.
\item
  Markdown:
  \passthrough{\lstinline!uv run python quadmath/scripts/validate\_markdown.py!}.
\end{itemize}

\subsection{Cross-references}\label{cross-references-3}

\begin{itemize}
\tightlist
\item
  Site geometry, normalization, and the embedding:
  \href{03_quadray_methods.md}{Quadray Methods}.
\item
  Shell enumeration and nearest-site search:
  \href{13_lattice_tooling.md}{Lattice Tooling}.
\item
  The sliced learner: \href{11_ivm_field_learning.md}{Static IVM Field
  Learning}.
\item
  The stripped dynamics: \href{12_ivm_dynamics.md}{Dynamics and Learning
  on the IVM Lattice}.
\end{itemize}

\newpage

\section{Benchmarks and Statistics}\label{benchmarks-and-statistics}

\subsection{Overview}\label{overview-10}

This section documents the measurement surface behind the preceding
chapters. It is split across three modules with one job each:
\passthrough{\lstinline!src/quadmath/stats/benchmarks.py!} measures (a
small \passthrough{\lstinline!perf\_counter!} timing harness over the
core numerical surfaces),
\passthrough{\lstinline!src/quadmath/stats/statistics.py!} analyzes
(resampling inference, effect sizes, multiplicity correction, and
scaling fits, on standalone numpy), and
\passthrough{\lstinline!src/quadmath/viz/vis\_stats.py!} renders
(input-agnostic matplotlib primitives; the figures they produce are
collected in \passthrough{\lstinline!18\_stats\_gallery.md!}). The
measured workloads are the four surfaces the manuscript already treats:
quadray conversions (\passthrough{\lstinline!quadray.py!}, conventions
in \passthrough{\lstinline!SPEC.md!}), shell enumeration
(\passthrough{\lstinline!omni\_numbering.py!},
\passthrough{\lstinline!13\_lattice\_tooling.md!}), nearest-site search
(\passthrough{\lstinline!lattice\_search.py!}), and field fitting
(\passthrough{\lstinline!ivm\_field.py!},
\passthrough{\lstinline!11\_ivm\_field\_learning.md!}). Everything here
is deterministic: every random draw comes from a seeded
\passthrough{\lstinline!numpy.random.default\_rng!}, and every timing
routine is fully specified by its arguments, so the numbers below
reproduce bit-for-bit.

\subsection{Timing methodology}\label{timing-methodology}

\passthrough{\lstinline!time\_callable(fn, *, trials=5, warmup=1)!}
measures wall-clock duration with
\passthrough{\lstinline!time.perf\_counter!}, the highest-resolution
monotonic clock available. It runs one warmup call first and discards it
(first-call effects: module imports, key caches), then times
\passthrough{\lstinline!trials!} calls and reports the summary
statistics

\begin{equation}
\label{eq:bench-mean}
\bar{t} \;=\; \frac{1}{T}\sum_{i=1}^{T} t_i ,
\qquad t_m \;=\; \operatorname{median}(t_1,\dots,t_T),
\end{equation}

\begin{equation}
\label{eq:bench-percentile}
t_{95} \;=\; Q_{0.95}(t_1,\dots,t_T),
\end{equation}

where \(T\) is the number of trials and \(Q_p\) the empirical
\(p\)-quantile. The median is robust to a single outlier trial; the 95th
percentile bounds the tail that a user would actually feel. Results are
returned as \passthrough{\lstinline!BenchRow!}, a frozen dataclass with
fields \passthrough{\lstinline!name!}, \passthrough{\lstinline!n!},
\passthrough{\lstinline!trials!}, \passthrough{\lstinline!total\_s!},
\passthrough{\lstinline!mean\_s!}, \passthrough{\lstinline!median\_s!},
\passthrough{\lstinline!p95\_s!}, and
\passthrough{\lstinline!ops\_per\_s!}, plus
\passthrough{\lstinline!as\_dict()!} for plain data. Here
\passthrough{\lstinline!n!} is the per-call workload size (conversions
per call, shell depth, sites or queries per call), and throughput is

\begin{equation}
\label{eq:bench-ops}
\text{ops\_per\_s} \;=\; \frac{n}{\bar{t}} .
\end{equation}

Four constructors wrap the core surfaces with these defaults:
\passthrough{\lstinline!bench\_conversions(n=200, trials=5)!} times
round-trip quadray/embedding conversions,
\passthrough{\lstinline!bench\_shell\_enumeration(k\_max=4, trials=5)!}
times shell enumeration through
\passthrough{\lstinline!omni\_numbering.generate\_shell!},
\passthrough{\lstinline!bench\_lattice\_search(n\_sites=200, queries=50, trials=5)!}
times \passthrough{\lstinline!lattice\_search!} nearest-site queries
over a random ball of sites, and
\passthrough{\lstinline!bench\_field\_fit(n\_sites=64, trials=3)!} times
a synthetic field fit through \passthrough{\lstinline!IVMField.learn!}.
\passthrough{\lstinline!run\_all()!} executes all four with the shared
defaults collected in the module constant
\passthrough{\lstinline!BENCH\_DEFAULTS!} and returns the list of
\passthrough{\lstinline!BenchRow!} records;
\passthrough{\lstinline!summary\_table(rows)!} renders them as a
fixed-width text table. The harness measures only; it asserts nothing
--- timing distributions belong to analysis, not to test assertions.

\subsection{Percentile bootstrap}\label{percentile-bootstrap}

\passthrough{\lstinline!bootstrap\_ci(x, stat=np.mean, *, iters=2000, seed=0, alpha=0.05)!}
gives a confidence interval for any statistic \(\hat{\theta}\) by
resampling the data with replacement. With \(n = |x|\) and a seeded
\passthrough{\lstinline!numpy.random.default\_rng(seed)!}, iteration
\(b\) draws resampling indices
\(i_{bj} \sim \mathrm{Uniform}\{0,\dots,n-1\}\) in one vectorized call
(\passthrough{\lstinline!rng.integers(0, n, size=(iters, n))!}) and
evaluates

\begin{equation}
\label{eq:stat-bootstrap}
\hat{\theta}^{*}_{b} \;=\; \hat{\theta}\big(x_{i_{b1}},\dots,x_{i_{bn}}\big),
\qquad
\mathrm{CI}_{1-\alpha}
\;=\;
\Big[\, Q_{\alpha/2}\big(\hat{\theta}^{*}_{1..B}\big),\;
Q_{1-\alpha/2}\big(\hat{\theta}^{*}_{1..B}\big) \Big],
\end{equation}

with \(B\) = \passthrough{\lstinline!iters!}. The percentile interval
needs no distributional assumption about \(\hat{\theta}\) --- only that
the empirical resampling distribution approximates its sampling
distribution. All randomness is consumed from the seeded generator, so
the interval is exactly reproducible.

\subsection{Pooled permutation tests}\label{pooled-permutation-tests}

\passthrough{\lstinline!permutation\_test(a, b, *, iters=2000, seed=0, alternative="two-sided")!}
tests whether two samples differ in location without assuming a
distribution. Both samples are pooled; the observed statistic is
\(o = \big|\bar{a} - \bar{b}\big|\); each of the
\passthrough{\lstinline!iters!} iterations draws
\passthrough{\lstinline!rng.permutation!} of the pool, splits it at
\(\lvert a\rvert\) into a fake \(a\) and \(b\), and recomputes the
difference. The two-sided p-value uses the add-one convention

\begin{equation}
\label{eq:stat-permutation}
p \;=\;
\frac{1 + \#\big\{ r :\; |d_r| \,\ge\, o \big\}}{\text{iters} + 1},
\end{equation}

where \(d_r\) is the difference of permutation \(r\). The \(+1\) in
numerator and denominator counts the observed arrangement itself and
guarantees \(p \ge 1/(\text{iters}+1) > 0\): a permutation test can
never report an impossible zero p-value, and the estimate is
conservative. \passthrough{\lstinline!alternative="greater"!} and
\passthrough{\lstinline!"less"!} use the same convention with signed
comparisons of \(d_r\) against the signed observed difference.

\subsection{Effect size, multiplicity, and
scaling}\label{effect-size-multiplicity-and-scaling}

\passthrough{\lstinline!cohens\_d(a, b)!} reports the standardized mean
difference

\begin{equation}
\label{eq:stat-cohens}
d \;=\; \frac{\bar{a} - \bar{b}}{s_p},
\qquad
s_p \;=\;
\sqrt{\frac{(n_a - 1)\,s_a^2 + (n_b - 1)\,s_b^2}{n_a + n_b - 2}},
\end{equation}

with \(s_a, s_b\) the unbiased sample standard deviations. A p-value
says whether a difference is distinguishable from noise; \(d\) says how
large the difference is, which is the question that matters when
comparing benchmark configurations.

\passthrough{\lstinline!p\_adjust\_bonferroni(pvals)!} controls the
family-wise error rate when \(m\) hypotheses are tested at once:

\begin{equation}
\label{eq:stat-bonferroni}
p'_i \;=\; \min\!\big(1,\; m\, p_i\big).
\end{equation}

Each \(p'_i\) can be read as a valid p-value at level \(\alpha\) for the
whole family --- conservative, but assumption-free.

\passthrough{\lstinline!scaling\_fit(sizes, times)!} fits ordinary least
squares on the log-log data,

\begin{equation}
\label{eq:stat-scaling}
\log t \;=\; \beta_1 \log n + \beta_0,
\end{equation}

and returns \passthrough{\lstinline!(slope, intercept, r2)!}. The slope
is the empirical complexity exponent: \(\beta_1 \approx 1\) indicates
linear work in the input size \(n\), \(\beta_1 \approx 2\) quadratic,
and the coefficient of determination \(r^2\) measures how well the power
law describes the measurements.

\subsection{Jackknife, FDR, Welch, and circular
statistics}\label{sec:jackknife_fdr_welch_circular}

\passthrough{\lstinline!jackknife\_ci(x, stat=np.mean, alpha=0.05)!}
produces both a confidence interval and a bias estimate for any scalar
statistic without consuming a single random draw. Where the percentile
bootstrap of \eqref{eq:stat-bootstrap} resamples with replacement, the
jackknife deletes one observation at a time: each of the \(n\)
leave-one-out samples is scored, and the spread of those scores
determines the interval. With \(\hat{\theta}\) the full-sample statistic
and \(\hat{\theta}_{(i)}\) the statistic on the sample with observation
\(i\) deleted, the returned triple
\passthrough{\lstinline!(low, high, bias)!} is

\begin{equation}
\label{eq:stat-jackknife}
\begin{aligned}
\mathrm{bias} &= (n-1)\big(\bar{\theta}_{(\cdot)} - \hat{\theta}\big), &
\hat{\theta}_{\text{corr}} &= \hat{\theta} - \mathrm{bias}, \\
\mathrm{se}_{\text{jack}} &=
\sqrt{\tfrac{n-1}{n} \sum_{i=1}^{n} \big(\hat{\theta}_{(i)} -
\bar{\theta}_{(\cdot)}\big)^{2}}, &
\mathrm{CI}_{1-\alpha} &=
\hat{\theta} \pm z_{1-\alpha/2}\,\mathrm{se}_{\text{jack}},
\end{aligned}
\end{equation}

where \(\bar{\theta}_{(\cdot)}\) is the mean of the leave-one-out
scores. The bias-corrected estimate \(\hat{\theta}_{\text{corr}}\)
removes the first-order bias that the jackknife detects in curved
statistics; the interval itself stays centered on the uncorrected
\(\hat{\theta}\). The normal quantile \(z\) comes from bisecting
\passthrough{\lstinline!math.erfc(z / sqrt(2))!}, which decreases
monotonically from \(2\) to \(0\) as \(z\) runs from \(-\infty\) to
\(\infty\), so one hundred halvings of the bracket \([-40, 40]\) pin
\(z\) to double precision. This deterministic erfc-bisection quantile
needs no seed and no lookup table, so repeated calls with identical
arguments return bit-identical results --- the same reproducibility
contract as the seeded routines above, but with no generator to seed at
all.

\passthrough{\lstinline!benjamini\_hochberg(pvals)!} controls a
different error rate from the Bonferroni map of
\eqref{eq:stat-bonferroni}: instead of the family-wise error rate it
controls the expected proportion of false discoveries among rejections,
the false discovery rate. The step-up procedure sorts the \(m\) p-values
ascending, forms the raw values \(m\,p_{(i)} / i\) for ranks
\(i = 1,\dots,m\), and then enforces monotonicity with a running minimum
taken from the largest rank backwards,

\begin{equation}
\label{eq:stat-bh}
p'_{(i)} \;=\;
\min\!\Big(1,\; \min_{j \ge i}\, \frac{m\,p_{(j)}}{j}\Big),
\qquad
p'_{(m)} \;=\; \min(1,\, p_{(m)}),
\end{equation}

so the adjusted values are monotone non-decreasing in the sorted
p-values and never exceed the largest raw p-value. The backward
cumulative minimum is the step that the naive \(m\,p_{(i)}/i\) mapping
misses: without it, a large early-rank raw value could exceed a smaller
later-rank one, breaking monotonicity and invalidating the FDR
guarantee. The stable argsort of the input is recorded so the adjusted
values can be mapped back to the original input order, and the final cap
at \(1.0\) keeps the output a valid p-value vector.

\passthrough{\lstinline!welch\_t\_test(a, b, alternative="two-sided")!}
compares two sample means without assuming equal variances, the
assumption that the pooled standard deviation \(s_p\) of
\passthrough{\lstinline!cohens\_d!} makes. Writing \(w_a = s_a^2/n_a\)
and \(w_b = s_b^2/n_b\), the statistic and its Welch--Satterthwaite
degrees of freedom are

\begin{equation}
\label{eq:stat-welch}
t \;=\; \frac{\bar{a} - \bar{b}}{\sqrt{w_a + w_b}},
\qquad
\nu \;=\;
\frac{(w_a + w_b)^{2}}
{\dfrac{w_a^{2}}{n_a - 1} + \dfrac{w_b^{2}}{n_b - 1}} .
\end{equation}

The p-value is evaluated from the exact Student-\(t\) tails rather than
a normal approximation, and it needs no \passthrough{\lstinline!scipy!}:
the two-sided tail uses the identity

\begin{equation}
\label{eq:stat-tail}
\mathrm{P}\big(|T| \ge |t|\big) \;=\; I_{x}\!\big(\tfrac{\nu}{2},
\tfrac{1}{2}\big),
\qquad
x \;=\; \frac{\nu}{\nu + t^{2}},
\end{equation}

where \(I_x(a, b)\) is the regularized incomplete beta function computed
by a Lentz continued fraction (with the standard symmetry swap above
\(x = (a+1)/(a+b+2)\), an \(\exp\)-\(\log\Gamma\) front factor, and
guard floors against vanishing denominators).
\passthrough{\lstinline!"greater"!} and \passthrough{\lstinline!"less"!}
reuse the same two-sided tail, halved on the appropriate side of
\(t = 0\). The degenerate branches matter for constant inputs: when both
samples are constant the denominator of \eqref{eq:stat-welch} vanishes,
and equal means give \(t = 0.0\) (p-value \(1.0\) two-sided) while
different means give a signed infinite \(t\) whose tail probabilities
are exactly \(0\) or \(1\), evaluated against the finite pooled fallback
\(\nu = n_a + n_b - 2\) so the tails stay well-defined.

\passthrough{\lstinline!rotation\_stats(angles)!} treats angles in
radians as points on a circle rather than points on a line, where
\(359°\) and \(1°\) are neighbors:

\begin{equation}
\label{eq:stat-circular}
\bar{\theta} \;=\; \operatorname{atan2}\!\Big(
\tfrac{1}{n}\textstyle\sum_{j} \sin\theta_{j},\;
\tfrac{1}{n}\textstyle\sum_{j} \cos\theta_{j}\Big),
\qquad
R \;=\; \Big|\tfrac{1}{n}\textstyle\sum_{j} e^{i\theta_{j}}\Big| .
\end{equation}

The circular mean is the \passthrough{\lstinline!atan2!} of the averaged
sine and cosine, wrapped into \((-\pi, \pi]\), and the mean resultant
length \(R \in [0, 1]\) measures concentration: \(R = 1\) only when
every angle agrees, \(R = 0\) when the unit vectors cancel completely.
The returned dictionary gives the circular variance \(1 - R\) in
\([0, 1]\), and \passthrough{\lstinline!variance\_2pi!} \(= 2\,(1 - R)\)
under the convention for angles on the full \([0, 2\pi)\) circle, which
ranges over \([0, 2]\) and approaches the familiar linear variance for
tightly clustered angles. Because sine and cosine are \(2\pi\)-periodic,
wrapping the angles into any full circle leaves every returned value
unchanged; when the resultant vanishes the mean degenerates to whatever
\passthrough{\lstinline!atan2!} returns for the cancelled components.
Like everything else in
\passthrough{\lstinline!src/quadmath/stats/statistics.py!}, all four
routines are pure
\passthrough{\lstinline!numpy!}-plus-\passthrough{\lstinline!math!}
computations: no randomness, no optional scientific stack, and
bit-identical repeats for identical arguments.

\subsection{A reproducible example}\label{a-reproducible-example}

Both workhorses above are fully determined by their arguments, so the
following numbers are exact for any reader. The first call computes the
95\% bootstrap confidence interval of the mean of the first seven
Fibonacci numbers; the second asks whether the shifted ranges \(1..10\)
and \(11..20\) differ in location:

\begin{lstlisting}[language=Python]
import numpy as np

from quadmath.stats.statistics import bootstrap_ci, permutation_test

x = [2, 3, 5, 8, 13, 21, 34]
lo, hi = bootstrap_ci(x, np.mean, iters=2000, seed=0, alpha=0.05)
# (5.142857142857143, 20.571428571428573)

a = list(range(1, 11))
b = list(range(11, 21))
p = permutation_test(a, b, iters=2000, seed=0, alternative="two-sided")
# 0.0014992503748125937
\end{lstlisting}

With
\passthrough{\lstinline!bootstrap\_ci(x, np.mean, iters=2000, seed=0, alpha=0.05)!}
the mean of \(x\) is \(86/7 \approx 12.285714\) and the 95\% percentile
bootstrap interval is \([5.142857142857143,\; 20.571428571428573]\).
Both endpoints are exact multiples of \(1/7\): the mean of any 7-value
resample with repetition is a multiple of \(1/7\), so the percentile
grid is discrete rather than dense. With
\passthrough{\lstinline!permutation\_test(a, b, iters=2000, seed=0, alternative="two-sided")!}
on \(a = 1,\dots,10\) and \(b = 11,\dots,20\), the observed difference
is \(o = |\bar{a} - \bar{b}| = 10\) exactly, and exactly 2 of the 2000
seeded permutations reach it, giving

\begin{equation}
\label{eq:stat-example-p}
p \;=\; \frac{2 + 1}{2000 + 1} \;=\; \frac{3}{2001}
\;\approx\; 0.00149925 .
\end{equation}

The count of 2 deserves a comment. Among the
\(\binom{20}{10} = 184{,}756\) possible splits of the pool, exactly two
achieve \(|d_r| = 10\) --- the observed arrangement (\(1..10\) against
\(11..20\)) and its complement --- so the expected hit count over 2000
random draws is only \(2000 \cdot 2/184756 \approx 0.02\). The add-one
floor of \eqref{eq:stat-permutation} is \(1/2001 \approx 0.0005\), which
is what a typical seed reports here; the stream seeded with
\passthrough{\lstinline!seed=0!} happens to draw both exact-maximum
arrangements (at iterations 70 and 1146 of the 2000), so the count is 2
and the reported p-value is \(3/2001\). The example illustrates the
convention rather than a typical run: with any other seed the same call
reproduces the floor value, and in either case the test correctly never
reports the impossible \(p = 0\).

The two interval families meet on common ground in a second example. Two
synthetic lattice-error samples with a planted mean shift are drawn once
from a seeded generator: routine A contributes \(n = 64\) error
residuals centered on \(0.0\) and routine B \(64\) residuals whose
population mean sits \(0.08\) higher, both with scale \(0.15\) --- the
unit being whatever the error residual measures. For each sample mean,
the 95\% percentile bootstrap interval of \eqref{eq:stat-bootstrap}
(\(B = 2000\) resamples, re-seeded at 17) and the jackknife interval of
\eqref{eq:stat-jackknife} (which consumes no randomness at all) are
computed, and the script
\passthrough{\lstinline!quadmath/scripts/stats\_diagnostics\_gallery.py!}
draws all four intervals in a single
\passthrough{\lstinline!plot\_ci\_bars!} call:

\begin{lstlisting}[language=Python]
import numpy as np

from quadmath.stats.statistics import bootstrap_ci, jackknife_ci

rng = np.random.default_rng(17)
a = rng.normal(0.0, 0.15, size=64)
b = rng.normal(0.08, 0.15, size=64)
bootstrap_ci(a, np.mean, iters=2000, seed=17, alpha=0.05)
# (-0.07609532416913227, -0.0004865391750512773)
jackknife_ci(a, np.mean, alpha=0.05)[:2]
# (-0.07682589122849623, 0.0013292867175937334)
\end{lstlisting}

What the two families assume is where they differ. The percentile
bootstrap needs no shape assumption at all --- only that the empirical
resampling distribution of \eqref{eq:stat-bootstrap} approximates the
sampling distribution of \(\hat{\theta}\) --- while the jackknife
interval is the normal approximation \(\hat{\theta} \pm
z_{1-\alpha/2}\,\mathrm{se}_{\text{jack}}\) built from the \(n\)
leave-one-out scores of \eqref{eq:stat-jackknife}. On the sample mean
with \(n = 64\) the two agree to within a whisker: routine A's bootstrap
interval is \([-0.076095,\,-0.000487]\) against the jackknife's
\([-0.076826,\,0.001329]\), and routine B's is \([0.056788,\,0.132418]\)
against \([0.056011,\,0.130639]\) (sample means \(-0.037748\) and
\(0.093325\); the jackknife bias estimates sit at floating-point zero,
as \eqref{eq:stat-jackknife} predicts for a linear statistic). The
figure records the comparison: routine B's planted shift shows up as
both of its intervals clearing the dashed zero line, while routine A
hugs the line so tightly that the bootstrap upper endpoint stops
\(0.0005\) short of it and the jackknife upper endpoint crosses it by
\(0.0013\) --- a hair-width disagreement that is the honest lesson of
the panel. Coverage verdicts this close to the boundary are
method-sensitive; a Welch test on the same samples reports
\(t \approx -4.75\) with \(p \approx 5.4 \times 10^{-6}\)
(\eqref{eq:stat-welch}), in agreement with the shift both interval
families detect.

\begin{figure}
\centering
\pandocbounded{\includegraphics[keepaspectratio,alt={Bootstrap versus jackknife 95\% confidence intervals for two lattice-error samples. Sample means (dots) of two seeded synthetic error samples --- n = 64 per group from numpy.random.default\_rng(17), scale 0.15, planted mean shift 0.08 in the B population --- each carrying a 95\% percentile bootstrap interval (bootstrap\_ci, 2000 resamples, seed 17, \textbackslash alpha = 0.05) and a deterministic normal-approximation jackknife interval (jackknife\_ci, \textbackslash alpha = 0.05), drawn together in one plot\_ci\_bars call; the dashed line marks the unbiased population mean 0.0. Routine B's intervals exclude 0 while routine A's straddle it, and the two CI families agree within whisker width on both samples. Regenerate: uv run python quadmath/scripts/stats\_diagnostics\_gallery.py.}]{../output/figures/stats_ci_comparison.png}}
\caption{\textbf{Bootstrap versus jackknife 95\% confidence intervals
for two lattice-error samples.} Sample means (dots) of two seeded
synthetic error samples --- \(n = 64\) per group from
\protect\passthrough{\lstinline!numpy.random.default\_rng(17)!}, scale
0.15, planted mean shift 0.08 in the B population --- each carrying a
95\% percentile bootstrap interval
(\protect\passthrough{\lstinline!bootstrap\_ci!}, 2000 resamples, seed
17, \(\alpha = 0.05\)) and a deterministic normal-approximation
jackknife interval (\protect\passthrough{\lstinline!jackknife\_ci!},
\(\alpha = 0.05\)), drawn together in one
\protect\passthrough{\lstinline!plot\_ci\_bars!} call; the dashed line
marks the unbiased population mean 0.0. Routine B's intervals exclude 0
while routine A's straddle it, and the two CI families agree within
whisker width on both samples. Regenerate:
\protect\passthrough{\lstinline!uv run python quadmath/scripts/stats\_diagnostics\_gallery.py!}.}
\end{figure}

\subsection{Cross-references}\label{cross-references-4}

\begin{itemize}
\tightlist
\item
  Timing harness, \passthrough{\lstinline!BenchRow!}, and the four
  benchmark constructors:
  \passthrough{\lstinline!src/quadmath/stats/benchmarks.py!}.
\item
  Bootstrap, permutation tests, \passthrough{\lstinline!cohens\_d!},
  \passthrough{\lstinline!p\_adjust\_bonferroni!},
  \passthrough{\lstinline!scaling\_fit!}:
  \passthrough{\lstinline!src/quadmath/stats/statistics.py!}.
\item
  The figure primitives built on these results:
  \passthrough{\lstinline!18\_stats\_gallery.md!}.
\item
  The interval-comparison script and figure behind the second example:
  \passthrough{\lstinline!quadmath/scripts/stats\_diagnostics\_gallery.py!}
  (\passthrough{\lstinline!quadmath/output/figures/stats\_ci\_comparison.png!}).
\item
  Measured surfaces: \passthrough{\lstinline!quadray.py!} (conventions
  in \passthrough{\lstinline!SPEC.md!}),
  \passthrough{\lstinline!omni\_numbering.py!} and
  \passthrough{\lstinline!lattice\_search.py!}
  (\passthrough{\lstinline!13\_lattice\_tooling.md!}),
  \passthrough{\lstinline!ivm\_field.py!}
  (\passthrough{\lstinline!11\_ivm\_field\_learning.md!}).
\item
  Sibling gallery of the lattice layer:
  \passthrough{\lstinline!16\_lattice\_gallery.md!}.
\end{itemize}

\newpage

\section{Statistics and Scaling
Gallery}\label{statistics-and-scaling-gallery}

\subsection{Overview}\label{overview-11}

This section is the figure surface of the benchmarks-and-statistics
layer: every rendering primitive used by section
\passthrough{\lstinline!17\_benchmarks\_statistics.md!} lives in
\passthrough{\lstinline!src/quadmath/viz/vis\_stats.py!}
(\passthrough{\lstinline!plot\_latency\_hist!},
\passthrough{\lstinline!plot\_scaling\_loglog!},
\passthrough{\lstinline!plot\_ci\_bars!},
\passthrough{\lstinline!plot\_ecdf!}). The primitives are input-agnostic
--- plain numpy arrays in, artists out; they import nothing from the
other \passthrough{\lstinline!src/!} modules and receive their
matplotlib axes explicitly, so panels compose inside caller-owned
figures and nothing renders at import time. The command-line entry is
\passthrough{\lstinline!quadmath/scripts/stats\_gallery.py!} --- a thin
orchestrator that sets a headless backend and a fixed seed
(\passthrough{\lstinline!seed=32!}), delegates to the gallery function
in \passthrough{\lstinline!src/quadmath/viz/vis\_stats.py!}
(thin-orchestrator contract of
\passthrough{\lstinline!quadmath/scripts/AGENTS.md!}), and prints each
written path on its own line: those stdout lines are the
\passthrough{\lstinline!make\_all\_figures!} manifest contract. The
gallery writes exactly four PNGs ---
\passthrough{\lstinline!stats\_gallery\_latency.png!},
\passthrough{\lstinline!stats\_gallery\_scaling.png!},
\passthrough{\lstinline!stats\_gallery\_ci.png!},
\passthrough{\lstinline!stats\_gallery\_ecdf.png!} --- and every random
draw comes from the fixed seed, so the whole set is deterministic.

\subsection{Latency histogram}\label{latency-histogram}

\passthrough{\lstinline!plot\_latency\_hist(times, ax=None, *, bins=20, title=...)!}
renders the distribution of per-call wall-clock durations measured by
\passthrough{\lstinline!src/quadmath/stats/benchmarks.py::time\_callable!}
(trials summary in
\passthrough{\lstinline!17\_benchmarks\_statistics.md!}) as a histogram
with \passthrough{\lstinline!bins!} bins on the given axes. The shape of
the distribution carries what the mean alone hides: a tight spike means
the harness saw a stable workload, while a heavy right tail is exactly
what the \passthrough{\lstinline!median\_s!} /
\passthrough{\lstinline!p95\_s!} fields of
\passthrough{\lstinline!BenchRow!} are there to expose. The gallery
panel histograms a deterministic synthetic stand-in for such a sample:
256 draws of a lognormal distribution (mean 0, \(\sigma = 0.6\), clipped
to \([0.05, 5.0]\) in seconds-scale units) taken from the single
fixed-seed \passthrough{\lstinline!numpy.random.default\_rng(32)!}
generator, so re-renders are byte-identical.

\begin{figure}
\centering
\pandocbounded{\includegraphics[keepaspectratio,alt={Latency distribution. Histogram (20 bins) of a synthetic 256-draw lognormal latency sample (mean 0, \textbackslash sigma = 0.6) clipped to {[}0.05, 5.0{]} in seconds-scale units, drawn from the fixed-seed numpy.random.default\_rng(32) generator of stats\_gallery.py; the dashed line marks the sample mean. Rendered by plot\_latency\_hist; reproduced by uv run python quadmath/scripts/stats\_gallery.py. The spread and right tail complement the mean\_s, median\_s, and p95\_s fields of BenchRow, whose per-call times time\_callable measures with time.perf\_counter after one discarded warmup call.}]{../output/figures/stats_gallery_latency.png}}
\caption{\textbf{Latency distribution.} Histogram (20 bins) of a
synthetic 256-draw lognormal latency sample (mean 0, \(\sigma = 0.6\))
clipped to {[}0.05, 5.0{]} in seconds-scale units, drawn from the
fixed-seed
\protect\passthrough{\lstinline!numpy.random.default\_rng(32)!}
generator of \protect\passthrough{\lstinline!stats\_gallery.py!}; the
dashed line marks the sample mean. Rendered by
\protect\passthrough{\lstinline!plot\_latency\_hist!}; reproduced by
\protect\passthrough{\lstinline!uv run python quadmath/scripts/stats\_gallery.py!}.
The spread and right tail complement the
\protect\passthrough{\lstinline!mean\_s!},
\protect\passthrough{\lstinline!median\_s!}, and
\protect\passthrough{\lstinline!p95\_s!} fields of
\protect\passthrough{\lstinline!BenchRow!}, whose per-call times
\protect\passthrough{\lstinline!time\_callable!} measures with
\protect\passthrough{\lstinline!time.perf\_counter!} after one discarded
warmup call.}
\end{figure}

\subsection{Scaling fit}\label{scaling-fit}

\passthrough{\lstinline!plot\_scaling\_loglog(sizes, times, ax=None, *, title=...)!}
plots measured workload sizes against durations on log-log axes and
overlays the least-squares power law of
\passthrough{\lstinline!src/quadmath/stats/statistics.py::scaling\_fit!}
--- the slope is the empirical complexity exponent \(\beta_1\) of
\eqref{eq:stat-scaling}. On log-log axes a power law is a straight line,
so the panel shows at a glance whether a lattice routine scales
linearly, quadratically, or worse, and how tightly the data follow the
law (\(r^2\)). The gallery panel fits a synthetic sweep with known
ground truth: durations \(t \approx 3\,n^{1.35}\,(1 + \varepsilon)\)
with \(\varepsilon \sim \mathcal{N}(0, 0.02^{2})\) relative noise over
sizes \(n \in \{8, 16, 32, 64, 128, 256\}\), drawn from the same
fixed-seed generator (times rounded to four decimals for byte
stability), so the fit should recover the planted exponent 1.35.

\begin{figure}
\centering
\pandocbounded{\includegraphics[keepaspectratio,alt={Log-log scaling fit. Workload sizes n \textbackslash in \textbackslash\{8, 16, 32, 64, 128, 256\textbackslash\} (dimensionless units) against synthetic durations t \textbackslash approx 3\textbackslash,n\^{}\{1.35\}\textbackslash,(1 + \textbackslash varepsilon) with \textbackslash varepsilon \textbackslash sim \textbackslash mathcal\{N\}(0, 0.02\^{}\{2\}) relative noise, times rounded to four decimals, all drawn from the fixed-seed numpy.random.default\_rng(32) generator; the least-squares power law from scaling\_fit is overlaid by plot\_scaling\_loglog, and its fitted slope --- annotated in the legend --- should recover the planted exponent 1.35, the empirical complexity exponent \textbackslash beta\_1 of .}]{../output/figures/stats_gallery_scaling.png}}
\caption{\textbf{Log-log scaling fit.} Workload sizes
\(n \in \{8, 16, 32, 64, 128, 256\}\) (dimensionless units) against
synthetic durations \(t \approx 3\,n^{1.35}\,(1 + \varepsilon)\) with
\(\varepsilon \sim \mathcal{N}(0, 0.02^{2})\) relative noise, times
rounded to four decimals, all drawn from the fixed-seed
\protect\passthrough{\lstinline!numpy.random.default\_rng(32)!}
generator; the least-squares power law from
\protect\passthrough{\lstinline!scaling\_fit!} is overlaid by
\protect\passthrough{\lstinline!plot\_scaling\_loglog!}, and its fitted
slope --- annotated in the legend --- should recover the planted
exponent 1.35, the empirical complexity exponent \(\beta_1\) of
\eqref{eq:stat-scaling}.}
\end{figure}

\subsection{Confidence intervals}\label{confidence-intervals}

\passthrough{\lstinline!plot\_ci\_bars(labels, means, lows, highs, ax=None, *, title=...)!}
renders paired estimates as a bar-and-whisker chart: one bar per label
at its point estimate, with error bars spanning the low and high ends of
a percentile bootstrap interval
\passthrough{\lstinline!src/quadmath/stats/statistics.py::bootstrap\_ci!}
(\eqref{eq:stat-bootstrap}). Overlapping intervals visually encode the
same information a permutation test quantifies
(\eqref{eq:stat-permutation}): two conditions whose intervals do not
overlap are the ones the pooled test flags. The gallery panel shows
three benchmark estimates (\passthrough{\lstinline!to\_xyz!},
\passthrough{\lstinline!quadray\_from\_xyz!},
\passthrough{\lstinline!shell\_enum!}) with symmetric intervals whose
half-widths are seeded uniform draws on \([0.05, 0.25]\) from the same
fixed-seed \passthrough{\lstinline!numpy.random.default\_rng(32)!}
generator: the primitive renders whatever \passthrough{\lstinline!lows!}
/ \passthrough{\lstinline!highs!} it is given, and in a measured
analysis those bounds are the percentile endpoints
\passthrough{\lstinline!bootstrap\_ci!} returns
(\eqref{eq:stat-bootstrap}).

\begin{figure}
\centering
\pandocbounded{\includegraphics[keepaspectratio,alt={Point estimates with confidence intervals by condition. Three benchmark estimates (to\_xyz, quadray\_from\_xyz, shell\_enum) in synthetic estimate units, rendered by plot\_ci\_bars as points with symmetric error bars whose half-widths are uniform draws on {[}0.05, 0.25{]} from the fixed-seed numpy.random.default\_rng(32) generator; in a measured analysis the bar ends would be the 95\% percentile bootstrap endpoints from bootstrap\_ci (2000 resamples, ), whose non-overlap is the visual cue the permutation test of  quantifies.}]{../output/figures/stats_gallery_ci.png}}
\caption{\textbf{Point estimates with confidence intervals by
condition.} Three benchmark estimates
(\protect\passthrough{\lstinline!to\_xyz!},
\protect\passthrough{\lstinline!quadray\_from\_xyz!},
\protect\passthrough{\lstinline!shell\_enum!}) in synthetic estimate
units, rendered by \protect\passthrough{\lstinline!plot\_ci\_bars!} as
points with symmetric error bars whose half-widths are uniform draws on
{[}0.05, 0.25{]} from the fixed-seed
\protect\passthrough{\lstinline!numpy.random.default\_rng(32)!}
generator; in a measured analysis the bar ends would be the 95\%
percentile bootstrap endpoints from
\protect\passthrough{\lstinline!bootstrap\_ci!} (2000 resamples,
\eqref{eq:stat-bootstrap}), whose non-overlap is the visual cue the
permutation test of \eqref{eq:stat-permutation} quantifies.}
\end{figure}

\subsection{Empirical CDF}\label{empirical-cdf}

\passthrough{\lstinline!plot\_ecdf(values, ax=None, *, title=...)!}
renders the empirical distribution function

\begin{equation}
\label{eq:vis-ecdf}
\hat{F}(t) \;=\; \frac{1}{n}\sum_{i=1}^{n}
\mathbf{1}\{t_i \le t\},
\end{equation}

a monotone step function that jumps by \(1/n\) at every observed value
--- the exact distribution of the sample, with no binning choices at
all. For timing data the ECDF answers directly what a percentile summary
rounds off: the curve passes through the empirical median at height
\(0.5\) and through the 95th percentile at height \(0.95\), so
\passthrough{\lstinline!median\_s!} and \passthrough{\lstinline!p95\_s!}
are readable off the plot. The gallery panel draws the ECDF of the same
256-draw clipped lognormal latency sample as the histogram panel above,
so the two panels read as one distribution in two renderings.

\begin{figure}
\centering
\pandocbounded{\includegraphics[keepaspectratio,alt={Empirical cumulative distribution of the latency sample. Exact step-function ECDF \textbackslash hat\{F\}(t) of the same 256-draw clipped lognormal sample as the histogram panel above (fixed seed 32, latency in seconds-scale units), rendered by plot\_ecdf; the empirical median (height 0.5) and 95th percentile (height 0.95) are read directly off the curve.}]{../output/figures/stats_gallery_ecdf.png}}
\caption{\textbf{Empirical cumulative distribution of the latency
sample.} Exact step-function ECDF \(\hat{F}(t)\) of the same 256-draw
clipped lognormal sample as the histogram panel above (fixed seed 32,
latency in seconds-scale units), rendered by
\protect\passthrough{\lstinline!plot\_ecdf!}; the empirical median
(height 0.5) and 95th percentile (height 0.95) are read directly off the
curve.}
\end{figure}

\subsection{Reproducibility and test
contract}\label{reproducibility-and-test-contract-2}

\begin{itemize}
\tightlist
\item
  The gallery function in
  \passthrough{\lstinline!src/quadmath/viz/vis\_stats.py!} writes
  exactly the four files listed above, in that order, into the output
  directory (creating it when missing) and returns their paths. Every
  panel is fully seeded --- the script fixes
  \passthrough{\lstinline!seed=32!} --- and the PNGs carry no timestamp
  metadata, so re-renders are byte-identical.
\item
  Tests: \passthrough{\lstinline!tests/test\_vis\_stats.py!} pins the
  contract with deterministic assertions, no mocks, and no pixel diffs
  --- artist placement on caller-provided axes, input-agnostic behavior
  over plain arrays, and the invariants of each primitive (monotone
  non-decreasing ECDF steps, CI whisker endpoints matching the
  lows/highs inputs). Consistent with
  \passthrough{\lstinline!17\_benchmarks\_statistics.md!}, no test
  asserts wall-clock timing values.
\item
  Figure regeneration:
  \passthrough{\lstinline!uv run python quadmath/scripts/stats\_gallery.py!}
  (the script sets \passthrough{\lstinline!MPLBACKEND=Agg!} itself);
  stdout is exactly the four written paths, the same manifest contract
  as \passthrough{\lstinline!quadmath/scripts/lattice\_gallery.py!}.
\item
  Markdown:
  \passthrough{\lstinline!uv run python quadmath/scripts/validate\_markdown.py!}.
\end{itemize}

\subsection{Cross-references}\label{cross-references-5}

\begin{itemize}
\tightlist
\item
  The measured quantities behind every panel:
  \passthrough{\lstinline!17\_benchmarks\_statistics.md!}
  (\passthrough{\lstinline!src/quadmath/stats/benchmarks.py!},
  \passthrough{\lstinline!src/quadmath/stats/statistics.py!}).
\item
  The sibling gallery of the lattice layer, same thin-orchestrator and
  determinism contract:
  \passthrough{\lstinline!16\_lattice\_gallery.md!}
  (\passthrough{\lstinline!src/quadmath/viz/vis\_lattice.py!}).
\end{itemize}

\end{document}
